\documentclass[11pt,oneside,openany]{book}
% ============================================================
%  Preamble -- Economics of Money and Banking (P. Mehrling)
% ============================================================
\usepackage[T1]{fontenc}
\usepackage[utf8]{inputenc}
\usepackage{amsmath,amsthm,mathtools}
\usepackage{newpxtext,newpxmath} % provides the AMS symbols: do not load amssymb
\usepackage{bm}
\usepackage{microtype}
% Screen layout (read on a laptop, not printed): two parallel columns (paracol), the text column
% with the width it had on A4 (15 cm) and a 6.5 cm right column for notes, answers, feedback,
% examples, proofs, course notes, history and asides. Every right-column item starts level with the
% paragraph it belongs to; if the right column is still busy, the text waits (white space).
% The old A4 print layout was: a4paper,inner=2.4cm,outer=3.6cm,marginparwidth=3.0cm,marginparsep=0.3cm (twoside).
\usepackage[paperwidth=25cm,paperheight=29.7cm,top=2.6cm,bottom=2.6cm,left=2cm,textwidth=22cm,
  marginparsep=0.5cm,marginparwidth=6.5cm,headheight=23pt]{geometry}
\usepackage{paracol}
\setcolumnwidth{15cm/0.5cm,6.5cm}
% counters shared by the two columns (equations in a proof, examples, answers ... in the right column)
% (sectioning counters cannot be global with this kernel: they are copied to the right column by
% \synccounter before each switch, see \side@emit)
\globalcounter{equation}\globalcounter{footnote}\globalcounter{figure}
\usepackage{booktabs,array,tabularx,longtable}
\usepackage[dvipsnames]{xcolor}
\usepackage[shortlabels,inline]{enumitem}
\usepackage{tikz}
\usetikzlibrary{arrows.meta,positioning,decorations.pathreplacing,calc}
\usepackage{pgfplots}
\pgfplotsset{compat=1.18}
\usepackage[most]{tcolorbox}
\usepackage{fancyhdr}
\usepackage[colorlinks=true,linkcolor=NavyBlue,urlcolor=NavyBlue,citecolor=NavyBlue]{hyperref}
\usepackage[nameinlink]{cleveref}

\setlist{itemsep=2pt,topsep=4pt}
\setlength{\parskip}{3pt}
\setlength{\emergencystretch}{2em}
\hfuzz=5pt

% ---------- headers ----------
\pagestyle{fancy}
\fancyhf{}
\fancyhead[L]{\small\leftmark}
\fancyhead[R]{\small\makebox[\marginparwidth][r]{\rightmark\quad\thepage}}
\renewcommand{\headrulewidth}{0.3pt}
\fancypagestyle{plain}{\fancyhf{}\fancyfoot[C]{\small\thepage}\renewcommand{\headrulewidth}{0pt}}

% ---------- colours (bright on purpose: yellow notes, blue definitions, red theorems, green examples) ----------
\definecolor{defcol}{HTML}{1565C0}  % blue: definitions
\definecolor{thmcol}{HTML}{D32F2F}  % red: theorems, propositions, lemmas, corollaries
\definecolor{intcol}{HTML}{2E7D32}  % green: examples (also the green of the figures)
\definecolor{notecol}{HTML}{FBC02D} % yellow: margin notes and their answers
\definecolor{intucol}{HTML}{8E24AA} % purple: intuition
\definecolor{fincol}{HTML}{EF6C00}  % orange: finance
\definecolor{cscol}{HTML}{00897B}   % teal: case studies
\definecolor{keycol}{HTML}{C2185B}  % magenta: key formulas
\definecolor{exercol}{HTML}{3949AB} % indigo: exercises
\definecolor{srccol}{HTML}{546E7A}  % grey-blue: sources, course notes, remarks
\definecolor{histcol}{HTML}{B8860B} % bronze: history boxes (classical columns)
\definecolor{doubtcol}{HTML}{C62828}
\definecolor{ideacol}{HTML}{1565C0}
\definecolor{okcol}{HTML}{388E3C}
\colorlet{anscol}{notecol!85!black}

% ---------- table rows: a little air between lines in every table ----------
\AddToHook{env/tabular/begin}{\renewcommand{\arraystretch}{1.35}}
\AddToHook{env/tabularx/begin}{\renewcommand{\arraystretch}{1.35}}
\AddToHook{env/longtable/begin}{\renewcommand{\arraystretch}{1.35}}

% ---------- theorem-like environments (numbered by chapter), coloured headings ----------
\newtheoremstyle{coldef}{}{}{\normalfont}{}{\bfseries\color{defcol}}{.}{ }{}
\newtheoremstyle{colex}{}{}{\normalfont}{}{\bfseries\color{intcol}}{.}{ }{}
\newtheoremstyle{colexer}{}{}{\normalfont}{}{\bfseries\color{exercol}}{.}{ }{}
\newtheoremstyle{colthm}{}{}{\itshape}{}{\bfseries\color{thmcol}}{.}{ }{}
\newtheoremstyle{colrem}{}{}{\normalfont}{}{\itshape\color{srccol}}{.}{ }{}
\theoremstyle{coldef}
\newtheorem{definition}{Definition}[chapter]
\theoremstyle{colex}
\newtheorem{example}[definition]{Example}
\theoremstyle{colexer}
\newtheorem{exercise}{Exercise}[chapter]
\theoremstyle{colthm}
\newtheorem{theorem}[definition]{Theorem}
\newtheorem{proposition}[definition]{Proposition}
\newtheorem{lemma}[definition]{Lemma}
\newtheorem{corollary}[definition]{Corollary}
\theoremstyle{colrem}
\newtheorem{remark}[definition]{Remark}
\crefname{exercise}{Exercise}{Exercises}
\globalcounter{definition}\globalcounter{exercise}

% Right-column items are queued and written into the right column at the next point where the
% text column is between paragraphs at top level (end of the paragraph, end of a box, end of a
% list): there the columns are synchronised, the queued items are typeset in the right column,
% and the text continues from where it was.
\makeatletter
\newif\ifnote@inbox
\newif\ifside@inright
\def\note@queue{}
\def\side@pcname{paracol}
\newcommand{\side@emit}{\ifx\note@queue\@empty\else
  \let\note@tmp\note@queue\gdef\note@queue{}%
  \synccounter{chapter}\synccounter{section}\synccounter{subsection}\synccounter{subsubsection}%
  \switchcolumn*\side@inrighttrue\note@tmp\par\side@inrightfalse\switchcolumn\fi}
\newcommand{\note@flush}{\ifvmode\ifnote@inbox\else\ifside@inright\else
  \ifx\@currenvir\side@pcname\side@emit\fi\fi\fi\fi}
\newcommand{\side@put}[1]{\gappto\note@queue{#1}\note@flush}
\AddToHook{para/after}{\note@flush}
\tcbset{note queue/.code={\appto\kvtcb@afterbox{\note@flush}}}
% every chapter/appendix file is one paracol environment: it starts right after the \chapter
% line of the file (\begin{paracol}{2}) and is closed here (appendix A, only long tables, has none:
% longtable does not work inside paracol)
\AddToHook{include/end}{\par\note@flush\ifx\@currenvir\side@pcname\end{paracol}\fi}
\newcommand{\flushside}{\par\note@flush}
\makeatother
\tcbset{thmbox/.style={enhanced,breakable,colback=#1!7,colframe=#1!7,borderline west={3pt}{0pt}{#1},
  sharp corners,left=6pt,right=4pt,top=2pt,bottom=2pt,before skip=8pt,after skip=8pt,
  before upper={\csname note@inboxtrue\endcsname},note queue}}
\tcolorboxenvironment{definition}{thmbox=defcol}
\tcolorboxenvironment{theorem}{thmbox=thmcol}
\tcolorboxenvironment{proposition}{thmbox=thmcol}
\tcolorboxenvironment{lemma}{thmbox=thmcol}
\tcolorboxenvironment{corollary}{thmbox=thmcol}
% (examples are boxed in the right column, see \begin{example} below)
\tcolorboxenvironment{exercise}{thmbox=exercol}
\tcolorboxenvironment{remark}{thmbox=srccol}

% ---------- content boxes ----------
\tcbset{notebox/.style={enhanced,breakable,colback=#1!7,colframe=#1,
  fonttitle=\bfseries,coltitle=white,colbacktitle=#1,boxrule=0.6pt,arc=2pt,
  left=6pt,right=6pt,top=4pt,bottom=4pt,before skip=10pt,after skip=10pt,
  before upper={\csname note@inboxtrue\endcsname},note queue}}
% Box title: the box type in capitals, then the topic
\newcommand{\boxtitle}[2]{\MakeUppercase{#1}\ifx&#2&\else: #2\fi}
% Intuition: the "physicist's reading" of a result
\newtcolorbox{intuition}[1][]{notebox=intucol,title={\boxtitle{Intuition}{#1}}}
% In the markets: how a mechanism shows up in practice (orange, applications)
\newtcolorbox{finance}[1][]{notebox=fincol,title={\boxtitle{In the markets}{#1}}}
% Case study: the Financial Times article that opens most lectures, or a historical episode
\newtcolorbox{casestudy}[1][]{notebox=cscol,title={\boxtitle{Case study}{#1}}}
% Summary / key formulas
\newtcolorbox{keyformulas}{notebox=keycol,fontupper=\small,title={\boxtitle{Key formulas}{}}}
% Sources: which lectures, videos and files a chapter is built on
\newtcolorbox{sources}{notebox=srccol,fontupper=\footnotesize,title={\boxtitle{Sources for this chapter}{}},before upper=\raggedright}
\newcolumntype{L}{>{\raggedright\arraybackslash}X}
% Reading: summary of an assigned reading (Allyn Young, Minsky, Bagehot, ...), sources colour
\newtcolorbox{reading}[1][]{notebox=srccol,title={\boxtitle{Reading}{#1}}}

% ---------- T-accounts (balance sheets), the basic tool of the course ----------
%  \tacct{Bank}{Reserves\\Loans}{Deposits\\Net worth}
%  \tacct[3.2cm]{...}  sets the width of each side (default 2.6cm). Rows are separated by \\;
%  changes are written with a sign, e.g. {$+$Loans}{$+$Deposits}; \amt{852} right-aligns an amount.
\newcommand{\amt}[1]{\hspace{0.6em plus 1filll}#1}
\newcommand{\tacct}[4][2.6cm]{\begin{tabular}[t]{@{}p{#1}|p{#1}@{}}
  \multicolumn{2}{@{}c@{}}{\textbf{#2}}\\\hline
  \begin{tabular}[t]{@{}>{\raggedright\arraybackslash}p{#1}@{}}#3\end{tabular} &
  \begin{tabular}[t]{@{\,}>{\raggedright\arraybackslash}p{\dimexpr#1-0.2em}@{}}#4\end{tabular}
  \end{tabular}}
% non-floating figure (floats inside paracol columns can leave a page almost empty): same caption
% and numbering as figure
\makeatletter
\newenvironment{nfigure}{\par\medskip\begin{center}\def\@captype{figure}}{\end{center}\medskip}
\makeatother
% several T-accounts side by side, centred
\newenvironment{taccts}{\par\smallskip\begin{center}\renewcommand{\arraystretch}{1.15}}{\end{center}\smallskip}
% Note on course material (errata, missing lectures, differences 2013/2024): right column, see below
% Secondary material ("beyond the core"): can be skipped on a first reading
\newtcolorbox{secondary}[1][]{enhanced,breakable,colback=white,colframe=white,boxrule=0pt,
  borderline west={1.5pt}{0pt}{srccol!60},sharp corners,left=8pt,right=0pt,top=1pt,bottom=1pt,
  before skip=6pt,after skip=6pt,fontupper=\small,
  title={\footnotesize\sffamily$\diamond$\ \boxtitle{Beyond the core}{#1}},
  coltitle=srccol,colbacktitle=white,fonttitle=\bfseries,titlerule=0pt,toptitle=0pt,bottomtitle=0pt,
  before upper={\csname note@inboxtrue\endcsname},note queue}
% ============================================================
%  The right column: the reader's notes/answers/feedback in \tiny; examples, proofs, course notes,
%  history and asides in \scriptsize
% ============================================================
%  \note{...}     the reader's comment/doubt, NEW (not yet read by Claude)
%  \note*{...}    a request that has been carried out (green DONE tag)
%  \note+{...}    a comment that has been read and needs no answer (blue SEEN tag)
%  \begin{answer}{question}[topic] ... \end{answer}
%                 a doubt and its answer in one box: the question on top, then a dark yellow
%                 bar "ANSWER n: topic" and the answer; \answerpart{topic} starts a further
%                 answer in the same box (one per question); \label{ans:...} makes it citable
%  \feedback{...} a comment on one of the reader's own solutions (indigo)
%  \begin{history}[topic], \begin{aside}[topic], \begin{coursenote}[topic]
%  example and proof environments are also printed here (sideexample = example)
%  Inside a coloured box of the main text, margin items appear right after the box.
\makeatletter
\tcbset{sidebox/.style={enhanced,colback=#1!25,colframe=#1,boxrule=0.5pt,arc=1pt,
    left=3pt,right=3pt,top=2pt,bottom=2pt,before skip=0pt,after skip=0pt,
    fontupper=\tiny\sffamily\mdseries\upshape\raggedright,fonttitle=\tiny\sffamily\bfseries,
    before upper={\note@inboxtrue}}}
\newcommand{\side@tag}[2]{\colorbox{#1}{\textcolor{white}{\bfseries\tiny #2}}\par}
% Notes and answered notes are numbered per chapter like footnotes: a small dark-yellow number in
% the text marks the point where the note was written, the same number starts its box.
\newcounter{sidenote}[chapter]
\globalcounter{sidenote}
\newcommand{\side@numput}[1]{\stepcounter{sidenote}%
  \ifhmode\ifdim\lastkern=0.01pt\textsuperscript{\sffamily\bfseries\color{anscol!75!black},}\fi
    \textsuperscript{\sffamily\bfseries\color{anscol!75!black}\thesidenote}\kern0.01pt\fi
  \begingroup\edef\x{\endgroup\noexpand\side@put{\def\noexpand\side@num{\thesidenote}\unexpanded{#1}}}\x}
\newcommand{\side@numtag}{{\bfseries\color{anscol!75!black}\side@num}\enspace}
\NewDocumentCommand{\note}{s t+ +m}{\side@numput{\begin{tcolorbox}[sidebox=notecol]%
  \IfBooleanT{#1}{\side@tag{okcol}{DONE}}\IfBooleanT{#2}{\side@tag{defcol}{SEEN}}\side@numtag#3\end{tcolorbox}}}
\newcommand{\feedback}[1]{\side@put{\begin{tcolorbox}[sidebox=exercol]\side@tag{exercol}{FEEDBACK}#1\end{tcolorbox}}}
\newcounter{answer}[chapter]
\renewcommand{\theanswer}{\thechapter.\arabic{answer}}
\crefname{answer}{answer}{answers}
\globalcounter{answer}
\newcommand{\answerpart}[1]{\refstepcounter{answer}%
  \tcbsubtitle[colback=anscol,colframe=anscol,coltext=black,fontupper=\tiny\sffamily\bfseries,
    left=3pt,top=0.5pt,bottom=0.5pt,before skip=2pt,after skip=2pt]{ANSWER \theanswer\ifx&#1&\else: #1\fi}}
\NewDocumentEnvironment{answer}{m O{} +b}{\side@numput{\begin{tcolorbox}[sidebox=notecol,colback=notecol!12]%
  \side@numtag#1\answerpart{#2}#3\end{tcolorbox}}}{}
% Content boxes of the right column (examples, proofs, course notes, history, asides): \scriptsize,
% one step above the \tiny of the reader's notes
\tcbset{sidecontent/.style={sidebox=#1,colback=#1!6,breakable,
    fontupper=\scriptsize\mdseries\upshape\raggedright,fonttitle=\scriptsize\sffamily\bfseries,coltitle=white,colbacktitle=#1}}
\NewDocumentEnvironment{aside}{O{} +b}{\side@put{\begin{tcolorbox}[sidecontent=histcol,
  title={\boxtitle{Aside}{#1}}]#2\end{tcolorbox}}}{}
\NewDocumentEnvironment{coursenote}{O{} +b}{\side@put{\begin{tcolorbox}[sidecontent=srccol,
  title={\boxtitle{Course note}{#1}}]#2\end{tcolorbox}}}{}
% Examples and proofs keep their theorem-like look (numbering, labels, QED) but go to the right column
\let\origexample\example \let\endorigexample\endexample
\RenewDocumentEnvironment{example}{o +b}{\side@put{\begin{tcolorbox}[thmbox=intcol,
  fontupper=\scriptsize,before skip=0pt,after skip=0pt]%
  \IfValueTF{#1}{\begin{origexample}[#1]}{\begin{origexample}}#2\end{origexample}\end{tcolorbox}}}{}
\NewDocumentEnvironment{sideexample}{o +b}{\IfValueTF{#1}{\begin{example}[#1]}{\begin{example}}#2\end{example}}{}
\let\origproof\proof \let\endorigproof\endproof
% A proof must start level with the top of its theorem: a theorem directly followed by \begin{proof}
% is held back (\thm@pend) and typeset by the proof, right after the proof has been queued.
\def\thm@pend{}
\newtoks\thm@tk
\newcommand{\thm@typeset}{{\edef\@currenvir{\side@pcname}\begin{tcolorbox}[thmbox=thmcol]\expandafter\IfValueTF\expandafter{\thm@opt}%
  {\begin{orig\thm@cur}[\thm@opt]\the\thm@tk\end{orig\thm@cur}}{\begin{orig\thm@cur}\the\thm@tk\end{orig\thm@cur}}\end{tcolorbox}}}
\newcommand{\thm@peek}[3]{\def\thm@a{proof}\def\thm@b{#3}%
  \ifx\thm@a\thm@b\gdef\thm@pend{\thm@typeset}\else\thm@typeset\fi#1#2{#3}}
\newcommand{\thm@look}[2]{\@ifnextchar\begin{\thm@peek{#1{#2}}}{\thm@typeset#1{#2}}}
\RenewDocumentEnvironment{proof}{o +b}{\par{\edef\@currenvir{\side@pcname}\side@put{\begin{tcolorbox}[thmbox=thmcol,colback=thmcol!3,
  fontupper=\scriptsize,before skip=0pt,after skip=0pt]%
  \IfValueTF{#1}{\begin{origproof}[#1]}{\begin{origproof}}#2\end{origproof}\end{tcolorbox}}%
  \thm@pend}\gdef\thm@pend{}}{}
\@for\thm@n:=theorem,proposition,lemma,corollary\do{%
  \RemoveFromHook{env/\thm@n/before}[tcolorbox]\RemoveFromHook{env/\thm@n/after}[tcolorbox]%
  \expandafter\let\csname orig\thm@n\expandafter\endcsname\csname\thm@n\endcsname
  \expandafter\let\csname endorig\thm@n\expandafter\endcsname\csname end\thm@n\endcsname
  \expandafter\edef\csname thm@def@\thm@n\endcsname{\noexpand\RenewDocumentEnvironment{\thm@n}{o +b}%
    {\noexpand\global\noexpand\thm@tk{##2}\noexpand\gdef\noexpand\thm@cur{\thm@n}\noexpand\gdef\noexpand\thm@opt{##1}\noexpand\thm@look}{}}%
  \csname thm@def@\thm@n\endcsname}
\makeatother
% History: a fact from the history of finance/mathematics, framed like a classical
% temple front: a dentilled entablature (the title) on two Ionic columns (not breakable)
\newcommand{\histcolumn}[2]{% #1 = foot of the column, #2 = top of the capital (coordinate names)
  \fill[histcol] ([xshift=-6pt]#1) rectangle ([xshift=6pt,yshift=2.5pt]#1);
  \fill[histcol] ([xshift=-4.5pt,yshift=2.5pt]#1) rectangle ([xshift=4.5pt,yshift=4pt]#1);
  \fill[histcol!15] ([xshift=-3.5pt,yshift=4pt]#1) rectangle ([xshift=3.5pt,yshift=-5pt]#2);
  \foreach \dx in {-3.5,-1.2,1.2,3.5}
    \draw[histcol,line width=0.4pt] ([xshift=\dx pt,yshift=4pt]#1) -- ([xshift=\dx pt,yshift=-5pt]#2);
  \fill[histcol] ([xshift=-4.5pt,yshift=-5pt]#2) rectangle ([xshift=4.5pt,yshift=-2.5pt]#2);
  \foreach \dx in {-4.5,4.5} {
    \draw[histcol,line width=0.8pt,fill=white] ([xshift=\dx pt,yshift=-4.5pt]#2) circle (1.9pt);
    \fill[histcol] ([xshift=\dx pt,yshift=-4.5pt]#2) circle (0.6pt);}
  \fill[histcol] ([xshift=-6.5pt,yshift=-2.5pt]#2) rectangle ([xshift=6.5pt]#2);}
\makeatletter
\NewDocumentEnvironment{history}{O{} +b}{\side@put{\begin{tcolorbox}[histbox,title={\boxtitle{History}{#1}}]#2\end{tcolorbox}}}{}
\makeatother
\tcbset{histbox/.style={enhanced,colback=histcol!6,colframe=histcol,boxrule=0.8pt,sharp corners,
  fonttitle=\scriptsize\sffamily\bfseries,fontupper=\normalfont\scriptsize\raggedright,coltitle=white,colbacktitle=histcol,
  left=18pt,right=18pt,top=6pt,bottom=3pt,before skip=0pt,after skip=0pt,
  before upper={\csname note@inboxtrue\endcsname},
  overlay={
    \draw[histcol,line width=2.5pt,dash pattern=on 2.5pt off 2pt]
      ([yshift=-1.25pt]title.south west) -- ([yshift=-1.25pt]title.south east);
    \coordinate (hfL) at ([xshift=9pt]frame.south west);
    \coordinate (htL) at ([xshift=9pt,yshift=-2.5pt]title.south west);
    \coordinate (hfR) at ([xshift=-9pt]frame.south east);
    \coordinate (htR) at ([xshift=-9pt,yshift=-2.5pt]title.south east);
    \histcolumn{hfL}{htL}\histcolumn{hfR}{htR}}}}
% The reader's own solution to an exercise, placed right below the exercise
%   \begin{mysolution}[optional topic] ... \end{mysolution}
\newtcolorbox{mysolution}[1][]{notebox=exercol,colback=exercol!5,fontupper=\small,
  title={\boxtitle{My solution}{#1}}}
% A request about the notes themselves (collected in the "Requests" chapter)
%   \request{done|open|noted}{where it was written}{the request}{what was done}
\makeatletter
\newcommand{\request}[4]{\ifstrequal{#1}{done}{\def\req@col{okcol}}{\ifstrequal{#1}{open}{\def\req@col{thmcol}}{\def\req@col{defcol}}}%
  \begin{tcolorbox}[sidebox=notecol,colback=notecol!12,fontupper=\small\sffamily\raggedright,
  before skip=5pt,after skip=5pt]%
  \side@tag{\req@col}{\MakeUppercase{#1}}%
  {\itshape #2.}\ #3%
  \tcbsubtitle[colback=anscol,colframe=anscol,coltext=black,fontupper=\scriptsize\sffamily\bfseries,
    left=3pt,top=0.5pt,bottom=0.5pt,before skip=2pt,after skip=2pt]{WHAT WAS DONE}#4\end{tcolorbox}}
\makeatother



% ---------- video timestamps ----------
%  \ts{2}{5}{3:10}  -> lecture 2, video segment (part) 5, at 3:10 of that segment
%  Each tag is a link to the segment's YouTube video at that second. The IDs (playlist
%  "MOOC | Perry Mehrling - Economics of Money and Banking", ColumbiaLearn) are in videos.tex,
%  generated by ../tools/transcripts.py; an unknown segment or a non-numeric time prints the tag
%  without link.
\ExplSyntaxOn
\cs_new_protected:Npn \tsvideo #1#2 { \tl_const:cx { c__ts_#1_tl } { \tl_to_str:n {#2} } }
\int_new:N \l__ts_int
\seq_new:N \l__ts_seq
\cs_new_protected:Npn \tslink #1#2#3#4
  {
    \tl_if_exist:cTF { c__ts_#1@#2_tl }
      { \regex_match:nnTF { \A\d+(:\d+)*\Z } {#3} { \__ts_href:nnnn {#1} {#2} {#3} {#4} } {#4} }
      {#4}
  }
\cs_new_protected:Npn \__ts_href:nnnn #1#2#3#4
  {
    \int_zero:N \l__ts_int
    \exp_args:NNV \seq_set_split:Nnn \l__ts_seq \c_colon_str {#3}
    \seq_map_inline:Nn \l__ts_seq { \int_set:Nn \l__ts_int { 60*\l__ts_int + ##1 } }
    \exp_args:Ne \href { \tl_to_str:n { https://youtu.be/ } \tl_use:c { c__ts_#1@#2_tl } \tl_to_str:n { ?t= } \int_use:N \l__ts_int } {#4}
  }
\ExplSyntaxOff
% Generated by tools/transcripts.py: YouTube ID of every video segment (lecture@part)
\tsvideo{1@1}{KNEouYM5wRE} % The Big Picture
\tsvideo{1@2}{OZPvJa9_hWI} % Prerequisites?
\tsvideo{1@3}{zz7kk7s-v4k} % What is a Bank, a Shadow Bank, a Central Bank
\tsvideo{1@4}{gPEnV3ZC708} % Central Themes
\tsvideo{1@5}{1VstsTPHaB4} % Readings: Allyn Young
\tsvideo{2@1}{w6FEujzITdc} % The Eurocrisis,  Liquidity vs. Solvency
\tsvideo{2@2}{XB2DFuAP8SY} % Hierarchy of Financial Instruments
\tsvideo{2@3}{TLMhtj1FYs8} % Hierarchy of Financial Institutions
\tsvideo{2@4}{YXVaY-gIahw} % Dynamics of the Hierarchy
\tsvideo{2@5}{Vuy2kh-_Pic} % Discipline and Elasticity
\tsvideo{2@6}{5NnlQIz2UOw} % Hierarchy of Market Makers
\tsvideo{2@7}{nVLB1dAIix0} % Managing the Hierarchy
\tsvideo{3@1}{2lYU2TejA4Q} % Quantitative Easing and the Fed
\tsvideo{3@2}{qTgTTf-gHKQ} % Allyn Young:  Money and Economic Orthodoxy
\tsvideo{3@3}{Ikdgs9latBU} % National Banking System before the Fed
\tsvideo{3@4}{AGkH0w9hpSQ} % Civil War Finance, bonds and loans
\tsvideo{3@5}{ymZAkB8QI28} % Civil War Finance, legal tenders
\tsvideo{3@6}{p09FurGoiM0} % National Banking System, origins
\tsvideo{3@7}{NB0FVBHzK3E} % National Banking System, instability
\tsvideo{3@8}{QW0pB44FqaM} % Federal Reserve System, plan
\tsvideo{3@9}{yKtiYyM_8oQ} % Federal Reserve System, actual
\tsvideo{4@1}{2dIsnwMxXTM} % FT: Dealer of Last Resort
\tsvideo{4@2}{f6qXHMUMfzQ} % Reading: Hyman Minsky
\tsvideo{4@3}{Kce8ST2g3aY} % Payments: Money and Credit
\tsvideo{4@4}{W0IGZW0QOr4} % Payments: Discipline and Elasticity
\tsvideo{4@5}{8fJAUlpixpY} % The Survival Constraint
\tsvideo{4@6}{5Fa-9RTK8zc} % Sources and Uses Accounts
\tsvideo{4@7}{80_z5IMo9mk} % Payment Example: Money and Credit
\tsvideo{4@8}{zbRwRasqTuA} % Flow of Funds Accounts
\tsvideo{4@9}{1_QLYX-99gQ} % The Survival Constraint, redux
\tsvideo{4@10}{sdHLuC8fNfc} % Liquidity, Long and Short
\tsvideo{4@11}{_XwuGgyYrGs} % Financial Fragility, Flows and Stocks
\tsvideo{5@1}{iH0e1jmhQmw} % Martin Wolf on QE3
\tsvideo{5@2}{PyTGohPHqy0} % One Big Bank
\tsvideo{5@3}{BZiQD75mpN4} % Multiple Banks, a challenge
\tsvideo{5@4}{_FrU0BVeCJg} % Reading: Charles F. Dunbar
\tsvideo{5@5}{ySeQT37belQ} % Correspondent banking
\tsvideo{5@6}{_BJEASfyXtc} % Correspondent banking, system network
\tsvideo{5@7}{Uok9KslQdrA} % Clearinghouse, normal operations
\tsvideo{5@8}{TKWrmleNog4} % Clearinghouse, private lender of last resort
\tsvideo{5@9}{6ocg8pO-hpc} % Central Bank Clearing
\tsvideo{5@10}{LG9X5x78g8g} % Central Bank Cooperation
\tsvideo{6@1}{3qD6VOUjbb4} % European Bank Deleveraging
\tsvideo{6@2}{DIZ8o0_PYcI} % What are Fed Funds?
\tsvideo{6@3}{mPd3pjmUxHM} % Payment settlement versus Required Reserves
\tsvideo{6@4}{j0tHZ_9LiWQ} % Payment elasticity/discipline
\tsvideo{6@5}{9xE-QFUxZ88} % The Function of the Fed Funds Market
\tsvideo{6@6}{aS1MVEThBnU} % Payment versus Funding:  an example
\tsvideo{6@7}{yG9WhU3aTWs} % Brokers versus Dealers
\tsvideo{6@8}{6XJQQk8K9Yc} % Payments Imbalances and the Fed Funds Rate
\tsvideo{6@9}{V1TN_TcLp_E} % Secured versus Unsecured Interbank Credit
\tsvideo{6@10}{Jq32n5K5B20} % Required Reserves, redux
\tsvideo{7@1}{MpIiGdYN174} % The impact of QE3
\tsvideo{7@2}{NmlWNEONCNU} % Money Market Interest Rate Patterns
\tsvideo{7@3}{7WXxQGfY4WI} % What is repo?
\tsvideo{7@4}{hF5WI1LoGoY} % Repo in balance sheets
\tsvideo{7@5}{_TgBBNLZ3e8} % Comparison with Fed Funds
\tsvideo{7@6}{N2xCU-S6fRs} % Legal construction of repo
\tsvideo{7@7}{SdKS8AfX1-c} % Security dealers balance sheet
\tsvideo{7@8}{dI9jj1OPjWw} % Repo, modern finance, and the Fed
\tsvideo{7@9}{jIzc7XIrzZM} % Interest rate spreads: before the crisis
\tsvideo{7@10}{_PlS_EOWPcE} % Interest rate spreads:  after the crisis
\tsvideo{8@1}{f5fuDRx3qc0} % FT:  Ring-fencing and the Volcker Rule
\tsvideo{8@2}{C8J92oXhnU8} % The Eurodollar Market in Crisis
\tsvideo{8@3}{Pm7LBbAgX2w} % What are Eurodollars?
\tsvideo{8@4}{G9Kz8VkJG4Y} % Why is there a Eurodollar market?
\tsvideo{8@5}{gbr-8rSBP-Q} % Eurodollar as global funding market
\tsvideo{8@6}{VmSH0JIxYuA} % Liquidity challenge of Eurodollar banks
\tsvideo{8@7}{8WBk34CHXKw} % FRA as implicit swap of IOUs
\tsvideo{8@8}{7JXyFsCl4D8} % Forward Parity, Interest Rates, EH
\tsvideo{8@9}{iLSolU5-CSY} % Forward Parity, Exchange Rates, UIP
\tsvideo{8@10}{2V0Y_TVbEuE} % Forward rates aren't expected spot rates
\tsvideo{9@1}{u7x5KMzwBks} % FT:  Depreciation of Iran's currency
\tsvideo{9@2}{VOXNm8Iwj5w} % Reading: John Hicks
\tsvideo{9@3}{Ae98xdlw0l0} % Bagehot's World, wholesale money market
\tsvideo{9@4}{wEBqtutMGJ0} % Economizing on notes:  deposits, acceptances
\tsvideo{9@5}{o4m7B1GgWcI} % Managing cash flow:  discount, rediscount
\tsvideo{9@6}{9AS-ABRspqM} % Market rate of interest
\tsvideo{9@7}{CLLdBIWbom8} % Central Bank and bank rate
\tsvideo{9@8}{NoQypZ3Mn8Q} % The Bagehot Rule, origin of monetary policy
\tsvideo{9@9}{kA_zvcqRiBs} % Limits on central banking
\tsvideo{10@1}{CSIsOeWMGpc} % FT:  Asymmetric Credit Growth in Europe
\tsvideo{10@2}{d1VOy-cACgQ} % Liquidity, dealers, and inventories
\tsvideo{10@3}{uS-upCmbYzc} % Two-sided dealer basics
\tsvideo{10@4}{GHK79gTisSo} % Economics of the dealer function
\tsvideo{10@5}{YshcXyGeszE} % Leveraged dealer basics
\tsvideo{10@6}{vesU2BtPgsg} % Real world dealers
\tsvideo{10@7}{aWvi6L5-xtA} % The assumption of perfect liquidity
\tsvideo{11@1}{qwlh_Tx_eA8} % FT:  Money Market Mutual Funds
\tsvideo{11@2}{b_2PCNCSG_I} % Banks as money dealers, a puzzle
\tsvideo{11@3}{oEYxykCngTE} % Security dealers as money dealers
\tsvideo{11@4}{r6swYplUcb0} % Adapting Treynor to liquidity risk
\tsvideo{11@5}{1sJu8XNKMik} % Evolution of American banking
\tsvideo{11@6}{GQSw9BSms5Q} % The Fed in the Fed Funds market
\tsvideo{11@7}{ZAL0WPROmOQ} % Return to the initial puzzle
\tsvideo{12@1}{reGRRtHWz80} % FT:  Citibank and the SIVs
\tsvideo{12@2}{5oriNHybOsc} % The Art of Central Banking
\tsvideo{12@3}{PZStGPgCWKk} % Evolution of Monetary Policy:  1951-1987
\tsvideo{12@4}{SJt4soq71HA} % The Taylor Rule: 1987-2007
\tsvideo{12@5}{N8r9LGsysjE} % Monetary Transmission Mechanism
\tsvideo{12@6}{8qGW896JlHg} % Anatomy of a normal crisis
\tsvideo{12@7}{808nhtHPTac} % Anatomy of a serious crisis
\tsvideo{12@8}{L_2FGkdSzgE} % Should the Fed intervene or not?
\tsvideo{12@9}{xQ0HUa7FQwE} % The Fed as Dealer of Last Resort: 2007-2009
\tsvideo{13@1}{qxRSkeY4R4U} % FT: Autonomy of Bank of Japan
\tsvideo{13@2}{z1cHjKdbbWc} % Key Currencies
\tsvideo{13@3}{xvHm3VFwivY} % What is Money?  Chartalism versus Metallism
\tsvideo{13@4}{jti2J0aCK0U} % What is Money? Chartalism as a theory of money
\tsvideo{13@5}{AEMloDPuk60} % What is Money? Quantity Theory of Money
\tsvideo{13@6}{n5c25B6Rns0} % Purchasing Power Parity
\tsvideo{13@7}{5-XCz1rxm0A} % Metallism as a theory of money
\tsvideo{13@8}{fB73U2syJzQ} % A Money View of International Payments, FX Dealers
\tsvideo{13@9}{i1_mHyCU98o} % Chartallism, Metallism, and the Money View compared
\tsvideo{13@10}{K11XZuuJX10} % Private and Public Money:  A Hybrid System
\tsvideo{13@11}{uwdP4RX5h5A} % Hybridity in FX Market-making
\tsvideo{14@1}{_7MD7izX294} % FT: Costs of Japan's Monetary Policy
\tsvideo{14@2}{8GCkSBgguKA} % Reading:  Robert Mundell
\tsvideo{14@3}{3uizoDN3C-I} % Confrontation
\tsvideo{14@4}{b17LLX9FMHI} % Contradiction
\tsvideo{14@5}{NpC0tGD0KzY} % The Dollar System
\tsvideo{14@6}{Jk42hn5ms6U} % Flexible exchange
\tsvideo{14@7}{wyPdLvCehHE} % Global Financial Crisis
\tsvideo{15@1}{UuFPZnVORWQ} % FT: European money market funds
\tsvideo{15@2}{BHhrmhZ_Dic} % International transactions
\tsvideo{15@3}{u1tIaPGNHXg} % Dealer model for foreign exchange
\tsvideo{15@4}{7zyTAA-Mqdw} % Central banking, defense of domestic exchange
\tsvideo{15@5}{ysMHukKPGjI} % Bank of England, defense against external drain
\tsvideo{15@6}{PE6npYTbwaE} % Toward a theory of exchange
\tsvideo{16@1}{aHh4yoeGOQ4} % FT: High frequency trading
\tsvideo{16@2}{wnQGKtSg0r8} % Uncovered interest parity
\tsvideo{16@3}{yQf8Pey75b4} % FX dealers under the gold standard, redux
\tsvideo{16@4}{lQh9I8LEm0w} % Private FX dealing system
\tsvideo{16@5}{1yNnLIFnbVw} % Economics of the dealer function, speculative dealer
\tsvideo{16@6}{q6p_inIHGG8} % Economics of the dealer function, matched-book dealer
\tsvideo{16@7}{tGWcH6wZljU} % Digression:  Why do UIP and EH fail?
\tsvideo{16@8}{PTptQUV-t2A} % Central bank as FX dealer of last resort
\tsvideo{16@9}{vIDHMEj49Rg} % Reading:  McCauley on internationalization of renminbi
\tsvideo{17@1}{VKMUAy9AIuA} % FT: Shadow banking
\tsvideo{17@2}{gk4G48jjCb4} % Bagehot's World
\tsvideo{17@3}{pf7uOvxBeak} % The New World
\tsvideo{17@4}{Rt9ELlcFKQ4} % Funding liquidity versus market liquidity
\tsvideo{17@5}{hGTVBK-teK0} % Digression:  Schumpeter on banking
\tsvideo{17@6}{sqqqk31z9mU} % Payment versus funding
\tsvideo{17@7}{fJS_A8_wuJU} % Reading:  Gurley and Shaw
\tsvideo{17@8}{_g1FG8YVAUA} % Financial evolution
\tsvideo{17@9}{PsU7Y5O0-GU} % Banking evolution
\tsvideo{17@10}{-iMW1R4Ds1w} % Preview:  Central banking and shadow banking
\tsvideo{18@1}{hDsHi3wAirE} % Preview: FT: Argentina in court to fight debt ruling
\tsvideo{18@2}{Zy4u9k7aw24} % Banking as advance clearing
\tsvideo{18@3}{4YVNeG8JHHM} % Forwards versus futures
\tsvideo{18@4}{sS7idgU2IoY} % Futures contracts, fluctuations II
\tsvideo{18@5}{mFlb0CS2hIU} % Futures contracts, fluctuations II
\tsvideo{18@6}{oo-kTCTcUEM} % Cash and carry arbitrage, defined
\tsvideo{18@7}{Q-N1s8zcBWQ} % Cash and carry arbitrage, explained as liquidity risk
\tsvideo{18@8}{UBbbt6ZrW3Y} % Cash and carry arbitrage as counterparty risk
\tsvideo{18@9}{ySY8hRUF34k} % Cash and carry arbitrage, as natural banking business
\tsvideo{19@1}{WXOw8n1YBZU} % FT: Sovereign debt crises
\tsvideo{19@2}{dlaIO_35Vp8} % Reading:  FOMC Report (1952)
\tsvideo{19@3}{snG05OFavjU} % Treasury-swap spread, a puzzle
\tsvideo{19@4}{nCrxDIuCZ2Y} % What is a swap
\tsvideo{19@5}{B57VC6qKmA0} % Why swap? An example from Stigum
\tsvideo{19@6}{HjeQRfq8WE4} % Market making in swaps
\tsvideo{19@7}{T9TfVmhp3EY} % Money market swaps, example
\tsvideo{19@8}{0WiQpgTIvRo} % Life in Arbitrage Land
\tsvideo{19@9}{RYsIpeGpQms} % Treasury-swap spread
\tsvideo{20@1}{mmgYLekyw3Y} % FT: Internationalization of the Euro
\tsvideo{20@2}{ypmwhTFWFfU} % Credit indices
\tsvideo{20@3}{47gfinIdtFo} % Fischer Black (1970), risk-free security
\tsvideo{20@4}{UD2e2FgC-vs} % What is a Credit Default Swap (CDS)
\tsvideo{20@5}{8I9IN6i0kOY} % Corporate bonds
\tsvideo{20@6}{t2o_EuZv43U} % CDS pricing
\tsvideo{20@7}{2TgKw3ulHYA} % Market making in CDS
\tsvideo{20@8}{qnC-_nLfSHE} % Example: Negative basis trade and Liquidity Risk
\tsvideo{20@9}{ynwQ5-gJUcQ} % Example: Private backstop of marketmaking in CDS
\tsvideo{20@10}{jIfYYYztlm4} % Example: Synthetic CDO as Collateral Prepayment
\tsvideo{21@1}{MkQ1Qb0d-pc} % FT: Regulation of shadow banking
\tsvideo{21@2}{UDkkRmqr2uc} % Shadow banking vs traditional banking
\tsvideo{21@3}{saW6Xo1rrIA} % Liquidity and solvency backstops
\tsvideo{21@4}{MsJ1GOM33OU} % Global dimension
\tsvideo{21@5}{1vIUlJM-uuY} % Evolution of modern finance
\tsvideo{21@6}{KRYE7SvR8gQ} % What is shadow banking?
\tsvideo{21@7}{T3iRxwBzwAM} % Backstopping the market makers
\tsvideo{21@8}{L1bDEFGQrCQ} % Regulation of systemic risk
\tsvideo{21@9}{dTIbBI9yn2Q} % Regulation of Collateral and Payment
\tsvideo{21@10}{UOM5uixfcsk} % Private backstop, and public
\tsvideo{22@1}{xhNA3VKGCYk} % FT: Future of banking
\tsvideo{22@2}{Wkq8PF0Zv1Q} % Three world views
\tsvideo{22@3}{PiNVe1Lf7I4} % Economics View: Commodity Exchange
\tsvideo{22@4}{HgUXoGWvq34} % Finance View: Risk
\tsvideo{22@5}{B0TlLwf2Mz4} % The education of Fischer Black
\tsvideo{22@6}{exdVFqH8tHM} % Steps from finance view to money view
\tsvideo{22@7}{NQ0a1q5isuo} % A money view of economics and finance

\makeatletter % consecutive tags may break between them, but never before punctuation
\DeclareRobustCommand{\ts}[3]{\mbox{\footnotesize\sffamily\color{srccol}\tslink{#1}{#2}{#3}{\color{srccol}[$\blacktriangleright$\,L#1.#2\,@\,#3]}}\@ifnextchar\ts{\allowbreak}{}}
% link to a whole segment (no time), e.g. in the Sources box: \seg{2}{5}{text}
\newcommand{\seg}[3]{\tslink{#1}{#2}{0}{#3}}
\makeatother

% ---------- Mehrling's written lecture notes ----------
%  \mn{19}{2} -> page 2 of Mehrling's notes for lecture 19, linked to the PDF on sites.bu.edu/perry
%  (local copies in materials/mehrling-notes/). \mnfile{19}{...} lists the file names.
\newcommand{\mnfile}[2]{\expandafter\def\csname mn@file@#1\endcsname{#2}}
\mnfile{1}{Lec-01-The-Four-Prices-of-Money}\mnfile{2}{Lec-02-The-Natural-Hierarchy-of-Money}
\mnfile{3}{Lec-03-Money-and-the-State-Domestic}\mnfile{4}{Lec-04-The-Money-View-Micro-and-Macro}
\mnfile{5}{Lec-05-The-Central-Bank-as-a-Clearinghouse}\mnfile{6}{Lec-06-Federal-Funds-Final-Settlement}
\mnfile{7}{Lec-07-Repos-Postponing-Settlement}\mnfile{8}{Lec-08-Eurodollars-Parallel-Settlement}
\mnfile{9}{Lec-09-The-World-that-Bagehot-Knew}\mnfile{10}{Lec-10-Dealers-and-Liquid-Security-Markets}
\mnfile{11}{Lec-11-Banks-and-The-Market-for-Liquidity}\mnfile{12}{Lec-12-Lender-or-Dealer-of-Last-Resort}
\mnfile{13}{Lec-13-Chartalism-Metallism-and-Key-Currencies}\mnfile{14}{Lec-14-Money-and-the-State-International}
\mnfile{15}{Lec-15-Banks-and-Global-Liquidity}\mnfile{16}{Lec-16-Foreign-Exchange}
\mnfile{17}{Lec-17-Direct-and-Indirect-Finance}\mnfile{18}{Lec-18-Forwards-and-Futures}
\mnfile{19}{Lec-19-Interest-Rate-Swaps}\mnfile{20}{Lec-20-Credit-Derivatives}
\mnfile{21}{Lec-21-Shadow-Banking-Central-Banking}\mnfile{22}{Lec-22-Touching-the-Elephant-three-views}
\DeclareRobustCommand{\mn}[2]{\mbox{\footnotesize\sffamily\href{https://sites.bu.edu/perry/files/2019/01/\csname mn@file@#1\endcsname.pdf\#page=#2}{\color{srccol}[Notes\,#1\,p.#2]}}}

% ---------- maths shortcuts ----------
\newcommand{\R}{\mathbb{R}}
\newcommand{\N}{\mathbb{N}}
\newcommand{\E}{\mathbb{E}}
\newcommand{\Prob}{\mathbb{P}}
\newcommand{\Q}{\mathbb{Q}}
\DeclareMathOperator{\Var}{Var}
\DeclareMathOperator{\Cov}{Cov}
\DeclareMathOperator{\Corr}{Corr}
\newcommand{\dd}{\,\mathrm{d}}
\newcommand{\ind}{\mathbf{1}}
\newcommand{\eqd}{\stackrel{d}{=}}

% ---------- file references ----------
% file paths: breakable at / . - _ (long reading file names overflowed the right column)
\DeclareRobustCommand{\file}[1]{{\small\nolinkurl{#1}}}
\appto\UrlBreaks{\do\-\do\_}


% Compile only some chapters while working on them, e.g.:
% \includeonly{chapters/ch01-four-prices}

\begin{document}

\frontmatter
\begin{titlepage}
  \centering
  \vspace*{3cm}
  {\Huge\bfseries Economics of Money and Banking\par}
  \vspace{1cm}
  {\Large Lecture notes\par}
  \vspace{0.5cm}
  {\large based on the course by Perry G.~Mehrling\\(Barnard College, Columbia University;
  Coursera and INET)\par}
  \vspace{2cm}
  {\large Personal study notes of Alessio Martini\par}
  \vfill
  \begin{minipage}{0.85\textwidth}\small
    These notes are a personal, non-commercial study aid derived from the video lectures of
    \emph{Economics of Money and Banking} by Perry G.~Mehrling, filmed at Barnard College
    (Columbia University) and published by ColumbiaLearn on YouTube, on Coursera, and by the
    Institute for New Economic Thinking. Lecture content is paraphrased and re-derived from the
    (automatic) transcripts of the videos, not transcribed; any mistakes are mine. All rights on
    the original lectures belong to their authors.
  \end{minipage}
  \vspace{1cm}
\end{titlepage}

\chapter*{How to use these notes}
\addcontentsline{toc}{chapter}{How to use these notes}
\begin{paracol}{2}

Each chapter covers one \emph{lecture} of the course (22 in all), and its sections follow the short
video segments into which every lecture was cut. Every chapter opens with a \textbf{Sources} box
listing the segments (each one a link to its video), the reading and the newspaper article
discussed in class. Part~\ref{part:intro} (lectures 1--12) is the ``Introduction'' of the course:
the payment system and the money market. Part~\ref{part:advanced} (lectures 13--22) extends the
same ``money view'' to foreign exchange and capital markets.

\paragraph{Boxes.}
\begin{itemize}
  \item \textcolor{defcol}{\textbf{Definitions}} (blue), \textcolor{thmcol}{\textbf{propositions}} (red),
        \textcolor{exercol}{\textbf{exercises}} and \textcolor{srccol}{\textbf{remarks}} are marked with a
        coloured bar on the left. A definition box titled \emph{New concepts} marks the first time a concept appears.
  \item \textcolor{keycol}{\textbf{Key formulas}} boxes collect what to remember.
  \item \textcolor{intucol}{\textbf{Intuition}} boxes give the reading of a mechanism in a physicist's language.
  \item \textcolor{fincol}{\textbf{In the markets}} boxes say where a mechanism shows up in practice.
  \item \textcolor{cscol}{\textbf{Case study}} boxes summarise the \emph{Financial Times} article that
        opens most lectures, or a historical episode.
  \item \textcolor{srccol}{\textbf{Reading}} boxes summarise the assigned reading of the week.
  \item \textbf{Beyond the core} blocks (grey bar on the left, smaller type) hold material that can be
        skipped on a first reading.
\end{itemize}

\paragraph{Balance sheets.} The basic tool of the course is the balance sheet, drawn as a ``T-account'':
assets on the left, liabilities (and net worth) on the right. A transaction is written as the
\emph{changes} it makes to the balance sheets of everyone involved:
\begin{taccts}
\tacct{Bank}{$+$Loan}{$+$Deposit}\qquad
\tacct{Borrower}{$+$Deposit}{$+$Loan}
\end{taccts}

\paragraph{The right column.} The page is laid out for a screen: the text is always on the left and a
column on the right holds \textcolor{intcol}{\textbf{examples}}, \textcolor{srccol}{\textbf{course notes}},
\textcolor{histcol}{\textbf{history}} boxes and \textcolor{histcol}{\textbf{asides}} (curiosities), your
\textcolor{notecol!80!black}{\textbf{notes}} with their \textcolor{anscol}{\textbf{answers}}, and
\textcolor{exercol}{\textbf{feedback}} on your solutions.

\paragraph{Video timestamps.} A tag such as \ts{2}{5}{3:10}\ means: lecture~2, video segment~5,
at minute 3:10 of that segment. Every tag is a link to the video at that second.

\paragraph{Your annotations.} While reading, write directly in the chapter files. The macro
\verb|\note{...}| prints a comment in the right column next to the text where it is written: a doubt,
an idea, a reminder. When a doubt is answered, the note and its answer become one box in the right
column. Requests about the notes themselves are collected in the chapter \emph{Requests about these
notes}, which also lists how read and unread notes are marked.

\flushside
\end{paracol}

% !TEX root = ../main.tex
\chapter*{Requests about these notes}\label{ch:requests}
\addcontentsline{toc}{chapter}{Requests about these notes}
\markboth{Requests about these notes}{}
\begin{paracol}{2}

Your margin notes are of two kinds: notes on the \emph{content} (a doubt, an idea about a mechanism),
which stay next to the text they refer to, and requests about the \emph{notes themselves} (layout,
boxes, conventions), which do not depend on where they were written. The second kind is collected here,
with what was done about it, so that the chapters read more smoothly. Write new ones anywhere with
\verb|\note{...}| and they will be moved here, or add them directly at the end of this list.

\paragraph{How the margin notes are marked.}
\begin{itemize}
  \item \verb|\note{...}|: a \emph{new} note, not yet read. ``Read my notes'' means these, and only these.
  \item A note that has been answered becomes one box with the question on top and a dark yellow
        \textbf{ANSWER} bar below it (\verb|\begin{answer}{question}[topic]| \dots\ \verb|\end{answer}|). A note with
        several questions gets one ANSWER bar per question. A question asked a second time gets a short
        answer that links to the first one.
  \item \verb|\note*{...}|: a request that has been carried out (green DONE tag).
  \item \verb|\note+{...}|: a comment that has been read and needs no answer (blue SEEN tag).
  \item \verb|\feedback{...}|: a comment on one of your own solutions (indigo).
\end{itemize}

\section*{The list}

\request{done}{All chapters, exercises}{esercizi brevi/a sé stanti lungo il testo, quelli collegati o complessi alla fine}{Standalone exercises now follow the section they practise (right before the next heading); connected or integrative exercises stay together at the end of the chapter, in the original order. In chapters where none were left, the final Exercises section was removed.}


\tableofcontents

\mainmatter
\part{Introduction: payments and the money market}\label{part:intro}
% !TEX root = ../main.tex
\chapter{The Four Prices of Money}\label{ch:fourprices}
\begin{paracol}{2}

\begin{sources}
\begin{tabularx}{\linewidth}{@{}l l L@{}}
\toprule
\textbf{Segment} & \textbf{Length} & \textbf{Content}\\
\midrule
\seg{1}{1}{L1.1 The Big Picture} & 19:26 & Where the course comes from: Stigum, the \emph{Financial Times}, the history of monetary and financial thought; the ``money view''; dialogue between central banking and modern finance.\\
\seg{1}{2}{L1.2 Prerequisites?} & 7:22 & IS/LM and the Fisher diagram: the two textbook ways of teaching money and banking, and why this course starts elsewhere.\\
\seg{1}{3}{L1.3 What is a Bank, a Shadow Bank, a Central Bank} & 12:10 & Balance sheets as the analytical tool; solvency versus liquidity; shadow banks; the Fed's balance sheet in 2012.\\
\seg{1}{4}{L1.4 Central Themes} & 13:09 & Plan of the course; the four prices of money; Bagehot's \emph{Lombard Street}; ``all banking is a swap of IOUs''.\\
\seg{1}{5}{L1.5 Reading: Allyn Young} & 3:10 & Presentation of the first reading.\\
\bottomrule
\end{tabularx}

\smallskip
\textbf{Files.} Mehrling's lecture notes \file{money-lecture_notes-Lec 01--The Four Prices of Money.pdf};
captions \file{materials/week1/srt/1 - *.srt}; transcripts \file{materials/transcripts/L01-P0*.txt}.
\textbf{Reading.} Allyn Young, four chapters from \emph{The Book of Popular Science} (1924), file
\file{Allyn Young.pdf}, with the study questions \file{Allyn Young(study).pdf}; discussed in
\cref{sec:fp-young} and again in lecture~3.
\end{sources}

\section*{Motivation}
The first lecture is an orientation. It says where the course stands among the ways money and banking
are usually taught, introduces the one analytical tool that will be used everywhere (the balance sheet),
and names the four prices that the whole course tries to explain:
\begin{itemize}
  \item \textbf{par}, the price of one kind of money in terms of another kind today;
  \item \textbf{interest}, the price of money today in terms of money tomorrow;
  \item the \textbf{exchange rate}, the price of domestic money in terms of foreign money;
  \item the \textbf{price level}, the price of money in terms of goods.
\end{itemize}
The thread of the lecture is that the standard textbooks describe the banking system of 1950, while
the system of today is \emph{global} and \emph{market-based}, and the problems it runs into are
problems of \emph{liquidity}, not only of solvency. For a physicist, the useful attitude is that of an
experimentalist entering a new lab: the course does not start from a model and apply it, but starts
from what practitioners actually do and tries to extract the regularities \ts{1}{2}{4:41}.

% ------------------------------------------------------------------
\section{The big picture}\label{sec:fp-bigpicture}

\paragraph{Where the knowledge comes from.} Mehrling has never traded in the money markets: his
knowledge is ``book learning'', but from practitioners' books, above all Marcia Stigum's
\emph{The Money Market}, and from the \emph{Financial Times}, rather than from academic journals
\ts{1}{1}{0:29}. His academic work is in the history of monetary and financial economics, in two
books: \emph{The Money Interest and the Public Interest}, a history of American monetary thought, and
\emph{Fischer Black and the Revolutionary Idea of Finance}, a history of the rise of modern finance
\ts{1}{1}{0:50}. These are the intellectual foundations of the course, and they explain why the
readings include authors such as Allyn Young, Walter Bagehot, Ralph Hawtrey and Edward Shaw
\ts{1}{1}{1:10}.\begin{aside}[Mehrling's books]
Perry Mehrling (Barnard College, Columbia University, at the time of the lectures; later Boston
University) wrote \emph{The Money Interest and the Public Interest: American Monetary Thought,
1920--1970} (Harvard University Press, 1997), \emph{Fischer Black and the Revolutionary Idea of
Finance} (Wiley, 2005) and \emph{The New Lombard Street: How the Fed Became the Dealer of Last Resort}
(Princeton University Press, 2011), the book that accompanies this course.
\end{aside}

\paragraph{A course about the present.} Although its foundations are historical, this is not a course
in the history of economic thought: it uses that history as a frame to understand the world as it is
now \ts{1}{1}{1:31}. The big fact to understand is \emph{financial globalization} (or ``global
financialization''). The lesson of the history of monetary thought is that every generation has to
work out for itself the world in front of it: some ideas carry over from the past, but the system is
not constant, and today it is ``a strange new world'' \ts{1}{1}{1:51}. This is why Mehrling does not use
the standard money and banking textbooks: they evolved from the textbooks of 1950, and a bank in 1950
was a different thing from a bank today \ts{1}{1}{2:53}. Better to put them in the library and start
again from the real world, accepting that the result is a little ``looser'' but attached to reality
\ts{1}{1}{3:14}.

\paragraph{Listening to the bankers.} Being attached to reality means listening to practitioners and
trying to understand their world \ts{1}{1}{3:35}. Practitioners know their own piece very well (repo,
Eurodollars, \dots) but have no picture of how the system as a whole works; an academic can add value
by noticing the patterns and the connections between the pieces \ts{1}{1}{3:55}. This is not easy:
there will be no ``take the first derivative and set it equal to zero'' in this course \ts{1}{1}{4:36}.

\begin{intuition}[Reading Stigum as an anthropologist]
Stigum's \emph{The Money Market} is the desk reference of people who trade the money markets, ``the
bible'' of the field \ts{1}{1}{4:56}. Mehrling's advice is to read it as an anthropologist reads the
sacred text of a foreign tribe \ts{1}{1}{5:58}: not to learn to trade, but to understand how the
members of the tribe think about their world. Read for every detail, it ``can blow your mind''; read for
the way of thinking, it is a window into a rather special world \ts{1}{1}{6:19}.
\end{intuition}\begin{aside}[Stigum]
Marcia Stigum (1934--2005) first published \emph{The Money Market: Myth, Reality, and Practice} in 1978;
the fourth edition (with Anthony Crescenzi, McGraw-Hill, 2007) is the one mentioned in class, with a
new chapter on Eurodollars \ts{1}{1}{5:16}. Mehrling recalls that when he started teaching the course,
Stigum was the whole text, with no lecture notes.
\end{aside}

\paragraph{Reading the Financial Times.} In a sense the whole course is about learning to read the
\emph{Financial Times} (FT) and to ``translate it into English'' \ts{1}{1}{7:21}; the final exam contains
exactly that exercise, an article to be explained to one's roommate \ts{1}{1}{7:42}. Every lecture
after this one starts with an article of the day. Three examples from the paper of that day show the
kind of question the course is about \ts{1}{1}{8:23}:
\begin{enumerate}[(a)]
  \item German banks split over the role of the European Central Bank (ECB): should the ECB buy the
        government debt of Spain or Italy to support its price? What does it even mean for a central
        bank to buy the debt of a state \ts{1}{1}{8:43}?
  \item Chinese steel mills brace for a slowdown in demand. The monetary angle: in China, inventories of
        iron ore had been used as \emph{collateral} for borrowing; as the banks call in the loans, the
        borrowers dump the collateral, the beginning of a financial crisis \ts{1}{1}{9:04}.
  \item Foreign exchange brokers back plans of the trading platform EBS to rein in ``predatory'' traders:
        a battle between algorithmic high-frequency traders (``cyborgs'') and human dealers, who lose
        money to them \ts{1}{1}{9:45}.
\end{enumerate}

\subsection{The money view and its ancestors}
Asked which school of economics he belongs to (saltwater or freshwater, Keynesian or monetarist),
Mehrling answers that he identifies with the \emph{money view}, a name he invented himself
\ts{1}{1}{10:47}. It has two lines of intellectual ancestors \ts{1}{1}{11:08}, shown in
\cref{fig:fp-ancestors}.
\begin{itemize}
  \item \textbf{The American line.} It starts with Allyn Young, because central banking in the United
        States starts in his time (the Fed was founded in 1913), when there was no tradition of central
        banking thought to draw on \ts{1}{1}{11:49}. Then Alvin Hansen and Edward Shaw (the three
        generations of Mehrling's first book), and Hyman Minsky, who would have been the missing last
        chapter of that book \ts{1}{1}{12:09}.
  \item \textbf{The British line.} Britain had a central bank for a long time and was the centre of the
        world economy for a long time, so it has a long tradition of thinking about central banking
        \ts{1}{1}{12:50}: Walter Bagehot, Ralph Hawtrey, Richard Sayers (who taught at the London School of
        Economics, LSE) and Charles Goodhart \ts{1}{1}{13:10}. Almost all of them talked to bankers and
        central bankers: Goodhart spent most of his career at the Bank of England \ts{1}{1}{13:30}, and to
        teach at the LSE is to be a mile from the City of London, as teaching in Manhattan is to be a mile
        from Wall Street \ts{1}{1}{13:51}.
\end{itemize}
Some of them called themselves Keynesians (Hansen, in a way Minsky), others leaned towards monetarism
(Shaw, Hawtrey); Mehrling reads them as one continuous stream of thought that kept its eye on the real
world of its own time \ts{1}{1}{14:11}.

\begin{nfigure}
\begin{tikzpicture}[x=1cm,y=1cm,font=\small,
  p/.style={draw=defcol,fill=defcol!8,rounded corners=2pt,inner sep=3pt,align=center,minimum width=2.5cm},
  q/.style={draw=thmcol,fill=thmcol!8,rounded corners=2pt,inner sep=3pt,align=center,minimum width=2.5cm},
  arr/.style={-{Stealth[length=4pt]},semithick}]
  \node[font=\bfseries\small,defcol] at (0,0.8) {American line};
  \node[p] (y) at (0,0) {Allyn Young\\\scriptsize 1876--1929};
  \node[p] (h) at (0,-1.2) {Alvin Hansen\\\scriptsize 1887--1975};
  \node[p] (s) at (0,-2.4) {Edward Shaw\\\scriptsize 1908--1994};
  \node[p] (m) at (0,-3.6) {Hyman Minsky\\\scriptsize 1919--1996};
  \draw[arr] (y)--(h); \draw[arr] (h)--(s); \draw[arr] (s)--(m);
  \node[font=\bfseries\small,thmcol] at (8,0.8) {British line};
  \node[q] (b) at (8,0) {Walter Bagehot\\\scriptsize 1826--1877};
  \node[q] (ha) at (8,-1.2) {Ralph Hawtrey\\\scriptsize 1879--1975};
  \node[q] (sa) at (8,-2.4) {Richard Sayers\\\scriptsize 1908--1989};
  \node[q] (g) at (8,-3.6) {Charles Goodhart\\\scriptsize b.~1936};
  \draw[arr] (b)--(ha); \draw[arr] (ha)--(sa); \draw[arr] (sa)--(g);
  \node[draw=intucol,fill=intucol!10,very thick,rounded corners=3pt,align=center,inner sep=5pt] (mv) at (4,-1.8)
    {\textbf{the money view}\\\scriptsize (Mehrling)};
  \draw[arr,intucol,dashed] (y.east) to[out=0,in=150] (mv);
  \draw[arr,intucol,dashed] (m.east) to[out=0,in=-150] (mv);
  \draw[arr,intucol,dashed] (b.west) to[out=180,in=30] (mv);
  \draw[arr,intucol,dashed] (g.west) to[out=180,in=-30] (mv);
  \node[draw=fincol,fill=fincol!10,rounded corners=2pt,align=center,inner sep=3pt] (fb) at (4,-4.4)
    {Fischer Black, modern finance\\\scriptsize ``no separate theory of money''};
  \draw[{Stealth[length=4pt]}-{Stealth[length=4pt]},semithick,fincol] (fb) -- node[right,font=\scriptsize]{dialogue} (mv);
\end{tikzpicture}
\caption{The intellectual ancestors of the money view \ts{1}{1}{11:28}, and its opponent in the
dialogue that the course sets up \ts{1}{1}{16:14}.}\label{fig:fp-ancestors}
\end{nfigure}

\subsection{The challenge of modern finance}
Mehrling also learned modern finance, through his biography of Fischer Black. Finance holds that there
is no such field as monetary economics \ts{1}{1}{14:53}: in an efficient market every asset trades at its
efficient price, so there is nothing for liquidity to distort. In ``straight-ahead'' finance there is no
\emph{fire sale}: if a bank had to sell its loans, a price just a tiny bit below their true value would be
enough for someone to snatch them up, so prices never move far from value and there are no runs on banks
\ts{1}{1}{15:13}. The logical end of this view is that there is no reason for central banks to exist
\ts{1}{1}{15:54}.

The course is, in Mehrling's words, a \emph{dialogue} between the classical central banking perspective and
the modern finance perspective, two ideas that seem contradictory, rubbed against each other to make
sparks for new economic thinking \ts{1}{1}{16:14}. The opposition is not Keynesian versus monetarist
\ts{1}{1}{16:55}.

\paragraph{Where Mehrling stands.} Pressed by the student, he states the position he will defend for the
rest of the course \ts{1}{1}{17:38}: there \emph{is} a role for central banks in setting bounds on the
system, a modern version of the role that Bagehot had in mind in 1873 (``lend freely at a high rate against
good security'') \ts{1}{1}{17:59}. That rule does not quite work in the modern world, and the subtitle of his
book, \emph{How the Fed Became the Dealer of Last Resort}, signals the update. Financial markets do not work
perfectly; they run into problems, and banking is all about these problems, which are \emph{liquidity}
problems rather than problems of asymmetric information \ts{1}{1}{18:20}. The most compelling work on them
comes not from academic journals, ``fighting old battles'', but from the Bank of England and the Bank for
International Settlements (BIS), the central bankers' bank \ts{1}{1}{19:01}.\begin{history}[Bagehot's \emph{Lombard Street} (1873)]
Walter Bagehot (1826--1877), editor of \emph{The Economist}, published \emph{Lombard Street: A Description of
the Money Market} in 1873, after the panic that followed the failure of the discount house Overend, Gurney \&
Co.\ in 1866. His prescription for the Bank of England in a panic is usually summarised as ``lend freely, at a
high rate, against good collateral''; in his own words, loans ``should only be made at a very high rate of
interest'' and ``on all good banking securities, and as largely as the public ask for them''. Lombard Street,
in the City of London, takes its name from the Lombard (northern Italian) merchants and bankers who settled
there in the Middle Ages.\par
\emph{References:} W.~Bagehot, \emph{Lombard Street}, London: Henry S.~King, 1873, ch.~VII; Mehrling,
\emph{The New Lombard Street}, 2011, ch.~1.
\end{history}

% ------------------------------------------------------------------
\section{Prerequisites: two textbook approaches}\label{sec:fp-prereq}

The formal prerequisites are intermediate macroeconomics and intermediate microeconomics
\ts{1}{2}{0:08}. The course makes little use of them, but other money and banking courses start from them,
and at the end of the course the money view will be connected to both \ts{1}{2}{0:29}. It is worth seeing what
they are, because they are what the course is \emph{not} doing.

\paragraph{Intermediate macro: IS/LM.} The macro story determines national income $Y$ and the
\emph{nominal} interest rate $i$ at the crossing of two curves \ts{1}{2}{1:15}:
\begin{itemize}
  \item the IS curve (``investment = saving''): the pairs $(Y,i)$ at which the market for goods clears;
        a higher rate depresses investment, so the curve slopes down;
  \item the LM curve (``liquidity preference = money supply''): the pairs at which the demand for money
        equals the supply; higher income raises the demand for money, so the rate must rise to keep it
        equal to a given supply, and the curve slopes up.
\end{itemize}
Monetary policy is then a shift of the LM curve. A whole course can be built on the foundations of the LM
curve, the Taylor rule, or the argument that with inflation targeting the LM curve should be horizontal
\ts{1}{2}{1:35}. This course instead questions the very notion of a stable money demand curve and money supply
curve, at a more fundamental level \ts{1}{2}{1:55}.

\paragraph{Intermediate micro: the Fisher diagram.} The micro story determines the \emph{real} interest rate in a
two-period general equilibrium \ts{1}{2}{2:16}, shown in \cref{fig:fp-prereq}. With wealth $W$ in units of
today's goods, consumption $c_1$ today and $c_2$ tomorrow, and real rate $r$, the budget line is
\begin{equation}\label{eq:fp-budget}
  c_1 + \frac{c_2}{1+r} = W, \qquad \frac{\dd c_2}{\dd c_1} = -(1+r),
\end{equation}
so its slope is $-(1+r)$ in the $(c_1,c_2)$ plane (Mehrling's notes write $-1/(1+r)$, which is the slope with
the axes exchanged). Maximising utility $u(c_1,c_2)$ on this line gives the first-order condition
\[
  \frac{\partial u/\partial c_1}{\partial u/\partial c_2} = 1+r,
\]
the tangency between an indifference curve and the budget line \ts{1}{2}{2:36}. Two people with the same
wealth and the same production possibilities but different tastes choose different points: one consumes more
than she produces today and borrows, the other consumes less and lends \ts{1}{2}{2:58}. Financial markets allow
this inter-temporal trade, and the real interest rate is the price that clears it \ts{1}{2}{3:19}. The diagram
is named after Irving Fisher \ts{1}{2}{3:39}. A money and banking course built on it treats banks as
intermediaries between savers and investors, needed because of some imperfection such as asymmetric
information \ts{1}{2}{4:00}.

\begin{nfigure}
\begin{tikzpicture}[font=\small]
  \begin{scope}
    \draw[-{Stealth[length=4pt]}] (0,0) -- (4.2,0) node[below left]{income $Y$};
    \draw[-{Stealth[length=4pt]}] (0,0) -- (0,3.4) node[left]{$i$};
    \draw[thick,defcol] (0.5,3.0) -- (3.8,0.5) node[right]{IS};
    \draw[thick,thmcol] (0.6,0.6) -- (3.8,2.2) node[right]{LM};
    \fill (2.16,1.39) circle (1.5pt);
    \draw[dashed,thin] (2.16,0) -- (2.16,1.39) -- (0,1.39);
    \node[font=\bfseries] at (2,3.5) {macro: nominal rate};
  \end{scope}
  \begin{scope}[xshift=6.2cm,scale=0.95]
    % PPF: quarter circle of radius 3; production at angle 40 deg, P=(2.298,1.928), slope -cot 40 = -1.192
    % budget line through P; indifference curves c1^a c2^(1-a) tangent to it (a=0.35 lender, a=0.75 borrower)
    \draw[-{Stealth[length=4pt]}] (0,0) -- (4.6,0) node[below]{$c_1$};
    \draw[-{Stealth[length=4pt]}] (0,0) -- (0,5.0) node[left]{$c_2$};
    \draw[thick,intcol] (0,3) arc (90:0:3);
    \node[intcol,font=\scriptsize] at (0.55,2.75) {PPF};
    \draw[thick,fincol] (0,4.667) -- (3.916,0);
    \node[fincol,font=\scriptsize,anchor=west] at (2.05,3.45) {slope $-(1+r)$};
    \draw[fincol,thin] (2.0,3.42) -- (0.95,3.53);
    \draw[defcol,domain=0.78:3.6,smooth,variable=\x] plot (\x,{3.031*(1.371/\x)^0.5385});
    \draw[thmcol,domain=1.92:4.4,smooth,variable=\x] plot (\x,{1.166*(2.937/\x)^3});
    \fill (2.298,1.928) circle (2pt);
    \node[font=\scriptsize,anchor=east] at (2.25,1.75) {production};
    \fill[defcol] (1.371,3.031) circle (2pt);
    \fill[thmcol] (2.937,1.166) circle (2pt);
    \draw[dashed,thin,defcol] (1.371,0) -- (1.371,3.031);
    \draw[dashed,thin] (2.298,0) -- (2.298,1.928);
    \draw[dashed,thin,thmcol] (2.937,0) -- (2.937,1.166);
    \draw[decorate,decoration={brace,amplitude=3pt,mirror},defcol] (1.40,-0.08) -- (2.27,-0.08)
      node[midway,below=3pt,font=\scriptsize]{lend};
    \draw[decorate,decoration={brace,amplitude=3pt,mirror},thmcol] (2.33,-0.08) -- (2.91,-0.08)
      node[midway,below=3pt,font=\scriptsize]{borrow};
    \node[font=\bfseries\small] at (2.3,5.5) {micro: real rate};
  \end{scope}
\end{tikzpicture}
\caption{The two textbook starting points \ts{1}{2}{1:15}\ts{1}{2}{2:16}. Left: IS/LM fixes income and the nominal
interest rate. Right: the Fisher diagram; two agents produce at the same point of the production possibility
frontier (PPF) and trade along the budget line to the tangency with their own indifference curve: the one who
consumes less today lends, the other borrows.}\label{fig:fp-prereq}
\end{nfigure}

\paragraph{The strategy of this course.} The course does not start from an economic apparatus and ask how to
apply it to banking: it starts from the concrete realities of banking and tries to draw from them, inductively,
how the system works and what regularities it has \ts{1}{2}{4:41}. Mehrling warns that the first two weeks
feel disorienting, because there is nothing familiar to grab onto; it comes together later \ts{1}{2}{5:02}. His own
way of reading the FT is to ask what Bagehot, Minsky or Fischer Black, whose biographies he wrote, would make of
it \ts{1}{2}{6:23}.\begin{aside}[The shelf in the library]
Monetary economics is a field of big books: to get anywhere one has to write a treatise. Mehrling chose it after
seeing the shelf of those books in the library and thinking that a life spent reading them from left to right
would be the best life ever \ts{1}{2}{7:04}.
\end{aside}

% ------------------------------------------------------------------
\section{What is a bank, a shadow bank, a central bank}\label{sec:fp-banks}

The analytical apparatus of the course is, largely, the balance sheet \ts{1}{3}{0:07}. Not balance sheets as
accounting law (what GAAP prescribes), but as a structure that forces rigorous thinking about the economy
\ts{1}{3}{0:27}.\begin{aside}[GAAP]
GAAP, \emph{Generally Accepted Accounting Principles}, are the accounting standards set in the United
States by the Financial Accounting Standards Board. The course ignores them: its balance sheets are a tool of
economic analysis, not of legal reporting.
\end{aside} We introduce three new concepts.

\begin{definition}[New concepts: balance sheet, solvency, liquidity]\label{def:fp-bs}
\leavevmode
\begin{enumerate}[(i)]
  \item A \emph{balance sheet} lists what an agent owns, its \emph{assets} (left side), and what it owes, its
        \emph{liabilities} (right side). \emph{Net worth} (equity) is the balancing item,
        \[ \text{net worth} = \text{assets} - \text{liabilities}, \]
        so that the two sides always have the same total. Drawn as a ``T-account'', the balance sheet has assets
        on the left and liabilities on the right \ts{1}{3}{0:49}.
  \item An agent is \emph{solvent} if the market value of its assets exceeds the market value of its
        liabilities, i.e.\ net worth is positive \ts{1}{3}{2:37}. (From Latin \emph{solvere}, to loosen,
        hence to pay off a debt.)
  \item An agent is \emph{liquid} if it can make the payments it has to make when they come due: it has cash, or
        can obtain it \ts{1}{3}{3:23}. (From Latin \emph{liquidus}, flowing: liquid assets flow easily into
        means of payment.)
\end{enumerate}
\end{definition}

\paragraph{A bank.} A generic bank holds loans, securities (for example Treasury bills) and cash reserves, and
owes deposit accounts and other borrowing (bonds, short-term money market borrowing); net worth balances the two
sides \ts{1}{3}{1:09}. ``Other borrowing'' will be a very big part of the course \ts{1}{3}{1:56}.
\begin{taccts}
\tacct[3cm]{Bank}{Loans\\Securities\\Cash reserves}{Deposits\\Other borrowing\\Net worth}
\end{taccts}
Two issues arise \ts{1}{3}{2:17}. The first is \emph{solvency}: are the assets worth more than the liabilities?
Much of the discussion after 2008, and in the euro crisis, is about this; perhaps entire national banking systems
are insolvent \ts{1}{3}{2:37}. The second, which is what makes banks different from other corporations and which
the course puts first, is \emph{liquidity} \ts{1}{3}{2:58}. Deposits are payable on demand: when a depositor pays
someone at another bank, the bank must deliver cash reserves. Does it have them? If not, can it acquire them, by
borrowing in the money market, by borrowing from the central bank (the Fed's ``discount window''), by selling or
\emph{repo}-ing a security \ts{1}{3}{3:44}? Or must it sell its loans at a \emph{fire-sale} price, below their
value?

\begin{intuition}[Liquidity can become solvency]
A fire sale turns a liquidity problem into a solvency problem: assets sold below their value reduce net worth,
and enough of them can make it negative \ts{1}{3}{4:04}. Most economists approach banking from solvency; the
course approaches it from liquidity, because that is how bankers think: they face the liquidity constraint every
day, when they have to clear their accounts \ts{1}{3}{4:25}. This is why the course starts with the payment system.
A rough analogy: solvency is a statement about stocks (values at one instant), liquidity a statement about flows
(cash in versus cash out, every day). A system can have enough ``energy'' in total and still fail because it
cannot deliver it at the rate required.
\end{intuition}\begin{example}[A fire sale]
A bank has loans 100, reserves 5; deposits 95, net worth 10. Depositors withdraw 20. Reserves cover 5; the bank
must raise 15 by selling loans, and in a panic buyers pay only 60\% of face value: it sells loans of face value
25. Now loans 75, reserves 0, deposits 75, net worth 0. A run of 20 out of 95 wiped out all the capital, although
before it the bank was solvent by 10.
\end{example}

\paragraph{A shadow bank.} Mehrling's research group had just finished a paper on shadow banking, discussed at the
end of the course (lecture 21) \ts{1}{3}{4:45}. A shadow bank, in the sense of the course, is an entity that holds,
for example, residential mortgage-backed securities (loans turned into a kind of bond), strips most of the risk
out of them with an interest rate swap and a credit default swap, and uses what is left, a short-term and
relatively risk-free asset, as collateral to borrow in the \emph{wholesale} money market \ts{1}{3}{5:06}. It
borrows not from households but from money market mutual funds and large corporate investors, through repurchase
agreements (repo), Eurodollar borrowing, asset-backed commercial paper \ts{1}{3}{6:15}.

\begin{definition}[New concepts: shadow bank, wholesale money market]\label{def:fp-shadow}
A \emph{shadow bank} is an entity that does what a bank does (holds long-term credit and finances it with
short-term liabilities that its creditors regard as money-like) outside the regulated deposit-taking banking
system. Its funding comes from the \emph{wholesale money market}, the market where institutions, not households,
lend to each other for short terms, often against collateral. The term ``shadow banking'' was coined by the
economist Paul McCulley (PIMCO) in 2007.
\end{definition}

\begin{taccts}
\tacct[4.2cm]{``Shadow bank''}{Mortgage-backed security\\\quad tranche\\Interest rate swap\\Credit default swap}%
{Money market borrowing:\\\quad repurchase agreements\\\quad asset-backed commercial paper\\\quad auction rate securities\\\quad Eurodollar borrowing}
\end{taccts}
The instruments on this balance sheet are only named here; each gets its own lecture (repo in lecture~7,
Eurodollars in lecture~8, swaps in lectures 19--20). ``It's not rocket science'' \ts{1}{3}{6:36}. The point now is
that this model of credit became dominant: the amount of credit in the US that ran through it became larger than
the amount that ran through traditional banks; the textbooks are all about banks, reality increasingly about
shadow banks \ts{1}{3}{6:56}. Again there is an issue of solvency and an issue of liquidity, but liquidity is
harder: it means being able to \emph{roll over} the short-term funding. If the money market freezes, the shadow
bank cannot roll over and has to liquidate its assets; much of the financial crisis of 2007--2009 was exactly
this \ts{1}{3}{7:38}. The regulatory system was built for traditional banks; the Dodd--Frank Act of 2010 ``doesn't
touch'' shadow banking \ts{1}{3}{8:20}.

\paragraph{The central bank is a bank.} The central bank is also a bank, at a higher level \ts{1}{3}{8:40}. Its
balance sheet is published every week (the Fed's release H.4.1, ``Factors Affecting Reserve Balances of
Depository Institutions''). The release shown in class, of 30 August 2012, gives, in round numbers
\ts{1}{3}{9:00}\ts{1}{3}{11:04}:
\begin{taccts}
\tacct[5.6cm]{Federal Reserve, 30 Aug 2012 (\$ billion)}{Treasuries (no bills)\amt{$\approx$1{,}600}\\Agency debt (Fannie Mae, Freddie Mac)\\Mortgage-backed securities\amt{852}\\Maiden Lane (AIG, Bear Stearns)\\Central bank liquidity swaps\\\textbf{Total}\amt{$\approx$2{,}800}}%
{Currency in circulation\amt{$\approx$1{,}000}\\Reserve balances of banks\amt{$\approx$1{,}500}\\Other\\\\\\\textbf{Total}\amt{$\approx$2{,}800}}
\end{taccts}
\begin{itemize}
  \item The Fed had sold all its short-term Treasury bills and held only notes and bonds: it was moving into
        long-term debt \ts{1}{3}{9:00}.
  \item It held mortgage-backed securities, \$852 billion of them, for the first time in its history: a legacy of
        the crisis \ts{1}{3}{9:21}\ts{1}{3}{9:42}. The ``Maiden Lane'' entries are what was left of the rescues of
        Bear Stearns and AIG.
  \item \emph{Central bank liquidity swaps} are loans of dollars from the Fed to other central banks: the ECB, the
        Swiss National Bank, the Bank of Japan \ts{1}{3}{10:03}.
  \item The total was about \$2.8 trillion, against less than \$1 trillion before the crisis: the balance sheet
        roughly tripled \ts{1}{3}{11:04}.
  \item On the liability side, currency (``green pieces of paper'') was about \$1 trillion, and the reserve
        balances that banks hold as deposits at the Fed were \$1.5 trillion, against about \$50 billion before the
        crisis (as said in class), an order of magnitude change \ts{1}{3}{11:24}\ts{1}{3}{11:45}.
\end{itemize}
This transformation of the Fed's balance sheet is what ``QE1, QE2, QE3'' (quantitative easing) refer to.

\begin{intuition}[The central bank does not ``print money'']
Textbooks write money demand $=$ money supply, $M^d=M^s$, as if the central bank decided a quantity. Behind the
symbols are balance sheets: the central bank is a bank with assets and liabilities, and when it ``expands the money
supply'' it expands \emph{both sides} of its balance sheet at once, for example buying a security (asset) and paying
for it by crediting a bank's reserve account (liability) \ts{1}{3}{10:23}\ts{1}{3}{10:44}. This is a classic banking
operation, and it happens at every level of the system, not only at the top.
\end{intuition}

\begin{taccts}
\tacct[3.2cm]{Fed}{$+$Security}{$+$Reserves of Bank A}\qquad
\tacct{Bank A}{$-$Security\\$+$Reserves}{}
\end{taccts}

% ------------------------------------------------------------------
\section{Central themes}\label{sec:fp-themes}

\subsection{The plan of the course}
The course is about banking \emph{phenomena}, not about banks as legally defined institutions (Mehrling's notes).
It is organised around four functions of banking \ts{1}{4}{0:08}\ts{1}{4}{1:29}, summarised in
\cref{fig:fp-map}:
\begin{enumerate}
  \item \textbf{Banking as a clearing system}: the payment system. Understanding payments is the best starting point
        to understand \emph{liquid assets}.
  \item \textbf{Banking as market making}: once liquid assets are understood, \emph{liquid markets}. These first two
        functions are the subject of the first half of the course, up to the midterm.
  \item \textbf{International money and banking}: the same two ideas extended to the global system \ts{1}{4}{0:28}.
  \item \textbf{Banking as advance clearing}: how finance (forwards, futures, swaps and other derivatives) enters the
        picture \ts{1}{4}{0:48}.
\end{enumerate}
At the end come regulation, what to do about shadow banking, and ``touching the elephant'', the connection of the
money view with what one learns in economics and in finance \ts{1}{4}{1:09}. The fourth classical function,
\emph{intermediation} (channelling savings to investment), is the one textbooks start from; here it comes last.

\subsection{The four prices of money}
In a sense the whole course is about understanding where four prices come from \ts{1}{4}{1:49}.

\begin{definition}[New concepts: the four prices of money]\label{def:fp-four}
\leavevmode
\begin{enumerate}[(i)]
  \item \emph{Par}: the price of one kind of money in terms of another kind of money, today, normally fixed at one
        for one (a deposit of \$107 at a bank exchanges for \$107 of currency) \ts{1}{4}{2:34}\ts{1}{4}{3:16}.
        (Latin \emph{par}, equal.)
  \item \emph{Interest}: the price of money today in terms of money tomorrow; an interest rate $i$ means that $1$
        today exchanges for $1+i$ in one period \ts{1}{4}{2:14}. Its markets: Fed funds, Eurodollars, repo.
        (Medieval Latin \emph{interesse}, ``to be between, to matter'': the compensation due to a lender for the
        loss of the use of the money.)
  \item \emph{Exchange rate}: the price of domestic money in terms of foreign money, or in terms of the world reserve
        money \ts{1}{4}{4:58}.
  \item \emph{Price level}: the price of money in terms of commodities, i.e.\ the inverse of the purchasing power of
        money \ts{1}{4}{5:59}.
\end{enumerate}
\end{definition}

Economics focuses on one of them, the interest rate, the price of money in the sense of a deposit transforming money
today into money tomorrow \ts{1}{4}{2:14}. The other three deserve more attention than they usually get.

\paragraph{Par, the forgotten price.} Par is a more primitive price than interest, and it is neglected
\ts{1}{4}{2:34}. In the bank balance sheet of \cref{sec:fp-banks} deposits (private money, a promise of the bank) sit
on one side and cash reserves (central bank money) on the other, and we take for granted that they trade at par
\ts{1}{4}{2:55}. In financial crises par is sometimes broken: bank money trades at a different rate from central bank
money \ts{1}{4}{3:16}. The system is \emph{hybrid}, part private money and part state money, and par masks the
hybridity: it makes the two look like the same thing \ts{1}{4}{3:36}. When the system works well they \emph{are} the
same; but why does it work well, how do we make it work well, and what do we do in a crisis \ts{1}{4}{3:56}? Par is
key to the payment system. If a Californian seller charged a buyer in New York an extra 5\% because there is no par
clearing between the two states, we would be outraged; yet this is how it was in the United States before the Fed
\ts{1}{4}{4:17}.

\begin{intuition}[Par is a fixed price, and fixed prices are hard]
Par clearing had to be created by institutions, and it is not easy: a price fixed at one for one, which must not
budge whatever the pressure, is like rent control, and economists know how hard such prices are to hold
\ts{1}{4}{4:37}. The rest of the course explains what holds it: the central bank, the clearing system, the
dealers who stand ready to trade at that price.
\end{intuition}\begin{history}[Exchange charges before the Fed]
Before the Federal Reserve, checks drawn on distant banks were often paid at less than face value: the paying bank
deducted an ``exchange charge'' to cover the cost of shipping currency. Universal par clearing was one of the aims of
the Federal Reserve Act (1913); the Fed's campaign for it lasted into the 1920s, and some small banks kept charging
non-par fees for decades.\par
\emph{References:} J.~Jessup, ``The Federal Reserve and Par Clearing'', Federal Reserve Bank of St.~Louis; Young,
``Mobilizing Banking Credits'' (the reading of this week).
\end{history}

\paragraph{The exchange rate.} It is particularly obscured for people living in the United States, because the dollar
is the world's best money \ts{1}{4}{5:18}. Americans never worry that ``the dollar'' could fall below par with the
\emph{international} dollar, the dollar actually used as world reserve \ts{1}{4}{5:39}. Yet there are two dollars, again
a private one and a state one: the international dollar consists largely of dollar deposits at banks outside the US,
mostly European, that have no account at the Fed. These are the \emph{Eurodollars} (lecture~8).

\paragraph{The price level.} Economists usually treat it with the quantity theory of money, $MV=PY$ (or $MV=PT$),
imagining that the price level moves one for one with the quantity of money \ts{1}{4}{5:59}. The course comes to this
only at the end \ts{1}{4}{6:20}; most of the attention goes to par, interest and the exchange rate.

\begin{finance}[The FT market data page]
Two or three pages from the back of the FT there is a ``horrific'' page of tiny numbers \ts{1}{4}{6:40}: market
interest rates (US dollar LIBOR, US dollar certificates of deposit, US overnight repo, Federal funds effective rate),
LIBOR in pounds and Swiss francs, exchange rates, interest rate swap rates \ts{1}{4}{7:00}. By the end of the course
every number on that page should be understandable, and explainable to a roommate: it is not an arcane world reserved
to partners of investment banks, but something the average citizen can and should understand \ts{1}{4}{7:20}.
\end{finance}

\subsection{Lombard Street, then and now}
The inspiration of the course is Bagehot's \emph{Lombard Street} (1873), and Mehrling's \emph{New Lombard Street} is an
homage to it \ts{1}{4}{7:41}. Bagehot wrote a bestseller to explain the money market to his fellow citizens after a
series of financial crises in which the Bank of England had done what it had promised not to do; he argued that it
had done the right thing: lend freely, at a high rate, against good collateral \ts{1}{4}{8:01}. The 2007--2009 crisis was
a crisis of the shadow banking system, and the Fed too did things it had sworn it would never do \ts{1}{4}{8:21}. To get
out of a global crisis, citizens need to understand how the system works, so that every proposal does not get stuck in
``pro-banker'' versus ``anti-banker'' politics; the course takes no side, and approaches the system as scientists and as
citizens \ts{1}{4}{8:42}\ts{1}{4}{9:02}.

\paragraph{Globalization: back to the 19th century.} The world has changed: it is a global financial world
\ts{1}{4}{9:23}. The world of the 1950s textbooks had hardly any private capital markets, and was not global, because
the global financial markets of the 19th century had broken down \ts{1}{4}{9:43}. In terms of globalization the
present resembles the world before the First World War, when the pound sterling was king and the Bank of England
acted as central bank of the whole world, more than the world after the Second (Mehrling's notes). Hence
\begin{center}
\emph{Bagehot is a better guide for today than Irving Fisher} \ts{1}{4}{10:03}:
\end{center}
Fisher lived in a simpler, closed world; Bagehot lived in a global financial system organised around the pound. He did
not know derivatives or high-frequency trading, but he would know how to start thinking about them, because he
understood how money markets work \ts{1}{4}{10:24}.

\paragraph{Finance: all banking is a swap of IOUs.} The second novelty is the integration of finance. Connecting money
and finance, the subjects of Mehrling's two books, which people regard as separate fields, occupied him for ten years
\ts{1}{4}{10:44}. The motto of the course is the connection \ts{1}{4}{11:04}:

\begin{keyformulas}
\textbf{All banking is a swap of IOUs.} \ts{1}{4}{11:25}

\smallskip
\begin{taccts}
\tacct[2.8cm]{Bank}{$+$Loan\\\quad(borrower's IOU)}{$+$Deposit\\\quad(bank's IOU)}\qquad
\tacct[2.8cm]{Borrower}{$+$Deposit\\\quad(bank's IOU)}{$+$Loan\\\quad(borrower's IOU)}
\end{taccts}
A loan is an exchange of promises: the borrower's promise to pay later against the bank's promise to pay on demand.
\end{keyformulas}

By ``swap'' Mehrling means interest rate swaps and credit default swaps, but also ordinary loans \ts{1}{4}{11:25}.
Finance believes it knows everything about swaps; the money view looks at them differently: every swap has a
\emph{funding} piece and a \emph{liquidity} piece, and the liquidity piece drives the action \ts{1}{4}{11:45}. The
course is about money markets and liquidity, connected to solvency, but starting from where Bagehot would start:
the money market as a global system of clearing \ts{1}{4}{12:06}. The aim is monetary theory for the real world of
financial globalization, using as guideposts the great texts of the past and the ``great texts of the present'', Stigum
and the FT \ts{1}{4}{12:26}\ts{1}{4}{12:47}.\begin{aside}[IOU]
\emph{IOU} is a phonetic abbreviation of ``I owe you'', used since the 17th century for an informal written
acknowledgement of a debt. In the course it means any promise to pay, from a handwritten note to a Treasury bond or a
deposit.
\end{aside}

\begin{nfigure}
\begin{tikzpicture}[font=\small,
  box/.style={draw,rounded corners=2pt,align=center,inner sep=3pt,text width=2.6cm,minimum height=1.2cm},
  arr/.style={-{Stealth[length=4pt]},semithick}]
  \node[box,draw=defcol,fill=defcol!8] (c) at (0,0) {\textbf{1. Clearing}\\\scriptsize payment system,\\liquid assets};
  \node[box,draw=defcol,fill=defcol!8] (m) at (3.5,0) {\textbf{2. Market making}\\\scriptsize dealers,\\liquid markets};
  \node[box,draw=intcol,fill=intcol!8] (i) at (7,0) {\textbf{3. International}\\\scriptsize the same, global:\\FX, key currencies};
  \node[box,draw=fincol,fill=fincol!8] (a) at (10.5,0) {\textbf{4. Advance clearing}\\\scriptsize forwards, futures,\\swaps};
  \draw[arr] (c)--(m); \draw[arr] (m)--(i); \draw[arr] (i)--(a);
  \node[draw=srccol,fill=srccol!8,rounded corners=2pt,align=center,inner sep=3pt] (e) at (8.7,-2.1)
    {\textbf{Shadow banking, regulation}\\\scriptsize ``touching the elephant'':\\\scriptsize money view vs economics vs finance};
  \draw[arr] (a.south) -- (a.south |- e.north);
  \draw[decorate,decoration={brace,amplitude=5pt,mirror},thick] (-1.4,-0.8) -- (4.9,-0.8)
    node[midway,below=6pt,font=\scriptsize]{before the midterm (lectures 1--12)};
  \node[draw=keycol,fill=keycol!8,rounded corners=2pt,align=center,inner sep=3pt] (p) at (3.5,1.9)
    {\textbf{four prices of money}\\\scriptsize par $\cdot$ interest $\cdot$ exchange rate $\cdot$ price level};
  \draw[arr,keycol,dashed] (p) -- (c.north); \draw[arr,keycol,dashed] (p) -- (m.north); \draw[arr,keycol,dashed] (p) -- (i.north);
  \node[draw=intucol,fill=intucol!8,rounded corners=2pt,align=center,inner sep=3pt] (s) at (10.5,1.9)
    {\textbf{all banking is}\\\textbf{a swap of IOUs}};
  \draw[arr,intucol,dashed] (s) -- (a);
\end{tikzpicture}
\caption{Map of the course \ts{1}{4}{0:08}: the four functions of banking in the order in which they are treated,
the four prices that the course explains, and the motto that connects money and finance.}\label{fig:fp-map}
\end{nfigure}

% ------------------------------------------------------------------
\section{Reading: Allyn Young}\label{sec:fp-young}

The first reading is by Allyn Young, one of the authors who started Mehrling on his path: the first three chapters of
\emph{The Money Interest and the Public Interest} are a biography of Young \ts{1}{5}{0:00}.

\begin{reading}[Allyn Young, four encyclopedia chapters (1924)]
Young wrote the chapters, unsigned, for an encyclopedia, \emph{The Book of Popular Science}, in 1923--1924, just after
the First World War \ts{1}{5}{0:00}: ``The Mystery of Money'', ``The Monetary System of the United States'',
``Mobilizing Banking Credits'', ``Dear and Cheap Money''. The Fed had been founded in 1913, and the United States had
no tradition of thinking about central banking; Americans were traditionally suspicious of it \ts{1}{5}{0:40}. Young
writes about the Fed when it is a new institution whose role is unclear: he is a friend of the president of the
Federal Reserve Bank of New York and advises him, while open market operations and stabilisation policy are being
invented \ts{1}{5}{1:15}. He compares the Fed with the Bank of England and predicts that London will remain the centre
of world finance; he was wrong: it was the end of sterling's hegemony, although the dollar became the world reserve
currency only after the Second World War, after the Great Depression \ts{1}{5}{1:15}\ts{1}{5}{1:48}. The chapters are
easy to read but deep: on a second reading one notices how carefully Young steers away from minefields
\ts{1}{5}{2:19}. Their spirit is the spirit of the course: having lived through a crisis we do not fully understand,
in a system that is changing, we have to figure it out, as he did in 1924 \ts{1}{5}{1:48}.

\smallskip
\emph{Study questions} (\file{Allyn Young(study).pdf}): where does Young draw the line between money and credit?
Represent with T-accounts the three ways to obtain a bank deposit (\cref{ex:fp-deposits}). How is Young's hostility
to government ``fiat'' money and his sympathy for bank ``credit money'' consistent with the hierarchy of money
(lecture~2)? The reading is analysed in detail in lecture~3.
\end{reading}\begin{history}[Allyn Young (1876--1929)]
Allyn Abbott Young taught at Cornell, Stanford and Harvard; in 1927 he moved to the London School of Economics, where
he died of pneumonia in March 1929, during an influenza epidemic, at 52. He is best known today for ``Increasing
Returns and Economic Progress'' (\emph{Economic Journal}, 1928). The friend at the New York Fed was Benjamin Strong,
its first governor (1914--1928), the dominant figure of the early Fed. Young's encyclopedia chapters are reprinted in
P.~Mehrling and R.~Sandilands (eds.), \emph{Money and Growth: Selected Papers of Allyn Abbott Young}, Routledge, 1999,
from which the course reading is taken.\par
\emph{References:} Mehrling, \emph{The Money Interest and the Public Interest}, 1997, ch.~1--3; Mehrling and
Sandilands, 1999.
\end{history}

% ------------------------------------------------------------------
\flushside
\section*{Key formulas and ideas}
\begin{keyformulas}
\begin{tabularx}{\linewidth}{@{}l L@{}}
Balance sheet & assets $=$ liabilities $+$ net worth (always balances)\\
Solvency & net worth $=$ market value of assets $-$ market value of liabilities $>0$ (a stock)\\
Liquidity & cash available or obtainable $\ge$ payments due, every day (a flow)\\
Four prices of money & par (money vs money, today); interest (money today vs money tomorrow, $1\mapsto 1+i$);
exchange rate (domestic vs foreign money); price level (money vs goods)\\
Fisher diagram & $c_1 + c_2/(1+r) = W$, slope $-(1+r)$; optimum: $u_1/u_2 = 1+r$ \eqref{eq:fp-budget}\\
Central bank & expanding the money supply $=$ expanding both sides of its balance sheet\\
Motto & all banking is a swap of IOUs\\
\end{tabularx}
\end{keyformulas}

\section*{Exercises}

\begin{exercise}[Three ways to get a deposit]\label{ex:fp-deposits}
Using T-accounts for you, your bank and (where needed) a second person and a second bank, record:
\begin{enumerate}[(a)]
  \item you deposit \$100 of currency;
  \item you deposit a check for \$100 written by someone with an account at another bank;
  \item the bank grants you a loan of \$100 by crediting your deposit account.
\end{enumerate}
In which case does the total quantity of deposits in the economy increase? (From the study questions on Young.)
\end{exercise}

\begin{exercise}[Liquidity and solvency]
A bank has loans 200, securities 50, reserves 10; deposits 230, net worth 30. Depositors withdraw 40.
\begin{enumerate}[(a)]
  \item The bank sells securities at their market value to pay. Write the new balance sheet. Is the bank solvent?
        Liquid?
  \item The securities can only be sold at 80\% of their value and the loans at 50\%. Which assets does the bank sell
        first, and what is its net worth afterwards?
  \item Name two ways the bank could meet the withdrawal without selling any asset.
\end{enumerate}
\end{exercise}

\begin{exercise}[Which price of money?]
For each of the following, say which of the four prices of money is involved:
\begin{enumerate}[(a)]
  \item a money market fund announces that each share is now worth \$0.97 instead of \$1;
  \item the Fed funds effective rate is 0.15\%;
  \item a tourist pays 1.30 dollars for one euro;
  \item the consumer price index rises by 2\% in a year;
  \item before 1914 a New York bank paid a check drawn on a Chicago bank at \$99.85 per \$100.
\end{enumerate}
\end{exercise}

\begin{exercise}[The Fisher diagram]
An agent has endowment $(y_1,y_2)=(100,0)$ and utility $u=\ln c_1+\beta\ln c_2$.
\begin{enumerate}[(a)]
  \item Find the optimal $(c_1,c_2)$ as a function of $r$ and $\beta$, and the amount lent.
  \item A second agent has endowment $(0,110)$ and the same utility. For $\beta=1$, find the real rate $r$ that clears
        the market for loans between the two.
\end{enumerate}
\end{exercise}

\begin{exercise}[Quantitative easing in T-accounts]
The Fed buys \$100 billion of mortgage-backed securities from a primary dealer, which holds its account at Bank B.
Write the changes on the balance sheets of the Fed, Bank B and the dealer. Which item of the Fed's liabilities
grows? Relate your answer to the reserve balances of \cref{sec:fp-banks}.
\end{exercise}

\begin{exercise}[A shadow bank in a freeze]
A shadow bank holds mortgage-backed securities worth 100, financed by overnight repo of 95 and capital of 5. One
morning lenders agree to roll over only 80.
\begin{enumerate}[(a)]
  \item How much must the shadow bank raise, and from what?
  \item If buyers pay only 90\% of value for the securities, is the shadow bank still solvent after the sale?
  \item Compare with the bank of \cref{def:fp-bs}: which liability of a traditional bank plays the role of the repo?
\end{enumerate}
\end{exercise}

% !TEX root = ../main.tex
\chapter{The Natural Hierarchy of Money}\label{ch:hierarchy}
\begin{paracol}{2}

\begin{sources}
\begin{tabularx}{\linewidth}{@{}l l L@{}}
\toprule
\textbf{Segment} & \textbf{Length} & \textbf{Content}\\
\midrule
\seg{2}{1}{L2.1 FT: The Eurocrisis, Liquidity vs.\ Solvency} & 10:05 & Draghi's bond purchases (a liquidity operation) versus Soros's European Fiscal Authority (a solvency operation), in T-accounts.\\
\seg{2}{2}{L2.2 Hierarchy of Financial Instruments} & 9:39 & Money versus credit; gold, currency, deposits, securities; where the line between money and credit is drawn depends on the point of view.\\
\seg{2}{3}{L2.3 Hierarchy of Financial Institutions} & 6:37 & The same hierarchy in balance sheets; inside and outside money.\\
\seg{2}{4}{L2.4 Dynamics of the Hierarchy} & 6:08 & The pyramid of credit; expansion and contraction at every time scale, in quantity and in quality.\\
\seg{2}{5}{L2.5 Discipline and Elasticity} & 8:49 & Scarcity of ultimate money and elasticity of derivative credit; currency principle versus banking principle.\\
\seg{2}{6}{L2.6 Hierarchy of Market Makers} & 9:17 & Central bank, banks and security dealers as market makers; the prices of money as the links between layers.\\
\seg{2}{7}{L2.7 Managing the Hierarchy} & 18:03 & Lender of last resort, origin of monetary policy, the term structure and the ``artificial hierarchy''; Operation Twist.\\
\bottomrule
\end{tabularx}

\smallskip
\textbf{Files.} Mehrling's lecture notes \file{money-lecture_notes-Lec 02--The Natural Hierarchy of Money.pdf}
(dated 14 Sept 2009); captions \file{materials/week1/srt/2 - *.srt}; transcripts \file{materials/transcripts/L02-P0*.txt}.
\textbf{Reading.} Allyn Young (\cref{sec:fp-young}); the study questions ask to place Young's ``The Mystery of
Money'' in the framework of this lecture.
\end{sources}

\section*{Motivation}
Mehrling calls this lecture ``the whole course'': everything that follows hangs on its picture, like
ornaments on a tree \ts{2}{2}{0:08}\ts{2}{4}{5:47}. Its first sentence is the thesis:
\begin{center}
\emph{Always and everywhere, monetary systems are hierarchical.}
\end{center}
The lecture builds the hierarchy in four steps, each adding a dimension to the previous one:
\begin{enumerate}
  \item a hierarchy of \emph{instruments} (gold, currency, deposits, securities), \cref{sec:nh-instruments};
  \item a hierarchy of \emph{institutions} that issue and hold them, \cref{sec:nh-institutions};
  \item the \emph{dynamics} of the hierarchy, which expands and contracts at every time scale, and the two
        principles behind it, discipline and elasticity, \cref{sec:nh-dynamics,sec:nh-discipline};
  \item a hierarchy of \emph{market makers} and of the \emph{prices of money} that connect the layers,
        \cref{sec:nh-marketmakers}, and the central bank's attempt to \emph{manage} it, \cref{sec:nh-managing}.
\end{enumerate}

\paragraph{Four mental obstacles.} The lecture comes so early because of four habits of thought, partly from
economics courses, that stand in the way \ts{2}{2}{0:48}:
\begin{enumerate}[(a)]
  \item banks expand both sides of the balance sheet at once, which looks like a free lunch, ``the alchemy of
        banking, creating something from nothing'', and good economics students resist it \ts{2}{2}{1:09};
  \item money is usually thought of as a positive quantity, a ``stuff''; seeing it as somebody's
        \emph{liability} is hard (this is the inside/outside money distinction) \ts{2}{2}{1:50};
  \item money is thought of as a creature of the nation-state; in this lecture the state does not appear, the
        central bank is just a bankers' bank (the state enters in lecture~3) \ts{2}{2}{2:11};
  \item general equilibrium puts all goods at the same level, linked by relative prices; monetary instruments
        are \emph{not} at the same level, they only look so, until a crisis shows that things that seemed
        equivalent are not \ts{2}{2}{2:31}.
\end{enumerate}
Students without economics background often find the start easier, because they have nothing to unlearn.
Nothing here contradicts the economics one knows; it is integrated with it at the end of the course
\ts{2}{2}{3:33}.

% ------------------------------------------------------------------
\section{FT: the eurocrisis, liquidity versus solvency}\label{sec:nh-ft}

\begin{casestudy}[FT, September 2012: Draghi and Soros]
The ECB under Mario Draghi has just announced that it is prepared to support the value of the short-term debt
of the ``peripheral'' euro countries (Spain, Italy, \dots), to the dismay of some German central bankers
\ts{2}{1}{0:27}. In the same issue of the FT, George Soros, under the title ``Lead or leave currency bloc,
Soros tells Germany'', argues that the ECB is not doing enough and makes his own proposal \ts{2}{1}{0:48}.
Mehrling uses the two plans to consolidate the distinction of lecture~1: Draghi's is a \emph{liquidity}
operation, Soros's a \emph{solvency} operation \ts{2}{1}{1:08}.
\end{casestudy}\begin{history}[``Whatever it takes'' and OMT (2012)]
On 26 July 2012, in London, Draghi said that ``within our mandate, the ECB is ready to do whatever it takes to
preserve the euro. And believe me, it will be enough.'' On 6 September 2012 the ECB Governing Council announced
\emph{Outright Monetary Transactions} (OMT): unlimited purchases of government bonds with 1 to 3 years to maturity,
on the secondary market, for countries under an EU assistance programme. OMT was never actually used, but
spreads fell sharply after the announcement. In 2015 the European Court of Justice ruled it compatible with
the EU treaties.\par
\emph{References:} ECB press release ``Technical features of Outright Monetary Transactions'', 6 Sept 2012;
CJEU, case C-62/14 \emph{Gauweiler}, 16 June 2015.
\end{history}

\paragraph{Draghi: a liquidity operation.} Investors, banks among them, hold short-term debt (\emph{bills}: short
term; \emph{bonds}: usually long term) of the peripheral countries, which pays much higher rates than German or
French debt \ts{2}{1}{1:54}. The ECB says it will cap those rates, without announcing an exact cap, by
entering the market as a buyer of the bills, and it pays for them, at least at first, by expanding its own balance
sheet: ``creating money'' \ts{2}{1}{2:34}.
\begin{taccts}
\tacct[2.9cm]{ECB}{$+$Peripheral bills}{$+$Money (reserves)}\qquad
\tacct[2.9cm]{Bondholders}{$-$Peripheral bills\\$+$Money}{}
\end{taccts}
This expansion of both sides of the balance sheet is what people mean by ``monetizing government debt''
\ts{2}{1}{2:56}. It supports the price of the bills and drags down the rates the periphery pays: a liquidity
operation \ts{2}{1}{3:18}.

\paragraph{Soros: a solvency operation.} Soros says that this only buys time: the fundamental problem is that some
peripheral countries have too much debt, a debt-to-GDP ratio above 60\% \ts{2}{1}{3:58}. His plan
\ts{2}{1}{4:40}:
\begin{enumerate}
  \item a \emph{European Fiscal Authority} (EFA), a fiscal unit of Europe as a whole, ``essentially like the US
        Treasury'', buys the ``excess'' long-term bonds of the periphery (the part above 60\% of GDP);
  \item it finances the purchase by issuing its own short-term debt, a liability of Europe as a whole:
        \emph{Euro bills};
  \item the ECB accepts Euro bills as the best collateral in the world, and may at first buy some of them by
        expanding its balance sheet, to get the operation going \ts{2}{1}{5:44};
  \item once bought, the excess bonds can be forgiven: the EFA's net worth takes the hit, but Europe as a whole is
        solvent where Portugal is not; the debt has moved to the larger, solvent entity \ts{2}{1}{6:50}.
\end{enumerate}
\begin{taccts}
\tacct[2.15cm]{EFA}{$+$Excess bonds}{$+$Euro bills}\quad
\tacct[2.0cm]{ECB}{$+$Euro bills}{$+$Money}\quad
\tacct[2.15cm]{Bondholders}{$-$Excess bonds\\$+$Money}{}
\end{taccts}
``You are going to get used to this pattern of thinking'': this minus and that plus are the same thing, every
transaction has two sides and both are booked \ts{2}{1}{6:29}.

\begin{intuition}[Money that appears and disappears]
Banks all over Europe do not trust each other and hold huge reserves at the ECB, earning nothing, as the US banks
held \$1.5 trillion at the Fed \ts{2}{1}{7:50}. Soros expects them to swap those reserves for Euro
bills. Then the ECB's balance sheet shrinks back: it was used only temporarily to distribute the bills, the money
disappears, and in the end no debt has been monetized; excess reserves have simply been replaced by Euro bills
\ts{2}{1}{8:31}. Money is created when a balance sheet expands and destroyed when it contracts.
\end{intuition}
\begin{taccts}
\tacct[2.15cm]{ECB}{$-$Euro bills}{$-$Reserves of banks}\qquad
\tacct[2.15cm]{Banks}{$-$Reserves\\$+$Euro bills}{}
\end{taccts}
There is no European treasury, so the job is left to the ECB. ``When central banks do very bold things, it's because
somebody else isn't doing their job'' \ts{2}{1}{9:12}. Understanding the mechanics allows a saner
political discussion than throwing around words like ``hyperinflation'' \ts{2}{1}{9:53}.

% ------------------------------------------------------------------
\section{The hierarchy of financial instruments}\label{sec:nh-instruments}

Start from the familiar distinction between money and credit \ts{2}{2}{4:18}.

\begin{definition}[New concepts: money and credit]\label{def:nh-money}
\emph{Money} is the means of final settlement: what debts are cancelled in. \emph{Credit} is a promise to pay money:
a means of \emph{delaying} final settlement. Money is not just ``more expensive'' than credit, it is
\emph{qualitatively} better \ts{2}{2}{4:18}. (``Credit'' from Latin \emph{credere}, to believe, to trust: a credit is a
promise one trusts.)
\end{definition}\begin{coursenote}[Hawtrey]
The money/credit distinction is the framework of Ralph Hawtrey's \emph{Currency and Credit} (1919; Mehrling's notes
cite the 1923 edition).
\end{coursenote}

The distinction is fine as far as it goes, but it goes both too far and not far enough (Mehrling's notes). Not far
enough, because it has only two layers; too far, because what counts as final settlement depends on the layer:
\emph{what looks like money at one level looks like credit from the level above.} To see this, take a world where
the ultimate money is unproblematic: a gold standard \ts{2}{2}{4:59}, with four layers
\ts{2}{2}{5:42}:
\begin{enumerate}
  \item \textbf{Gold}, the ultimate money, the international means of payment.
  \item \textbf{National currency}, a liability of the central bank and a promise to pay gold at the mint parity.
  \item \textbf{Bank deposits}, liabilities of private banks, promises to pay currency on demand: twice-removed
        promises to pay gold.
  \item \textbf{Securities} (and loans of every kind), promises to pay currency at some date in the future: even
        more attenuated promises.
\end{enumerate}

\begin{nfigure}
\begin{tikzpicture}[font=\small]
  \fill[keycol!35] (0,4) -- (-0.8,3) -- (0.8,3) -- cycle;
  \fill[defcol!30] (-0.8,3) -- (-1.6,2) -- (1.6,2) -- (0.8,3) -- cycle;
  \fill[intcol!25] (-1.6,2) -- (-2.4,1) -- (2.4,1) -- (1.6,2) -- cycle;
  \fill[fincol!20] (-2.4,1) -- (-3.2,0) -- (3.2,0) -- (2.4,1) -- cycle;
  \draw[thick] (0,4) -- (-3.2,0) -- (3.2,0) -- cycle;
  \foreach \y in {1,2,3} \draw (-3.2+0.8*\y,\y) -- (3.2-0.8*\y,\y);
  \node at (0,3.35) {gold};
  \node at (0,2.5) {currency};
  \node at (0,1.5) {deposits};
  \node at (0,0.5) {securities};
  \draw[-{Stealth[length=4pt]},thick] (-4.6,0.2) -- (-4.6,3.8) node[midway,left,align=right,font=\scriptsize]{more\\``moneyness''};
  \node[left,font=\scriptsize] at (-4.6,3.9) {money};
  \node[left,font=\scriptsize] at (-4.6,0.1) {credit};
  % three dividing lines
  \draw[dashed,thmcol,thick] (0.8,3) -- (5.2,3) node[right,align=left,font=\scriptsize]{international\\settlement};
  \draw[dashed,defcol,thick] (1.6,2) -- (5.2,2) node[right,align=left,font=\scriptsize]{bank settlement\\(19th-century view)};
  \draw[dashed,intcol,thick] (2.4,1) -- (5.2,1) node[right,align=left,font=\scriptsize]{retail settlement\\(M1, M2, \dots)};
  \node[right,font=\scriptsize,align=left] at (5.2,3.9) {where the line between\\money and credit falls:};
\end{tikzpicture}
\caption{The simple hierarchy of instruments under a gold standard \ts{2}{2}{5:42}. Each layer is a promise to pay the
layer above. The line between money and credit depends on who is settling \ts{2}{2}{7:05}.}\label{fig:nh-instruments}
\end{nfigure}

\paragraph{Where is the line between money and credit?} It depends on the point of view \ts{2}{2}{7:05}:
\begin{itemize}
  \item \textbf{Banks.} Banks clear payments among themselves in national currency, so for a bank settling its
        accounts at the end of the day currency (``high-powered money'') is money and deposits are credit. This was
        the typical 19th-century view \ts{2}{2}{7:26}.
  \item \textbf{Households and firms.} They pay with deposits, so deposits and everything above are money and
        securities are credit. This is what textbooks mean by the money supply, with the ambiguity reflected in the
        many definitions M1, M2, M3 \ts{2}{2}{7:46}.
  \item \textbf{Countries.} Your own currency is no use in paying other countries: they want their own currency, or
        the international means of settlement, gold under a gold standard, perhaps SDRs (the IMF's Special Drawing
        Rights) today. Then only gold is money \ts{2}{2}{8:07}. (The US is the exception, because of the international
        role of the dollar.)
\end{itemize}
So what counts as money depends on who you are and on the layer you are talking about. The course starts from
concepts, not from definitions: in each situation, ask which layer we are at, and what is money and what is credit
there \ts{2}{2}{8:27}. ``Once you get used to thinking about a hierarchy, you'll see it everywhere''
\ts{2}{2}{9:08}.

\begin{secondary}[More layers than four]
Even with four layers the picture is much simpler than the real world (Mehrling's notes). Currency and high-powered
money are not the same: high-powered money includes the deposits of banks at the central bank. Deposits are not
homogeneous, and there are deposit-like instruments, such as money market mutual fund (MMMF) accounts, that are not
the liability of any bank. Securities are even more heterogeneous, by maturity and credit quality. All this
reinforces the point: avoid sterile debates about what is money and what is credit, and stand on the hierarchical
character of the system.
\end{secondary}

\begin{aside}[Not a cloakroom ticket]
A currency backed by gold reserves does not \emph{represent} gold: it is not a cloakroom ticket for gold held
somewhere on the holder's behalf. Reserves only make the promise to pay more credible. This holds even with 100\%
reserves (a currency board): there is still a promise, and a promise can be broken (Mehrling's notes).
\end{aside}

% ------------------------------------------------------------------
\section{The hierarchy of financial institutions}\label{sec:nh-institutions}

The same hierarchy can be seen through the institutions that issue and hold the instruments, i.e.\ through
balance sheets \ts{2}{3}{0:08}:
\begin{itemize}
  \item the \textbf{central bank}, a bankers' bank, holds reserves of the international money (gold) and issues the
        national money (currency) as its liability \ts{2}{3}{0:35};
  \item the \textbf{banking system} holds currency (or deposits at the central bank) as reserves and issues deposits,
        the private money supply \ts{2}{3}{1:15};
  \item the \textbf{private sector} holds deposits, and also securities (through pension funds, insurance companies),
        and issues securities \ts{2}{3}{1:57}.
\end{itemize}
\begin{taccts}
\tacct[2.0cm]{Central bank}{Gold}{Currency}\qquad
\tacct[2.0cm]{Banking system}{Currency}{Deposits}\qquad
\tacct[2.0cm]{Private sector}{Deposits\\Securities}{Securities}
\end{taccts}
Much is left out: government debt, the main asset of a modern central bank; loans, the main asset of banks; and the
shadow banking system altogether. This is 19th-century, gold-standard banking \ts{2}{3}{2:40}.

\paragraph{Every item appears twice, except gold} \ts{2}{3}{3:01}. Each instrument is an asset of one balance sheet and
a liability of another, except gold, which is no one's liability.

\begin{definition}[New concepts: outside and inside money]\label{def:nh-inside}
\emph{Outside money} is a financial asset that is no one's liability (gold). \emph{Inside money} is money that is
also someone's liability, a form of credit: currency, deposits, the liabilities of shadow banks \ts{2}{3}{3:21}.
If all balance sheets are consolidated into one, inside money cancels (asset of one, liability of another) and only
outside money remains.
\end{definition}

Almost all money is inside money: ``it's not stuff, it's not a heap of anything, it's somebody's promise''
\ts{2}{3}{4:02}. The terms come from John Gurley and Edward Shaw, \emph{Money in a Theory of Finance} (1960), a book
about banking as intermediation, of which the course reads some chapters \ts{2}{3}{4:24}. Surprisingly,
Gurley and Shaw treat \emph{currency} as outside money: they think of a single nation-state with fiat currency, and
consolidate only the private sector, so currency has no counterpart in their accounts \ts{2}{3}{5:06}. Off the gold
standard, the ``ontological status'' of currency is indeed less clear \ts{2}{3}{5:26}. The course takes a global view:
every national money must be convertible, at some exchange rate, into an international money, so currency is
always treated as \emph{inside} money \ts{2}{3}{5:46}.

\begin{intuition}[Consolidation hides what matters]
Consolidating all balance sheets makes the whole edifice of credit vanish, leaving a tiny quantity of gold. But the
consolidation ``misses the entire point'' (Mehrling's notes): interest rates and GDP depend on the gross quantity of
credit, on who issues it, who holds it, and where it lies in the hierarchy. Standard macro models keep only a
quantity of money, usually treated as an outside asset. A physicist would say: the total charge is zero, but the
field depends on where the charges are.
\end{intuition}

% ------------------------------------------------------------------
\section{The dynamics of the hierarchy}\label{sec:nh-dynamics}

\paragraph{The pyramid.} Take every balance sheet in the economy and record each instrument as a liability on the left
and as an asset on the right, at its level of the hierarchy \ts{2}{4}{0:07}. The result is a symmetric
pyramid centred on zero (\cref{fig:nh-pyramid}): the vertical axis is quality, from credit at the bottom to money at
the top; the horizontal axis is quantity \ts{2}{4}{1:56}. The tip is drawn at zero because the quantity of gold is
vanishingly small compared with the credit below it, ``gold is really even just in the mind'' \ts{2}{4}{1:15}: a
pyramid of promises to pay, all the way up \ts{2}{4}{1:35}.

\begin{nfigure}
\begin{tikzpicture}[font=\small]
  % contracted (steep) and expanded (flat) pyramids
  \fill[defcol!10] (0,3.6) -- (-4.8,0) -- (4.8,0) -- cycle;
  \draw[defcol,thick,dashed] (0,3.6) -- (-4.8,0) -- (4.8,0) -- cycle;
  \fill[thmcol!15] (0,3.6) -- (-2.2,0) -- (2.2,0) -- cycle;
  \draw[thmcol,thick] (0,3.6) -- (-2.2,0) -- (2.2,0) -- cycle;
  \draw[-{Stealth[length=4pt]}] (-5.4,0) -- (5.6,0);
  \draw[-{Stealth[length=4pt]}] (0,0) -- (0,4.2) node[above]{money};
  \node[below,font=\scriptsize] at (0,0) {credit};
  \node[below,font=\scriptsize] at (-3.2,-0.05) {quantity of liabilities};
  \node[below,font=\scriptsize] at (3.2,-0.05) {quantity of assets};
  \node[thmcol,font=\scriptsize,align=center] (ct) at (-3.2,2.7) {contraction:\\steep, scarce};
  \draw[thmcol,thin] (ct.east) -- (-0.45,2.2);
  \node[defcol,font=\scriptsize,align=center] at (3.9,1.6) {expansion:\\flat, elastic};
  \draw[-{Stealth[length=4pt]},thick,defcol] (2.2,0.6) -- (3.6,0.6);
  \draw[-{Stealth[length=4pt]},thick,defcol] (-2.2,0.6) -- (-3.6,0.6);
\end{tikzpicture}
\caption{The hierarchy as a pyramid centred on zero \ts{2}{4}{0:07}: liabilities to the left, assets to the right, quality on
the vertical axis. In an expansion the pyramid flattens (more credit per unit of money); in a contraction it becomes
steep again \ts{2}{4}{3:21}.}\label{fig:nh-pyramid}
\end{nfigure}

\paragraph{The system breathes.} The pyramid adds a third dimension: the system is dynamic, at every time scale
\ts{2}{4}{2:17}:
\begin{itemize}
  \item within the day, payments are cleared with expansion and contraction of credit (daylight overdrafts);
  \item over the business cycle, credit booms and busts;
  \item over the rise and fall of nations, wars and depressions.
\end{itemize}
``At all time scales, this system is breathing'': its expansions and contractions are ``the beating heart of the
financial system'' \ts{2}{4}{2:59}. In an expansion the pyramid becomes flatter and wider: the amount of credit for a
given amount of money (a ``branching ratio'') increases; this is \emph{irrational exuberance}. In a contraction it
shrinks and becomes steeper: \emph{financial crisis} \ts{2}{4}{3:21}. The system expands until it hits
some limit, collapses until it hits another, and goes back and forth \ts{2}{4}{4:04}.

\paragraph{Quantity and quality.} The fluctuation has the two dimensions of the picture \ts{2}{4}{4:24}:
\begin{enumerate}
  \item \textbf{quantity}: credit grows in a boom and collapses in a crisis (``deleveraging'');
  \item \textbf{quality}: in a boom credit starts to look like money, moves up the hierarchy, becomes more liquid and
        usable for payments; in a contraction you discover that what you hold is credit, not money: that gold and
        currency are not the same thing, that deposits and currency are not the same thing \ts{2}{4}{4:44}.
\end{enumerate}
The quality dimension shows up in prices: \emph{credit spreads} (the excess of a borrower's interest rate over the
rate on the best credit) become very flat in a boom and blow out in a crisis \ts{2}{4}{5:26}.

\begin{definition}[New concepts: moneyness, credit spread, deleveraging]
The \emph{moneyness} of an instrument is how close it is to money, i.e.\ how readily it is accepted in settlement at
par. A \emph{credit spread} is the difference between the interest rate paid by a borrower and the rate on a
reference of the highest quality. \emph{Deleveraging} is the reduction of debt relative to assets or capital (from
\emph{leverage}, the lever that multiplies the effect of a given capital).
\end{definition}

\begin{finance}[Daily, cyclical, secular]
A single trading day, a business cycle and a century of history show the same pattern of expansion and contraction:
the hierarchy flattens as everyone accepts everyone's promises, then steepens as the promises are tested. The
2007--2009 crisis is the reassertion phase at business-cycle frequency; the end of sterling as world money, which
Young did not foresee (\cref{sec:fp-young}), is the same phenomenon at the secular scale.
\end{finance}

% ------------------------------------------------------------------
\section{Discipline and elasticity}\label{sec:nh-discipline}

Two principles drive the fluctuation \ts{2}{5}{0:08}:

\begin{keyformulas}
\begin{tabularx}{\linewidth}{@{}L c L@{}}
\textbf{scarcity of ultimate money} (discipline) & $\leftrightarrow$ & \textbf{currency principle}\\
\textbf{elasticity of derivative credit} (elasticity) & $\leftrightarrow$ & \textbf{banking principle}\\
\end{tabularx}
\end{keyformulas}

\begin{enumerate}
  \item \textbf{Scarcity (discipline).} There is only so much gold at the top of the pyramid \ts{2}{5}{0:36}. More
        generally, no one can, by their own action, increase the quantity of money at a level above their own:
        governments cannot create gold, banks cannot create government currency, you and I cannot create
        currency, because our liabilities are not currency \ts{2}{5}{2:41}.
  \item \textbf{Elasticity.} Credit is called \emph{derivative} because it is a promise to pay money, denominated in
        money but not money itself; the word should get us used to derivatives before swaps appear
        \ts{2}{5}{0:58}. Agents at any level can expand credit at their own level: if I buy something
        from you and you accept my IOU, the trade happens, and the scarcity of gold has nothing to do with it
        \ts{2}{5}{2:00}. If banks accommodate us, our credit expansion becomes an expansion of
        deposits; if the central bank accommodates the banks, of currency too \ts{2}{5}{3:01}.
\end{enumerate}
In a crisis discipline ``smacks you'': this is money and that is not; in a boom elasticity says everything is money
\ts{2}{5}{1:39}. The system is always some balance between discipline at the top and elasticity at the bottom, and the
shifting balance shows up in prices and in output: it is a way to understand the fundamental instability of credit
\ts{2}{5}{3:22}. Assuming it away because it makes the mathematics easier takes us farther from reality,
not closer \ts{2}{5}{4:03}.

\paragraph{The same swing in monetary theory.} One can build a monetary theory on either principle \ts{2}{5}{4:44}:
\begin{itemize}
  \item from the top, on the scarcity of ultimate money: the \emph{currency principle}; monetarism can be read this way
        \ts{2}{5}{5:04};
  \item from the bottom, on the elasticity of credit: the \emph{banking principle} \ts{2}{5}{5:28}.
\end{itemize}
Internationally the same opposition is \emph{metallism} versus \emph{chartalism}, which come later in the course
\ts{2}{5}{6:33}. When scarcity of ultimate money is the issue, monetarism looks right; when elasticity is the issue,
the Keynesians look right; monetarists are in the ascendant when we need more discipline, Keynesians when we need more
elasticity \ts{2}{5}{6:09}. This makes economists look fickle, but Mehrling's claim is that \emph{both}
theories are right: each captures something true of the system, neither the whole truth \ts{2}{5}{7:36}.
This is his answer to ``which school are you'' (\cref{sec:fp-bigpicture}): embrace the fluctuation, of practice and of
theory, and ask where in the system there is too much scarcity or too much elasticity right now; that makes for a
healthier policy dialogue \ts{2}{5}{8:17}.\begin{history}[Currency School versus Banking School]
The names come from the British debate of the 1830s--1840s on the Bank of England's note issue. The Currency School
(Samuel Jones Loyd, later Lord Overstone; Robert Torrens; George Warde Norman) held that bank notes should vary
exactly as a fully metallic currency would, one for one with gold. The Banking School (Thomas Tooke, John Fullarton)
held that the note issue adapts to the ``needs of trade'' and cannot be overissued if notes are convertible. The
Currency School won: Peel's Bank Charter Act of 1844 split the Bank into an Issue Department, whose notes beyond a
fixed ``fiduciary issue'' of \pounds14 million had to be backed by gold, and a Banking Department. The Act had to be
suspended in the crises of 1847, 1857 and 1866.\par
\emph{References:} F.~W.~Fetter, \emph{Development of British Monetary Orthodoxy 1797--1875}, Harvard UP, 1965;
Young, ``Dear and Cheap Money'' (course reading).
\end{history}

% ------------------------------------------------------------------
\section{The hierarchy of market makers}\label{sec:nh-marketmakers}

Put the institutions back on the hierarchy of instruments: each of them \emph{straddles} two layers, and each is, in
its own way, a market maker \ts{2}{6}{0:07}. In the private sector, focus on security dealers.

\begin{definition}[New concepts: market maker, dealer, bid--ask spread]\label{def:nh-mm}
A \emph{market maker} (or \emph{dealer}) stands ready to buy or to sell an asset at prices it quotes, and absorbs the
difference between purchases and sales in its own inventory. It quotes two prices: the \emph{bid}, at which it buys,
and the \emph{ask} (or offer), at which it sells; the \emph{bid--ask spread} is the difference, its compensation
\ts{2}{6}{1:08}. A dealer trades for its own account, unlike a \emph{broker}, who only matches buyers and sellers.
\end{definition}

\begin{itemize}
  \item \textbf{Security dealers.} If you want to buy Apple stock a dealer sells it to you, if you want to sell it a
        dealer buys it; there are liquid markets in Apple stock because dealers quote prices
        \ts{2}{6}{1:28}. They knit together deposits and securities, and determine the price of securities, for which
        the interest rate stands in \ts{2}{6}{1:49}. In Mehrling's notes, they hold inventories of
        securities and money (in practice, of repos and reverse repos, which give access to securities and deposits
        held elsewhere).
  \item \textbf{Banks} are dealers too, quoting a two-way price in money: give Citibank cash and it gives you a
        deposit, and the other way round \ts{2}{6}{2:51}. The price they maintain is \emph{par}, one for one, a single
        price rather than two; and keeping a price fixed is a challenge \ts{2}{6}{3:12}.
  \item \textbf{The central bank} is a dealer at a higher level, between international and domestic money: if you need
        international money to buy abroad, it sells it to you; if you want to change it into domestic money, it buys
        it. The price made there is the \emph{exchange rate} \ts{2}{6}{3:33}.
\end{itemize}

\begin{nfigure}
\begin{tikzpicture}[font=\small,
  lay/.style={draw,rounded corners=2pt,minimum width=2.6cm,minimum height=0.7cm,fill=#1!15},
  mm/.style={draw,thick,rounded corners=2pt,minimum width=3cm,minimum height=0.6cm,fill=white,align=center}]
  \node[lay=keycol] (g) at (0,4.5) {gold};
  \node[lay=defcol] (c) at (0,3) {currency};
  \node[lay=intcol] (d) at (0,1.5) {deposits};
  \node[lay=fincol] (s) at (0,0) {securities};
  \node[mm] (cb) at (4.3,3.75) {central bank};
  \node[mm] (bk) at (4.3,2.25) {banking system};
  \node[mm] (sd) at (4.3,0.75) {security dealers};
  \foreach \m/\a/\b in {cb/g/c,bk/c/d,sd/d/s} {
    \draw[semithick] (\m.west) -- ++(-0.6,0) |- (\a.east);
    \draw[semithick] (\m.west) -- ++(-0.6,0) |- (\b.east);}
  \node[right,font=\small\bfseries,keycol] at (6.1,3.75) {exchange rate};
  \node[right,font=\small\bfseries,keycol] at (6.1,2.25) {par};
  \node[right,font=\small\bfseries,keycol] at (6.1,0.75) {interest rate};
  \node[font=\scriptsize\bfseries] at (0,5.3) {asset};
  \node[font=\scriptsize\bfseries] at (4.3,5.3) {market maker};
  \node[font=\scriptsize\bfseries,anchor=west] at (6.1,5.3) {price of money};
\end{tikzpicture}
\caption{The simple hierarchy of market makers \ts{2}{6}{0:07}: each market maker has one leg in each of two adjacent
layers (the two sides of its balance sheet), and the price at which it trades them is one of the prices of money of
\cref{def:fp-four}.}\label{fig:nh-marketmakers}
\end{nfigure}

\paragraph{From quality to quantity.} The prices of money are themselves hierarchical, one between each pair of layers
\ts{2}{6}{4:15}. As long as these prices exist, the hierarchy is hidden: there is a price at which securities exchange for
deposits, one at which deposits exchange for currency, and the qualitative difference looks like a quantitative
similarity \ts{2}{6}{4:56}. Economics abstracts from money to see good A and good B exchanging at a relative
price; applied to money, the same impulse loses the hierarchy \ts{2}{6}{4:35}. Economics prices a kidney and
an army tank in the same way; the hierarchy of money is a reminder that some things are different in kind
\ts{2}{6}{5:57}. In normal times the approximation is not so bad; in a crisis it fails.

\paragraph{Market makers under stress.} As the system swings between elasticity and scarcity, the market makers stand
in the middle, one leg on each level \ts{2}{6}{6:38}. When the hierarchy reasserts itself they are pulled
in two directions: long the lower layer and short the upper one, if the upper one becomes scarce it rises in value, and
they can become insolvent, but ``liquidity kills you quick'': if you cannot settle your books at the end of the day, you
are done \ts{2}{6}{7:19}. All three face this constraint, which is where the course starts next
\ts{2}{6}{8:00}.
\begin{itemize}
  \item Security dealers have it easier, because in a crisis they can change their price: asset prices plunge
        \ts{2}{6}{8:00}.
  \item Banks cannot: a bank offering ``80 cents on the dollar'' for its deposits would not be a bank any more
        \ts{2}{6}{8:20}.
  \item The same holds for a fixed exchange rate or the gold parity \ts{2}{6}{8:41}.
\end{itemize}
In the worst case the market makers lose all their money and stop making markets, and one cannot move from one level
of the hierarchy to another at any price, because there is no dealer: that is a financial crisis \ts{2}{6}{9:01}.

\begin{intuition}[The hierarchy is natural, not imposed]
Mehrling calls the hierarchy \emph{natural} (notes, section IV): the institutions are hierarchical \emph{because} the
underlying system is, not the other way round; the hierarchy is not imposed from the top by the state or the central
bank, it grows from the ground up. (He adds that he does not mean self-organisation in the sense of Austrian economists
such as Carl Menger.) The market makers are what turn qualitatively different layers into a continuum of prices.
\end{intuition}

% ------------------------------------------------------------------
\section{Managing the hierarchy}\label{sec:nh-managing}

So far the market makers have been profit maximizers: dealers and banks do not create liquidity to be nice, they do it to
make money, and the course will analyse the source of that profit \ts{2}{7}{0:08}. One entity is different.
Once the central bank becomes an agent of the state rather than a private bankers' bank, it need not maximize (short-term)
profit \ts{2}{7}{0:49}. Sitting at the top, it notices that when things get really bad what everybody wants is its own
liabilities, the best in the system, and that it can help \ts{2}{7}{1:10}. This happened gradually: the Bank of England was
a private bank for a long time; in the US, J.~P.~Morgan acted as a private central bank in 1907 \ts{2}{7}{1:30}.

\begin{definition}[New concepts: lender of last resort]\label{def:nh-lolr}
A \emph{lender of last resort} is an institution that, when the scarcity of money threatens to make the market makers stop
making markets, introduces elasticity from the top: it lends its own liabilities (expands the money supply) against
securities taken as collateral, and so rejoins the layers of the hierarchy that have come apart \ts{2}{7}{2:35}.
(``Last resort'': the lender to whom one turns when all other sources of funding are gone.)
\end{definition}

\paragraph{The widening remit of central banks} \ts{2}{7}{3:17}:
\begin{enumerate}
  \item \textbf{defend the exchange rate} (the mint parity), the most primitive job; central banks that did only this and
        ``let the rest of the economy go to hell'' learned that they would be sucked down with it \ts{2}{7}{3:59};
  \item \textbf{help banks maintain par}: in a run, people want currency instead of deposits, and the central bank can create
        currency: financial stability, the lender of last resort \ts{2}{7}{4:21};
  \item \textbf{monetary policy}: once the central bank has taken responsibility for the crisis, it would rather prevent it,
        leaning against the wind: raising rates to impose scarcity when the economy is all elasticity, lowering them to
        inject elasticity when it is all discipline. Counter-cyclical policy comes from taking responsibility for the
        crisis \ts{2}{7}{6:07};
  \item \textbf{regulation}: making sure banks do not get overextended or ``buy crap assets'' comes from the same source
        \ts{2}{7}{7:29};
  \item at the extreme, the ambition to eliminate the business cycle altogether by creating an artificial hierarchy
        \ts{2}{7}{14:16}.
\end{enumerate}

\begin{aside}[A copy of \emph{Lombard Street}]
Mehrling shows his used copy of Bagehot's \emph{Lombard Street}, inscribed ``Merry Christmas to father from Ed, 1907'': the
year of a financial crisis, and Ed was giving his father the key to understanding it \ts{2}{7}{4:44}.
\end{aside}
Bagehot's point was that the Bank of England, a private central bank, had in past crises been acting as a lender of last
resort without intending to, and should bring this to consciousness and embrace it: in a crisis, lend freely, at a high
rate, against good collateral. It was very controversial when published, then became orthodoxy \ts{2}{7}{5:27}.

\paragraph{Policy rate and market rates.} Most central banks intervene by setting a very short-term interest rate (in the
US the Fed funds rate, an overnight rate) \ts{2}{7}{8:18}. Rates at different maturities form the \emph{term structure}
(\cref{fig:nh-term}): at the short end, overnight money, set by policy; at the long end, 30-year Treasury bonds, set by
private security markets. ``So who wins?'' \ts{2}{7}{8:39}

\begin{definition}[New concepts: term structure, expectations hypothesis]\label{def:nh-eh}
The \emph{term structure of interest rates} (the yield curve) is the interest rate as a function of maturity. The
\emph{expectations hypothesis} says that a long rate is the expected average of the future short rates over its life:
with $R_n$ the yield on an $n$-period bond and $r_t$ the one-period rate at $t$,
\[
  (1+R_n)^n = \E\Bigl[\,\prod_{t=0}^{n-1}(1+r_t)\Bigr], \qquad\text{approximately}\qquad R_n \approx \frac1n\sum_{t=0}^{n-1}\E[r_t].
\]
It looks like an arbitrage condition (rolling over overnight loans for 30 years versus lending for 30 years), yet it is
rejected empirically ``basically all the time'' \ts{2}{7}{10:01}. Mehrling will argue that the reason is
liquidity.
\end{definition}

\begin{nfigure}
\begin{tikzpicture}[font=\small]
  \draw[-{Stealth[length=4pt]}] (0,0) -- (7.2,0) node[right]{maturity};
  \draw[-{Stealth[length=4pt]}] (0,0) -- (0,4) node[left]{rate};
  \draw[very thick,defcol] (0,1.2) .. controls (1.2,2.2) and (3,2.6) .. (7,2.7) node[right,font=\scriptsize]{normal};
  \draw[very thick,intcol] (0,0.3) .. controls (1.2,1.9) and (3,3.1) .. (7,3.4) node[right,font=\scriptsize]{elastic (steep)};
  \draw[very thick,thmcol] (0,3.3) .. controls (1.2,2.4) and (3,2.0) .. (7,1.9) node[right,font=\scriptsize]{discipline (inverted)};
  \node[left,font=\scriptsize,align=right] at (-0.1,1.2) {policy rate\\(Fed funds)};
  \node[below,font=\scriptsize] at (0.3,0) {overnight};
  \node[below,font=\scriptsize] at (6.6,0) {30 years};
  \node[font=\scriptsize,align=center] at (5.2,0.8) {market rates\\(security markets)};
\end{tikzpicture}
\caption{The term structure of interest rates \ts{2}{7}{8:39}: the central bank sets the overnight end, security markets the
long end. To impose discipline it pushes short rates above long rates (inverted curve); to make the pyramid more elastic
it pushes them far below \ts{2}{7}{12:09}.}\label{fig:nh-term}
\end{nfigure}

\paragraph{The artificial hierarchy.} Policy tries to move the market by moving the short rate, but the market constrains
what policy can do; and if policy simply follows the market it accepts the hierarchy as it is. ``Monetary policy is all
about creating an artificial hierarchy of money and credit'' \ts{2}{7}{11:03}. In the term structure overnight money at the
Fed plays the role of currency, long bonds the role of securities \ts{2}{7}{11:28}. To impose discipline the
central bank makes short-term money ``ridiculously expensive'', more than long-term money: an inverted (negatively sloped)
curve. To make the pyramid more elastic it lowers short rates far below long rates \ts{2}{7}{12:09}. Choosing
the point between elasticity and discipline is monetary policy, with its politics \ts{2}{7}{12:52}.

\begin{intuition}[A thin reed]
The whole of monetary policy hangs on an overnight interest rate, ``a very thin reed'', while the credit pyramid is
enormous. How can the ``puny little Fed'' have any influence on it? Where exactly is the transmission mechanism? This is a
serious question that the course will be able to ask precisely once inside the system \ts{2}{7}{13:12}.
Mehrling's notes add: the main policy instrument is the overnight rate, and there can be considerable ``slippage''
between it and the long rates on which investment and spending depend.
\end{intuition}\begin{history}[Punch bowl and string]
William McChesney Martin, Fed chairman 1951--1970, said in a speech of October 1955 that the Fed is ``in the position of
the chaperone who has ordered the punch bowl removed just when the party was really warming up'': raising rates in a boom.
``Pushing on a string'', for the impotence of lower rates in a slump, is usually attributed to Keynes, but the first
recorded use is in the 1935 House hearings on the Banking Act, in an exchange between Congressman T.~Alan Goldsborough and
Fed chairman Marriner Eccles. Mehrling's notes quote both, attributing the second to Keynes as is customary.\par
\emph{References:} W.~McC.~Martin, address to the New York Group of the Investment Bankers Association, 19 Oct 1955;
US House, \emph{Banking Act of 1935}, Hearings, March 1935.
\end{history}

\paragraph{A brave new world.} The central question of the course is the proper role of the central bank today
\ts{2}{7}{14:36}. Draghi, putting a floor under the price of peripheral bills by buying them himself, was acting as a
\emph{security dealer}, a role nobody had anticipated \ts{2}{7}{14:57}. The Fed, which had never held a
mortgage-backed security since 1913, bought about a trillion dollars of them in a few months
\ts{2}{7}{15:19}. Central banks are reacting, making it up: we live in a ``Bagehot moment''. Bagehot looked
back at what the Bank had done and told it to embrace it; we have no simple rule like his for today's global, shadow-banking
system, and the course is part of the search for one \ts{2}{7}{16:00}.

\begin{finance}[Operation Twist]
Asked about \emph{Operation Twist} (2011--2012: the Fed sold short-term Treasuries and bought long-term ones, to lower
long rates without expanding its balance sheet), Mehrling answers that by buying long-term bonds instead of only fixing the
overnight rate the Fed is trying to influence the shape of the pyramid directly; whether it works, and why, is much debated
\ts{2}{7}{17:28}. The name recalls the 1961 operation of the same kind, which ``twisted'' the yield curve.
\end{finance}

% ------------------------------------------------------------------
\section*{Key formulas and ideas}
\begin{keyformulas}
\begin{tabularx}{\linewidth}{@{}l L@{}}
Hierarchy & gold $>$ currency $>$ deposits $>$ securities: each layer a promise to pay the layer above\\
Money vs credit & means of final settlement vs promise to pay; the line depends on the layer\\
Inside / outside & inside money $=$ someone's liability (cancels on consolidation); outside $=$ no one's (gold)\\
Dynamics & expansion: flatter pyramid, more credit per unit of money, credit looks like money; contraction: steeper, spreads widen\\
Two principles & scarcity of ultimate money (currency principle) vs elasticity of derivative credit (banking principle)\\
Market makers & central bank: exchange rate; banks: par; security dealers: interest rate\\
Expectations hypothesis & $(1+R_n)^n = \E\prod_{t}(1+r_t)$, i.e.\ $R_n\approx \frac1n\sum_t\E r_t$ (empirically rejected)\\
Monetary policy & creating an artificial hierarchy by setting the overnight rate\\
\end{tabularx}
\end{keyformulas}

\section*{Exercises}

\begin{exercise}[Where is the money?]
For each agent, say which instruments are money and which are credit from its point of view, and why:
\begin{enumerate}[(a)]
  \item a household paying its rent;
  \item a bank settling its payments with other banks at the end of the day;
  \item the central bank of a small country paying for imports under a gold standard;
  \item the same central bank today.
\end{enumerate}
\end{exercise}

\begin{exercise}[Consolidation]
Write the balance sheets of \cref{sec:nh-institutions} with numbers: gold 10; currency 100, held by banks; deposits 1000,
held by the private sector; securities 5000, issued by the private sector and held half by the private sector itself.
\begin{enumerate}[(a)]
  \item Consolidate the three sectors. What remains?
  \item Consolidate only the private sector and the banking system, as Gurley and Shaw do. What is ``outside'' money now?
  \item Compute the ``branching ratios'' deposits/currency and currency/gold. What happens to them in an expansion?
\end{enumerate}
\end{exercise}

\begin{exercise}[Draghi versus Soros in T-accounts]
Peripheral governments have issued 100 of excess long-term bonds, held by banks.
\begin{enumerate}[(a)]
  \item Record Soros's plan step by step: the EFA buys the bonds with Euro bills; the ECB buys 100 of Euro bills from the
        EFA's buyers with new reserves; banks then swap 100 of excess reserves for Euro bills from the ECB.
  \item What is the final change in the ECB's balance sheet? In the money supply?
  \item The EFA forgives the bonds. Whose net worth falls, and by how much?
\end{enumerate}
\end{exercise}

\begin{exercise}[The expectations hypothesis]
The overnight rate is 1\% today and is expected to rise by 1 percentage point a year for the next 4 years, then stay
constant.
\begin{enumerate}[(a)]
  \item Using the approximate expectations hypothesis, compute the 2-, 5- and 10-year yields and sketch the term
        structure.
  \item The observed 10-year yield is 1 percentage point higher than your answer. Give two possible explanations.
\end{enumerate}
\end{exercise}

\begin{exercise}[Market makers in a crisis]
For each of the three market makers of \cref{fig:nh-marketmakers}, describe what happens to its balance sheet when the
layer above becomes scarce, and what options it has: change its price, use reserves, borrow from the layer above, stop
making the market.
\end{exercise}

% !TEX root = ../main.tex
\chapter{Money and the State: Domestic}\label{ch:state}
\begin{paracol}{2}

\begin{sources}
\begin{tabularx}{\linewidth}{@{}l l L@{}}
\toprule
\textbf{Segment} & \textbf{Length} & \textbf{Content}\\
\midrule
\seg{3}{1}{L3.1 FT: Quantitative Easing and the Fed} & 7:47 & The Fed's balance sheet 2007--2010; QE; the expected QE3.\\
\seg{3}{2}{L3.2 Allyn Young: Money and Economic Orthodoxy} & 9:09 & Where Young stands against the orthodoxy of economists and against outside opinions.\\
\seg{3}{3}{L3.3 National Banking System before the Fed} & 3:22 & The pyramid of reserves: country banks, reserve cities, New York.\\
\seg{3}{4}{L3.4 Civil War Finance, bonds and loans} & 8:49 & Bond sales, the \$150 million bank loan of 1861, suspension of gold payments.\\
\seg{3}{5}{L3.5 Civil War Finance, legal tenders} & 7:07 & The greenbacks: state money inserted into the hierarchy.\\
\seg{3}{6}{L3.6 National Banking System, origins} & 6:30 & Bonds sold to banks, bank notes backed by bonds; a fixed quantity of base money.\\
\seg{3}{7}{L3.7 National Banking System, instability} & 5:57 & Seasonal demand for currency, call loans, crises; 1907 and clearinghouse certificates.\\
\seg{3}{8}{L3.8 Federal Reserve System, plan} & 6:32 & Discounts and rediscounts: elastic reserves and elastic currency backed by loans to Main Street.\\
\seg{3}{9}{L3.9 Federal Reserve System, actual} & 6:48 & War finance fills the Fed with Treasury bills; open market operations; state money versus private money.\\
\bottomrule
\end{tabularx}

\smallskip
\textbf{Files.} Transcripts \file{materials/transcripts/L03-P0*.txt}.
\textbf{Reading.} Allyn Young's four chapters (\file{Allyn Young.pdf}, see \cref{sec:fp-young}); this lecture translates
their history of American money into balance sheets.
\end{sources}

\section*{Motivation}
Lecture~2 left the state out: the central bank was just a bankers' bank \ts{2}{2}{2:11}. Here the state comes in, through
the history that Allyn Young tells: how the United States financed the Civil War, how that war finance produced the
National Banking System, why that system was unstable, and how the Federal Reserve was designed to fix it and what it became
instead. The history serves two purposes \ts{3}{9}{3:37}:
\begin{enumerate}
  \item practice with balance sheets: every episode is a set of T-accounts;
  \item a conceptual point: the monetary system is a \emph{hybrid} of state money and private money, and a theory built on only
        one of the two misses half of it.
\end{enumerate}
A recurring pattern: \emph{the history of American central banking is a history of war finance} \ts{3}{8}{0:07}.

% ------------------------------------------------------------------
\section{FT: quantitative easing and the Fed}\label{sec:st-ft}

\begin{casestudy}[FT, September 2012: ``Fed's expected QE move weighs on the dollar'']
Global stocks rallied and the dollar was sold amid expectations that the Fed would announce a third round of quantitative
easing, and that Germany's constitutional court would allow Europe's new bailout fund (the court did, with reservations)
\ts{3}{1}{0:28}. ``QE'' is a word not yet in the textbooks: it was invented during the crisis
\ts{3}{1}{1:09}.
\end{casestudy}

\paragraph{The Fed's balance sheet in the crisis.} The chart shown in class plots the Fed's balance sheet from January 2007 to
March 2010, assets as positive numbers and liabilities as negative ones, with two vertical lines: the collapse of Bear Stearns
(March 2008) and of Lehman Brothers and AIG (September 2008) \ts{3}{1}{1:52}. It shows three phases.
\begin{enumerate}
  \item \textbf{Interest rates.} The crisis began in August 2007, and the Fed first lowered its rate from 5\% to 2\% in steps,
        trying to make the system more elastic \ts{3}{1}{2:33}. It was not enough: Bear Stearns failed \ts{3}{1}{2:53}.
  \item \textbf{Changing the composition of assets.} Before the crisis the Fed held essentially Treasury bills, against currency
        \ts{3}{1}{3:13}. After Bear Stearns it sold about half of its bills and lent the proceeds directly to banks and
        broker-dealers to support financial markets; liabilities did not change \ts{3}{1}{3:34}.
  \item \textbf{Expanding both sides.} After Lehman and AIG this was not enough either, and the Fed expanded its own liabilities:
        the balance sheet roughly doubled in a few weeks \ts{3}{1}{3:54}. First came short-term loans, to domestic
        and to international banks: the \emph{central bank liquidity swaps} with the ECB, the Bank of Japan and others reached
        about \$600 billion, because the problem had become global \ts{3}{1}{4:34}. As these facilities ran off
        around the turn of the year, the Fed started buying mortgage-backed securities, about a trillion dollars of them
        \ts{3}{1}{5:15}.
\end{enumerate}

\begin{nfigure}
\begin{tikzpicture}[font=\scriptsize,x=0.42cm,y=0.42cm]
  % stylised, not data: time on x (2007-2010), assets up, liabilities down
  \draw[-{Stealth[length=3pt]}] (0,0) -- (14.5,0) node[right]{time};
  \draw[-{Stealth[length=3pt]}] (0,-6.5) -- (0,6.5);
  \node[left] at (0,5) {assets}; \node[left] at (0,-5) {liabilities};
  % assets
  \fill[defcol!45] (0,0) -- (0,3.6) -- (5,3.6) -- (6,1.8) -- (14,1.8) -- (14,0) -- cycle;
  \fill[fincol!45] (5,3.6) -- (6,1.8) -- (8,1.8) -- (9,1.8) -- (8.5,5.5) -- (7,5.8) -- cycle;
  \fill[cscol!45] (7,5.8) -- (8.5,5.5) -- (9,4) -- (14,6) -- (14,1.8) -- (9,1.8) -- (8,1.8) -- cycle;
  % liabilities
  \fill[defcol!25] (0,0) -- (0,-3.6) -- (14,-4.2) -- (14,0) -- cycle;
  \fill[thmcol!30] (6.5,-3.7) -- (7.2,-6.2) -- (14,-6.0) -- (14,-4.2) -- cycle;
  \draw[dashed] (5,-6.5) -- (5,6.5) node[above]{Bear Stearns};
  \draw[dashed] (7,-6.5) -- (7,6.5) node[above right]{Lehman, AIG};
  \node[align=center] at (2.5,1.8) {Treasury\\bills};
  \node[align=center] at (7.2,3.3) {loans,\\swaps};
  \node[align=center] at (11.6,3.8) {MBS};
  \node[align=center] at (3,-1.8) {currency};
  \node[align=center] at (10.5,-5.1) {bank reserves};
  \node[below] at (0.5,-6.5) {2007}; \node[below] at (13.5,-6.5) {2010};
\end{tikzpicture}
\caption{Schematic of the Fed's balance sheet in the crisis, after the chart shown in class \ts{3}{1}{2:12} (shapes only
indicative, not data): first a change of composition (bills sold, loans made), then an expansion of both sides (loans, swaps,
then mortgage-backed securities, against reserves).}\label{fig:st-fedbs}
\end{nfigure}

\begin{coursenote}[QE1, QE2, QE3]
In class Mehrling calls the expansion after Lehman ``QE1'' and the purchases of mortgage-backed securities ``QE2''
\ts{3}{1}{4:14}. In the usual numbering the first program of large-scale asset purchases (QE1, from
November 2008, extended in March 2009) included the MBS purchases, about \$1.25 trillion by March 2010; QE2 (November
2010) was \$600 billion of long-term Treasuries; QE3, announced on 13 September 2012, the day after this lecture's FOMC
meeting, was open-ended MBS purchases of \$40 billion a month. Mehrling also says ``March of 07'' for Bear Stearns, which
failed in March 2008.
\end{coursenote}\begin{aside}[Why ``QE1'']
For someone ``of a certain age'', QE1 and QE2 are ocean liners, the \emph{Queen Elizabeth} and the \emph{Queen Elizabeth 2};
Mehrling guesses the joke was coined by the \emph{Financial Times}, a British paper \ts{3}{1}{4:14}.
\end{aside}

\paragraph{QE3: two ways to lower long rates.} At Jackson Hole the Fed had signalled that the economy was softening. Short
rates were already near zero; the question was how to lower the long end of the term structure further
\ts{3}{1}{5:36}:
\begin{enumerate}
  \item \textbf{Announce} that short rates will stay low for a very long time: if the market believes it, long rates fall, by the
        expectations hypothesis (\cref{def:nh-eh}) \ts{3}{1}{6:38};
  \item \textbf{buy long bonds}, expanding the balance sheet on both sides, which pushes up their prices: a capital gain for
        whoever owns them, hence the excitement of the market \ts{3}{1}{6:59}.
\end{enumerate}
In practice the Fed would probably do both \ts{3}{1}{6:59}.

\begin{definition}[New concepts: quantitative easing]
\emph{Quantitative easing} (QE) is the purchase by a central bank of large quantities of assets, especially long-term ones,
paid for by creating reserves, i.e.\ by expanding both sides of its balance sheet, when the short-term policy rate is already
near zero. ``Quantitative'' because the policy is stated as a quantity of purchases rather than as a rate.
\end{definition}

% ------------------------------------------------------------------
\section{Allyn Young against orthodoxy}\label{sec:st-young}

Young wrote in 1924, after the First World War; he had been the economist of the American delegation at Versailles, a
contemporary of Keynes but a much quieter personality \ts{3}{2}{0:50}. His encyclopedia chapters read
smoothly, and do not cite the literature they answer, so it is easy to miss that he takes position on many controversial
issues \ts{3}{2}{1:31}. Mehrling lists six.

\paragraph{Against economic orthodoxy} \ts{3}{2}{1:54}:
\begin{enumerate}
  \item \textbf{Barter.} Economics tells the story that money was invented to escape the inconvenience of barter (a trope Mehrling
        attributes to Jevons). Young: there was never a period of barter; money has been there ever since trade, it is
        ``baked into the cake'' \ts{3}{2}{1:54}.
  \item \textbf{Growth.} A developing economy needs a proper financial system. Orthodoxy then and now disagrees: in the Solow growth
        model growth comes from capital accumulation, population and productivity, technical change is a residual, and any
        growth due to finance is dropped into that residual; money is a veil to look through \ts{3}{2}{2:35}.
  \item \textbf{The currency principle.} Most economists think money gets its value from scarcity relative to demand, as in the
        neoclassical theory of value, and fear over-issue \ts{3}{2}{3:59}. Young criticises the currency principle:
        the National Banking System was built on it, limiting both note and deposit issue, and the lack of elasticity caused
        periodic crises; the Fed arose to give elasticity \ts{3}{2}{4:40}.
\end{enumerate}
Young was professor of monetary economics at Harvard and adviser to Benjamin Strong of the New York Fed, but his view was a
minority one among economists, then and now, though not inside central banks \ts{3}{2}{5:22}.

\paragraph{Against outside opinions} \ts{3}{2}{5:43}:
\begin{enumerate}[resume]
  \item \textbf{Chartalism.} Against the idea that the dollar has value because the state declares it legal tender
        \ts{3}{2}{6:04}.
  \item \textbf{Efficient markets.} Young points out how the inefficiency of the monetary system made asset prices swing apart from
        fundamental values; he is against speculation and for public policy to control bubbles \ts{3}{2}{6:45}.
  \item \textbf{Populism.} Above all he defends the Federal Reserve against those who saw it as a behemoth, and, against William
        Jennings Bryan, he wants the gold standard: the dollar is valuable because it is convertible into gold
        \ts{3}{2}{7:47}. This fits the hierarchy of lecture~2, with gold at the top. But the gold
        standard does not work perfectly by itself: it needs better financial institutions \ts{3}{2}{8:48}.
\end{enumerate}\begin{history}[Bryan's ``Cross of Gold'' (1896)]
At the Democratic national convention in Chicago, on 9 July 1896, William Jennings Bryan closed his speech for the free coinage
of silver with: ``you shall not crucify mankind upon a cross of gold''. He won the nomination but lost the election to William
McKinley; the Gold Standard Act of 1900 put the US formally on gold.\par
\emph{References:} \emph{Official Proceedings of the Democratic National Convention}, 1896; Gold Standard Act, 14 March 1900.
\end{history}\begin{history}[Young at Versailles]
In 1918--1919 Allyn Young headed the economic division of the Inquiry and then of the American Commission to Negotiate Peace
in Paris; Keynes was there for the British Treasury.\par
\emph{References:} C.~P.~Blitch, \emph{Allyn Young: The Peripatetic Economist}, Macmillan, 1995.
\end{history}

% ------------------------------------------------------------------
\section{The National Banking System before the Fed}\label{sec:st-nbs}

The issue of gold versus silver versus fiat money is left to the lectures on international money \ts{3}{3}{0:07}. Young's
chapters contain a table of the reserves of national banks before the Fed, in three groups \ts{3}{3}{0:30}:
\begin{enumerate}
  \item \textbf{Central reserve cities}: New York, Chicago, St.~Louis, required to hold 25\% of their deposits in ``lawful money''
        (currency, or gold) \ts{3}{3}{0:50}. Deposits in the table: \$825 million in New York, \$262 million in
        Chicago, \$116 million in St.~Louis; New York held \$218 million of lawful money, 26.8\% \ts{3}{3}{1:52}. To some extent
        the central bank of the US was the collection of New York banks, ``J.~P.~Morgan and his buddies'', which also faced
        Europe in international transactions \ts{3}{3}{1:11}.
  \item \textbf{Other reserve cities} (Boston, San Francisco, Kansas City, \dots): also 25\%, but partly held as deposits in New
        York, ``due from reserve agents'' \ts{3}{3}{2:12}.
  \item \textbf{Country banks}, about six thousand small banks: 15\% reserves, partly cash (\$199 million) and partly deposits in
        reserve cities or New York (\$226 million) \ts{3}{3}{2:32}.
\end{enumerate}

\begin{nfigure}
\begin{tikzpicture}[font=\small,
  b/.style={draw,rounded corners=2pt,align=center,minimum width=3.4cm,minimum height=0.8cm}]
  \node[b,fill=keycol!15] (g) at (0,3.6) {lawful money (gold, currency)};
  \node[b,fill=defcol!15] (ny) at (0,2.4) {New York banks (25\%)};
  \node[b,fill=intcol!15] (rc) at (0,1.2) {reserve city banks (25\%)};
  \node[b,fill=fincol!15] (cb) at (0,0) {country banks (15\%)};
  \draw[-{Stealth[length=4pt]}] (cb.east) to[out=0,in=0] node[right,font=\scriptsize]{deposits} (rc.east);
  \draw[-{Stealth[length=4pt]}] (rc.east) to[out=0,in=0] node[right,font=\scriptsize]{deposits} (ny.east);
  \draw[-{Stealth[length=4pt]}] (ny.east) to[out=0,in=0] node[right,font=\scriptsize]{reserves} (g.east);
  \draw[-{Stealth[length=4pt]},dashed] (cb.west) to[out=180,in=180] node[left,font=\scriptsize]{cash} (g.west);
\end{tikzpicture}
\caption{The layering of reserves in the National Banking System \ts{3}{3}{2:53}: country banks keep reserves as deposits in a
reserve city, reserve city banks keep them in New York, New York banks keep lawful money. ``Hierarchy, you can see that.''}
\label{fig:st-pyramid}
\end{nfigure}

% ------------------------------------------------------------------
\section{Civil War finance}\label{sec:st-civilwar}

The problem of war finance: the government must mobilise as many resources as it can to win the war \ts{3}{4}{0:07}.
It taxes as much as it can, it conscripts labour, and then it asks for volunteers: it sells bonds \ts{3}{4}{0:49}.
Salmon P.~Chase, Lincoln's Secretary of the Treasury, tried three methods in turn.

\subsection{Stage 1: bonds and loans}
\paragraph{Version 1: bonds sold to the public.} The government issues bonds and receives deposits, which it then spends on war
goods; within the banking system only the ownership of deposits changes, and the total quantity of deposits does not
\ts{3}{4}{1:33}.
\begin{taccts}
\tacct[2.0cm]{Government}{$+$Deposits (G)}{$+$Bonds}\quad
\tacct[2.0cm]{Private sector}{$+$Bonds\\$-$Deposits}{}\quad
\tacct[2.0cm]{Banks}{}{$-$Deposits (private)\\$+$Deposits (G)}
\end{taccts}
Spending the deposits on war goods transfers them back to the private sector \ts{3}{4}{2:56}.

\paragraph{Version 2: a loan from the banks.} If the public will not buy the bonds, or only at a high rate, the government can
borrow directly from the banking system \ts{3}{4}{3:16}. A loan is like a bond, a promise to pay at a future date,
but not a security traded in the open capital market \ts{3}{4}{3:57}. It is a swap of IOUs: the banks say to the government ``I
owe you \$150 million'' (a deposit), the government says ``I owe you \$150 million'' (the loan) \ts{3}{4}{4:18}. This
is what Chase did in August 1861 with the big New York bankers \ts{3}{4}{4:59}.
\begin{taccts}
\tacct[3cm]{Government}{$+$Deposits \$150m}{$+$Loan \$150m}\qquad
\tacct[3cm]{Banking system}{$+$Loan \$150m}{$+$Deposits (G) \$150m}
\end{taccts}
The deposits are bank money ``created from thin air'', an expansion of both sides of the balance sheet; ``you might call this
quantitative easing''. There is no central bank and no private sector in this operation \ts{3}{4}{5:20}.

\paragraph{Chase takes the gold.} Instead of spending the deposits at home, Chase withdrew them: deposits are promises to pay on
demand, and he demanded gold \ts{3}{4}{6:21}. He wanted the gold to buy war material abroad, where dollars were not
wanted; this turned out to be key, because the South had to barter cotton for European war material, while the North could pay
in gold \ts{3}{4}{7:23}.
\begin{taccts}
\tacct[3cm]{Government}{$-$Deposits\\$+$Gold}{}\qquad
\tacct[3cm]{Banking system}{$-$Gold}{$-$Deposits (G)}
\end{taccts}
With their gold reserves gone, the banks could not pay any other deposit in gold and suspended specie payments: the United States
left the gold standard \ts{3}{4}{8:03}.

\begin{history}[The loans of 1861 and the suspension]
In August 1861 the associated banks of New York, Boston and Philadelphia agreed to lend the Treasury \$150 million in three
instalments, paying in coin as the government drew on it. Chase insisted on withdrawing the coin rather than leaving it on deposit.
After the \emph{Trent} affair raised fears of war with Britain, the New York banks suspended specie payments on 30 December 1861,
followed by the Treasury.\par
\emph{References:} W.~C.~Mitchell, \emph{A History of the Greenbacks}, Chicago UP, 1903, part I ch.~2; B.~Hammond,
\emph{Sovereignty and an Empty Purse: Banks and Politics in the Civil War}, Princeton UP, 1970.
\end{history}

\begin{definition}[New concepts: specie, suspension, legal tender]
\emph{Specie} is coined metal (Latin \emph{in specie}, ``in kind''); \emph{specie payments} are the conversion of notes and deposits into
coin on demand, and their \emph{suspension} is the refusal to convert. \emph{Legal tender} is money that, by law, a creditor cannot
refuse in payment of a debt (Latin \emph{tendere}, to offer).
\end{definition}

\subsection{Stage 2: the greenbacks}
Banks would not fall for it twice \ts{3}{5}{0:08}. Chase's second method was legal tender notes, the \emph{greenbacks}. Mehrling
shows a dollar authorized by the act of 3 March 1863 (the third Greenback act), with the words ``this note is a legal tender for one
dollar'': you cannot get gold for it, only another dollar bill, and within the US you cannot refuse it in payment of a debt
\ts{3}{5}{0:29}.

The government buys war goods by paying with notes it prints \ts{3}{5}{2:56}; the private sector deposits
some of them (say 100) in banks, where they become the banks' reserves \ts{3}{5}{4:21}:
\begin{taccts}
\tacct[2.0cm]{Government}{$+$War goods}{$+$Legal tenders 400}\quad
\tacct[2.0cm]{Private sector}{$-$War goods\\$+$Legal tenders 300\\$+$Deposits 100}{}\quad
\tacct[2.0cm]{Banks}{$+$Legal tenders 100}{$+$Deposits 100}
\end{taccts}
Depositors now withdraw legal tenders, not gold. In the language of lecture~2, Chase inserted a layer of state money between gold and
deposits, at the top of the US hierarchy \ts{3}{5}{5:09}.

\paragraph{Consequences.} Legal tenders could not buy foreign goods: foreigners ask their gold value \ts{3}{5}{5:51}. Throughout the
war their value kept falling against gold, to about 50 cents on the dollar (a fall of the exchange rate), and domestic prices rose,
because the government was adding purchasing power instead of transferring it from the private sector as a bond sale does
\ts{3}{5}{6:11}. Young is ``aghast'' at this, while understanding that in war you do what it takes; the
better way is the third one \ts{3}{5}{4:01}.

\begin{coursenote}[Numbers]
In class: \$400 million authorized, ``almost 400 million'' issued, falling to ``50 cents on the dollar'' \ts{3}{5}{0:50}.
Historically the three Legal Tender Acts (25 February 1862, 11 July 1862, 3 March 1863) authorized \$450 million, of which about
\$431 million were outstanding at the peak (1864), and the greenback reached its low in July 1864, at about 35 cents in gold
(Mitchell, 1903).
\end{coursenote}\begin{aside}[Wesley Clair Mitchell]
In these passages Young draws on his close friend Wesley Clair Mitchell, author of \emph{A History of the Greenbacks} (1903),
a long-time professor at Columbia and one of the founders of the National Bureau of Economic Research (1920), the great American
institutionalist and quantitative economist \ts{3}{5}{1:55}.
\end{aside}

\subsection{Stage 3: the National Banking System}\label{sec:st-nbsorigin}
In 1863 Chase went back to bonds \ts{3}{6}{0:09}. He asked the banks to buy, at full price, bonds with a low interest
rate, paying with deposits (no gold withdrawals this time: the banks had none) \ts{3}{6}{0:58}. To sweeten the deal,
banks could \emph{issue bank notes using the bonds as collateral} \ts{3}{6}{1:39}. The government spends the deposits
on war goods; then, when the public wants to withdraw its deposits in currency, the banks print their own notes instead of paying
out their reserves \ts{3}{6}{2:20}:
\begin{taccts}
\textit{Step 1: bonds for deposits, deposits spent on war goods.}\par
\tacct[2.0cm]{Government}{$+$War goods}{$+$Bonds}\quad
\tacct[2.0cm]{Banks}{$+$Bonds}{$+$Deposits}\quad
\tacct[2.0cm]{Private sector}{$-$War goods\\$+$Deposits}{}

\medskip
\textit{Step 2: the public withdraws its deposits in bank notes.}\par
\tacct[2.0cm]{Banks}{}{$-$Deposits\\$+$Bank notes}\quad
\tacct[2.0cm]{Private sector}{$-$Deposits\\$+$Bank notes}{}
\end{taccts}
The bonds were deposited with a central authority as collateral: if the bank failed, a note holder did not need to go to the
bankrupt bank, he could get a government bond for his note \ts{3}{6}{4:48}. In effect the government borrowed from the
banks by allowing them to create money: notes one step removed from legal tenders, promises to pay legal tender, private money
\ts{3}{6}{5:30}.

\paragraph{A fixed base.} The number of such bonds was fixed, so at the end of the war they became the money supply: together with the
legal tenders and gold (after resumption), a \emph{fixed supply of base money} after 1863, which did not expand or contract with the
business cycle \ts{3}{6}{5:50}.

\begin{coursenote}[Which bonds?]
Mehrling describes ``special two percent bonds'' \ts{3}{6}{0:58}. The National Currency Act (25 February 1863) and the National Bank
Act (3 June 1864) required national banks to deposit US bonds with the Comptroller of the Currency to secure their notes (notes up to
90\% of the bonds' value); in the 1860s these were mostly the war bonds paying 5--6\%. The 2\% ``consols of 1930'', used for note
issue until 1935, date from the Gold Standard Act of 1900. Also, the 25\%/15\% reserve requirements were set in 1863--64; Chicago
and St.~Louis became central reserve cities only in 1887.
\end{coursenote}

% ------------------------------------------------------------------
\section{The instability of the National Banking System}\label{sec:st-instability}

Most deposits were in the country banks, \$2,627 million in Young's table, whose reserves were mostly deposits in the money
centres \ts{3}{7}{0:07}. The United States was an agricultural country, and the demand for currency had a strong seasonal cycle: at
harvest time farmers withdrew notes to move the crops; Young mentions \$50 million \ts{3}{7}{0:29}.

\paragraph{The seasonal squeeze} \ts{3}{7}{1:09}:
\begin{enumerate}
  \item farmers withdraw \$50 million in cash from country banks;
  \item country banks replenish their cash by drawing \$50 million on their deposits in New York;
  \item New York banks held total reserves of \$221 million: they lose 25\% of them and fall far below the required 25\%
        \ts{3}{7}{1:30};
  \item to rebuild reserves they call in their \emph{call loans} to stock-market speculators: in slack seasons they had lent idle
        reserves overnight to brokers, and at harvest they want the money back \ts{3}{7}{3:09};
  \item the speculators sell what they bought: stock prices fall, interest rates rise \ts{3}{7}{3:50};
  \item the banks also try to draw reserves from the world money market, London, disrupting the world gold market
        \ts{3}{7}{4:13}.
\end{enumerate}
``The United States is this big child in the world monetary system, and every fall it throws a little fit'' \ts{3}{7}{4:41}.

\begin{taccts}
\tacct[2.4cm]{Country banks}{$-$Deposits in NY 50}{$-$Deposits 50}\quad
\tacct[2.4cm]{New York banks}{$-$Reserves 50\\later: $-$Call loans}{$-$Deposits of country banks 50}
\end{taccts}

\begin{definition}[New concepts: call loan, clearinghouse]
A \emph{call loan} is a loan repayable on demand (``at call''), typically to a stockbroker against securities. A \emph{clearinghouse} is an
association of banks that settles the claims of its members on each other in one place, netting them (from \emph{clearing}, settling
accounts so that they are ``clear'').
\end{definition}

\paragraph{1907.} Europe eventually ``got sick of this'' and sent help, Paul Warburg, to help create a central bank; there was also
domestic agitation for one \ts{3}{7}{5:02}. The key moment was the crisis of 1907, when J.~P.~Morgan showed how to get out of a crisis:
the clearinghouse of the New York banks issued \emph{clearinghouse loan certificates}, a substitute for money, creating elastic reserves,
but privately and only for its members. The Fed's idea was to do the same for everyone, with a central bank \ts{3}{7}{5:23}.
(Clearinghouses are the subject of lecture~5.)\begin{history}[From 1907 to 1913]
The panic of October 1907 started with the failed attempt to corner United Copper and the run on the Knickerbocker Trust; J.~P.~Morgan
organised the rescue, and clearinghouses across the country issued loan certificates. The Aldrich--Vreeland Act (1908) created a National
Monetary Commission and emergency currency; Paul Warburg, a German banker who had settled in New York in 1902, became the main advocate of a
central bank. The Federal Reserve Act was signed by President Woodrow Wilson on 23 December 1913; the twelve Federal Reserve Banks opened in
November 1914.\par
\emph{References:} R.~F.~Bruner and S.~D.~Carr, \emph{The Panic of 1907}, Wiley, 2007; Federal Reserve Act, 1913.
\end{history}

% ------------------------------------------------------------------
\section{The Federal Reserve System: plan and reality}\label{sec:st-fed}

\subsection{The plan}
Attempts by the bankers themselves to create a central bank were not trusted; President Wilson got it through, just in time to use it
for First World War finance \ts{3}{8}{0:07}. Young describes three levels \ts{3}{8}{0:50}:
\begin{enumerate}
  \item \textbf{member banks}, ordinary banks making loans to farmers and merchants, ``Main Street'', by expanding their own deposits
        \ts{3}{8}{1:32};
  \item the \textbf{Federal Reserve Banks}, located in the same cities as the old reserve-city banks (Boston, San Francisco, \dots);
  \item the \textbf{Federal Reserve} itself at the top.
\end{enumerate}
What holds member banks back from lending is the fear that deposits are withdrawn and reserves run out \ts{3}{8}{2:13}. The plan removes
the fear: a bank short of reserves can get more from its Reserve Bank by using its loans as collateral (a \emph{discount loan}), or by
selling them to the Reserve Bank if they are in the right form \ts{3}{8}{2:33}.

\begin{taccts}
\tacct[2.6cm]{Member bank}{$+$Loans to Main Street\\$+$Reserves}{$+$Deposits\\$+$Discount loan}\quad
\tacct[2.6cm]{Reserve bank}{$+$Discount loan}{$+$Reserves of member bank}
\end{taccts}
``The alchemy of banking once again'': the member bank lends without fear, because its loans can be turned into reserves, and the Reserve
Bank creates the reserves from thin air by expanding its own balance sheet: \emph{elastic reserves}, unlike the National Banking System
\ts{3}{8}{3:41}. If depositors want cash instead, Reserve Banks can \emph{rediscount} at the Federal Reserve, which prints
notes: \emph{elastic currency} \ts{3}{8}{4:21}. (This is Young's description, not how it works today.)

\begin{definition}[New concepts: discount, rediscount]
To \emph{discount} a bill is to buy it before maturity for less than its face value; the difference is the interest. A central bank
\emph{discount loan} is a loan to a bank against eligible paper, and to \emph{rediscount} is to discount again paper that has already been
discounted once (by the member bank, then by the Reserve Bank). The rate charged is the \emph{discount rate}.
\end{definition}

The whole idea was a seasonal one: in an agricultural country the demand for currency and deposits fluctuates with the seasons, so the
supply of currency and of reserves must fluctuate too, backed by loans to Main Street \ts{3}{8}{5:26}. The Federal Reserve
Act says the assets of the Reserve Banks are to be these discounts; that is how Wilson sold it \ts{3}{8}{6:08}.

\subsection{The reality}
The Fed's balance sheet before the 2007 crisis held hardly any loans to Main Street: its main asset was Treasury bills \ts{3}{9}{0:08}.
This was true almost from the start, because the US went to war soon after the Fed was created, and at the end of the war the Reserve
Banks, and the member banks too, were ``stuffed full of Treasury bills'' \ts{3}{9}{0:28}. Young, writing in the 1920s,
treats this as temporary, and imagines \emph{open market operations} as a temporary way to provide elasticity \ts{3}{9}{0:48}.

\begin{definition}[New concepts: open market operations]
\emph{Open market operations} are purchases and sales of securities (classically Treasury bills) by the central bank in the open market,
paid for by creating or destroying reserves. They were developed by the New York Fed under Benjamin Strong in the 1920s \ts{3}{9}{1:34}.
\end{definition}

The importance of open market operations is a historical accident: the Fed was never supposed to lend to the government
\ts{3}{9}{2:15}. The framers had the greenbacks in mind: give the government the power to print money and it will print money; so only
loans to Main Street were to back the expansion of reserves and notes. Then came the First World War, ``the cat was out of the bag'', and
everybody got used to a Fed full of Treasury bills, so much so that in 2008 people found it weird when the Fed sold them and lent to banks,
to the private sector, in its commercial paper operations \ts{3}{9}{2:35}. (Mehrling adds that the textbook
account of open market operations changing the money supply is not how they actually work: repo is involved, lecture~7
\ts{3}{9}{1:54}.)

\subsection{State money and private money}
The conceptual point \ts{3}{9}{3:58}: the system is a \emph{hybrid}. The central bank is partly a bankers' bank; in war it is the
government's bank, helping it sell its bonds, as in both world wars, and that is its job, not a perversion of it; in normal times it
facilitates private capital markets; and it moves back and forth between these roles \ts{3}{9}{4:18}.

\begin{intuition}[Two theories, both partly true]
One can build a monetary theory in which the state prints money and everything else is a promise to pay it (chartalism), or one in which
everything is private credit and the state is just a borrower (metallism, starting from gold and private pledges). Neither is any good
alone, because the system is a hybrid that moves back and forth between the two \ts{3}{9}{4:40}. As with
discipline and elasticity (\cref{sec:nh-discipline}), Mehrling is neither chartalist nor metallist, and uses both
\ts{3}{9}{5:40}. The balance sheet approach is the bridge: debits and credits work the same way for state money and for
private money \ts{3}{9}{6:22}.
\end{intuition}

% ------------------------------------------------------------------
\flushside
\section*{Key formulas and ideas}
\begin{keyformulas}
\begin{tabularx}{\linewidth}{@{}l L@{}}
Bond sale to public & government $+$deposits/$+$bonds; only ownership of deposits changes, total deposits unchanged\\
Loan from banks & swap of IOUs: bank $+$loan/$+$deposit; both sides expand (new bank money)\\
Legal tenders & state money inserted between gold and deposits; adds purchasing power: depreciation and inflation\\
National bank notes & bank notes backed by government bonds; fixed base money, inelastic\\
Seasonal squeeze & farmers' cash demand $\to$ drain on New York reserves $\to$ call loans called $\to$ stocks fall, rates rise\\
Fed plan & discount and rediscount of loans to Main Street: elastic reserves, elastic currency\\
Fed reality & war finance: Treasury bills as main asset; open market operations\\
Hybrid & state money and private money; chartalism and metallism each partly true\\
\end{tabularx}
\end{keyformulas}

\section*{Exercises}

\begin{exercise}[Three ways to finance a war]
The government needs 100 of war goods. Write the T-accounts of government, banks and private sector when it
\begin{enumerate}[(a)]
  \item sells 100 of bonds to the public;
  \item borrows 100 from the banks and spends the deposits;
  \item prints 100 of legal tenders, of which the public deposits 40 in banks.
\end{enumerate}
In each case, what happens to the total of deposits and of currency? Which method adds purchasing power?
\end{exercise}

\begin{exercise}[The seasonal squeeze]
Using the numbers of \cref{sec:st-instability}, suppose New York banks have deposits of 825 and reserves of 221, and country banks draw 50.
\begin{enumerate}[(a)]
  \item Compute the New York reserve ratio before and after.
  \item How many call loans must be called (or gold imported) to restore 25\%?
  \item Redo (b) when a Reserve Bank discounts loans for the New York banks.
\end{enumerate}
\end{exercise}

\begin{exercise}[Bank notes backed by bonds]
A national bank buys 100 of government bonds with deposits and may issue notes up to 90\% of their value. Write its balance sheet after the
public withdraws 90 of deposits in notes. Why could the total of notes not respond to the seasonal demand for currency?
\end{exercise}

\begin{exercise}[QE in T-accounts]
The Fed buys 100 of long-term Treasury bonds from a pension fund that banks with Bank A. Write the changes in the three balance sheets and
explain why the market price of bonds is expected to rise.
\end{exercise}

% !TEX root = ../main.tex
\chapter{The Money View, Micro and Macro}\label{ch:moneyview}
\begin{paracol}{2}

\begin{sources}
\begin{tabularx}{\linewidth}{@{}l l L@{}}
\toprule
\textbf{Segment} & \textbf{Length} & \textbf{Content}\\
\midrule
\seg{4}{1}{L4.1 FT: Dealer of Last Resort} & 5:23 & Bernanke's QE3 and Draghi's conditional promise as dealer-of-last-resort actions.\\
\seg{4}{2}{L4.2 Reading: Hyman Minsky} & 3:07 & Minsky and the financial instability hypothesis; Copeland's flow of funds.\\
\seg{4}{3}{L4.3 Payments: Money and Credit} & 5:36 & Paying with money, with pure credit, with a bank's IOU.\\
\seg{4}{4}{L4.4 Payments: Discipline and Elasticity} & 4:13 & The payment system is a credit system.\\
\seg{4}{5}{L4.5 The Survival Constraint} & 3:35 & Every agent as a bank: cash inflow must cover cash commitments.\\
\seg{4}{6}{L4.6 Sources and Uses Accounts} & 6:55 & The sources-and-uses matrix and its two rules.\\
\seg{4}{7}{L4.7 Payment Example: Money and Credit} & 10:05 & A coffee paid in cash, a dinner paid by credit card.\\
\seg{4}{8}{L4.8 Flow of Funds Accounts} & 10:10 & Copeland's macro accounts; statistical discrepancies; what they miss.\\
\seg{4}{9}{L4.9 The Survival Constraint, redux} & 2:40 & Three ways to meet a cash shortfall.\\
\seg{4}{10}{L4.10 Liquidity, Long and Short} & 9:42 & Time pattern of cash flows and commitments; price-insensitive borrowers.\\
\seg{4}{11}{L4.11 Financial Fragility, Flows and Stocks} & 6:23 & The scales of cash flows and commitments; the money market as stress detector.\\
\bottomrule
\end{tabularx}

\smallskip
\textbf{Files.} Transcripts \file{materials/transcripts/L04-P*.txt}.
\textbf{Reading.} Mehrling's article on Hyman Minsky (not among the course files); Minsky, \emph{Stabilizing an Unstable
Economy} (1986), quoted in class.
\end{sources}

\section*{Motivation}
Lectures 1--3 drew the hierarchy with balance sheets, which are \emph{stocks}. This lecture introduces the other half of the
method: \emph{flows} of cash. Its central idea, taken from Hyman Minsky, is to look at every economic unit ``as if it were a
bank'', facing every day a \emph{survival constraint}: cash inflows must cover cash commitments. The lecture develops it at two
levels:
\begin{enumerate}
  \item \textbf{micro}: payments with money and with credit, the sources-and-uses accounts of a single agent, the survival
        constraint and the ways to meet it;
  \item \textbf{macro}: Morris Copeland's flow of funds accounts, and the money market interest rate as the place where the
        balance of the whole economy's cash flows and commitments shows up.
\end{enumerate}

% ------------------------------------------------------------------
\section{FT: dealer of last resort}\label{sec:mv-ft}

\begin{casestudy}[FT, September 2012: QE3 and Münchau on the ECB]
The FOMC has announced that it will keep buying mortgage-backed securities until employment improves: ``quite astonishing'' open-ended
support, prepared by Bernanke's remarks at Jackson Hole \ts{4}{1}{0:07}\ts{4}{1}{0:27}. In the same week Draghi said he would do
whatever is necessary: central bankers are determined to use their balance sheets to prevent anything bad \ts{4}{1}{0:48}. In the FT,
Wolfgang Münchau argues that QE would be the right policy for Europe too \ts{4}{1}{1:09}. He points out that Draghi will buy Spanish or
Italian bonds only if those countries \emph{ask}, which a political leader may be reluctant to do, while Bernanke simply announced
that he would buy \ts{4}{1}{1:30}\ts{4}{1}{1:50}\ts{4}{1}{2:11}.
\end{casestudy}

\paragraph{Promises and purchases.} If Spain and Italy do not ask, why should the price of their bonds move? Only through
expectations, and Münchau, with Mehrling, considers that a weak reed \ts{4}{1}{2:32}. The course looks at what central banks actually
do \ts{4}{1}{2:53}. Mehrling will call both actions \emph{dealer of last resort} actions: the central bank promises to buy bonds at a
particular price, so that no one will ever sell them for less; buying at a high price lowers their yield \ts{4}{1}{3:16}\ts{4}{1}{3:37}.
For this to work the promise must be credible, i.e.\ the central bank must actually buy. Bernanke is buying, and the market price will
move to his price at once; Mehrling predicts that the price of Spanish and Italian bonds will not move all the way to Draghi's level
until Draghi actually buys \ts{4}{1}{3:57}\ts{4}{1}{4:18}.

\begin{intuition}[Why central banks are so bold]
Central banks are doing things they have never done, and boldly, because the fiscal authorities are not: in Europe they are caught up in
austerity and in blaming Greece, in the US Republicans and Democrats cannot agree. The central banks are trying to hold the system
together through the election cycle until fiscal authorities can sort things out, ``not a job that any central banker ever signed up
for'' \ts{4}{1}{4:38}\ts{4}{1}{4:58}.
\end{intuition}\begin{coursenote}[OMT conditionality]
Outright Monetary Transactions (\cref{sec:nh-ft}) require that the country first request assistance from the European Stability
Mechanism and accept its conditions; this is the ``asking'' that Münchau worried about.
\end{coursenote}

% ------------------------------------------------------------------
\section{Reading: Hyman Minsky}\label{sec:mv-minsky}

Mehrling knew Hyman Minsky, who belongs to the American tradition he calls his own (Wesley Clair Mitchell, Edward Shaw, Morris Copeland)
\ts{4}{2}{0:07}. When the crisis started, the \emph{New Yorker} ran a ``Talk of the Town'' piece calling it a ``Minsky moment''; for a man
who grew up in New York, after decades of neglect by economists, that would have been immortality \ts{4}{2}{0:28}\ts{4}{2}{0:49}.

\begin{reading}[Mehrling on Minsky]
Minsky's \emph{financial instability hypothesis}: there is something inherently unstable about financial systems, and the role of the
central bank is to put bounds on that instability, to keep it from flying off on the upside or on the downside; this is not easy
\ts{4}{2}{1:10}. It is a very different vision from the technocratic one of moving the Fed funds rate by another 25 basis points to get
it just right; Minsky's was a mentality prepared for financial crisis \ts{4}{2}{1:30}\ts{4}{2}{1:51}. A common misreading, Mehrling warns,
takes Minsky's hedge, speculative and Ponzi units to be about leverage; read through Copeland, they are about the \emph{alignment in time}
of cash flows and commitments (\cref{sec:mv-longshort}) \ts{4}{10}{0:07}.
\end{reading}\begin{history}[The Minsky moment]
Hyman P.~Minsky (1919--1996), born in Chicago, taught at Washington University in St.~Louis and ended his career at the Levy Economics
Institute of Bard College. The phrase ``Minsky moment'' was coined by Paul McCulley in 1998 for the Russian crisis; John Cassidy's
\emph{New Yorker} piece ``The Minsky Moment'' appeared on 4 February 2008. Mehrling's ``grew up in New York'' refers to Minsky's
youth: he attended high school in New York City.\par
\emph{References:} H.~P.~Minsky, \emph{Stabilizing an Unstable Economy}, Yale UP (Twentieth Century Fund), 1986; J.~Cassidy,
\emph{The New Yorker}, 4 Feb 2008.
\end{history}

\paragraph{Three images of the macroeconomy.} Macroeconomics courses give two images of how money enters the economy
\ts{4}{2}{2:31}:
\begin{enumerate}
  \item the national income and product accounts (NIPA), $C+I+G+NX = Y$;
  \item the quantity theory, $MV = PT$.
\end{enumerate}
A third, proposed by Morris Copeland in 1952, is the \emph{flow of funds}: used on Wall Street and inside the Fed, but never turned into a
proper macroeconomic theory, which many think is what should be attempted \ts{4}{2}{2:52}.

% ------------------------------------------------------------------
\section{Payments: money and credit}\label{sec:mv-payments}

Suppose there are only two people, you and I; each produces what the other wants, but not at the same time, so barter does not work
\ts{4}{3}{0:07}\ts{4}{3}{0:28}. There are three ways to arrange the trade (I buy goods from you):
\begin{enumerate}
  \item \textbf{Money}: you give me goods, I give you a token; only existing money changes pocket \ts{4}{3}{0:56}\ts{4}{3}{1:18}.
  \item \textbf{Pure credit}: I write you a note, ``I owe you''; this requires trust, as between friends who take turns paying lunch
        \ts{4}{3}{1:39}\ts{4}{3}{2:00}\ts{4}{3}{2:41}.
  \item \textbf{Bank credit}: if we do not trust each other but both trust a third party, the bank accepts my IOU and issues its own IOU,
        which you accept \ts{4}{3}{3:42}\ts{4}{3}{4:02}\ts{4}{3}{4:24}.
\end{enumerate}
\begin{taccts}
\textit{(1) Money:}\quad
\tacct[1.9cm]{Me}{$+$Goods\\$-\Delta M$}{}\quad
\tacct[1.9cm]{You}{$-$Goods\\$+\Delta M$}{}

\medskip
\textit{(2) Pure credit:}\quad
\tacct[1.9cm]{Me}{$+$Goods}{$+\Delta$IOU}\quad
\tacct[1.9cm]{You}{$-$Goods\\$+\Delta$IOU}{}

\medskip
\textit{(3) Bank credit:}\quad
\tacct[1.6cm]{Me}{$+$Goods}{$+$IOU}\quad
\tacct[1.6cm]{Bank}{$+$IOU (mine)}{$+M$}\quad
\tacct[1.6cm]{You}{$-$Goods\\$+M$}{}
\end{taccts}
The difference between (1) and the other two: with money the quantity of money does not change, existing pieces of paper move; with
credit the quantity of outstanding IOUs \emph{increases} with the transaction \ts{4}{3}{3:01}\ts{4}{3}{3:21}. In (3) the bank's IOU is
called $M$ because it is used as money, but it is a liability of the bank created by expanding its balance sheet: money is a form of
credit, and it expands and contracts with the pattern of trade \ts{4}{3}{4:45}\ts{4}{3}{5:06}. This stripped-down picture is how the
modern economy works \ts{4}{3}{5:27}.

\section{Discipline and elasticity in payments}\label{sec:mv-de}

\begin{nfigure}
\begin{tikzpicture}[font=\scriptsize,x=0.55cm,y=0.55cm]
  \begin{scope}
    \draw[-{Stealth[length=3pt]}] (0,0) -- (9,0) node[right]{time};
    \draw[-{Stealth[length=3pt]}] (0,0) -- (0,5);
    \draw[dashed] (0,4) -- (8.5,4) node[right]{total $M$};
    \draw[thick,defcol] plot[smooth] coordinates {(0,2) (1,2.8) (2,3.6) (3,2.6) (4,1.2) (5,0.3) (6,1.5) (7,2.9) (8,2.2)};
    \draw[thick,thmcol] plot[smooth] coordinates {(0,2) (1,1.2) (2,0.4) (3,1.4) (4,2.8) (5,3.7) (6,2.5) (7,1.1) (8,1.8)};
    \node[defcol] at (8.8,2.4) {mine}; \node[thmcol] at (8.8,1.6) {yours};
    \node[font=\small\bfseries] at (4,5.5) {money: fixed total};
  \end{scope}
  \begin{scope}[xshift=7.4cm,yshift=1.4cm]
    \draw[-{Stealth[length=3pt]}] (0,0) -- (9,0) node[right]{time};
    \draw[-{Stealth[length=3pt]}] (0,-2.5) -- (0,2.7);
    \draw[dashed,intcol] (0,1.8) -- (8.5,1.8) node[right]{your limit};
    \draw[dashed,intcol] (0,-1.8) -- (8.5,-1.8) node[right]{my limit};
    \draw[thick,keycol] plot[smooth] coordinates {(0,0) (1,-0.8) (2,-1.5) (3,-0.6) (4,0) (5,1.2) (6,1.6) (7,0.5) (8,-0.4)};
    \node[left] at (0,0) {0};
    \node[font=\small\bfseries] at (4,3.0) {credit: net position};
  \end{scope}
\end{tikzpicture}
\caption{Payments with a fixed quantity of money (left: balances only move between pockets, and a trade is impossible when the buyer has
run out) and with credit (right: the net position starts at zero and expands and contracts within credit limits)
\ts{4}{4}{0:28}\ts{4}{4}{1:35}.}\label{fig:mv-de}
\end{nfigure}

This is where discipline and elasticity (\cref{sec:nh-discipline}) become concrete \ts{4}{4}{0:07}.
\begin{itemize}
  \item With a \textbf{fixed quantity of money}, my balance goes up when I sell to you and down when I buy, while the total does not
        change. If there is little money, people run out, and some mutually advantageous trades are impossible: that is
        \emph{discipline} \ts{4}{4}{0:28}\ts{4}{4}{0:49}\ts{4}{4}{1:13}.
  \item With \textbf{pure credit}, positions start at zero and expand and contract with the pattern of payments, within each party's credit
        limit; there is no money, only credit that comes and goes \ts{4}{4}{1:35}\ts{4}{4}{1:56}\ts{4}{4}{2:16}.
  \item With a \textbf{bank}, I can borrow up to my limit at the bank, and I am willing to hold large positive balances with the bank, much
        more than I would lend to you; the bank does this with many people at once: a multilateral credit system
        \ts{4}{4}{2:58}\ts{4}{4}{3:18}.
\end{itemize}

\begin{proposition}[The payment system is a credit system]\label{prop:mv-credit}
The payment system is a credit system, not a money system, and it has to be one in order to have enough elasticity: people who want to
make mutually advantageous trades must be able to make them even if none of them has money (gold, the thing fixed in quantity) at the
moment \ts{4}{4}{3:39}. When looking at the actual system, expect to find credit, even where it does not look like credit
\ts{4}{4}{4:00}.
\end{proposition}

% ------------------------------------------------------------------
\section{The survival constraint}\label{sec:mv-survival}

Mehrling reads from Minsky's \emph{Stabilizing an Unstable Economy} (1986), p.~198 \ts{4}{5}{0:07}: to analyse how financial commitments
affect the economy one must look at economic units in terms of their cash flows; the cash flow approach looks at all units, households,
corporations, state and municipal governments, even national governments, \emph{as if they were banks} \ts{4}{5}{0:28}. That is: every unit
is a money-flow operation, cash in and cash out, and the goods side is secondary \ts{4}{5}{0:49}. The macro side of this view comes from
Copeland (\cref{sec:mv-fof}); the micro side is the survival constraint \ts{4}{5}{1:09}\ts{4}{5}{1:30}.

\begin{definition}[New concepts: survival (reserve) constraint]\label{def:mv-survival}
The \emph{survival constraint} (Minsky), or \emph{reserve constraint}, of an agent is the requirement that in every period
\[
  \text{cash inflow} \;\ge\; \text{cash outflow (cash commitments)} ,
\]
where cash commitments include promises already made, such as debt repayments. If it fails, the agent defaults: ``it doesn't matter if
you have lots of wealth, if you have no cash'' \ts{4}{5}{2:10}\ts{4}{5}{2:31}\ts{4}{5}{2:51}.
\end{definition}

The binding constraint is not the budget constraint of economics, and certainly not an \emph{intertemporal} budget constraint: it is a
liquidity constraint \ts{4}{5}{1:50}\ts{4}{5}{2:10}. ``Liquidity kills you quick'': a bank can go on doing business for a long time while
insolvent, many do, but cannot survive one day illiquid. If everyone is a bank, this is what focuses the mind \ts{4}{5}{2:51}\ts{4}{5}{3:12}.

\begin{intuition}[Budget constraint versus survival constraint]
The intertemporal budget constraint, $\sum_t c_t/(1+r)^t \le W$ (\cref{eq:fp-budget}), restricts present values: any timing of payments
is allowed as long as the totals balance, because one can always borrow against future income at the rate $r$. The survival constraint is
a constraint at \emph{each} date: a pointwise condition rather than an integral one. Two agents with the same present values can
differ completely in whether they survive.
\end{intuition}

\paragraph{Three ways to meet a shortfall.} If cash inflow does not cover a commitment, there are three choices \ts{4}{9}{0:07}\ts{4}{9}{0:29}:
\begin{enumerate}
  \item \textbf{dishoard}: run down cash reserves;
  \item \textbf{borrow}: find someone to lend more cash;
  \item \textbf{sell assets}, at whatever price one can get: a fire sale.
\end{enumerate}
(One could also sell goods; the course focuses on financial assets \ts{4}{9}{0:49}.) Only the first is dependable in a crisis: the others
need someone to take the other side of the trade, at a good price, and in a crisis they may not want to \ts{4}{9}{1:09}\ts{4}{9}{1:31}. This is
why reserves matter: at the end of the day it is cash that meets commitments. But nobody wants to hold only money; credit lets us escape the
constraint by pushing it into the future: we live to fight another day, but that day the borrowing comes due and we face the constraint
again \ts{4}{9}{1:52}\ts{4}{9}{2:12}. \emph{Credit gives us elasticity today but discipline tomorrow} \ts{4}{9}{2:32}.

% ------------------------------------------------------------------
\section{Sources and uses of funds}\label{sec:mv-su}

The \emph{flow of funds} accounts are real accounts, collected by the Fed and watched closely by financial advisers; start from their
micro foundation \ts{4}{6}{0:06}\ts{4}{6}{0:27}. For one agent, distinguish \emph{uses} from \emph{sources} of funds, and the goods-and-services
account (``above the line'') from the financial account (``below the line''), which has three categories: financial assets, financial
liabilities and money, a hierarchy of credit and money in which money is the ultimate means of payment \ts{4}{6}{0:48}\ts{4}{6}{1:08}\ts{4}{6}{1:29}.

\begin{nfigure}
\renewcommand{\arraystretch}{1.35}
\begin{tabular}{l|c|c}
 & \textbf{Uses} (cash out) & \textbf{Sources} (cash in)\\\hline
Goods and services (incl.\ labour) & expenditure & receipts\\\hline
Financial assets & accumulation & decumulation\\
Financial liabilities & repayment & borrowing\\
Money & hoarding & dishoarding\\
\end{tabular}
\caption{The sources-and-uses framework \ts{4}{6}{1:53}\ts{4}{6}{4:01}. Every entry is a \emph{flow} over a period, not a stock: this is not a
T-account. Above the line, goods; below the line, the financial account.}\label{fig:mv-su}
\end{nfigure}

\begin{itemize}
  \item Selling anything, including labour, is a receipt, a source of funds; buying is an expenditure, a use. Unlike the national
        accounts, which count only final sales (value added), here every sale counts: it is a source of cash for the seller even if it
        adds no value for the economy \ts{4}{6}{2:14}\ts{4}{6}{2:38}.
  \item Selling a financial asset is a source (decumulation), buying one a use (accumulation) \ts{4}{6}{2:58}.
  \item Borrowing is a source, repaying a use \ts{4}{6}{3:19}.
  \item Taking cash out of one's pile is dishoarding (a source), adding to it hoarding (a use) \ts{4}{6}{3:41}\ts{4}{6}{4:01}.
\end{itemize}
The framework is complete: every transaction fits in one of the categories, which lets one sort out complicated transactions
\ts{4}{6}{4:22}.

\begin{proposition}[The two rules of the flow of funds]\label{prop:mv-rules}
\leavevmode
\begin{enumerate}[(i)]
  \item For each agent, every use has a corresponding source and vice versa: nothing can be bought, no asset accumulated, no debt
        repaid, no cash hoarded, without the money coming from somewhere \ts{4}{6}{4:42}\ts{4}{6}{5:03}\ts{4}{6}{5:24}.
  \item Every agent's use is some other agent's source and vice versa: my receipt is somebody's expenditure, my purchase of an asset is
        somebody's sale \ts{4}{6}{5:45}\ts{4}{6}{6:08}.
\end{enumerate}
Consequently each transaction requires at least four entries, two for each of two agents \ts{4}{7}{1:56}. Rule (ii) is what turns the
framework into macroeconomics: the economy is a web in which everything is connected \ts{4}{6}{6:28}.
\end{proposition}

\begin{intuition}[Double entry, twice]
Rule (i) is the conservation law of each agent's cash (inflow $=$ outflow, with hoarding as the ``storage'' term); rule (ii) is the
conservation law between agents (every outflow is someone's inflow). Together they make the matrix of all agents' flows have zero row sums
(each type of flow balances across agents) and equal column sums for each agent (uses $=$ sources), as Kirchhoff's laws do for currents at
nodes and around a circuit.
\end{intuition}

\subsection{Two examples}\label{sec:mv-examples}
\paragraph{A coffee paid in cash.} Before class Mehrling bought a \$2 coffee at Oren's Daily Roast, in cash \ts{4}{7}{0:08}\ts{4}{7}{0:28}:
\begin{center}
\begin{tabular}{l|cc|cc}
 & \multicolumn{2}{c|}{\textbf{Mehrling}} & \multicolumn{2}{c}{\textbf{Oren's}}\\
 & uses & sources & uses & sources\\\hline
goods & coffee 2 & & & coffee 2\\\hline
money & & dishoard 2 & hoard 2 & \\
\end{tabular}
\end{center}
Four entries, both rules satisfied \ts{4}{7}{0:48}\ts{4}{7}{1:09}\ts{4}{7}{1:35}. The point is to build intuition for the inside character of money,
for the interlocking of the monetary system, for its money flows; these intuitions are not natural for economists, trained to see money as a
veil \ts{4}{7}{1:56}\ts{4}{7}{2:17}.

\paragraph{A dinner paid by credit card.} After class, dinner at Virelli, \$20, charged to MasterCard \ts{4}{7}{2:38}\ts{4}{7}{3:00}. Three parties,
MasterCard in the middle as the bank of \cref{sec:mv-payments} \ts{4}{7}{3:25}. In stages \ts{4}{7}{3:46}--\ts{4}{7}{9:04}:
\begin{enumerate}
  \item \textbf{Dinner.} Mehrling's expenditure, Virelli's receipt. Signing the slip is not paying but \emph{borrowing}: an IOU to MasterCard
        (Mehrling borrows, MasterCard accumulates). Virelli lets him go because it gets MasterCard's IOU (MasterCard borrows, Virelli
        accumulates) \ts{4}{7}{4:33}\ts{4}{7}{4:53}\ts{4}{7}{5:14}\ts{4}{7}{5:36}. Two new credits, out of thin air.
  \item \textbf{Merchant settlement.} The next day or week Virelli presents its claim: Virelli decumulates and hoards cash, MasterCard dishoards
        and repays \ts{4}{7}{7:00}\ts{4}{7}{7:21}\ts{4}{7}{8:03}. For Virelli it is as if it had been paid in cash.
  \item \textbf{Cardholder settlement.} At the end of the month Mehrling repays (dishoarding cash, repaying debt), and MasterCard decumulates the
        IOU and hoards cash \ts{4}{7}{8:24}\ts{4}{7}{8:44}\ts{4}{7}{9:04}.
\end{enumerate}
\begin{center}
\begin{tabular}{l|cc|cc|cc}
 & \multicolumn{2}{c|}{\textbf{Mehrling}} & \multicolumn{2}{c|}{\textbf{MasterCard}} & \multicolumn{2}{c}{\textbf{Virelli}}\\
 & uses & sources & uses & sources & uses & sources\\\hline
goods & 1: dinner & & & & & 1: dinner\\\hline
assets & & & 1: IOU$_\text{M}$ & 3: IOU$_\text{M}$ & 1: IOU$_\text{MC}$ & 2: IOU$_\text{MC}$\\
liabilities & 3: repay & 1: IOU$_\text{M}$ & 2: repay & 1: IOU$_\text{MC}$ & & \\
money & & 3: dishoard & 3: hoard & 2: dishoard & 2: hoard & \\
\end{tabular}
\end{center}
Credit expanded, two credits at once, and then collapsed in two phases; at the end there is no credit left, and cash settled the debts: the
elasticity came from credit, the discipline from the debts having to be paid \ts{4}{7}{9:26}\ts{4}{7}{9:47}.

% ------------------------------------------------------------------
\section{The flow of funds accounts}\label{sec:mv-fof}

Morris Copeland, \emph{A Study of Moneyflows in the United States} (NBER, 1952), proposed these accounts as a framework for
macroeconomic analysis \ts{4}{8}{0:07}. They apply the sources-and-uses framework to whole sectors, with aggregation and netting, which is
what makes them hard to read for people who do not have the framework in mind \ts{4}{8}{0:27}.
\begin{itemize}
  \item \textbf{Columns}: sectors (households and nonprofit organizations, nonfinancial business, state and local governments, federal
        government, domestic financial sectors, rest of the world), each with a ``uses'' and a ``sources'' column \ts{4}{8}{0:48}\ts{4}{8}{2:31}.
  \item \textbf{Rows}: first the goods and services part (about 12 lines), then the financial instruments: various kinds of money first
        (official reserves, SDRs, Treasury currency, deposits; lines 14--22), then forms of credit (from line 23), in the reverse order
        of the hierarchy \ts{4}{8}{1:09}\ts{4}{8}{1:29}.
  \item \textbf{Checks}: by rule (i) each sector's uses and sources columns must have the same sum; by rule (ii) each row must sum to zero.
        They do not quite: the \emph{statistical discrepancies} show where the net has missed something \ts{4}{8}{1:50}\ts{4}{8}{2:10}.
  \item \textbf{Negative entries}: the published accounts record repayment as a negative source (a reduction of borrowing) and
        decumulation as a negative use; this moves entries to the other side but keeps the rules \ts{4}{8}{3:12}\ts{4}{8}{3:33}\ts{4}{8}{3:53}.
\end{itemize}

\paragraph{Copeland against Keynes and the monetarists} \ts{4}{8}{4:54}:
\begin{itemize}
  \item $C+I+G=Y$ ignores purchases and sales of financial assets and of existing goods (old houses, a very big market), and much of what
        happens in the monetary and financial system \ts{4}{8}{5:15}\ts{4}{8}{5:35};
  \item $MV=PT$ treats all transactions as made with money, ``as if the whole economy were Oren's Daily Roast'', while most transactions
        are credit transactions, absorbed into velocity \ts{4}{8}{5:56}.
\end{itemize}
Mehrling thinks Copeland was right, but he did not win: the accounts were compiled by a small division of the Board of Governors, but not
used for policy, because no macroeconomic theory was attached to them; Keynesians and monetarists each had theirs. ``So there's a Nobel
prize for you'' \ts{4}{8}{6:17}\ts{4}{8}{6:37}\ts{4}{8}{6:58}.

\paragraph{What the accounts miss.} In the table shown, the largest discrepancy in the last column is \$115 billion on line 22, ``Fed funds
and security RPs'': the wholesale money market escapes the statistical net \ts{4}{8}{7:18}\ts{4}{8}{7:38}. Sector discrepancies are large too
(\$561 billion for households and nonprofits, \$440 billion for domestic nonfinancial sectors), partly because ``nonfinancial'' firms are
also financial: GM makes cars, but it is also a kind of bank with international operations \ts{4}{8}{7:59}\ts{4}{8}{8:19}\ts{4}{8}{8:40}. And there
are no derivatives at all, no swaps, no options: they did not exist in 1952, and they are off balance sheet \ts{4}{8}{9:22}\ts{4}{8}{9:42}. After the
crisis the Office of Financial Research was created to build new accounting structures; Mehrling proposed to ``soup up'' the accounts to
include swaps and derivatives, as the course will do \ts{4}{8}{8:59}\ts{4}{8}{9:42}.\begin{coursenote}[Flow of funds today]
The flow of funds accounts have been published by the Fed since 1955, now as the \emph{Financial Accounts of the United States} (release
Z.1). The Office of Financial Research, created by the Dodd--Frank Act in 2010, is an office of the US Treasury Department, not of the Fed
(in class it is called ``a division of the Fed'' \ts{4}{8}{8:59}).
\end{coursenote}

% ------------------------------------------------------------------
\section{Liquidity, long and short}\label{sec:mv-longshort}

The survival constraint applies not only today but tomorrow and after: pushing a problem into the future lets you fight another day, but
you should push it where you expect cash inflows, not where you have other commitments \ts{4}{10}{0:28}\ts{4}{10}{0:48}. Example, over four
periods $t,\dots,t+3$, with cash inflow 10 every period (40 in all) and total commitments 20: no solvency problem, a ``profit'' of 20
\ts{4}{10}{1:31}\ts{4}{10}{1:53}\ts{4}{10}{2:14}.

\begin{nfigure}
\renewcommand{\arraystretch}{1.35}
\begin{tabular}{l|cccc|l}
 & $t$ & $t+1$ & $t+2$ & $t+3$ & \\\hline
cash inflow (all cases) & 10 & 10 & 10 & 10 & \\\hline
commitments, hedge & 5 & 5 & 5 & 5 & never binds\\
commitments, short-term & 20 & 0 & 0 & 0 & binds today\\
\end{tabular}
\caption{Same stocks, different flows \ts{4}{10}{2:35}\ts{4}{10}{3:48}: in both cases inflows total 40 and commitments 20 (ignoring
discounting), so the balance sheet is the same; only the time pattern differs.}\label{fig:mv-patterns}
\end{nfigure}

\begin{enumerate}
  \item \textbf{Hedge (liquid) financing}: commitments of 5 every period; the constraint never binds, the question ``how am I going to
        meet it?'' never arises \ts{4}{10}{2:35}\ts{4}{10}{2:56}.
  \item \textbf{Short-term financing}: the business is fine (10 every period) but financed with short-term debt, all due today: 20 owed,
        10 coming in. One of the three choices is forced now: use reserves, roll over the debt, or sell the business
        \ts{4}{10}{3:48}\ts{4}{10}{4:09}. Short-term borrowing is the sign of vulnerability, and real banks, with demand deposits and overnight
        repo funding, look like this all the time; hence the importance of the overnight market, enormous and invisible to households
        \ts{4}{10}{4:30}\ts{4}{10}{4:50}\ts{4}{10}{5:11}.
  \item \textbf{A problem coming later} (a third case, drawn on the board with ``a little leeway''): no problem today, but a future misalignment of flows and commitments that the agent, and its creditors,
        can see. While still liquid, the agent refinances, spreading the debt over time or prepaying part of it: \emph{anticipated} liquidity
        constraints have consequences for liquidity today, because tomorrow the markets may be frozen and a perfectly solvent agent could be
        ``dead'' \ts{4}{10}{6:00}\ts{4}{10}{6:21}\ts{4}{10}{6:41}\ts{4}{10}{7:02}.
\end{enumerate}

\paragraph{Price-insensitive borrowers.} The economy is a crowd of agents with different time patterns of flows and commitments, all trading
in the same market \ts{4}{10}{7:22}\ts{4}{10}{7:43}. An agent under stress must borrow now, whether the short rate is 10\% or 20\%: its demand
is price-inelastic, and the more such agents there are, the further rates can move, for reasons that have nothing to do with time preference
\ts{4}{10}{8:03}\ts{4}{10}{8:23}. Think of Bear Stearns or Lehman before their collapse: this is the stuff of financial crises, and all banks, with
short-term liabilities that can vanish tomorrow, look like this \ts{4}{10}{8:44}. Agents with surplus cash lend it into the market. Over the
business cycle the balance swings between easy and tight credit, and it shows up in the money market \ts{4}{10}{9:04}\ts{4}{10}{9:25}.

\begin{definition}[New concepts: hedge financing, rollover]
\emph{Hedge financing} (Minsky): commitments in every period are covered by expected cash inflows of the same period. To \emph{roll over}
a debt is to repay it with the proceeds of a new debt, pushing the commitment into the future. (Minsky also defines \emph{speculative}
units, whose near-term commitments exceed near-term inflows and must be rolled over, and \emph{Ponzi} units, which must borrow even to pay
interest.)
\end{definition}

% ------------------------------------------------------------------
\section{Financial fragility: flows and stocks}\label{sec:mv-fragility}

\begin{nfigure}
\begin{tikzpicture}[font=\small]
  \draw[thick] (0,0) -- (0,2.6);
  \fill (0,2.6) circle (2pt);
  \draw[thick] (-2.6,2.3) -- (2.6,2.9);
  \draw (-2.6,2.3) -- (-3.3,1.2) (-2.6,2.3) -- (-1.9,1.2);
  \draw[thick,intcol,fill=intcol!15] (-3.4,1.2) arc (180:360:0.75);
  \node[intcol,align=center,font=\scriptsize] at (-2.65,0.1) {cash\\flows};
  \draw (2.6,2.9) -- (1.9,1.8) (2.6,2.9) -- (3.3,1.8);
  \draw[thick,thmcol,fill=thmcol!15] (1.8,1.8) arc (180:360:0.75);
  \node[thmcol,align=center,font=\scriptsize] at (2.55,0.7) {cash\\commitments};
  \draw[-{Stealth[length=4pt]},very thick,keycol] (0,2.6) -- (0.9,3.6);
  \draw[keycol] (-1.2,3.5) arc (135:45:1.7);
  \node[keycol,font=\scriptsize,align=center] at (0,4.3) {money market rate of interest};
  \draw[thick] (-0.8,0) -- (0.8,0);
\end{tikzpicture}
\caption{The scales \ts{4}{11}{0:07}\ts{4}{11}{0:30}: the balance between the pattern of cash flows and the pattern of cash commitments, across the
whole economy, sets the money market interest rate.}\label{fig:mv-scales}
\end{nfigure}

The image, ``because we're in New York'', is the scales of Lady Liberty: one pan holds cash flows, the other cash commitments, and the
pointer is the money market interest rate \ts{4}{11}{0:07}\ts{4}{11}{0:30}\ts{4}{11}{0:50}. The Fed does something too, but behind the scenes there are
balance sheet facts: promises to pay at future dates that either are or are not lined up with the prospects of whoever made them; if they are
not, that person is forced to act. ``It's not about preference, it's about being forced to do something by the promises you've made in the
past'' \ts{4}{11}{1:10}\ts{4}{11}{1:31}. You buy a house, then lose your job: the cash flows go away, the commitments stay \ts{4}{11}{1:52}.

\begin{proposition}[Minsky's fragility, in the money view]\label{prop:mv-fragility}
Financial fragility is the balance of cash commitments (promises to pay) relative to the cash flows actually arising from activity
(wages, profits, sales). When most agents' cash flows cover their commitments there is no crisis; when many agents' do not, the system
becomes fragile, and the crisis shows up in the money market rate of interest \ts{4}{11}{2:12}\ts{4}{11}{2:33}\ts{4}{11}{2:53}.
\end{proposition}

The money market is thus the ``crisis detector'', the ``stress detector'': the bouncing of its rate reflects the shifting balance between cash
flows and commitments, a ``sufficient statistic'' for the state of the macroeconomic balance \ts{4}{11}{3:15}\ts{4}{11}{3:36}.

\paragraph{Stocks are residues of flows.} In a T-account, financial assets and money on the left, financial liabilities and net worth on the
right; net worth is solvency \ts{4}{11}{3:57}\ts{4}{11}{4:17}. Each stock is the accumulation of past flows: assets bought, cash hoarded, debts
incurred \ts{4}{11}{4:38}\ts{4}{11}{4:58}. And each looks forward too: liabilities and assets have dates, the dates of future payments; the stock
value adds them up, as a discounted present value (40 and 20 in the examples) \ts{4}{11}{5:18}. The examples of \cref{fig:mv-patterns} have
identical stocks and very different prospects: the balance sheet shows solvency; to see fragility one must look behind it, at the time
patterns \ts{4}{11}{5:38}\ts{4}{11}{5:59}.

% ------------------------------------------------------------------
\flushside
\section*{Key formulas and ideas}
\begin{keyformulas}
\begin{tabularx}{\linewidth}{@{}l L@{}}
Money vs credit payments & money: existing balances move, total fixed; credit: outstanding IOUs increase with the transaction\\
Payment system & a credit system (elasticity), settled in money (discipline)\\
Survival constraint & every period: cash inflow $\ge$ cash commitments (not an intertemporal budget constraint)\\
Shortfall & dishoard, borrow, or sell assets; only reserves are dependable in a crisis\\
Sources and uses & receipts/expenditure; decumulation/accumulation; borrowing/repayment; dishoarding/hoarding\\
Two rules & (i) each agent: uses $=$ sources; (ii) each use is another agent's source $\Rightarrow$ $\ge4$ entries per transaction\\
Fragility & commitments vs cash flows in time; shows up in the money market rate\\
Stocks & accumulated past flows and discounted future flows; same stocks can hide different time patterns\\
\end{tabularx}
\end{keyformulas}

\section*{Exercises}

\begin{exercise}[Sources and uses]
Record in the sources-and-uses framework, for all agents involved:
\begin{enumerate}[(a)]
  \item you receive a \$3000 salary paid into your bank account;
  \item you buy \$1000 of shares from another investor, paying from your account;
  \item you repay \$500 of a bank loan;
  \item a firm sells a used machine to another firm on 30-day trade credit.
\end{enumerate}
Check both rules in each case.
\end{exercise}

\begin{exercise}[The credit card, completed]
Write the three stages of the dinner example as changes in T-accounts (stocks) for Mehrling, MasterCard and Virelli, assuming each pays
through a bank. What is the total quantity of credit outstanding after each stage?
\end{exercise}

\begin{exercise}[Time patterns]
A firm expects cash inflows of 10, 10, 10, 10 and has debt of 24.
\begin{enumerate}[(a)]
  \item Give a repayment schedule that never violates the survival constraint and one that violates it at $t$.
  \item In the second case, how much must it roll over at $t$? At what rate does it become insolvent if it must pay interest $i$ on the
        rolled-over amount for three periods?
  \item Explain why its demand for funds at $t$ is price-inelastic.
\end{enumerate}
\end{exercise}

\begin{exercise}[Statistical discrepancy]
In a two-sector economy households lend 100 to firms through repo, and firms report borrowing 85. Explain how the discrepancy shows up in the
row and column checks of the flow of funds matrix, and where you would look for the missing flow.
\end{exercise}

% !TEX root = ../main.tex
\chapter{The Central Bank as a Clearinghouse}\label{ch:clearinghouse}
\begin{paracol}{2}

\begin{sources}
\begin{tabularx}{\linewidth}{@{}l l L@{}}
\toprule
\textbf{Segment} & \textbf{Length} & \textbf{Content}\\
\midrule
\seg{5}{1}{L5.1 Martin Wolf on QE3} & 3:16 & Wolf: QE3 more likely to do too little than too much; the nominal GDP gap.\\
\seg{5}{2}{L5.2 One Big Bank} & 8:20 & A payment system as one big bank: money version and credit version; the payment matrix.\\
\seg{5}{3}{L5.3 Multiple Banks, a challenge} & 3:56 & Many banks: paying in reserves is more constraining.\\
\seg{5}{4}{L5.4 Reading: Charles F.~Dunbar} & 2:05 & Dunbar's \emph{Chapters on the Theory and History of Banking} (1891).\\
\seg{5}{5}{L5.5 Correspondent banking} & 10:24 & Netting, correspondent balances as a swap of IOUs; who expands.\\
\seg{5}{6}{L5.6 Correspondent banking, system network} & 3:55 & The hierarchical network of correspondents up to New York.\\
\seg{5}{7}{L5.7 Clearinghouse, normal operations} & 8:28 & Clearinghouse certificates, multilateral netting, settlement.\\
\seg{5}{8}{L5.8 Clearinghouse, private lender of last resort} & 10:44 & Clearinghouse loan certificates in crises; why the Fed replaced them.\\
\seg{5}{9}{L5.9 Central Bank Clearing} & 4:54 & The Fed as a public clearinghouse: Fedwire, Fed funds, discount loans.\\
\seg{5}{10}{L5.10 Central Bank Cooperation} & 5:41 & Internal and external drain; the major central banks acting as a clearinghouse.\\
\bottomrule
\end{tabularx}

\smallskip
\textbf{Files.} Transcripts \file{materials/transcripts/L05-P*.txt}; Mehrling's notes \mn{5}{1}.
\textbf{Reading.} Charles F.~Dunbar, \emph{Chapters on the Theory and History of Banking} (1891; 2nd ed.\ 1901), the chapter on the
checking system (not among the course files).
\end{sources}

\section*{Motivation}
This lecture opens the first block of the course, \emph{banking as a system of payments}, and \emph{central banking as a clearinghouse}
\ts{5}{2}{0:08}. Starting from payments is unusual, but it builds the intuition for what follows: why the Fed funds market is so central,
what goes on in repo and in the Eurodollar market \ts{5}{2}{0:28}. The guiding idea \ts{5}{2}{0:48}:
\begin{center}
\emph{all the institutions and practices of the monetary system are attempts, worked out over a long time, to make a system of many banks
work as if it were one big bank.}
\end{center}
The lecture goes up the hierarchy: one big bank (\cref{sec:ch-onebank}); many banks (\cref{sec:ch-many}) linked by correspondent balances
(\cref{sec:ch-corr}); the clearinghouse at the top, in normal times and in crises (\cref{sec:ch-chouse,sec:ch-lolr}); the Fed as a public
clearinghouse (\cref{sec:ch-fed}); and the central banks of the world as a clearinghouse one level higher (\cref{sec:ch-coop}).

% ------------------------------------------------------------------
\section{FT: Martin Wolf on QE3}\label{sec:ch-ft}

\begin{casestudy}[FT, September 2012: ``Bernanke makes a historic choice'']
Martin Wolf, the FT's chief economics commentator, writes every Wednesday \ts{5}{1}{0:09}. He argues that critics of the new Fed policy
are right that it will probably fail to work as hoped, but wrong that it will do vast damage: the fear of hyperinflation is misguided, and
the costs of deficient demand are far greater; the Fed ``is in truth more likely to achieve too little of what it seeks than too much''
\ts{5}{1}{0:50}. Mehrling agrees: it is not clear how QE translates into aggregate demand \ts{5}{1}{1:10}. Wolf's chart of nominal GDP
against its pre-crisis trend, the goal that Bernanke and Michael Woodford discussed at Jackson Hole, shows a gap that is widening, not
closing; perhaps because fiscal policy is stuck, too much is being asked of the central bank, in the US and in Europe
\ts{5}{1}{1:50}.
\end{casestudy}\begin{aside}[The Money View blog]
Mehrling and his co-blogger Daniel Neilson posted on QE3 on their blog, together with a short video on Draghi and the Soros plan
(\cref{sec:nh-ft}) \ts{5}{1}{1:30}.
\end{aside}

% ------------------------------------------------------------------
\section{One big bank}\label{sec:ch-onebank}

``One big bank'' can mean two things, following the money/credit distinction of lecture~4 \ts{5}{2}{1:29}.

\paragraph{The money version.} The bank holds reserves (say gold) and owes deposits to people, $\alpha$ and $\beta$. A payment from
$\alpha$ to $\beta$ is a book entry from one account to another; no reserves move at all \ts{5}{2}{2:11}. The limit: when a
deposit is run down to zero, no more purchases are possible, unless one can borrow, bilaterally, from someone who has a deposit
\ts{5}{2}{2:54}. The IMF, with its Special Drawing Rights, is built on this discipline model: a fixed balance sheet, reallocated,
expanded only now and then ``in one big tranche'' \ts{5}{2}{3:14}.

\paragraph{The credit version.} The bank has deposits and also \emph{overdrafts} of people $\gamma$ and $\lambda$, on the asset side: it
uses the fact that no reserves are really needed \ts{5}{2}{3:35}. Now the balance sheet may expand or contract, depending on the
pattern of payments \ts{5}{2}{4:16}.
\begin{taccts}
\tacct[2.4cm]{One big bank (money)}{Reserves}{Deposit $\alpha$\\Deposit $\beta$}\qquad
\tacct[2.4cm]{One big bank (credit)}{Overdraft $\gamma$\\Overdraft $\lambda$}{Deposit $\alpha$\\Deposit $\beta$}
\end{taccts}

\begin{nfigure}
\renewcommand{\arraystretch}{1.35}
\begin{tabular}{c|cc|cc}
from $\backslash$ to & $\alpha$ & $\beta$ & $\gamma$ & $\lambda$\\\hline
$\alpha$ & 0 & 0 & $-$ & $-$\\
$\beta$ & 0 & 0 & $-$ & $-$\\\hline
$\gamma$ & $+$ & $+$ & 0 & 0\\
$\lambda$ & $+$ & $+$ & 0 & 0\\
\end{tabular}
\caption{The payment matrix of the credit one-big-bank \ts{5}{2}{4:37}: change in the size of the balance sheet when the row agent pays the
column agent. $\alpha,\beta$ have deposits, $\gamma,\lambda$ overdrafts.}\label{fig:ch-matrix}
\end{nfigure}

Reading the matrix \ts{5}{2}{4:58}:
\begin{itemize}
  \item paying oneself, or a depositor paying a depositor: no change, exactly as in the money version;
  \item an overdrawn agent paying another overdrawn agent: one overdraft grows, the other shrinks, total unchanged;
  \item a depositor ($\alpha$) paying an overdrawn agent ($\gamma$): $\alpha$'s deposit falls and $\gamma$'s overdraft is partly repaid:
        the balance sheet \emph{contracts} \ts{5}{2}{6:21};
  \item an overdrawn agent paying a depositor: the overdraft grows and the deposit grows: the balance sheet \emph{expands} \ts{5}{2}{6:41}.
\end{itemize}
The aim of a payment system is that whenever there is a willing buyer and a willing seller, the payment system does not stand in the way.
The credit system can do everything the money system does and more; the actual US domestic payment system has this credit character, and
the Federal Reserve System is an example of it \ts{5}{2}{7:22}.

\begin{secondary}[What is the quantity of money?]
In the money version it is natural to call the deposits the quantity of money. With overdrafts the measure becomes ambiguous; the notes list three
candidates \mn{5}{1}:
\begin{enumerate}
  \item deposits, which counts only the positive balances;
  \item deposits minus overdrafts;
  \item deposits minus credit limits (the unused borrowing power counts as means of payment too).
\end{enumerate}
This ambiguity led Fischer Black to propose not trying to measure money at all, in ``Banking in a World without Money'' (lecture~\ref{ch:threeviews})
\mn{5}{2}.
\end{secondary}

\begin{definition}[New concepts: overdraft]
An \emph{overdraft} is a negative balance on a deposit account, i.e.\ a loan from the bank created automatically when the holder pays more
than he has (one ``draws over'' the balance). It is an asset of the bank.
\end{definition}

% ------------------------------------------------------------------
\section{Many banks}\label{sec:ch-many}

There is no one big bank. In the United States there are thousands of banks, fewer than the six or seven thousand of Young's time, some very
big (Citibank, Bank of America), many small or regional; in some countries, such as the UK or Canada, there are a handful of big banks
\ts{5}{3}{0:07}. The problem is how to knit them into one country-wide system in which anybody can pay anybody
\ts{5}{3}{0:48}.

Take bank A with depositors $\alpha,\beta$ and bank B with depositors $\gamma,\lambda$, all with positive deposits; the credit element now
is \emph{interbank} credit. This is the wholesale money market, behind the scenes, and its elasticity is what provides elasticity in the retail
market \ts{5}{3}{1:09}. One could pay in reserves: when $\alpha$ pays $\gamma$ 10, bank A sends 10 of reserves to
bank B \ts{5}{3}{2:31}:
\begin{taccts}
\tacct[2.2cm]{Bank A}{$-$Reserves 10}{$-$Deposit $\alpha$ 10}\qquad
\tacct[2.2cm]{Bank B}{$+$Reserves 10}{$+$Deposit $\gamma$ 10}
\end{taccts}
This is \emph{more} constraining than one big bank: not only must $\alpha$ have a deposit, bank A must have reserves. So banks do not
generally do this within a country \ts{5}{3}{3:12}.

% ------------------------------------------------------------------
\section{Reading: Charles F.~Dunbar}\label{sec:ch-dunbar}

\begin{reading}[Dunbar, \emph{Chapters on the Theory and History of Banking} (1891)]
This was the textbook used at Harvard College; the course reads its chapter on how the checking system works in the United States
\ts{5}{4}{0:08}. It explains a payment system without a central bank (there was no Fed in 1891), which helps to understand why one was
eventually created; and it is ``crystal clear'': no four-colour diagrams, no workbook, no slides \ts{5}{4}{0:29}. Dunbar was
the first professor of money and banking at Harvard, a former journalist who had written on banking for a Boston paper; the second edition
was enlarged by O.~M.~W.~Sprague, of the Harvard Business School, a good friend of Allyn Young \ts{5}{4}{1:09}.
\end{reading}\begin{history}[Dunbar and Sprague]
Charles Franklin Dunbar (1830--1900) edited the \emph{Boston Daily Advertiser} before becoming, in 1871, the first professor of political economy
at Harvard; he founded the \emph{Quarterly Journal of Economics} in 1886. The second edition of his \emph{Chapters} (1901) was edited by Oliver
M.~W.~Sprague (1873--1953), later the author of \emph{History of Crises under the National Banking System} (1910) for the National Monetary
Commission. In class the paper is called ``the Boston Globe'' \ts{5}{4}{1:09}.\par
\emph{References:} F.~W.~Taussig, ``Charles Franklin Dunbar'', \emph{QJE} 14 (1900); Dunbar, \emph{Chapters on the Theory and History of
Banking}, 2nd ed., Putnam, 1901.
\end{history}

% ------------------------------------------------------------------
\section{Correspondent banking}\label{sec:ch-corr}

\paragraph{Collecting and netting.} During the day bank A receives checks, orders to pay. It does not pay them one by one: it collects them
and sorts them by \emph{bank} (the routing number at the bottom of a check), not by customer; which account at bank B gets credited is bank B's
problem \ts{5}{5}{0:27}. Checks between customers of bank A itself need no reserve flows; for each other bank it
sums what it owes (``due to B'') and what it is owed (``due from B''), and at the end of the day computes the net, which may be positive or
negative \ts{5}{5}{1:29}. Netting already saves moving reserves back and forth across the street every
day \ts{5}{5}{3:43}.

\paragraph{Correspondent balances: a swap of IOUs.} Banks that do business regularly can do better: bank A holds a deposit at B in its own
name (a \emph{banker's balance}), and B holds one at A. Instead of paying for these deposits with reserves, they swap them: each says to the
other ``I owe you \$1000''. Both balance sheets expand, deposits ``created from thin air'' \ts{5}{5}{4:04}.
\begin{taccts}
\tacct[2.3cm]{Bank A}{$+$Deposit at B}{$+$Deposit of B}\qquad
\tacct[2.3cm]{Bank B}{$+$Deposit at A}{$+$Deposit of A}
\end{taccts}
It looks like a free lunch that cannot do anything, but it does: now net payments can be made without moving reserves
\ts{5}{5}{5:30}. For $\alpha$ (at A) to pay $\lambda$ (at B) 10, there are two ways \ts{5}{5}{5:51}--\ts{5}{5}{8:16}:
\begin{enumerate}
  \item A draws down its deposit at B: A's balance sheet \emph{shrinks}, B's only changes owner (B swaps an IOU to A for an IOU to $\lambda$);
  \item B increases its deposit at A (A's liabilities to B grow): A's balance sheet does not change in size, B's \emph{expands}.
\end{enumerate}
\begin{taccts}
\textit{(1)}\quad\tacct[2.2cm]{Bank A}{$-$Deposit at B 10}{$-$Deposit $\alpha$ 10}\quad
\tacct[2.2cm]{Bank B}{}{$-$Deposit of A 10\\$+$Deposit $\lambda$ 10}

\medskip
\textit{(2)}\quad\tacct[2.2cm]{Bank A}{}{$-$Deposit $\alpha$ 10\\$+$Deposit of B 10}\quad
\tacct[2.2cm]{Bank B}{$+$Deposit at A 10}{$+$Deposit $\lambda$ 10}
\end{taccts}
Either way it is credit: either the debtor bank's balance sheet shrinks or the creditor bank's expands \ts{5}{5}{8:16}. In the first case total
deposits, counting interbank balances, fall; in the second they rise \mn{5}{3}. The invention of correspondent banking moves payments \emph{between
banks} from a money system to a credit system: they are no longer constrained by the quantity of gold, only by bilateral credit limits \mn{5}{3}.

\begin{definition}[New concepts: correspondent balances]
A \emph{correspondent balance} (banker's balance) is a deposit that one bank holds at another to settle payments between them. The bank that
holds the account is the \emph{respondent}, the bank that keeps it the \emph{correspondent} (the two ``correspond'', i.e.\ exchange letters and
orders, in the old sense of the word).
\end{definition}

\paragraph{Which one gets picked? Hierarchy.} In practice the big city bank says to the small country bank: I have no reason to hold deposits
with you; you hold a deposit with me, paid in reserves up front, or I lend it to you for a fee, and my deposits are the clearing balances
\ts{5}{5}{8:37}. This is discipline for the debtor bank: to join the payment system a country bank must keep balances at an urban
bank; if it runs out there is no guarantee the city bank will expand to help it, and it may charge extra. ``This is the hierarchy of money
showing up in the payment system'' \ts{5}{5}{9:17}. Correspondent balances make the system work much more like the ideal one big
bank on a credit basis; the retail customer sees none of this \ts{5}{5}{9:58}. The notes: the more central bank always takes deposits from the
less central; in the US country banks held deposits at city banks, and those reserve cities survive as the cities of the twelve Federal Reserve Banks
\mn{5}{3}.

\paragraph{The network.} Before the central bank, money-centre banks (Cleveland, Chicago, \dots) had correspondent relations with banks in their
hinterland; payments within a region just moved deposits at the centre bank \ts{5}{6}{0:07}. Payments between regions need a payment
one layer up, between centre banks; in practice this was not done bilaterally but hierarchically, through a few New York banks at which all the
money centres held accounts, with further regional layers \ts{5}{6}{1:10}. The whole structure nets as
much as it can, so that reserves do not have to flow, and builds credit relations locally and higher up, balancing elasticity and discipline from
the smallest bank ``in the middle of nowhere'' up to J.~P.~Morgan in New York \ts{5}{6}{2:53}. At the top, the biggest New York banks
must clear with each other: that is where the clearinghouse comes in \ts{5}{6}{3:34}.

\begin{nfigure}
\begin{tikzpicture}[font=\scriptsize,
  ny/.style={draw,thick,fill=keycol!20,rounded corners=2pt,minimum width=1.6cm,minimum height=0.6cm},
  mc/.style={draw,fill=defcol!15,rounded corners=2pt,minimum width=1.3cm,minimum height=0.5cm},
  cb/.style={draw,fill=fincol!15,circle,inner sep=1.6pt}]
  \node[ny] (ny) at (0,3) {New York bank};
  \foreach \x/\n [count=\i] in {-4.2/Cleveland,0/Chicago,4.2/St.\ Louis} {
    \node[mc] (m\i) at (\x,1.5) {\n};
    \draw[semithick] (m\i) -- (ny);
    \foreach \d in {-1.2,-0.6,0,0.6,1.2} {
      \draw[thin] (\x+\d,0) -- (m\i);
      \node[cb] at (\x+\d,0) {};}}
  \node[right,font=\scriptsize] at (5.6,0) {country banks};
  \node[right,font=\scriptsize] at (5.6,1.5) {money-centre banks};
  \node[right,font=\scriptsize] at (5.6,3) {New York};
\end{tikzpicture}
\caption{The correspondent network before the Fed \ts{5}{6}{0:07}: country banks hold balances at their money-centre bank, money-centre
banks at a New York bank; a payment between two country banks in different regions moves balances at every level up to the common correspondent.}
\label{fig:ch-network}
\end{nfigure}

% ------------------------------------------------------------------
\section{The clearinghouse in normal times}\label{sec:ch-chouse}

The New York Clearing House Association (NYCHA) is a mutual organisation of the big New York banks, ``first among equals'' J.~P.~Morgan
\ts{5}{7}{0:08}. It was set up on the discipline model: every member deposits gold or currency, and the clearinghouse issues
\emph{clearinghouse certificates} to its credit, physical notes standing for the reserves \ts{5}{7}{0:55}.
\begin{taccts}
\tacct[2.6cm]{Clearinghouse}{Gold (reserves)}{Clearinghouse certificates (to A, B, \dots)}
\end{taccts}

\paragraph{During the day: credit.} The members agree to treat every amount due to or from another member as due to or from the clearinghouse
itself \ts{5}{7}{1:58}. Each bank then has a single net position against the clearinghouse, smaller than the sum of its bilateral nets,
because positives and negatives offset: \emph{multilateral netting} \ts{5}{7}{2:40}. In Dunbar's description these are
actual checks: at the end of the day the banks add them up and meet at the clearinghouse, and those in net deficit hand clearinghouse certificates to
those in net surplus; then they go home \ts{5}{7}{3:42}.

\begin{intuition}[Bilateral versus multilateral netting]
With $n$ banks and gross obligations $x_{ij}$ (from $i$ to $j$), bilateral netting settles $|x_{ij}-x_{ji}|$ for each pair, while multilateral
netting settles only each bank's net position $\sum_j(x_{ij}-x_{ji})$. The sum of the multilateral nets is never larger than the sum of the
bilateral ones (triangle inequality), and is usually much smaller: a ring of payments $A\to B\to C\to A$ of equal size nets to zero.
\end{intuition}

``There is a credit expansion during the day that creates elasticity, and then there's a settlement at the end of the day'': elasticity and
discipline at the very highest level of the system \ts{5}{7}{4:23}. The members know they have this elasticity and can sell it to the banks
below them, and so on down: the whole system operates as a credit system because the clearinghouse does, during the day. Sometimes settlement is
every three days, which is why the length of the \emph{float} matters \ts{5}{7}{5:05}.

\paragraph{Failing to settle.} A bank that has run down its certificates over several days of deficits has three options
\ts{5}{7}{5:48}:
\begin{enumerate}
  \item pay in clearinghouse certificates, if it has them;
  \item borrow certificates bilaterally from another member, e.g.\ from J.~P.~Morgan, or from a member in net surplus;
  \item if nobody will lend, because it is thought insolvent, default: the clearinghouse's gold goes to pay the creditors pro rata, the losses are
        shared among all the members, and the defaulter is expelled, ``banished to the outer ring of hell''.
\end{enumerate}

\begin{proposition}[Joint credit]
The credit of a group of banks that jointly and severally guarantee an obligation is better than the credit of any individual bank: clearinghouse
credit is better than the bilateral credit of its members \ts{5}{7}{7:32}.
\end{proposition}

\begin{definition}[New concepts: clearing, settlement, netting, float]
\emph{Clearing} is the process of collecting and offsetting payment orders between banks; \emph{settlement} is the final transfer of the means of
payment (reserves) that discharges the net obligations. \emph{Netting} replaces gross obligations by net ones. The \emph{float} is the credit that
exists between the moment a payment is credited to the payee and the moment it is settled.
\end{definition}

% ------------------------------------------------------------------
\section{The clearinghouse as private lender of last resort}\label{sec:ch-lolr}

In times of stress all the members need to make payments, to their correspondents or abroad, in gold; they all want to withdraw their accounts at
the clearinghouse, and even the biggest banks can run out of certificates: the country banks hold their reserves in New York, and the New York banks
do not have enough gold to cover them \ts{5}{8}{0:08}. Then the clearinghouse creates money, not by bilateral borrowing but
by joint borrowing and lending \ts{5}{8}{0:50}: it lends to bank A (at 6\% in the example) by printing a \emph{clearinghouse loan
certificate}, a liability of the clearinghouse, guaranteed by all its members, against whatever collateral A can pledge, for example government
securities \ts{5}{8}{1:32}.
\begin{taccts}
\tacct[2.6cm]{Clearinghouse}{$+$Loan to A (6\%)}{$+$Loan certificates}\qquad
\tacct[2.6cm]{Bank A}{$+$Loan certificates}{$+$Loan from clearinghouse}
\end{taccts}
Money has been created at the apex of the system: more of the best money, which circulates among banks as if it were gold. The certificates say
``loan certificate'', but they are treated exactly like the ordinary certificates, since both are liabilities of the clearinghouse
\ts{5}{8}{2:38}. This is a lender of last resort operation, managed by the clearinghouse; the New York
clearinghouse still exists, still clears for its members, and clears much of the Eurodollar market \ts{5}{8}{3:40}.

\paragraph{Bagehot's rule in private hands.} A committee decides \ts{5}{8}{4:21}. Six percent in 1907 was a high rate, ``lend freely at a high rate
against good security'': high, so that the borrower wants to repay. So high, in fact, that banks liked to keep these certificates, the safest asset in
the system at 6\%, and after the crisis the clearinghouse had trouble getting them back \ts{5}{8}{4:42}. Good security: the
bankers know each other and would not take ``garbage''; a bank without acceptable collateral had to borrow some from another bank first, unlike the Fed,
which can decide to let it pass \ts{5}{8}{5:43}.

\paragraph{Playing favourites.} Young describes how, in crises, New York banks refused to pay the deposits of country banks in reserves, offering loan
certificates, take it or leave it, and kept their gold to pay abroad, where the Bank of England would accept nothing but proper international money
\ts{5}{8}{6:46}. It was never clear that the certificates were legal, since they were used as money (unlike the bank notes
backed by government bonds of \cref{sec:st-nbsorigin}) \ts{5}{8}{7:47}; as long as they stayed among bankers nobody noticed, but sometimes
they circulated. After 1907 a law made the practice legal, with rules on acceptable collateral \ts{5}{8}{8:29}. Some argue that this is all
that should have been done, and there is some truth in it: it was a decentralised lender of last resort, discipline in good times and elasticity in
crises \ts{5}{8}{9:11}. But the New York bankers were playing favourites; the banks of the hinterland, refused payment when
everybody needed it, sometimes went bankrupt to save the New York banks, and pressed for democratic control. That is how the Fed came about: ``one way to
understand the Fed is that it is basically a clearinghouse'' \ts{5}{8}{10:13}.

\begin{coursenote}[History of the NYCHA]
The New York Clearing House Association was founded in 1853, long before Morgan's prominence; Morgan's role as ``first among equals'' refers to the
crises of the 1890s and 1907. Clearinghouse loan certificates were first issued in New York in 1857 and 1860 and in most crises thereafter (1873, 1884,
1890, 1893, 1907). The law Mehrling refers to is presumably the Aldrich--Vreeland Act (1908), which authorised groups of national banks
(``national currency associations'') to issue emergency currency against commercial paper and other securities; it was used once, in 1914. Today the
association's successor, The Clearing House, operates CHIPS.
\end{coursenote}\begin{history}[The panic of 1907]
In October 1907 the run on the Knickerbocker Trust Company spread to the trust companies, which were not members of the New York Clearing House. J.~P.~Morgan
assembled bankers in his library to organise loans to the weaker institutions, and the clearinghouses of New York and of many other cities issued loan
certificates; currency went to a premium over certified checks, i.e.\ par between currency and deposits was broken.\par
\emph{References:} Bruner and Carr, \emph{The Panic of 1907}, 2007; Sprague, \emph{History of Crises under the National Banking System}, 1910.
\end{history}

% ------------------------------------------------------------------
\section{The Fed as a public clearinghouse}\label{sec:ch-fed}

The Federal Reserve, when on the gold standard, held gold and issued Federal Reserve notes and bank reserves (deposits of banks at the Fed)
\ts{5}{9}{0:07}. It can run intraday credit on its own books: \emph{Fedwire}, now electronic \ts{5}{9}{0:49}. The private clearing system
of the New York banks is today called \emph{CHIPS}. CHIPS clears first, then Fedwire, and Fedwire is ``real money'', deposits at the Fed, the ultimate
settlement asset \ts{5}{9}{1:12}.
\begin{taccts}
\tacct[2.6cm]{Federal Reserve}{Gold\\Discount loans}{Federal Reserve notes\\Bank reserves}
\end{taccts}
Overdrafts and balances build up during the day. At the end of the day a bank without a positive balance at the Fed must \ts{5}{9}{1:33}:
\begin{enumerate}
  \item borrow from another bank that has a surplus: the \emph{Fed funds market} (next lecture);
  \item or borrow from the Fed, which makes \emph{discount loans} by expanding its own balance sheet, if the banks cannot find each other or B is
        unwilling: nobody forces them \ts{5}{9}{2:16};
  \item or default.
\end{enumerate}
Discount loans to troubled banks appear on the Fed's balance sheet and are small: with thousands of banks, one fails every week or so, and whole
divisions deal with them; they matter for the aggregate only in system-wide stress \ts{5}{9}{2:58}. Then the Fed does what the
clearinghouse did, but legally, as monetary policy: it expands both sides of its balance sheet to create more of the best money. Bernanke goes further than
any clearinghouse would: he does not only lend to banks, he buys mortgage-backed securities from whoever holds them, and lends to broker-dealers; the
clearinghouse was a club, the Fed has more members and in a crisis deals with the whole economy \ts{5}{9}{3:18}.
Private and public clearinghouse differ in membership, rules and priorities, but the mechanics are very much the same \ts{5}{9}{4:41}.

\begin{intuition}[One step beyond the clearinghouse]
In the notes, central banking is ``nothing more than one step beyond the clearinghouse'': it regularizes and strengthens the clearinghouse and
erases the difference between clearinghouse certificates (backed by gold) and clearinghouse \emph{loan} certificates (backed by members' loans), both
becoming deposits or currency of the central bank \mn{5}{5}. A clearinghouse can be broken by any single member who insists on redeeming its
certificates in gold; a central bank cannot be broken by internal forces, because its liabilities \emph{are} the member banks' reserves. All payments
inside the system net out on its books: it makes the system of banks work as if it were a single bank. What remains is managing relations with the
rest of the world, internal and external drains, and that is the origin of the ``art of central banking'' \mn{5}{5}.
\end{intuition}

\begin{definition}[New concepts: Fedwire, CHIPS]
\emph{Fedwire} is the Federal Reserve's real-time gross settlement system, which transfers reserve balances between banks' accounts at the Fed.
\emph{CHIPS} (Clearing House Interbank Payments System), operated by The Clearing House since 1970, is the private system on which large-value, mostly
international dollar payments are netted among member banks and then settled through their accounts at the Fed.
\end{definition}

% ------------------------------------------------------------------
\section{Central bank cooperation}\label{sec:ch-coop}

What can kill a clearinghouse is members withdrawing their gold because they need it \ts{5}{10}{0:08}. Why would they? Distinguish
\ts{5}{10}{0:28}:

\begin{definition}[New concepts: internal and external drain]
An \emph{internal drain} is a loss of reserves to the domestic public, which wants cash; an \emph{external drain} is a loss of reserves to foreigners,
who want the international money (gold). An internal drain can be met by expanding the balance sheet at the top of the national hierarchy; an
external drain comes from a level higher than the national system, and can kill a clearinghouse or a central bank.
\end{definition}

The same holds for central banks: if foreign central banks bring Federal Reserve notes and ask for gold, the Fed can run out, because its notes are
not fully covered \ts{5}{10}{1:09}. Then it \emph{suspends} payments: ``I know these were promises to pay gold, I changed my mind; I will pay
some time later'': it abandons the fixed mint parity, and all central banks have done it, the Bank of England periodically
\ts{5}{10}{1:51}. An internal crisis can be handled by expanding the balance sheet; when the whole country loses reserves, the only way to
get elasticity is to eliminate the discipline by going off gold \ts{5}{10}{2:32}, unless all the central banks agree to act as a clearinghouse
and expand together, one layer higher \ts{5}{10}{2:53}.

\paragraph{The world clearinghouse.} That is what is happening now: Draghi and Bernanke announce QE; the Bank of England and the Swiss National Bank
already do it; the Bank of Japan was in the FT the day before. ``There are only five central banks that matter'', and whether or not they met in a room
like the New York bankers, they act in concert: all expanding at the same time, so there is no external drain, ``the Martians are not yet doing business
with us'' \ts{5}{10}{3:15}. They do what the IMF cannot, since the IMF needs all its members to agree to expand
SDRs; the ``C5'' can do it privately, leaving the rest of the world out, like the New York bankers. Expect complaints from the rest of the world that the
big central banks take care of each other and not of them \ts{5}{10}{4:37}.

\begin{nfigure}
\begin{tikzpicture}[font=\scriptsize,
  l/.style={draw,rounded corners=2pt,align=center,minimum width=3.3cm,minimum height=0.8cm}]
  \node[l,fill=keycol!15] (w) at (0,3.6) {central banks acting together\\(external drain disappears)};
  \node[l,fill=defcol!15] (c) at (0,2.4) {central bank / clearinghouse\\(internal drain, lender of last resort)};
  \node[l,fill=intcol!15] (m) at (0,1.2) {money-centre banks\\(correspondent balances)};
  \node[l,fill=fincol!15] (b) at (0,0) {banks and their customers\\(deposits, overdrafts)};
  \foreach \u/\d in {w/c,c/m,m/b} \draw[-{Stealth[length=4pt]}] (\u) -- (\d);
  \node[right,align=left] at (2.4,1.8) {each level: credit during the day (elasticity),\\settlement in the asset of the level above (discipline)};
\end{tikzpicture}
\caption{Clearing all the way up: the same mechanism of intraday credit and final settlement repeats at each level of the hierarchy, from bank customers to
the central banks of the world.}\label{fig:ch-levels}
\end{nfigure}

% ------------------------------------------------------------------
\flushside
\section*{Key formulas and ideas}
\begin{keyformulas}
\begin{tabularx}{\linewidth}{@{}l L@{}}
One big bank & money version: payments are book entries, balance sheet fixed; credit version: depositor$\to$overdrawn contracts, overdrawn$\to$depositor expands\\
Many banks & paying in reserves is more constraining than one big bank\\
Correspondent balances & a swap of IOUs; a payment either shrinks the debtor bank or expands the creditor bank; hierarchy decides which\\
Clearinghouse & certificates against gold; intraday multilateral netting (elasticity); end-of-day settlement in certificates (discipline)\\
Failing to settle & pay, borrow from a member, or default (losses shared pro rata)\\
Loan certificates & clearinghouse lends its own liabilities against collateral at a high rate: private lender of last resort\\
Fed & public clearinghouse: Fedwire, Fed funds, discount loans\\
Drains & internal: met by expansion; external: suspension, unless central banks expand together\\
\end{tabularx}
\end{keyformulas}

\section*{Exercises}

\begin{exercise}[The payment matrix]
In the credit one-big-bank, $\alpha$ and $\beta$ have deposits of 50 each, $\gamma$ and $\lambda$ overdrafts of 30 each. Compute the size of the balance sheet
after each of the following payments, applied in sequence: $\gamma$ pays $\alpha$ 20; $\alpha$ pays $\lambda$ 40; $\lambda$ pays $\gamma$ 10.
\end{exercise}

\begin{exercise}[Netting]
Four banks have the following gross obligations at the end of the day (row pays column): A$\to$B 30, B$\to$C 40, C$\to$A 35, A$\to$C 10, D$\to$A 5, B$\to$D 15.
\begin{enumerate}[(a)]
  \item Compute the total gross payments, the total of bilateral nets, and the total of multilateral nets.
  \item Which banks must deliver certificates to the clearinghouse, and how many?
\end{enumerate}
\end{exercise}

\begin{exercise}[Correspondent balances]
Country bank C holds a balance of 100 at city bank N. Customers of C pay customers of N 30 more than they receive.
\begin{enumerate}[(a)]
  \item Write the T-accounts when the net is settled by drawing down C's balance.
  \item The balance is now too low; N offers to lend C 50 by crediting its account. Write the T-accounts and explain why this is ``discipline for the debtor
        bank''.
\end{enumerate}
\end{exercise}

\begin{exercise}[Loan certificates]
Bank A owes the clearinghouse 20 at settlement and has no certificates. The clearinghouse issues 20 of loan certificates at 6\% against A's government bonds.
Write the T-accounts of the clearinghouse, A, and the creditor bank B. What happens a week later when A repays? Why might B prefer not to be repaid?
\end{exercise}

% !TEX root = ../main.tex
\chapter{Federal Funds, Final Settlement}\label{ch:fedfunds}
\begin{paracol}{2}

\begin{sources}
\begin{tabularx}{\linewidth}{@{}l l L@{}}
\toprule
\textbf{Segment} & \textbf{Length} & \textbf{Content}\\
\midrule
\seg{6}{1}{L6.1 FT: European Bank Deleveraging} & 5:43 & European banks sell property loans to meet capital ratios; three overnight rates in the FT.\\
\seg{6}{2}{L6.2 What are Fed Funds?} & 5:13 & An interbank market in reserves; the Fed funds rate over the cycle; excess reserves.\\
\seg{6}{3}{L6.3 Payment settlement versus Required Reserves} & 1:11 & The two jobs of a Fed funds desk (Stigum).\\
\seg{6}{4}{L6.4 Payment elasticity/discipline, public and private} & 9:03 & Daylight overdrafts at the Fed; CHIPS; real-time gross settlement.\\
\seg{6}{5}{L6.5 The Function of the Fed Funds Market} & 9:15 & Bilateral overnight borrowing replaces central overnight lending; paying today by promising tomorrow.\\
\seg{6}{6}{L6.6 Payment versus Funding: an example} & 11:24 & Buying an apartment with a mortgage: the Fed funds market and a matched-book dealer.\\
\seg{6}{7}{L6.7 Brokers versus Dealers} & 2:03 & The difference.\\
\seg{6}{8}{L6.8 Payments Imbalances and the Fed Funds Rate} & 7:06 & Imbalances, dealer prices, and how the Fed influences the rate through repo.\\
\seg{6}{9}{L6.9 Secured versus Unsecured Interbank Credit} & 5:05 & Fed funds are unsecured; credit lines; stress in unsecured markets.\\
\seg{6}{10}{L6.10 Required Reserves, redux} & 7:20 & Reserve requirements, lagged and contemporaneous reserve accounting, and why they are ineffective.\\
\bottomrule
\end{tabularx}

\smallskip
\textbf{Files.} Transcripts \file{materials/transcripts/L06-P*.txt}.
\textbf{Reading.} Stigum, \emph{The Money Market}, the chapter on the Fed funds market (quoted in class).
\end{sources}

\section*{Motivation}
The next three lectures cover the three main dollar money markets \ts{6}{1}{3:57}\ts{6}{1}{4:18}:
\begin{enumerate}
  \item \textbf{Fed funds} (this lecture): unsecured interbank loans of reserves at the Fed, the market for ``the best money in the world'';
  \item \textbf{repo} (lecture~7): secured loans;
  \item \textbf{Eurodollars} (lecture~8): unsecured interbank loans of dollars offshore.
\end{enumerate}
In the FT of the day their overnight rates are 0.15\% (US dollar LIBOR, the Eurodollar rate), 0.29\% (US overnight repo) and 0.16\% (Fed funds
effective): all overnight, all different. Why are they what they are? \ts{6}{1}{4:38}\ts{6}{1}{5:00}\ts{6}{1}{5:23} The thread of this lecture: the
Fed funds market exists to meet the \emph{survival constraint} of banks at the end of each day, i.e.\ \emph{final settlement}.

% ------------------------------------------------------------------
\section{FT: European bank deleveraging}\label{sec:ff-ft}

\begin{casestudy}[FT, September 2012: ``Europe's banks race to offload property loans'']
European banks hold property loans left over from the boom in Spain, Ireland and elsewhere, many no longer worth face value
\ts{6}{1}{0:28}\ts{6}{1}{0:50}. A new round of regulation will require more capital against such risky loans, so that losses fall on private
capital rather than the taxpayer \ts{6}{1}{1:10}\ts{6}{1}{1:31}. Capital requirements are ratios of capital to risky assets; rather than raise
capital, the banks shrink the denominator: they package the loans, securitise some, and sell them, sometimes ``at pennies on the dollar'', to
private equity funds, and use the cash to pay down short-term borrowing or borrowing from the ECB \ts{6}{1}{1:52}\ts{6}{1}{2:14}\ts{6}{1}{2:56}.
\end{casestudy}
\begin{taccts}
\tacct[2.4cm]{European bank}{$-$Property loans\\$+$Cash, then\\$-$Cash}{then: $-$Short-term borrowing}\qquad
\tacct[2.4cm]{Private equity fund}{$+$Property loans\\$-$Cash}{}
\end{taccts}
This is \emph{deleveraging}: shrinking the balance sheet on both sides to reach the desired capital ratio \ts{6}{1}{2:56}. The loans are sold at
fire-sale prices, and the banks make no new loans: a credit crunch in Europe. The article is about the connection between banks and other
financial institutions, and between banking and the real economy \ts{6}{1}{3:16}\ts{6}{1}{3:37}.

\begin{definition}[New concepts: capital requirement, securitisation]
A \emph{capital requirement} sets a minimum ratio of a bank's capital (net worth) to its (risk-weighted) assets. \emph{Securitisation} is the pooling
of loans into a package whose cash flows back securities that can be sold to investors (the loans are turned into ``securities'').
\end{definition}

% ------------------------------------------------------------------
\section{What are Fed funds?}\label{sec:ff-what}

In the flow of funds matrix for 2011 (here in stock form, assets and liabilities), line 12, ``Fed funds and security repos'', shows for the
domestic financial sectors \$823 billion of assets and \$1.1 trillion of liabilities: an \emph{interbank} market, in which the financial sector
borrows from and lends to itself; the numbers differ because some non-banks participate too \ts{6}{2}{0:07}\ts{6}{2}{0:28}\ts{6}{2}{0:52}\ts{6}{2}{1:12}.

\begin{definition}[New concepts: Fed funds]\label{def:ff}
\emph{Federal funds} are unsecured loans, typically overnight, of reserves (deposits at the Federal Reserve) between banks: a promise to pay reserves
tomorrow morning in exchange for reserves today \ts{6}{5}{5:17}. Fed funds are \emph{not} a liability or an asset of the Fed, despite the name: they are
assets and liabilities of banks \ts{6}{2}{2:34}.
\end{definition}

In the hierarchy, from the banks' point of view reserves are money and Fed funds are credit: interbank promises to pay money, a credit expansion on the
base of the money on the Fed's balance sheet \ts{6}{2}{1:53}\ts{6}{2}{2:13}. ``If you get off on the wrong foot there, you can be halfway through the lecture
before realizing where you are'' \ts{6}{2}{2:34}.

\paragraph{The rate.} The Fed funds effective rate (FRED, St.~Louis Fed) fluctuates with the business cycle; in the crisis the Fed lowered its target from 5\%
to 2\% and then to nearly zero, where it has stayed, in a range of 0 to 0.25\%, with the effective rate at 0.16\% \ts{6}{2}{2:54}\ts{6}{2}{3:16}\ts{6}{2}{3:36}. The
Fed used to steer the economy by moving this rate; at the zero bound it uses instead the \emph{size} of its balance sheet, quantitative easing
\ts{6}{2}{3:57}\ts{6}{2}{4:17}. Reserves used to be about \$50 billion; now there are over a trillion, excess reserves: reserves are not scarce any more, and yet the
Fed funds market is still active, because it serves the payment system \ts{6}{2}{4:38}\ts{6}{2}{4:58}.

\begin{finance}[The two jobs of a Fed funds desk]
Stigum: the primary job of the manager of a bank's Fed funds desk is to ensure
\begin{enumerate*}[(i)]
  \item that the bank settles with the Fed, every day, and
  \item that in doing so it holds no more excess reserves than it can carry into the next week
\end{enumerate*} \ts{6}{3}{0:09}. The second job, tied to reserve requirements, is now basically defunct: with so many reserves, missing a reserve target is
``ridiculous''. Only the first, the payment job, matters \ts{6}{3}{0:30}\ts{6}{3}{0:50}.
\end{finance}

% ------------------------------------------------------------------
\section{Intraday elasticity: daylight overdrafts and CHIPS}\label{sec:ff-intraday}

The survival constraint, in Stigum's words that the bank must settle with the Fed, takes the precise institutional form of a requirement to \emph{end
the day with a non-negative balance at the Fed} \ts{6}{4}{0:08}. It does not bind \emph{during} the day \ts{6}{4}{0:29}.

\paragraph{Daylight overdrafts at the Fed.} Mehrling spent the summer of 1986 at the Board of Governors as a graduate student, when the Fed was worried
about its credit exposure from daylight overdrafts \ts{6}{4}{0:49}\ts{6}{4}{1:10}. Bank A must pay bank B but has no reserves: it simply runs an overdraft,
which is a loan from the Fed; the Fed creates the reserves and credits them to B, which feels paid and knows nothing of the arrangement
\ts{6}{4}{1:33}\ts{6}{4}{1:53}\ts{6}{4}{2:14}\ts{6}{4}{2:34}:
\begin{taccts}
\tacct[2.2cm]{Fed}{$+$Overdraft of A}{$+$Reserves of B}\quad
\tacct[2.2cm]{Bank A}{}{$-$Deposits\\$+$Overdraft at Fed}\quad
\tacct[2.2cm]{Bank B}{$+$Reserves}{$+$Deposits}
\end{taccts}
This is ``the credit grease that makes the payment system go during the day'' \ts{6}{4}{2:54}. The overdraft is a promise to come up with the reserves by
the end of the day; if A does not, the Fed is forced to lend overnight, but at a penalty, 100 basis points over the Fed funds rate as said in class, to give
A an incentive to find someone who will lend cheaper. The Fed cannot say ``pay me or you're dead'', because it is protecting the payment system; it says
``pay me or you owe me more than you want to''. Sometimes a bank cannot find a counterparty and pays the penalty \ts{6}{4}{3:15}\ts{6}{4}{3:36}\ts{6}{4}{3:56}\ts{6}{4}{4:16}.
Credit expands during the day and collapses at the end of it: the intraday version of the pattern at every level and every time scale \ts{6}{4}{4:36}\ts{6}{4}{4:57}.

\paragraph{CHIPS.} The same happens on the balance sheet of a private clearinghouse, CHIPS, a mutual organisation of banks that collects dues to and from
during the day and settles at the end of the day on Fedwire \ts{6}{4}{4:57}\ts{6}{4}{5:23}\ts{6}{4}{5:44}. Every morning members deposit collateral; during the day
A can owe CHIPS while B is owed by CHIPS \ts{6}{4}{6:05}\ts{6}{4}{6:25}\ts{6}{4}{6:46}\ts{6}{4}{7:07}:
\begin{taccts}
\tacct[2.2cm]{CHIPS}{Due from A}{Due to B}\quad
\tacct[2.2cm]{Bank A}{}{Due to CHIPS}\quad
\tacct[2.2cm]{Bank B}{Due from CHIPS}{}
\end{taccts}
Two layers of intraday credit expansion, a public one on the Fed's balance sheet and a private one on CHIPS's \ts{6}{4}{6:05}. If A fails, all members of CHIPS
are supposed to pay B: B holds an obligation of the clearinghouse, better than one of A, though not as good as one of the Fed \ts{6}{4}{7:28}\ts{6}{4}{7:49}.

\begin{definition}[New concepts: real-time gross settlement]
A \emph{real-time gross settlement} (RTGS) system settles each payment individually (gross) and immediately (real time) in central bank money: once
received, the reserves are money \emph{now}, not to be netted at the end of the day. Fedwire is RTGS; CHIPS nets during the day and settles at the end,
so a CHIPS credit is money only ``hopefully'', after settlement on Fedwire \ts{6}{4}{8:09}\ts{6}{4}{8:30}\ts{6}{4}{8:50}.
\end{definition}

\begin{coursenote}[Daylight overdraft policy]
The Fed and CHIPS have tightened up over time: the size of overdrafts allowed depends on a bank's capital, and the Fed charges for them, on the average
overdraft during the day, so a trillion-dollar overdraft for a few microseconds costs essentially nothing, but still disciplines the bank, which must know
where the money is coming from \ts{6}{5}{2:52}\ts{6}{5}{3:13}\ts{6}{5}{3:33}. (Fees on daylight overdrafts were introduced in 1994, with caps related to capital.
The overnight penalty rate quoted in class may have changed since.) Large security and international transactions cluster at the end of the day, when CHIPS
settles through Fedwire \ts{6}{5}{3:53}.
\end{coursenote}

% ------------------------------------------------------------------
\section{The function of the Fed funds market}\label{sec:ff-function}

At the end of the day, in a closed system, the overdrafts and surpluses sum to zero: if I owe the Fed, somebody out there has excess reserves, and if we
find each other we can both end flat and neither needs the Fed; the same for CHIPS \ts{6}{5}{0:48}\ts{6}{5}{1:09}\ts{6}{5}{1:30}. \emph{Interbank markets exist to let
banks find each other} \ts{6}{5}{1:51}. For CHIPS the relevant interbank market is the Eurodollar market: anyone can be a member of CHIPS, or have a
correspondent that is, including many foreign banks; Fedwire is a much more exclusive club, and its members borrow and lend in the Fed funds market.
Roughly, Fed funds are domestic, public, ``Fed'' dollars; Eurodollars international, private dollars; both serve the same function
\ts{6}{5}{2:11}\ts{6}{5}{2:32}. The Fed funds market creates bilateral overnight borrowing to substitute for the centralised overnight lending the Fed would
otherwise be forced to do \ts{6}{5}{4:14}.

\paragraph{A Fed funds trade.} Bank A, short at the end of the day, borrows from bank B, long: it promises reserves tomorrow morning and gets reserves today,
which pay off its overdraft; the Fed's balance sheet expands during the day and contracts back \ts{6}{5}{4:35}\ts{6}{5}{4:57}\ts{6}{5}{5:37}\ts{6}{5}{5:58}\ts{6}{5}{6:19}.
\begin{taccts}
\tacct[2.4cm]{Bank A}{$+$Reserves (repays overdraft)}{$+$Fed funds borrowed}\qquad
\tacct[2.4cm]{Bank B}{$-$Reserves\\$+$Fed funds lent}{}
\end{taccts}
Banks do not wait for the end of the day: a desk anticipates needs and trades all day long \ts{6}{5}{6:40}\ts{6}{5}{7:00}. Without any overdraft: A borrows reserves
from B and then uses them to pay B; the reserves go back and forth and what is left is a Fed funds loan \ts{6}{5}{7:21}\ts{6}{5}{8:02}\ts{6}{5}{8:30}:
\begin{taccts}
\tacct[2.4cm]{Bank A}{}{$-$Deposits of payer\\$+$Fed funds borrowed}\qquad
\tacct[2.4cm]{Bank B}{$+$Fed funds lent}{$+$Deposits of payee}
\end{taccts}

\begin{proposition}[Paying today by promising tomorrow]
A Fed funds loan pushes final settlement into the future: ``I'm going to pay you today by promising to pay you tomorrow'' \ts{6}{5}{8:50}\ts{6}{5}{9:11}.
\end{proposition}

% ------------------------------------------------------------------
\section{Payment versus funding: an example}\label{sec:ff-example}

Mehrling buys your apartment in Manhattan for \$1 million, with a mortgage from his bank, Citibank; you bank at Chase. Credit checks are ``interesting but not
monetary economics'' \ts{6}{6}{0:08}\ts{6}{6}{0:28}\ts{6}{6}{0:49}. Step by step:
\begin{enumerate}
  \item \textbf{The loan.} Citibank grants the mortgage by crediting a deposit: a swap of IOUs, but with different timing and different signers: Mehrling
        owes \$1 million over 30 years, Citibank owes \$1 million on demand; when Citibank signs an IOU it is money, when Mehrling does it is credit
        \ts{6}{6}{1:09}\ts{6}{6}{1:29}\ts{6}{6}{1:49}\ts{6}{6}{2:10}.
  \item \textbf{Funding the payment.} The deposit must go to Chase, and Citibank does not keep \$1 million of idle reserves. Its credit department knows the
        closing is coming, so it borrows \$1 million in the Fed funds market, not from Chase but from a third party, HSBC, which lends reserves it has (or
        runs an overdraft) \ts{6}{6}{2:30}\ts{6}{6}{2:51}\ts{6}{6}{3:34}\ts{6}{6}{3:54}. Eventually the mortgage will be packaged and securitised, but that takes weeks; the
        payment must happen now \ts{6}{6}{3:11}.
  \item \textbf{The payment.} Citibank pays Chase in reserves and Chase credits your deposit: it is happy to, because it is settled, it has money
        \ts{6}{6}{4:14}\ts{6}{6}{4:35}\ts{6}{6}{4:57}\ts{6}{6}{5:17}.
  \item \textbf{Chase lends on.} Reserves pay nothing, so Chase lends them in the Fed funds market, to HSBC \ts{6}{6}{5:38}\ts{6}{6}{5:59}.
\end{enumerate}
\begin{taccts}
\tacct[2.0cm]{Citibank}{$+$Mortgage}{$+$Fed funds (HSBC)}\quad
\tacct[2.0cm]{HSBC}{$+$Fed funds (Citi)}{$+$Fed funds (Chase)}\quad
\tacct[2.0cm]{Chase}{$+$Fed funds (HSBC)}{$+$Deposit (you)}

\medskip
\tacct[2.2cm]{Mehrling}{$+$Apartment}{$+$Mortgage}\qquad
\tacct[2.2cm]{You}{$-$Apartment\\$+$Deposit}{}
\end{taccts}

\paragraph{Who lent the money?} Follow the chain: Citibank borrowed from HSBC, HSBC from Chase, Chase from you. ``In a certain sense I bought it from you and you
lent me the money for it'', though Mehrling owes the intermediaries, not you \ts{6}{6}{6:19}\ts{6}{6}{6:40}\ts{6}{6}{7:00}. Until the mortgage is financed by some
pension fund, it has been funded through the \emph{payment system}, by an expansion of balance sheets: nobody saved extra money so that Mehrling could borrow;
you transformed a house into a deposit, he acquired a house and a mortgage \ts{6}{6}{7:00}\ts{6}{6}{7:21}.

\paragraph{HSBC as a dealer.} HSBC lends and borrows Fed funds at the same time, flat in exposure: it does so only because it lends at a higher rate than it
borrows. It is a \emph{dealer} in Fed funds, quoting rates at which it will lend and borrow and setting them so that over the day inflows match outflows; when
they do not, it must find funds itself, or place the surplus \ts{6}{6}{7:41}\ts{6}{6}{8:01}\ts{6}{6}{8:22}\ts{6}{6}{8:44}.

\begin{definition}[New concepts: matched book]
A dealer runs a \emph{matched book} when its borrowing and its lending in a market have the same size (and maturity), so that it earns the spread without
net exposure to the price \ts{6}{6}{8:44}\ts{6}{6}{9:04}.
\end{definition}

\paragraph{An internal drain.} If you withdraw the \$1 million in cash (``1MM'', in Stigum's notation), Chase must get currency: it asks the Fed to convert
reserves into cash, which is physically ``downtown at 33 Liberty Street''. The Fed's liabilities change form, one million more of currency and one million less
of reserves; an internal drain, which causes no monetary problem, since both are liabilities of the Fed \ts{6}{6}{9:25}\ts{6}{6}{10:06}\ts{6}{6}{10:46}\ts{6}{6}{11:07}.

% ------------------------------------------------------------------
\section{Brokers versus dealers}\label{sec:ff-brokers}

A \emph{broker} puts buyer and seller together and takes a fee, like a real-estate broker, who does not buy your house and resell it. In a brokered Fed funds
trade the loan sits on the balance sheets of A and B only \ts{6}{7}{0:07}\ts{6}{7}{0:28}\ts{6}{7}{0:48}. A \emph{dealer} like HSBC stands in between, and the loans sit on
its balance sheet: Chase and Citibank probably do not even know that they are, in effect, lending to and borrowing from each other \ts{6}{7}{1:08}\ts{6}{7}{1:29}. The
dealer is ``the source of liquidity in these markets'', a theme for the rest of the course \ts{6}{7}{1:50}. (Brokers and dealers were defined in
\cref{def:nh-mm}.)

% ------------------------------------------------------------------
\section{Payment imbalances and the Fed funds rate}\label{sec:ff-rate}

During the day the Fed and CHIPS expand to make payments possible; at the end both shrink back, and the Fed funds market is one way to do it. But the credit does
not vanish: it moves to the bilateral balance sheets of borrowers and lenders and onto the balance sheets of the Fed funds dealers \ts{6}{8}{0:07}\ts{6}{8}{0:28}\ts{6}{8}{0:50}.
The centre protects itself (``we won't extend this overnight, you do it''); daily payment imbalances that could not be cleared show up as expansion of Fed funds credit
and of dealer balance sheets, pushed into the future in the hope that tomorrow the flow reverses \ts{6}{8}{1:11}\ts{6}{8}{1:32}\ts{6}{8}{1:52}. These quantities are highly
sensitive measures of payment patterns, and so are the prices: a dealer with many more borrowers than lenders raises its rate to attract lenders, and lowers it in the
opposite case. The Fed funds rate is a \emph{market rate}, and it fluctuates within the day too \ts{6}{8}{1:52}\ts{6}{8}{2:14}\ts{6}{8}{2:34}\ts{6}{8}{2:56}.

\paragraph{How does the Fed ``set'' the rate?} Not as in intermediate macro: the Fed does not participate in the Fed funds market at all, neither borrowing nor
lending \ts{6}{8}{3:16}\ts{6}{8}{3:37}. It has a \emph{target}, and it influences the rate by changing the quantity of reserves, by entering other markets, in particular the
repo market \ts{6}{8}{3:57}. Almost every day it anticipates payment imbalances that would push the rate away from target, and adjusts the quantity of reserves: it
changes money in order to influence credit \ts{6}{8}{4:18}\ts{6}{8}{4:38}\ts{6}{8}{4:59}. To add reserves it makes an overnight repurchase agreement loan, for example to a
security dealer funding its securities, which ends up with reserves it does not want and lends them in the Fed funds market, putting downward pressure on the rate
\ts{6}{8}{5:20}\ts{6}{8}{5:41}\ts{6}{8}{6:01}\ts{6}{8}{6:21}:
\begin{taccts}
\tacct[2.0cm]{Fed}{$+$Repo loan}{$+$Reserves}\quad
\tacct[2.0cm]{Security dealer}{$+$Reserves, then $-$Reserves\\$+$Fed funds lent}{$+$Repo from Fed}
\end{taccts}
The \emph{Fed funds effective} rate is the average of the rates at which funds trade during the day; good borrowers pay less, bad ones more \ts{6}{8}{6:45}.

\begin{definition}[New concepts: target rate, effective rate]
The \emph{Fed funds target} is the rate (since 2008 a range) that the Federal Open Market Committee (FOMC) aims at; the \emph{effective} Fed funds rate is the
volume-weighted average of the rates actually paid in the market on a given day.
\end{definition}

% ------------------------------------------------------------------
\section{Secured versus unsecured interbank credit}\label{sec:ff-unsecured}

Fed funds are \emph{unsecured} overnight loans: no collateral. If Citibank does not pay HSBC, HSBC can sue, but there is no asset to seize; a mortgage, by contrast,
is secured by the house \ts{6}{9}{0:08}\ts{6}{9}{0:29}\ts{6}{9}{0:49}\ts{6}{9}{1:10}. Eurodollars too are unsecured interbank loans \ts{6}{9}{1:10}. Why lend a million dollars
overnight on a promise? Because banks are in a network, it is their business, and tomorrow they will be on the other side; they watch their exposure to each
counterparty very closely, with \emph{credit lines}: ``we're already in them for \$10 million, tell them your credit line to us is full'', and they diversify
their lending \ts{6}{9}{1:30}\ts{6}{9}{1:50}\ts{6}{9}{2:11}\ts{6}{9}{2:32}.

The market, especially the Eurodollar market, is under great stress: with many insolvent banks in Europe, nobody wants to lend unsecured, and even a dealer in the
middle worries, because its business is highly leveraged on both sides: if Citibank does not pay HSBC, HSBC must still pay Chase; it is not a pass-through, it faces its
own survival constraint \ts{6}{9}{2:53}\ts{6}{9}{3:13}\ts{6}{9}{3:33}\ts{6}{9}{3:54}. Repurchase agreements, next lecture, are \emph{secured}, by a Treasury bill for example
\ts{6}{9}{3:54}. This is where large amounts are transferred, where risk is taken and liquidity created: the Fed funds market is the market for the best money in
the world, Eurodollars another unsecured market, repo the secured one; understanding them puts us at the very core of the international payment system
\ts{6}{9}{4:14}\ts{6}{9}{4:35}\ts{6}{9}{4:56}.

\begin{definition}[New concepts: secured and unsecured credit, credit line]
A \emph{secured} loan is backed by collateral that the lender can take if the borrower defaults; an \emph{unsecured} loan rests only on the borrower's general
promise. A \emph{credit line} is the maximum exposure a lender is willing to have to a given counterparty.
\end{definition}

% ------------------------------------------------------------------
\section{Required reserves, redux}\label{sec:ff-required}

``The least important part of this lecture'': reserve requirements \ts{6}{10}{0:07}. Stigum describes two-week \emph{maintenance periods} over which a bank must hold
an average quantity of reserves, so a bank had two tasks: settle every day (easy, with the Fed funds market) and hit the average without overshooting
\ts{6}{10}{0:27}\ts{6}{10}{0:48}.

\paragraph{Monetarism and reserve accounting.} Under Volcker, monetarists wanted to use reserve requirements to control the money supply: if deposits require
reserves, a shortage of reserves limits the expansion of deposits and of bank credit \ts{6}{10}{1:08}\ts{6}{10}{1:29}.
\begin{itemize}
  \item With \emph{lagged reserve accounting}, required reserves depended on deposits of two weeks earlier: the loan and the deposit are already there, so the
        Fed has to supply the reserves, or somebody will miss the requirement; the Fed only accommodates demand \ts{6}{10}{1:50}\ts{6}{10}{2:12}\ts{6}{10}{2:32}.
  \item With \emph{contemporaneous reserve accounting}, required reserves depended on deposits in the same period: harder to hit, with both deposits and
        reserves moving; it caused a lot of chaos in the Fed funds market and was abandoned \ts{6}{10}{2:53}\ts{6}{10}{3:15}\ts{6}{10}{3:36}.
\end{itemize}
It was an attempt to impose discipline on a system that seemed out of control, in a period of double-digit inflation; according to Stigum it did not really work
\ts{6}{10}{3:36}\ts{6}{10}{3:56}. Today, with trillions of reserves, requirements are a dead letter, yet the Fed funds market exists for payments \ts{6}{10}{4:17}.

\begin{coursenote}[Dates]
The Fed's ``monetarist experiment'' of targeting reserves ran from October 1979 to 1982. Contemporaneous reserve accounting was introduced in February 1984 and
replaced by lagged accounting again in 1998. (In March 2020, after these lectures, the Fed set reserve requirement ratios to zero.)
\end{coursenote}

\paragraph{Were reserve requirements ever important?} Asked whether they will matter again after the crisis, Mehrling answers that a trillion of excess reserves will
not go away soon, but the deeper question is whether the quantity of reserves ever significantly restricted credit or money \ts{6}{10}{4:38}\ts{6}{10}{4:59}\ts{6}{10}{5:19}.
A parallel banking system, shadow banking, creates close substitutes for money and credit with no reserve or capital requirement, possibly offshore; marginal credit can flow
there. Squeezing banks, as with the European capital requirements of the FT article, does not mean there is no credit, only that it does not come from banks: ``like a
balloon, you punch it here and it blows out there'' \ts{6}{10}{5:39}\ts{6}{10}{5:59}\ts{6}{10}{6:20}. Mehrling sides with those who think the reserve constraint is ineffective in
modern financial markets; it probably was not in 1950, when it was invented and the textbooks were written, because there were no alternatives \ts{6}{10}{6:40}\ts{6}{10}{7:00}.

% ------------------------------------------------------------------
\flushside
\section*{Key formulas and ideas}
\begin{keyformulas}
\begin{tabularx}{\linewidth}{@{}l L@{}}
Fed funds & unsecured overnight interbank loans of reserves; assets/liabilities of banks, not of the Fed\\
Survival constraint & end the day with a non-negative balance at the Fed\\
Intraday & daylight overdrafts at the Fed (RTGS) and at CHIPS (net): credit expansion during the day\\
End of day & surplus and deficit banks find each other in the Fed funds (or Eurodollar) market; else penalty borrowing from the Fed\\
Payment vs funding & a payment funded by interbank borrowing: paying today by promising tomorrow\\
Dealer & borrows and lends at a spread, matched book; source of market liquidity\\
Fed funds rate & a market rate; the Fed targets it by adjusting reserves (repo), not by trading Fed funds\\
Reserve requirements & lagged vs contemporaneous accounting; ineffective with shadow banking\\
\end{tabularx}
\end{keyformulas}

\section*{Exercises}

\begin{exercise}[A day at the Fed]
Bank A starts the day with 0 reserves, B with 50. During the day A pays B 80 and B pays A 30.
\begin{enumerate}[(a)]
  \item Track A's balance at the Fed; what is its maximum daylight overdraft?
  \item What must A do by the end of the day? Write the T-accounts if it borrows from B in the Fed funds market.
  \item Repeat when B is unwilling to lend to A, and a dealer D borrows from B and lends to A.
\end{enumerate}
\end{exercise}

\begin{exercise}[The apartment]
Redo the apartment example when Chase itself lends Citibank the reserves in the Fed funds market. Then suppose the buyer's bank and the seller's bank are the same:
what happens to reserves and to the Fed funds market?
\end{exercise}

\begin{exercise}[Dealer pricing]
A Fed funds dealer quotes 0.15\% to borrow and 0.20\% to lend. During the morning it receives requests to borrow 300 and offers to lend 100.
\begin{enumerate}[(a)]
  \item What is its net position if it accepts everything?
  \item Which way should it move its quotes, and why does this move the effective rate?
\end{enumerate}
\end{exercise}

\begin{exercise}[Reserve accounting]
A bank with deposits of 1000 and a 10\% reserve requirement grants a loan of 100 that is spent within the bank. Compare what it must do under lagged and under
contemporaneous reserve accounting, and explain why the Fed must accommodate in the first case.
\end{exercise}

% !TEX root = ../main.tex
\chapter{Repos, Postponing Settlement}\label{ch:repo}
\begin{paracol}{2}

\begin{sources}
\begin{tabularx}{\linewidth}{@{}l l L@{}}
\toprule
\textbf{Segment} & \textbf{Length} & \textbf{Content}\\
\midrule
\seg{7}{1}{L7.1 FT: The Impact of QE3} & 2:45 & Second thoughts on QE3: buying time, and hedge funds betting against it.\\
\seg{7}{2}{L7.2 Money Market Interest Rate Patterns} & 3:34 & LIBOR, Fed funds, repo: the usual order, and its reversal.\\
\seg{7}{3}{L7.3 What is Repo?} & 3:19 & A collateralized loan organised as a sale and repurchase.\\
\seg{7}{4}{L7.4 Repo in Balance Sheets} & 7:34 & Repo and reverse repo around a security dealer; the confusing language.\\
\seg{7}{5}{L7.5 Comparison with Fed Funds} & 5:09 & Repo as an open, dealer-made market.\\
\seg{7}{6}{L7.6 Legal Construction of Repo} & 9:46 & Symmetry, margin, haircut, repo interest (Stigum's example).\\
\seg{7}{7}{L7.7 Security Dealers Balance Sheet} & 11:07 & Primary dealer statistics; dealers as banks; haircuts and liquidity spirals.\\
\seg{7}{8}{L7.8 Repo, Modern Finance, and the Fed} & 8:49 & Repo as money; Fed open market operations as repos with dealers.\\
\seg{7}{9}{L7.9 Interest Rate Spreads: Before the Crisis} & 5:33 & Why repo $<$ Fed funds: Bagehot's market rate and bank rate.\\
\seg{7}{10}{L7.10 Interest Rate Spreads: After the Crisis} & 8:07 & Repo $>$ Fed funds: stress in dollar funding; reading spreads as clues.\\
\bottomrule
\end{tabularx}

\smallskip
\textbf{Files.} Transcripts \file{materials/transcripts/L07-P*.txt}. Segment~5 has no YouTube captions: its transcript
\file{L07-P05.txt} is extracted from the original \file{transcript-lecture-7to12.txt}, so its times may be off by a few seconds.
\textbf{Reading.} Stigum, \emph{The Money Market}, chapter on repo (pages 533--534 are quoted in class).
\end{sources}

\section*{Motivation}
The Fed funds market (lecture~6) lets banks meet at the end of the day and push settlement to tomorrow. The repo market does the same, but for
almost everybody and against collateral: it is the core funding market of the security dealers, and through them of modern finance. The lecture
explains what repo is (\cref{sec:rp-what}), how it appears in balance sheets (\cref{sec:rp-bs}), how it is priced (\cref{sec:rp-legal}), why
dealers that fund themselves this way are banks (\cref{sec:rp-dealers}), how the Fed uses it (\cref{sec:rp-fed}), and what the spread between repo
and Fed funds tells us about the system (\cref{sec:rp-spreads}).

% ------------------------------------------------------------------
\section{FT: the impact of QE3}\label{sec:rp-ft}

\begin{casestudy}[FT, September 2012: second thoughts on QE3]
John Plender, ``Central bank firepower risks creating false sense of security'': QE3 puts a floor under certain security prices but does nothing
fundamental; if the problem is fiscal, central banks only buy time, and if nobody uses the time, nothing has been done. His last paragraph:
``throughout history monetary experiments have shown a nasty tendency to blow up\dots\ this experiment is vast; it must not breed complacency''
\ts{7}{1}{0:06}. On the same page, ``Hedge fund sceptics take on the bulls over QE infinity'': some funds argue that QE1
worked well, QE2 less, QE3 may not even lift asset prices, and trade against those who expect a boom \ts{7}{1}{1:08}.
\end{casestudy}
If private investors trade against it, QE buys even less time \ts{7}{1}{2:09}. Central banks act not because they are optimistic, but because
nobody else is doing anything; and in repo rates we can see symptoms of the stresses in the system \ts{7}{1}{2:09}.

% ------------------------------------------------------------------
\section{Money market interest rate patterns}\label{sec:rp-patterns}

In the FT: US dollar LIBOR overnight 0.1509\%, US overnight repo 0.31\%, Fed funds effective 0.16\% \ts{7}{2}{0:07}. Before the crisis
the typical order was
\[
  \text{repo} \;<\; \text{Fed funds} \;<\; \text{LIBOR (Eurodollar)} ,
\]
and Mehrling had a whole story for it, from a paper he wrote before the crisis \ts{7}{2}{0:56}. Now repo is \emph{above} Fed funds:
odd, since repo is secured and Fed funds are not \ts{7}{2}{1:38}.

\begin{nfigure}
\begin{tikzpicture}[font=\scriptsize,x=1cm,y=1cm]
  \draw[-{Stealth[length=3pt]}] (0,0) -- (8,0) node[right]{time (mid-2005 to mid-2006)};
  \draw[-{Stealth[length=3pt]}] (0,-1.8) -- (0,1.6) node[above]{spread vs target};
  \node[left] at (0,0) {0};
  \draw[thick,intcol] plot[smooth] coordinates {(0,0.45) (1,0.5) (1.5,-0.6) (1.7,0.45) (3,0.5) (3.5,-0.5) (3.7,0.4) (5,0.45) (5.5,-0.7) (5.7,0.4) (7.5,0.5)};
  \draw[thick,thmcol] plot[smooth] coordinates {(0,-0.4) (1,-0.35) (1.5,-1.5) (1.7,-0.4) (3,-0.35) (3.5,-1.3) (3.7,-0.4) (5,-0.35) (5.5,-1.6) (5.7,-0.4) (7.5,-0.35)};
  \node[intcol,right] at (7.5,0.5) {target $-$ repo $>0$};
  \node[thmcol,right] at (7.5,-0.35) {target $-$ Eurodollar $<0$};
\end{tikzpicture}
\caption{Schematic of the pre-crisis chart shown in class \ts{7}{2}{1:58}: Fed funds \emph{target} minus repo was typically positive, target
minus Eurodollar typically negative. The downward spikes occur before each step of the 2005--2006 tightening, as market rates rise in anticipation
\ts{7}{2}{2:39}. (Shape illustrative.)}\label{fig:rp-spreads}
\end{nfigure}

% ------------------------------------------------------------------
\section{What is repo?}\label{sec:rp-what}

\begin{definition}[New concepts: repurchase agreement]\label{def:rp}
A \emph{repurchase agreement} (``repo'', ``RP'') is a loan, overnight or for a term, collateralized by a security, organised as a
\emph{sale and repurchase}: the borrower of money sells a security today and promises to buy it back tomorrow at a fixed price; the securities go
from the borrower to the lender and back, the money from the lender to the borrower and back \ts{7}{3}{0:08}.
(Stigum draws the ``first leg'' and the ``second leg'' of the repo \ts{7}{3}{1:34}.)
\end{definition}

The securities are collateral in a strong sense: if the borrower does not repay, the lender \emph{already holds} the securities and can sell them;
there is no need to depend on the borrower's credit \ts{7}{3}{1:55}. And symmetrically: if the lender does not return the securities, the
borrower need not repay the money. Such a \emph{fail} turns into a free overnight continuation of the loan until the securities are delivered
\ts{7}{3}{2:37}. The money collateralizes the loan of securities, and the securities collateralize the loan of money.

\begin{nfigure}
\begin{tikzpicture}[font=\small,
  ag/.style={draw,rounded corners=2pt,minimum width=2.6cm,minimum height=0.8cm,align=center}]
  \node[ag,fill=defcol!10] (b) at (0,0) {borrower of money\\(cash taker)};
  \node[ag,fill=intcol!10] (l) at (6,0) {lender of money\\(cash provider)};
  \draw[-{Stealth[length=4pt]},thick,keycol] ([yshift=4pt]b.east) -- node[above,font=\scriptsize]{today: securities} ([yshift=4pt]l.west);
  \draw[-{Stealth[length=4pt]},thick,fincol] ([yshift=-4pt]l.west) -- node[below,font=\scriptsize]{today: money} ([yshift=-4pt]b.east);
  \draw[-{Stealth[length=4pt]},dashed,keycol] (l.south) to[out=-110,in=-70] node[below,font=\scriptsize]{tomorrow: securities back} (b.south);
  \draw[-{Stealth[length=4pt]},dashed,fincol] (b.north) to[out=70,in=110] node[above,font=\scriptsize]{tomorrow: money $+$ interest} (l.north);
\end{tikzpicture}
\caption{The two legs of a repo \ts{7}{3}{1:14}: a sale of securities today and their repurchase tomorrow at a slightly higher price.}\label{fig:rp-legs}
\end{nfigure}

% ------------------------------------------------------------------
\section{Repo in balance sheets}\label{sec:rp-bs}

There are trillions of these overnight loans \ts{7}{4}{0:07}. The language is not standardised: ``doing repo'', ``reversing in securities'', ``repoing
securities''; and, Mehrling warns, the same phrase can mean opposite sides of the trade in New York and in London. Balance sheets keep you straight \ts{7}{4}{0:27}.
Two rules:
\begin{enumerate}
  \item almost every repo has a \emph{security dealer} on one side \ts{7}{4}{0:48};
  \item ``follow the money'': a repo booked as a liability means money borrowed, as an asset money lent; money is better than securities, so its side tells
        you where to book \ts{7}{4}{1:54}.
\end{enumerate}

\begin{taccts}
\tacct[2.0cm]{Pension fund}{Repo loan to dealer}{}\quad
\tacct[2.0cm]{Security dealer}{Reverse repo (loan to bank)}{Repo (from pension fund)}\quad
\tacct[2.0cm]{Bank}{Securities}{Repo (from dealer)}
\end{taccts}

\begin{itemize}
  \item \textbf{Pension fund}: lends money overnight, against a Treasury, for interest. For it this is practically a deposit, a way to hold cash, and
        better than a deposit, because large deposits are not insured: if the bank fails the cash is gone, while in repo it holds a Treasury
        \ts{7}{4}{2:15}.
  \item \textbf{Dealer lending money}: takes in securities; Stigum calls it a \emph{reverse} \ts{7}{4}{3:16}. Securities move one way across the
        diagram, money the other \ts{7}{4}{3:57}.
  \item \textbf{Bank}: holds securities and funds them by repoing them out overnight, instead of with deposits: repo as the liability that funds an asset
        \ts{7}{4}{4:17}.
\end{itemize}
The dealer's ``repo'' and ``reverse'' are the same instrument seen from the two sides; they get different names because the dealer is in the middle,
where one is a liability and the other an asset \ts{7}{4}{5:40}. To add to the confusion, the Fed uses the words the other way round: when it
borrows from dealers it calls it a ``reverse repo'', when it lends (open market operations) a ``repo'' \ts{7}{4}{6:21}. Always translate
into balance sheets: who is borrowing money? \ts{7}{4}{7:23}

\begin{history}[Repo before the Fed]
Repo is ``very ancient'': in the United States it existed before the Fed, as the way banks traded with each other before the Fed funds market. With a
correspondent you could rely on the relationship; with a stranger you wanted collateral: lend against a bond, and you only need to check the bond, not the
borrower. Mehrling links the rise of bond ratings (Moody's) to this need to know how much to lend against which bond \ts{7}{4}{4:38}.
(John Moody published his first bond ratings in 1909.)\par
\emph{References:} R.~Sylla, ``An historical primer on the business of credit rating'', in R.~Levich et al.\ (eds.), \emph{Ratings, Rating Agencies and the
Global Financial System}, Kluwer, 2002.
\end{history}

\paragraph{Comparison with Fed funds.} The repo market is like the Fed funds market: it lets deficit and surplus agents meet and push their liabilities and
assets to another day, with a dealer in the middle \ts{7}{5}{0:03}. Two differences \ts{7}{5}{1:25}:
\begin{enumerate}
  \item \textbf{Openness}: anyone with a Treasury bond can do repo, without being a bank or a member of the Federal Reserve System: it is ``much more
        democratic'' \ts{7}{5}{1:33};
  \item \textbf{Dealing versus brokering}: in Fed funds there is little dealing, it is mostly brokered; in repo it is almost all dealing, on the balance
        sheets of security dealers \ts{7}{5}{0:52}.
\end{enumerate}
Why a dealer market? Repo is heterogeneous: any security can be repoed (corporate bonds, mortgage-backed securities, \dots) and anyone can participate.
The \emph{name of the dealer}, a primary dealer, a big bank, on the other side of every trade, is what makes the market homogeneous: you only worry about
the dealer \ts{7}{5}{2:12}. There is even \emph{general collateral} repo, in which the dealer may substitute one
security for another of the same class, to keep things smooth; the dealer is a liability on both sides, to return securities and to return money
\ts{7}{5}{3:53}. Repo defaults usually happen not because someone is insolvent but because the securities were sold on and not returned
\ts{7}{5}{4:42}.

% ------------------------------------------------------------------
\section{The legal construction of repo: haircut and interest}\label{sec:rp-legal}

Repo is like a collateralized loan, but much more symmetric \ts{7}{6}{0:07}. With a mortgage, the bank has a claim on the house, but foreclosure takes months
and legal fees, and nobody intends to use it \ts{7}{6}{0:28}. A repo is an overnight loan of \$10 million: the lender \emph{owns} the security, can sell it or
re-pledge it (\emph{rehypothecate}) and trade with it all day, and is only liable for returning it the next day; the borrower can spend the money as it likes
\ts{7}{6}{0:48}.

\paragraph{Who gives margin to whom?} The lender of money fears that the security falls below the value of the loan, which gives the borrower an incentive not
to repay; the lender of securities fears the opposite \ts{7}{6}{1:51}. Stigum's answer: it is always the \emph{lender of money} who gets margin
\ts{7}{6}{2:31}. ``This is part of the hierarchy of money and credit'': money is the better asset, and its holder can demand extra security; and the demand for
liquidity is the short side of the market, people pay up to get money, not to get collateral \ts{7}{6}{2:31}. (In swaps, by contrast,
typically both sides post margin \ts{7}{6}{9:09}.)

\begin{definition}[New concepts: haircut, repo rate]
The \emph{haircut} is the percentage by which the amount lent is smaller than the market value of the collateral; it protects the lender of money against a fall
in the collateral's price. The \emph{repo rate} is the interest rate on the loan, quoted on a simple-interest, actual/360 basis.
\end{definition}

\paragraph{Stigum's example} \ts{7}{6}{2:52}--\ts{7}{6}{7:47}. A 10-year Treasury bond trades a bit above par (its coupon is higher than the current 10-year
rate). It is offered as collateral for an overnight loan with a haircut of 2\%, so the amount lent per 100 of face value is 99.228. The repo rate is 4.92\%
(annualised). The next day the borrower repays principal plus simple interest for one day on a 360-day year:
\begin{keyformulas}
\[
  \text{repayment} = L\left(1 + r_{\text{repo}}\,\frac{d}{360}\right),\qquad L = (1-h)\,P ,
\]
with $P$ the market value of the collateral, $h$ the haircut, $d$ the number of days (1 overnight, 3 over a weekend, 7 for a week).
\end{keyformulas}
The conventions (simple, not compound interest; a 360-day year) date ``from time immemorial'': the market prices them in, and overnight they change the result
only at the fourth decimal, but traders must get them right \ts{7}{6}{6:04}. On a \$1 million bond the overnight interest
is about \$137: the cost of borrowing a million dollars overnight \ts{7}{6}{7:26}.

\begin{coursenote}[Checking the numbers]
The price of the bond is garbled in the captions (``129 30 seconds''); the other numbers are consistent with a price of about $101\tfrac{8}{32}=101.25$:
$0.98\times101.253=99.228$. Interest on \$1 million for one day at 4.92\%/360 is \$136.67, the ``137'' of the lecture; on the amount actually lent,
\$992{,}280, it is \$135.61. Bond prices were quoted in 32nds of a point (one point $=$ 1\% of face value).
\end{coursenote}

\paragraph{Haircuts make collateral homogeneous.} Different securities get different haircuts (2\% for a Treasury, 5, 10 or 50\% for a mortgage-backed
security, which in a crisis may not be accepted at all) and then the same repo rate \ts{7}{6}{4:38}. The haircut is negotiated first,
according to the volatility of the collateral's price, then the rate \ts{7}{6}{4:59}.

% ------------------------------------------------------------------
\section{The security dealer's balance sheet}\label{sec:rp-dealers}

The New York Fed publishes weekly statistics on the primary dealers \ts{7}{7}{0:07}:
\begin{itemize}
  \item \textbf{Transactions} by type of security: Treasury coupon securities, agencies (Fannie Mae, Freddie Mac), mortgage-backed, corporate
        \ts{7}{7}{0:48}.
  \item \textbf{Net positions} (inventory, hence risk exposure): about \$26 billion of Treasury bills, \$68 billion of coupon securities due in three years or
        less, negligible longer bonds \ts{7}{7}{1:08}. This is unusual: normally dealers are \emph{short} bills and \emph{long} bonds, picking up the term
        spread. With Operation Twist the Fed buys bonds from the dealers and sells them bills, sucking out their bond inventory and wiping out their short in
        bills; the dealers absorb the mismatch until they can distribute it \ts{7}{7}{1:49}. Analysts from
        Barclays, in a Reuters article, explained the high repo rate this way; Mehrling wonders why it should affect the repo rate \ts{7}{7}{2:30}.
        (Net positions say nothing about hedges with derivatives \ts{7}{7}{4:14}.)
  \item \textbf{Financing} (gross): securities in, about \$1.9 trillion, financed overnight or at term, and securities out (short positions)
        \ts{7}{7}{4:35}.
\end{itemize}
The memorandum items give the repo book, in billions \ts{7}{7}{5:38}:
\begin{taccts}
\tacct[3.0cm]{Primary dealers (\$ billion)}{Reverse repo, overnight\amt{878}\\Reverse repo, term\amt{1{,}316}}{Repo, overnight\amt{1{,}842}\\Repo, term\amt{931}}
\end{taccts}
Dealers fund their long positions in the repo market and their short positions in the reverse market \ts{7}{7}{6:44}. And look at the maturities: the big number
on the asset side is \emph{term}, on the liability side \emph{overnight}. Dealers borrow short and lend long, ``just like a bank'' (term may be a week, but it is
more than a day) \ts{7}{7}{7:25}.

\begin{proposition}[Security dealers are banks]
Overnight repo to dealers is, for institutional investors, a deposit, even better than one because secured; at \$1.8 trillion it is larger than M1. Dealers,
trading securities for profit, create an enormous volume of monetary assets: the wholesale money supply \ts{7}{7}{7:04}.
\end{proposition}

\paragraph{Haircuts and liquidity spirals.} Repo must be rolled over, and everything is ``up for grabs'' at each renewal: haircut and rate
\ts{7}{7}{9:28}. If the haircut on a security you finance with repo jumps overnight from 2\% to 10\%, you have a \emph{funding gap}: yesterday you
could borrow 98 cents on the dollar, today 90; if you cannot find the other eight, you must sell the security at a fire-sale price, which pushes prices down
further. In the crisis, haircuts on ``triple-A'' mortgage-backed securities went from 2\% to 10\% to 50\%, and the scramble for funding distorted every other price
in the world \ts{7}{7}{9:48}.

\begin{intuition}[Leverage and the haircut]
With haircut $h$, a dealer with capital $E$ can hold securities worth at most $E/h$: leverage $1/h$ (50 at $h=2\%$, 10 at $10\%$, 2 at $50\%$). A rise of the
haircut from 2\% to 10\% forces assets to fall by a factor of five for the same capital: a positive-feedback loop if the sales lower prices and raise haircuts
further.
\end{intuition}

% ------------------------------------------------------------------
\section{Repo, modern finance and the Fed}\label{sec:rp-fed}

Most money and banking textbooks do not mention repo: they are about retail money and banks, and a security dealer ``is not part of our subject''. This is where
modern finance changed the world: not being called a bank does not mean not operating like one, and producing liquid liabilities is producing money supply for
someone \ts{7}{8}{0:07}. The lender in repo may be a money market mutual fund, which holds repo and issues shares that function as deposits, to
corporations or even foreign central banks: money, directly or passed through a fund \ts{7}{8}{0:49}.

\paragraph{Open market operations are repos with dealers.} The textbook says the Fed buys or sells Treasury securities to change the money supply; in reality it
trades with the primary dealers, by auction: it asks them to bid \ts{7}{8}{1:10}. The New York Fed's web page lists the last temporary
open market operations \ts{7}{8}{2:10}. Before the crisis the Fed was in the repo market almost every day; with excess reserves, rarely \ts{7}{8}{2:31}:
\begin{itemize}
  \item the latest is a \emph{reverse repo}: the Fed borrows money from dealers against Treasuries, agencies or MBS; dealers submitted \$4.3 billion, the Fed took
        \$2.25 billion, from Thursday the 13th to Monday the 17th, ``part of an operational readiness program'': the Fed is practising a way to drain its excess
        reserves \ts{7}{8}{2:51};
  \item for an actual \emph{repo}, which adds reserves, one must go back to August \ts{7}{8}{5:37}.
\end{itemize}
For a long time the Fed refused to call its borrowing ``reverse repos'' and called them \emph{matched sale--purchases} (MSPs): ``I'm the Fed, I print this stuff,
why do I need to provide collateral?'' Now it calls them reverses, acting like a bank \ts{7}{8}{4:56}.

\paragraph{One more level.} A Fed repo lends reserves to a dealer against a Treasury bill; the dealer has no account at the Fed, so the reserves end up at the
dealer's bank, credited to its deposit (or repaying a bank loan); if they are excess, the bank lends them in the Fed funds market. This is how open market
operations influence the Fed funds rate \ts{7}{8}{6:42}:
\begin{taccts}
\tacct[2.0cm]{Fed}{$+$Repo loan to dealer}{$+$Reserves of bank}\quad
\tacct[2.0cm]{Dealer}{$+$Deposit at bank}{$+$Repo from Fed}\quad
\tacct[2.0cm]{Dealer's bank}{$+$Reserves}{$+$Deposit of dealer}
\end{taccts}

% ------------------------------------------------------------------
\section{Interest rate spreads}\label{sec:rp-spreads}

\subsection{Before the crisis: repo $<$ Fed funds}
Stigum puzzles over why Fed funds traded above repo, and gives two reasons \ts{7}{9}{0:30}:
\begin{enumerate}
  \item \textbf{a credit spread}: repo is secured, Fed funds are not. Mehrling was never convinced, and the present reversal disproves it: nobody lends unsecured
        with any thought of default; if you fear default you cut the credit line, you do not ask a higher rate. ``Who risks a million dollars to make 137 of
        interest overnight?'' And the borrowers are banks with access to the Fed \ts{7}{9}{1:10};
  \item \textbf{segmentation}: some must borrow Fed funds and have no access to repo. But many can arbitrage, certainly any bank, by borrowing in repo and lending
        Fed funds \ts{7}{9}{2:32}.
\end{enumerate}

\paragraph{Mehrling's explanation: Bagehot's two rates.} In Bagehot's world (lecture~9) there was a \emph{market rate} of interest, made by bill brokers in the
London money market, and the official \emph{bank rate}, the Bank of England's discount rate, normally above it; nobody came to the Bank until stress pushed the market
rate up to bank rate, and then everybody flooded to it. That was monetary policy then \ts{7}{9}{2:53}. Before the crisis the US
worked the same way: the analogue of bank rate was the Fed funds target, the analogue of the market rate was repo \ts{7}{9}{4:19}. Repo below Fed funds is
\emph{discipline}: repo is the cheapest funding, so dealers repo out all their securities first; if they need more, they must go to more expensive money, borrowing
from their bank not at the Fed funds rate but at Fed funds plus 50 basis points \ts{7}{9}{4:40}.

\subsection{After the crisis: repo $>$ Fed funds}
Since January repo has moved in a band of about 16 to 30 basis points, always above Fed funds effective, sometimes by a lot \ts{7}{10}{0:09}. With the
same Bagehot logic: a market rate above the official rate means the Fed is trying to be very elastic, to pull the market rate down by pegging Fed funds below it. That
offers an arbitrage: borrow Fed funds, lend in repo, pick up about 15 basis points, and only banks (mostly) can do it \ts{7}{10}{0:51}. That
repo \emph{stays} above Fed funds despite the incentive is worrying: there are people who need liquidity and keep bidding for it, and the arbitrage is not enough to meet
the needs of the dollar funding markets, which are under stress \ts{7}{10}{1:54}. This is what the Fed would have been watching when deciding
on another QE: pegging Fed funds low was not enough, so now it buys MBS outright, \$40 billion a month, injecting reserves directly and not temporarily
\ts{7}{10}{2:36}. Will it work? Keynes said monetary policy is better at slowing the economy than at speeding it up: ``pushing on a string''
(\cref{sec:nh-managing}) \ts{7}{10}{3:37}.

\begin{finance}[Reading spreads as clues]
The method: look at patterns of interest rates and of quantities and balance sheets, and ask what must be happening behind the scenes; spreads are ``detective clues'' to
where the strains are \ts{7}{10}{4:37}. In the fall of 2007 the spread of LIBOR over Fed funds reached 100 basis points, never seen in 15 years of teaching;
knowing the two markets, it pointed to Europe \ts{7}{10}{4:58}. Prices are visible every day, balance sheets only later; ``the spreads are the key'', while the
intermediate macro textbook has just one interest rate \ts{7}{10}{5:39}. A 15 basis point reversal looks tiny, but when spreads reverse the
incentives and the flows reverse: big information in small numbers \ts{7}{10}{6:42}. Mehrling's hypothesis was formed in real time, from the FT alone, to be
tested against what happens next \ts{7}{10}{7:23}.
\end{finance}

\begin{definition}[New concepts: LIBOR]
\emph{LIBOR} (London Interbank Offered Rate) was the average rate at which a panel of banks said they could borrow unsecured from other banks in London, for each
currency and maturity; US dollar LIBOR is the Eurodollar rate (lecture~8). (After the manipulation scandal that broke in 2012, LIBOR was phased out; US dollar LIBOR
ceased in 2023, replaced by SOFR, a rate computed from overnight Treasury repo transactions.)
\end{definition}

% ------------------------------------------------------------------
\flushside
\section*{Key formulas and ideas}
\begin{keyformulas}
\begin{tabularx}{\linewidth}{@{}l L@{}}
Repo & collateralized loan organised as sale and repurchase; fails continue the loan for free\\
Balance sheets & follow the money; dealer: reverse repo (asset), repo (liability); the Fed's language is the reverse\\
Margin & always to the lender of money (hierarchy): haircut $h$, amount lent $L=(1-h)P$\\
Repo interest & $L\,r\,d/360$, simple interest, 360-day year\\
Dealers & borrow overnight (repo), lend term (reverse): banks; overnight repo $\approx$ \$1.8 tn, larger than M1\\
Liquidity spiral & haircuts up $\Rightarrow$ funding gap $\Rightarrow$ fire sales $\Rightarrow$ prices down, haircuts up\\
Fed & open market operations = repos/reverse repos with primary dealers\\
Spreads & normal: repo $<$ FF $<$ LIBOR (repo as Bagehot's market rate below bank rate); repo $>$ FF signals funding stress\\
\end{tabularx}
\end{keyformulas}

\section*{Exercises}

\begin{exercise}[Repo arithmetic]
A Treasury note with market value 102.50 per 100 face is repoed for 3 days (a weekend), haircut 2\%, repo rate 0.30\%.
\begin{enumerate}[(a)]
  \item How much is lent on \$10 million face value, and how much is repaid?
  \item What is the same interest with compound interest on a 365-day year? Compare.
\end{enumerate}
\end{exercise}

\begin{exercise}[Funding gap]
A dealer has capital 5 and holds MBS worth 100, financed by repo with a 5\% haircut.
\begin{enumerate}[(a)]
  \item Write its balance sheet.
  \item The haircut rises to 20\%. How much new funding does it need? If it must sell MBS at 90\% of value, what is its capital after selling enough to restore the
        haircut constraint?
\end{enumerate}
\end{exercise}

\begin{exercise}[Who is doing repo?]
For each statement, write the T-account entry of every party: (a) ``the Fed did \$5 billion of reverse repos with primary dealers''; (b) ``the money fund reversed in
\$1 billion of Treasuries from a dealer''; (c) ``the dealer repoed out its inventory of corporate bonds''.
\end{exercise}

\begin{exercise}[The arbitrage]
Repo trades at 0.31\%, Fed funds at 0.16\%. Describe the arbitrage a bank could do, write its T-accounts, and give two reasons why it might not close the spread.
\end{exercise}

% !TEX root = ../main.tex
\chapter{Eurodollars, Parallel Settlement}\label{ch:eurodollar}
\begin{paracol}{2}

\begin{sources}
\begin{tabularx}{\linewidth}{@{}l l L@{}}
\toprule
\textbf{Segment} & \textbf{Length} & \textbf{Content}\\
\midrule
\seg{8}{1}{L8.1 FT: Ring-fencing and the Volcker Rule} & 9:45 & Ring-fencing retail banking; matched book versus proprietary trading.\\
\seg{8}{2}{L8.2 The Eurodollar Market in Crisis} & 4:51 & LIBOR--OIS in 2007--2009.\\
\seg{8}{3}{L8.3 What are Eurodollars?} & 7:11 & Stigum's Exxon example; dollar deposits outside the US with reserves at a New York correspondent.\\
\seg{8}{4}{L8.4 Why is there a Eurodollar Market?} & 4:43 & Demand for dollars abroad; capital controls; tail and dog.\\
\seg{8}{5}{L8.5 Eurodollar as Global Funding Market} & 11:33 & Eurobanks as money dealers; the interbank market and LIBOR.\\
\seg{8}{6}{L8.6 Liquidity Challenge of Eurodollar Banks} & 10:59 & Lining up cash flows in time; a forward deposit as a swap of IOUs.\\
\seg{8}{7}{L8.7 FRA as Implicit Swap of IOUs} & 4:34 & Forward forwards and forward rate agreements, off balance sheet.\\
\seg{8}{8}{L8.8 Forward Parity, Interest Rates, EH} & 8:58 & Forward interest parity; covered interest parity via a swap of deposits.\\
\seg{8}{9}{L8.9 Forward Parity, Exchange Rates, UIP} & 2:48 & Forward rates are not expected spot rates: EH and UIP fail.\\
\seg{8}{10}{L8.10 Forward Rates are NOT Expected Spot Rates} & 2:49 & (Repeats L8.9.)\\
\bottomrule
\end{tabularx}

\smallskip
\textbf{Files.} Transcripts \file{materials/transcripts/L08-P*.txt}.
\textbf{Reading.} Stigum, \emph{The Money Market}, chapter on the Eurodollar market (from p.~212), and on FRAs.
\end{sources}

\section*{Motivation}
The third of the three wholesale dollar money markets, after Fed funds and repo \ts{8}{2}{0:07}. Eurodollars are dollar deposits at banks outside
the United States, which hold their dollar reserves not at the Fed but at a correspondent bank in New York: a \emph{parallel} settlement system, one
level below the Fed funds market in the hierarchy. The lecture explains what Eurodollars are and why they exist, how Eurobanks act as dealers in money
and line up their cash flows in time, how this leads to the first derivatives of the course, forward rate agreements, understood as implicit swaps of
IOUs, and to two parity conditions, one of which fails in practice.

% ------------------------------------------------------------------
\section{FT: ring-fencing and the Volcker rule}\label{sec:ed-ft}

\begin{casestudy}[FT, autumn 2012: ``Banks braced for proposals to ring-fence trading'']
A full-page FT feature, ``Leaner and meaner'', speculates on the business model of banking after the crisis: banks have learned some lessons, and the
regulatory landscape has yet to firm up \ts{8}{1}{0:07}\ts{8}{1}{0:28}. The article on page five refers to two ideas: the UK Vickers Commission's
recommendation that banks' retail operations be \emph{ring-fenced}, and the US \emph{Volcker rule} limiting proprietary trading \ts{8}{1}{0:48}\ts{8}{1}{1:08}.
\end{casestudy}

\paragraph{Ring-fencing.} Take the retail side of a bank, retail deposits and retail loans (households, and businesses: the distinction is retail versus
\emph{wholesale}, such as Eurodollar borrowing or derivatives), allocate to it part of the bank's capital and reserves, and fence it off. Everything else,
trading assets, wholesale borrowing, must have its own capital and reserves, so that if it loses a lot of money it can fail while the retail part is carved
off and kept safe \ts{8}{1}{1:29}\ts{8}{1}{1:50}\ts{8}{1}{2:14}\ts{8}{1}{2:34}\ts{8}{1}{2:56}\ts{8}{1}{3:19}\ts{8}{1}{3:39}.

\paragraph{The Volcker rule.} A bank may not do \emph{proprietary trading}, i.e.\ speculate on the direction of a price; it may do \emph{market making},
a matched book \ts{8}{1}{4:20}\ts{8}{1}{4:41}. Example: a pension fund holds a risky asset and wants to hedge it with a \emph{credit default swap}, a kind of
insurance, though not regulated as insurance and not a large-numbers game, that pays it if the asset fails \ts{8}{1}{5:01}\ts{8}{1}{5:22}\ts{8}{1}{5:45}\ts{8}{1}{6:06}.
The bank may sell it, but then must find another client and buy an offsetting CDS: on both sides, earning the spread, like the repo dealer of lecture~7 but with a
derivative \ts{8}{1}{6:26}\ts{8}{1}{6:47}\ts{8}{1}{7:08}. What is outlawed is the naked position, betting that the assets will fail \ts{8}{1}{7:28}\ts{8}{1}{7:48}\ts{8}{1}{8:09}.
Must the hedge be the same instrument? What matters is the same risk exposure; but hedging with something only correlated leaves \emph{basis risk}: is that
proprietary trading? If it is outlawed, banks cannot hedge, and cannot make the market at all. ``The devil is in the details'' \ts{8}{1}{8:30}\ts{8}{1}{8:51}\ts{8}{1}{9:12}\ts{8}{1}{9:32}.
\begin{taccts}
\tacct[2.0cm]{Client 1 (pension fund)}{Risky asset\\CDS bought}{}\quad
\tacct[2.0cm]{Bank (matched book)}{CDS bought from client 2}{CDS sold to client 1}\quad
\tacct[2.0cm]{Client 2}{}{CDS sold}
\end{taccts}\begin{history}[Vickers and Volcker]
The UK Independent Commission on Banking, chaired by Sir John Vickers, published its final report in September 2011, recommending ring-fencing of retail
banking; it was enacted in the Financial Services (Banking Reform) Act 2013 and in force from 2019. The Volcker rule, proposed by former Fed chairman Paul
Volcker, is section 619 of the Dodd--Frank Act (2010); the final implementing rule was adopted in December 2013.\par
\emph{References:} Independent Commission on Banking, \emph{Final Report}, Sept 2011; Dodd--Frank Act, sec.~619.
\end{history}

% ------------------------------------------------------------------
\section{The Eurodollar market in crisis}\label{sec:ed-crisis}

The chart shown plots the Eurodollar rate (LIBOR) minus OIS, the \emph{overnight index swap} rate, ``a sort of term Fed funds'', at one and three months, through the
crisis \ts{8}{2}{0:28}\ts{8}{2}{0:48}\ts{8}{2}{1:08}. Normally the spread is a few basis points, and one expects it: banks can borrow and lend in either market, and any gap is an
arbitrage \ts{8}{2}{1:50}\ts{8}{2}{2:11}. In August 2007 it jumped from almost zero to almost 100 basis points: a measure of the disorganisation of the world money market
\ts{8}{2}{2:31}\ts{8}{2}{2:51}. Mehrling was teaching the course that fall, and told the class something very bad was happening: banks without access to the Fed funds market
needed to roll their one- or three-month funding, and banks with access would not lend to them; the rate rose to offer a huge incentive to borrow Fed funds and lend at LIBOR
\ts{8}{2}{3:12}\ts{8}{2}{3:32}\ts{8}{2}{3:53}. After Bear Stearns the Fed seemed to stabilise the spread, high but steady; it lost control at Lehman, regained it, and the crisis
subsided \ts{8}{2}{3:53}\ts{8}{2}{4:13}.

\begin{definition}[New concepts: basis point, overnight index swap]
A \emph{basis point} is one hundredth of a percentage point (2\% $+$ 1 bp $=$ 2.01\%); in money markets, with sums of millions, even the fourth decimal matters
\ts{8}{2}{1:08}\ts{8}{2}{1:29}. An \emph{overnight index swap} (OIS) exchanges a fixed rate for the average of the overnight rate (in dollars, Fed funds effective) over a
term; its fixed rate is the market's term version of Fed funds. LIBOR $-$ OIS measures the premium for term unsecured interbank funding.
\end{definition}

\begin{nfigure}
\begin{tikzpicture}[font=\scriptsize,x=0.9cm,y=0.9cm]
  \draw[-{Stealth[length=3pt]}] (0,0) -- (9,0) node[right]{time};
  \draw[-{Stealth[length=3pt]}] (0,0) -- (0,4) node[above]{LIBOR $-$ OIS (bp)};
  \foreach \y/\l in {1/90,2/180,3/270} {\draw (-0.08,\y) -- (0.08,\y); \node[left] at (0,\y) {\l};}
  \draw[thick,thmcol] plot[smooth] coordinates {(0,0.05) (1.8,0.06) (2,0.9) (2.4,0.6) (2.8,1.0) (3.2,0.7) (3.6,0.85) (4.2,0.8) (4.8,0.85) (5.2,3.6) (5.6,2.4) (6.2,1.2) (7,0.4) (8.5,0.12)};
  \draw[dashed] (2,0) -- (2,3.8) node[above]{Aug 2007};
  \draw[dashed] (3.6,0) -- (3.6,3.8) node[above]{Bear Stearns};
  \draw[dashed] (5.2,0) -- (5.2,3.8) node[above]{Lehman};
\end{tikzpicture}
\caption{Schematic of the LIBOR--OIS spread through the crisis \ts{8}{2}{0:28}: near zero before August 2007, around 100 basis points after it, steadied after Bear
Stearns, a peak at Lehman, then back to normal. (Shape illustrative; the three-month spread peaked at about 365 bp in October 2008.)}\label{fig:ed-libor}
\end{nfigure}

% ------------------------------------------------------------------
\section{What are Eurodollars?}\label{sec:ed-what}

Stigum's chapter is a recent addition and ``still doesn't go far enough'' \ts{8}{3}{0:07}. Her example (p.~212): Exxon moves money from its account at Chase New
York to Citibank London, a branch outside the country \ts{8}{3}{0:28}\ts{8}{3}{0:49}\ts{8}{3}{1:09}. Behind the scenes Chase transfers reserves at the New York Fed to
Citi New York, which acts as correspondent of Citi London \ts{8}{3}{1:55}\ts{8}{3}{2:15}\ts{8}{3}{2:37}:
\begin{taccts}
\tacct[2.2cm]{Exxon}{$-$Deposit at Chase\\$+$Deposit at Citi London}{}\quad
\tacct[2.2cm]{Chase NY}{$-$Reserves}{$-$Deposit of Exxon}

\medskip
\tacct[2.2cm]{Citi NY}{$+$Reserves}{$+$Deposit of Citi London}\quad
\tacct[2.2cm]{Citi London}{$+$Deposit at Citi NY}{$+$Deposit of Exxon}
\end{taccts}
Stigum stresses that ``no money ever leaves the United States'': reserves at the New York Fed only change owner, from Chase to Citi \ts{8}{3}{3:12}\ts{8}{3}{3:32}. True,
but a \emph{deposit} leaves the country, and that is money; and the deposit of Citi London at Citi New York is, for Citi London, its \emph{reserves}: it cannot hold an
account at the Fed, but it can hold one at a bank that does \ts{8}{3}{3:53}\ts{8}{3}{4:14}\ts{8}{3}{4:34}.

\paragraph{A clearer example.} A branch makes it look like internal accounting of one bank. Better: Exxon moves its money to Crédit Lyonnais in Paris, which holds its
dollar reserves as a deposit at Citi New York, a different bank: a correspondent relationship, and Citi New York certainly considers its reserves to back \emph{its own}
deposits \ts{8}{3}{4:55}\ts{8}{3}{5:16}\ts{8}{3}{5:36}\ts{8}{3}{5:57}. Crédit Lyonnais can now go into the dollar business: make dollar loans by creating dollar deposits, and
pay withdrawals with its deposit in New York, even though it is in France \ts{8}{3}{6:18}\ts{8}{3}{6:38}\ts{8}{3}{6:59}.

\begin{definition}[New concepts: Eurodollar]\label{def:ed}
\emph{Eurodollars} are dollar-denominated deposits (and loans) at banks outside the United States, including foreign branches of US banks. Their reserves are deposits at
banks in the United States, not at the Fed. (``Euro'' because the market started in Europe; the name is used for offshore dollars anywhere, and has nothing to do with the
euro currency.)
\end{definition}\begin{history}[Origins of the Eurodollar market]
The market grew in London in the late 1950s. Commonly cited causes: dollar balances held in Europe by Soviet-bloc banks wary of keeping them in the US (the Paris bank BCEN,
whose telex address was ``Eurobank'', is often credited with the name); the 1957 sterling crisis, when the UK restricted sterling credit to non-residents and London banks
turned to dollars; US Regulation Q ceilings on deposit rates; and later US capital controls (the Interest Equalization Tax of 1963, the voluntary credit restraints of 1965).\par
\emph{References:} C.~R.~Schenk, ``The origins of the Eurodollar market in London: 1955--1963'', \emph{Explorations in Economic History} 35 (1998).
\end{history}

\section{Why is there a Eurodollar market?}\label{sec:ed-why}
There is demand: the dollar is the world reserve currency, and traders in Japan and Argentina may invoice and pay in dollars without any office in the US
\ts{8}{4}{0:07}\ts{8}{4}{0:28}. They could do it through accounts in New York, but the Fed prefers that this happens offshore: it wants to focus on domestic credit
\ts{8}{4}{0:49}. And a payment system is a credit system: deficit agents must be able to borrow from surplus agents, so people want not only dollar deposits but dollar loans
\ts{8}{4}{1:09}\ts{8}{4}{1:30}. A bank like Crédit Lyonnais wants a matched book in dollars, as many dollar assets as dollar liabilities, so that it has no currency exposure in
its euro accounts \ts{8}{4}{1:50}.

\paragraph{History and size.} The market arose because of controls on international capital flows after the Second World War: large corporations did their dollar business
with each other in London, evading US controls, and the US let them, even encouraged them; the controls broke down, then Bretton Woods, and the market kept thriving, growing
an order of magnitude larger than the Fed funds market \ts{8}{4}{2:13}\ts{8}{4}{2:33}\ts{8}{4}{2:54}.

\paragraph{Tail and dog.} Before the crisis it was unclear which was the market for dollars, Fed funds or Eurodollars: they traded at the same rate. People at the New York Fed,
and Mehrling too, worried that with so much volume offshore, and big banks able to borrow in either market, the Fed could lose control of the rate
\ts{8}{4}{3:14}\ts{8}{4}{3:55}\ts{8}{4}{4:15}. The crisis settled it: the Eurodollar market froze, the Fed funds market did not, because the Fed supported it. Fed funds are the
``real'' dollars, Eurodollars a credit extension of them, below Fed funds in the hierarchy \ts{8}{4}{3:34}\ts{8}{4}{4:36}.

% ------------------------------------------------------------------
\section{The Eurodollar market as global funding market}\label{sec:ed-funding}

The Eurodollar market is not only a payment system: it is \emph{the} world funding market, and the world funding market is a dollar market \ts{8}{5}{0:08}\ts{8}{5}{0:29}.
Ultimate lenders and borrowers anywhere want to lend and borrow in dollars, for periods of time, and the Eurodollar system intermediates the flow: an overnight deposit (a money
balance, perhaps petrodollars) funds a six-month term loan, perhaps to a Korean manufacturer \ts{8}{5}{0:50}\ts{8}{5}{1:10}\ts{8}{5}{1:30}\ts{8}{5}{1:54}.
\begin{taccts}
\tacct[2.4cm]{Eurodollar system}{Term loan (6 months)}{Overnight deposit}
\end{taccts}
This is the repo dealer's balance sheet of lecture~7, overnight funding against term lending: here a bank acts as a dealer, quoting a deposit rate and a loan rate, a bid--ask
spread. ``Banks are dealers in money'' \ts{8}{5}{2:14}\ts{8}{5}{2:34}\ts{8}{5}{2:54}. Unlike repo, Eurodollars are \emph{unsecured}: the security is the whole balance sheet of the
bank, its capital, and its backstop, the national government; unlike Fed funds, a brokered market, Eurodollar funding shows up on the balance sheets of the Eurobanks
\ts{8}{5}{3:15}\ts{8}{5}{3:36}\ts{8}{5}{3:57}. A small country borrowing large amounts is most likely funded, ultimately, in the world dollar market; borrowers and depositors outside
the US face exchange risk, a topic for later \ts{8}{5}{4:17}\ts{8}{5}{4:38}\ts{8}{5}{4:58}.

\paragraph{Outside US regulation.} By the home-country rule a bank is supervised where it resides: a French bank dealing in dollars is supervised by the Banque de France and to
some extent the ECB, not by the Fed. It can expand its balance sheet in dollars and create what is essentially money, not counted in the US money supply, since it is neither
the liability of an onshore bank nor the asset of a US resident. Yet for Exxon a deposit in Europe is ``more or less the same thing'' as one in New York: they trade at par
\ts{8}{5}{5:19}\ts{8}{5}{5:39}\ts{8}{5}{5:59}\ts{8}{5}{6:20}. The crisis chart shows the stress on that par: to keep it, Eurodollars had to pay a full percentage point more, to keep
people from ``breaking the buck'' \ts{8}{5}{6:20}\ts{8}{5}{6:41}.

\paragraph{The interbank market.} Individual banks are not matched. Say Crédit Lyonnais has many customers who want dollar deposits, more than it can lend: excess dollar
deposits; Citi London can lend any amount but has a thin European deposit base: excess dollar loans. They do business with each other through an interbank loan
\ts{8}{5}{7:01}\ts{8}{5}{7:22}\ts{8}{5}{8:08}\ts{8}{5}{8:30}\ts{8}{5}{9:13}:
\begin{taccts}
\tacct[2.4cm]{Crédit Lyonnais}{Interbank loan to Citi London}{Dollar deposits}\qquad
\tacct[2.4cm]{Citi London}{Dollar loans}{Interbank deposit from Crédit Lyonnais}
\end{taccts}
LIBOR, the London interbank offered rate, is the rate on this interbank lending \ts{8}{5}{9:47}\ts{8}{5}{10:08}. People ask why banks do all this lending to each other instead of
lending to ``real productive activity'': this is why, to move money from where it is to where it is needed \ts{8}{5}{10:28}\ts{8}{5}{10:50}. It is this market that is breaking down
in Europe now: the loans are unsecured, and a bank that is not sure it can lend on may not want to take the deposit in the first place \ts{8}{5}{11:11}\ts{8}{5}{11:31}.

% ------------------------------------------------------------------
\section{The liquidity challenge of Eurodollar banks}\label{sec:ed-liquidity}

Eurobanks hold their dollar reserves at a New York bank, not at the Fed: they cannot borrow Fed funds, nor borrow from the Fed; their backstop is their own central bank, and the
ECB ``can't print dollars''. So they face more discipline than onshore banks \ts{8}{6}{0:07}\ts{8}{6}{0:27}\ts{8}{6}{0:48}. They respond by lining up their cash flows in time much
more carefully: deposits are \emph{term} deposits to a specific date (one month, three months), loans run to specific dates (six months), and the banks match not only the
amounts but the \emph{dates}, so that when they must pay, someone is paying them \ts{8}{6}{1:09}\ts{8}{6}{1:29}\ts{8}{6}{1:50}\ts{8}{6}{2:10}. Payments do not always come as promised,
so reserves are still needed, but one can try; and this is the motivation for forward rate agreements and foreign exchange agreements: instruments to line up cash inflows and
outflows in time \ts{8}{6}{2:30}\ts{8}{6}{2:50}\ts{8}{6}{3:10}.

\paragraph{A forward deposit as a swap of IOUs.} Bank X expects to make a three-month loan two months from now (a corporate customer has announced its needs) and wants to lock in
the funding now; bank Y expects a deposit inflow in two months and wants to lock in a use for it \ts{8}{6}{4:55}\ts{8}{6}{5:15}\ts{8}{6}{5:36}\ts{8}{6}{5:56}\ts{8}{6}{6:17}. ``All banking is a swap
of IOUs'': X makes a two-month deposit at Y of \$1 million, Y makes a five-month deposit at X of \$1 million \ts{8}{6}{6:37}\ts{8}{6}{7:02}:
\begin{taccts}
\textit{$t=0$:}\quad\tacct[2.4cm]{Bank X}{Deposit at Y (2 months)}{Deposit of Y (5 months)}\quad
\tacct[2.4cm]{Bank Y}{Deposit at X (5 months)}{Deposit of X (2 months)}

\medskip
\textit{$t=2$:}\quad\tacct[2.4cm]{Bank X}{$-$Deposit at Y\\$+$3-month loan}{Deposit of Y (3 months left)}\quad
\tacct[2.4cm]{Bank Y}{Deposit at X (3 months left)}{$-$Deposit of X\\$+$Customer deposit}
\end{taccts}
No cash moves at time 0 \ts{8}{6}{7:23}\ts{8}{6}{7:44}. At two months X's deposit at Y matures: Y pays X, which uses the money to make its loan; Y pays with the deposit inflow it
expected \ts{8}{6}{7:44}\ts{8}{6}{8:04}\ts{8}{6}{8:26}\ts{8}{6}{8:47}. The swap has locked in a flow of funds from Y to X in two months, for three months; Eurodollar deposits exist to
any date, so the timing can be fine-tuned with balance sheet entries alone \ts{8}{6}{9:08}\ts{8}{6}{9:28}. If the two deposits pay the same rate over the last three months, call it
$f$, X has locked in a funding rate and Y an investment rate, even if three-month LIBOR is much lower by then \ts{8}{6}{9:49}\ts{8}{6}{10:12}\ts{8}{6}{10:32}\ts{8}{6}{10:53}.

\section{FRAs as implicit swaps of IOUs}\label{sec:ed-fra}

The literal swap takes up balance sheet space, against which regulators require capital and reserves; so it is done off balance sheet \ts{8}{7}{0:08}. Stigum describes the
evolution \ts{8}{7}{0:28}\ts{8}{7}{0:50}:
\begin{enumerate}
  \item the \emph{forward forward}: Y promises to make a deposit at X at rate $f$ two months from now;
  \item the \emph{forward rate agreement}: Y promises nothing of the sort; X funds its loan at three-month LIBOR when the time comes, and the two banks act as if there were a
        three-month deposit paying $f-\text{LIBOR}$ between them \ts{8}{7}{1:10}\ts{8}{7}{1:31}.
\end{enumerate}

\begin{definition}[New concepts: forward rate agreement]\label{def:ed-fra}
A \emph{forward rate agreement} (FRA) between two banks fixes today a rate $f$ for a future period (here from month 2 to month 5); at the start of the period they settle the
difference between the reference rate then observed (LIBOR) and $f$, on a notional amount: if LIBOR $>f$ the forward lender pays the forward borrower, if LIBOR $<f$ the
forward borrower pays \ts{8}{6}{3:31}\ts{8}{6}{3:52}\ts{8}{7}{1:55}\ts{8}{7}{2:15}.
\end{definition}

Market people describe it as a bet on LIBOR, and it can be used to speculate; but that is not why it exists \ts{8}{6}{4:13}\ts{8}{6}{4:33}. X funds at LIBOR plus or minus the FRA payment,
i.e.\ at $f$ whatever happens: the FRA implements exactly what the swap of IOUs did, locking in funding, without balance sheet space. ``It looks like a side bet on LIBOR, but it's used to
line up payments in time'' \ts{8}{7}{2:37}\ts{8}{7}{2:58}\ts{8}{7}{3:18}.

\begin{proposition}[Derivatives as implicit swaps of IOUs]
All the derivatives of the course, interest rate swaps, foreign exchange swaps, credit default swaps, can be viewed as implicit, behind-the-scenes swaps of IOUs: they are integrated in
the balance sheet structure, and are ``not really off balance sheet''. The method: ask how one would build the same payoff with a swap of IOUs on bank balance sheets; that tells what
the instrument is for \ts{8}{7}{3:18}\ts{8}{7}{3:39}\ts{8}{7}{3:59}\ts{8}{7}{4:20}.
\end{proposition}

% ------------------------------------------------------------------
\section{Forward parity}\label{sec:ed-parity}

\paragraph{Forward interest parity.} Y would like $f$ high, X low; but they do not haggle: $f$ is fixed by an arbitrage condition \ts{8}{8}{0:07}\ts{8}{8}{0:28}. There are market rates
for two-month and five-month Eurodollar deposits, $R_{0,2}$ and $R_{0,5}$. If the swapped deposits were real deposits at market rates, X would have manufactured a three-month funding
rate and Y an investment rate; the only rate both accept is the one implied by the market rates \ts{8}{8}{0:48}\ts{8}{8}{1:12}\ts{8}{8}{1:37}\ts{8}{8}{2:17}\ts{8}{8}{2:37}:
\begin{keyformulas}
\textbf{Forward interest parity} \ts{8}{8}{1:37}\ts{8}{8}{2:58}:
\[
  (1+R_{0,2})(1+f_{2,5}) = 1+R_{0,5}
\]
(with each rate expressed per its own period; in money-market convention $(1+R_{0,2}\tfrac{60}{360})(1+f_{2,5}\tfrac{90}{360}) = 1+R_{0,5}\tfrac{150}{360}$).

\medskip
\textbf{Covered interest parity} \ts{8}{8}{7:52}\ts{8}{8}{8:12}: with $S$ the spot and $F_6$ the six-month forward exchange rate, both in yen per dollar, and $R$, $R^\ast$ the six-month
dollar and yen deposit rates,
\[
  F_6 = S\,\frac{1+R^\ast}{1+R}.
\]
\end{keyformulas}
FRAs use this $f$ \ts{8}{8}{2:58}.

\paragraph{Covered interest parity.} A bank is asked for a yen loan but has only dollar funding, six-month Eurodollars: matched in time, not in currency \ts{8}{8}{2:58}\ts{8}{8}{3:21}\ts{8}{8}{3:42}.
Forward exchange contracts exist for this; but with a swap of IOUs, it swaps a six-month dollar deposit for a six-month yen deposit with another bank, at today's spot rate
\ts{8}{8}{4:02}\ts{8}{8}{4:22}\ts{8}{8}{4:43}\ts{8}{8}{5:05}. For one dollar: it receives $S$ yen (say 100) to lend; after six months it owes the counterparty $S(1+R^\ast)$ yen and is owed
$1+R$ dollars \ts{8}{8}{5:25}\ts{8}{8}{6:28}\ts{8}{8}{6:48}\ts{8}{8}{7:08}\ts{8}{8}{7:31}. Paying yen and receiving dollars at a ratio fixed today is an implicit forward exchange rate, which must equal the
market forward $F_6$: covered interest parity \ts{8}{8}{7:52}\ts{8}{8}{8:12}. Both parities are almost definitions: since anyone can ``roll their own'' forward from existing instruments, the outright
forward must have the same price, or there is an arbitrage \ts{8}{8}{8:33}\ts{8}{8}{8:53}. (In the yen--dollar example the dollar has no star, ``because the dollar is special'', the world funding
currency \ts{8}{8}{5:46}.)

\begin{taccts}
\tacct[2.4cm]{Bank (yen lender)}{Yen loan $S$\\Dollar deposit at B 1}{Dollar Eurodollar funding 1\\Yen deposit of B $S$}\qquad
\tacct[2.4cm]{Bank B}{Yen deposit at bank $S$}{Dollar deposit of bank 1}
\end{taccts}

\section{Forward rates are not expected spot rates}\label{sec:ed-eh}

Two natural guesses \ts{8}{9}{0:08}\ts{8}{9}{0:29}\ts{8}{9}{0:49}\ts{8}{9}{1:31}:
\begin{itemize}
  \item the forward interest rate $f_{2,5}$ is an unbiased forecast of three-month LIBOR two months from now: the \emph{expectations hypothesis} of the term structure
        (\cref{def:nh-eh}), $f = \E[\text{future spot rate}]$;
  \item the forward exchange rate $F_6$ is an unbiased forecast of the spot rate six months from now: \emph{uncovered interest parity} (UIP), $F_6 = \E[S_6]$, equivalently
        $\E[S_6]/S = (1+R^\ast)/(1+R)$.
\end{itemize}
Both fail empirically, all the time: a puzzle for economics \ts{8}{9}{1:10}\ts{8}{9}{1:52}. Mehrling's answer, to be developed, is that it has to do with the hierarchy of money and credit, with
liquidity: payments today are worth more than promised payments in the future, and not only because of time preference, in foreign exchange as in the term structure \ts{8}{9}{2:12}\ts{8}{9}{2:32}.

\begin{coursenote}[Segment 8.10]
The video ``Forward Rates are NOT Expected Spot Rates'' (\seg{8}{10}{L8.10}), filed in the playlist under lecture~21, repeats almost word for word the second half of L8.9.
\end{coursenote}

% ------------------------------------------------------------------
\flushside
\section*{Key formulas and ideas}
\begin{keyformulas}
\begin{tabularx}{\linewidth}{@{}l L@{}}
Eurodollars & dollar deposits outside the US; reserves $=$ deposits at a New York correspondent; below Fed funds in the hierarchy\\
Global funding & Eurobanks as money dealers: overnight/term deposits fund term loans; LIBOR $=$ interbank rate\\
Volcker rule & matched book (market making) allowed, proprietary positions not; basis risk blurs the line\\
Liquidity & no Fed access $\Rightarrow$ match amounts \emph{and} dates of cash flows\\
FRA & implicit swap of term deposits; locks in a future funding/investment rate $f$ off balance sheet\\
Forward interest parity & $(1+R_{0,2})(1+f_{2,5}) = 1+R_{0,5}$\\
Covered interest parity & $F = S(1+R^\ast)/(1+R)$\\
EH, UIP & $f=\E[\text{future rate}]$, $F=\E[S_T]$: both fail empirically (liquidity)\\
\end{tabularx}
\end{keyformulas}

\section*{Exercises}

\begin{exercise}[Forward interest parity]
Two-month Eurodollar deposits pay 0.30\% and five-month deposits 0.45\% (annualised, actual/360, 60 and 150 days). Compute the implied three-month forward rate two months ahead. If an
FRA is quoted at 0.70\%, describe the arbitrage.
\end{exercise}

\begin{exercise}[FRA settlement]
Bank X buys an FRA at $f=0.60\%$ on \$10 million for the period from month 2 to month 5 (90 days). At month 2 three-month LIBOR is 0.90\%.
\begin{enumerate}[(a)]
  \item Who pays whom, and how much (paid at the end of the period)?
  \item Show that X's all-in funding cost for its loan is 0.60\%.
\end{enumerate}
\end{exercise}

\begin{exercise}[Covered interest parity]
The spot rate is 80 yen per dollar, six-month dollar deposits pay 0.5\% and yen deposits 0.1\% (for six months).
\begin{enumerate}[(a)]
  \item Compute the six-month forward rate.
  \item Under UIP, what is the expected spot rate in six months? What does the failure of UIP mean for a bank that funds in yen and lends in dollars?
\end{enumerate}
\end{exercise}

\begin{exercise}[Exxon's transfer]
Redo Stigum's example with Crédit Lyonnais in place of Citi London, and then let Crédit Lyonnais lend \$5 million to a Korean firm that pays a US supplier banking at Chase. Write all the
T-accounts, including the Fed.
\end{exercise}

% !TEX root = ../main.tex
\chapter{The World that Bagehot Knew}\label{ch:bagehot}
\begin{paracol}{2}

\begin{sources}
\begin{tabularx}{\linewidth}{@{}l l L@{}}
\toprule
\textbf{Segment} & \textbf{Length} & \textbf{Content}\\
\midrule
\seg{9}{1}{L9.1 FT: Depreciation of Iran's Currency} & 3:37 & Official and black-market rial; sanctions on the payment system.\\
\seg{9}{2}{L9.2 Reading: John Hicks} & 3:38 & Hicks, \emph{A Market Theory of Money} (1989): banks as dealers.\\
\seg{9}{3}{L9.3 Bagehot's World, Wholesale Money Market} & 7:57 & Bills of exchange and their discount.\\
\seg{9}{4}{L9.4 Economizing on Notes: Deposits, Acceptances} & 8:28 & Discounting for deposits; acceptances as guarantees.\\
\seg{9}{5}{L9.5 Managing Cash Flow: Discount, Rediscount} & 7:29 & The discount rate as a tool of cash management; rediscount between banks.\\
\seg{9}{6}{L9.6 Market Rate of Interest} & 2:44 & Decentralised management produces one market rate, the world rate.\\
\seg{9}{7}{L9.7 Central Bank and Bank Rate} & 8:13 & The Bank of England's balance sheet; bank rate.\\
\seg{9}{8}{L9.8 The Bagehot Rule, Origin of Monetary Policy} & 4:38 & Lend freely at a high rate; from lender of last resort to monetary policy.\\
\seg{9}{9}{L9.9 Limits on Central Banking: Internal vs.\ External Drain} & 10:20 & Discipline from above: gold, cooperation among central banks.\\
\bottomrule
\end{tabularx}

\smallskip
\textbf{Files.} Transcripts \file{materials/transcripts/L09-P*.txt}.
\textbf{Readings.} Bagehot, \emph{Lombard Street} (1873), two chapters; John Hicks, \emph{A Market Theory of Money} (1989), excerpt (not among the
course files).
\end{sources}

\section*{Motivation}
So far the course has dealt with \emph{par}: payments, in which a deposit or reserves are delivered to settle a debt; the interest rate stayed in the
background \ts{9}{2}{1:51}. The next two weeks ask where the interest rate comes from and what the central bank has to do with it; after the
midterm the course moves up to the exchange rate and out to asset prices, the price of risk \ts{9}{2}{2:11}. The shift is from payments to
\emph{market making}: banks, central banks and security dealers as dealers who create liquidity by standing ready to buy and sell \ts{9}{2}{2:53}.
The lecture starts in the 19th century, where things are simpler: the London money market of Bagehot \ts{9}{3}{0:30}.

% ------------------------------------------------------------------
\section{FT: depreciation of Iran's currency}\label{sec:bg-ft}

\begin{casestudy}[FT, October 2012: ``Iran uses force to strengthen rial'']
The rial has depreciated dramatically against the dollar in Tehran's bazaar, where private currency dealers trade. The government steps in and buys dollars at
28{,}000 rials; on the illegal black market a dollar costs 30{,}000, a depreciation of about 6\% \ts{9}{1}{0:23}. The
FT attributes it to Western sanctions: EU oil sanctions shrink Iran's oil revenue, and US financial sanctions make it ``slow and expensive'' for Iran to receive
payments, which are in dollars; investors stay away from the rial and citizens dump it for foreign notes \ts{9}{1}{2:07}.
\end{casestudy}
Two points for the course \ts{9}{1}{1:46}:
\begin{enumerate}
  \item the \emph{foreign exchange dealer}, a first example of the dealers studied from now on (foreign exchange itself comes after the midterm);
  \item the cause: an attempt to keep Iran out of the world \emph{payment system}, so that it cannot get dollars.
\end{enumerate}

% ------------------------------------------------------------------
\section{Reading: John Hicks}\label{sec:bg-hicks}

\begin{reading}[Hicks, \emph{A Market Theory of Money} (1989)]
Hicks's last book, written as an old man, influenced Mehrling deeply: he read it as a young professor (he started teaching at Barnard in the fall of 1987) and has
had it on the reading list since he first taught money and banking in 1996 \ts{9}{2}{0:07}. As a young man, Hicks had sent monetary economics down what
Mehrling sees as a misleading path, with an article of 1933; in 1989 he proposed another way, a market theory of money in which banks are \emph{dealers}, which is
exactly what the course does in the next two weeks. ``I think of him as the Moses of modern economics'': he saw the promised land but did not live to enter it
\ts{9}{2}{0:49}.
\end{reading}\begin{coursenote}[Which article?]
The early article is presumably Hicks's ``A Suggestion for Simplifying the Theory of Money'' (\emph{Economica}, 1935), which treated the demand for money as a choice
among assets, an approach that led to portfolio theories of money; in class the date is given as 1933 \ts{9}{2}{0:49}. Hicks received the Nobel prize in 1972.
\end{coursenote}

% ------------------------------------------------------------------
\section{Bagehot's world: bills of exchange}\label{sec:bg-bills}

Firms in a production chain (in Britain, say, a wool producer and a clothing firm) buy and sell intermediate goods from each other: business to business, the
wholesale money market \ts{9}{3}{0:51}. Firm A buys goods from firm B on credit, issuing a \emph{bill of exchange}: a promise to
pay in 90 days that refers specifically to these goods, which will have been sold by then \ts{9}{3}{2:25}. Business-to-business
trade still works this way \ts{9}{3}{3:28}.

\begin{definition}[New concepts: bill of exchange, discount]\label{def:bg-bill}
A \emph{bill of exchange} is a written order by which a debtor (here firm A) promises to pay a sum on a given date, typically 90 days, for goods received: a form of trade
credit, transferable, ``self-liquidating'' because the sale of the goods provides the means to pay. To \emph{discount} a bill is to buy it before maturity for less than
its face value; the difference is the interest (\cref{sec:st-fed}).
\end{definition}

\paragraph{The four steps} \ts{9}{3}{3:49}--\ts{9}{3}{6:23}:
\begin{enumerate}
  \item A buys goods from B with a bill;
  \item B, needing cash, discounts the bill at a bank, which pays in bank notes, say 95 for a bill of 100 (an exaggeration: it would be closer to 100);
  \item in 90 days A sells the goods to its customers for notes;
  \item A uses the notes to redeem the bill at the bank.
\end{enumerate}
\begin{taccts}
\textit{(1)}\quad\tacct[2.0cm]{Firm A}{$+$Goods}{$+$Bill}\quad
\tacct[2.0cm]{Firm B}{$-$Goods\\$+$Bill}{}

\medskip
\textit{(2)}\quad\tacct[2.0cm]{Firm B}{$-$Bill\\$+$Notes}{}\quad
\tacct[2.0cm]{Bank}{$+$Bill\\$-$Notes}{}

\medskip
\textit{(3--4)}\quad\tacct[2.0cm]{Firm A}{$-$Goods\\$+$Notes, then $-$Notes}{$-$Bill}\quad
\tacct[2.0cm]{Bank}{$-$Bill\\$+$Notes}{}
\end{taccts}
After step 2, B has been paid and is out; A owes the bank, not B \ts{9}{3}{4:31}. This is the basic discount mechanism by which firms pay each
other with the help of a bank \ts{9}{3}{6:23}.

\paragraph{The bank's opportunity and danger.} The bank gives up notes, which pay no interest, for a bill, which does: good. But it gives up a means of payment for an asset
that turns into cash only in 90 days, hopefully: a liquidity risk. Its notes are its reserves against its deposits, which can be withdrawn or transferred at any time
\ts{9}{3}{7:04}.

\section{Economizing on notes: deposits and acceptances}\label{sec:bg-economizing}

Two alternatives to step 2 save notes \ts{9}{4}{0:09}:
\begin{enumerate}[(a)]
  \item \textbf{Discount for deposits.} If B accepts a deposit instead of notes, the bank expands its balance sheet and keeps its notes. Still a liquidity risk: B discounted
        because it wants to spend, and if it pays a customer of the same bank, fine, but if it pays someone in France the bank must find notes, and behind them gold
        \ts{9}{4}{0:53}.
  \item \textbf{Acceptance.} B may prefer to hold the bill to maturity and earn the full 100 rather than 95, but wants to be able to sell it if needed. So it has the bank
        \emph{endorse} the bill: the bank promises to pay if A does not. Then the bill is ``a 90-day bill on this bank'', almost a time deposit, and easy to sell
        \ts{9}{4}{2:57}.
\end{enumerate}
\begin{taccts}
\textit{(a)}\quad\tacct[2.0cm]{Firm B}{$-$Bill\\$+$Deposit}{}\quad
\tacct[2.0cm]{Bank}{$+$Bill}{$+$Deposit of B}

\medskip
\textit{(b)}\quad\tacct[2.3cm]{Firm B}{Bill\\$+$Acceptance\\\quad(contingent)}{}\quad
\tacct[2.3cm]{Bank}{}{$+$Acceptance\\\quad(contingent)}
\end{taccts}

\begin{definition}[New concepts: acceptance, contingent liability]
A \emph{banker's acceptance} is a bank's guarantee written on a bill (physically, the word ``accepted'' and the banker's signature across it): the bank will pay if the drawee does
not. It is a \emph{contingent} liability of the bank and a contingent asset of the holder: it becomes an actual payment only if the event (the drawee's failure to pay)
occurs \ts{9}{4}{4:22}.
\end{definition}

Mehrling writes the acceptance as a separate balance sheet entry, because in the modern world it is a separate instrument: ``an early form of a credit default swap: if they
don't pay, I'll pay'' \ts{9}{4}{5:26}. The bank charges a fee, like an insurance premium, smaller than the discount, say 1\% instead of 5\%, since it is a guarantee and
not a 90-day investment; if A pays, the bank has ``money for nothing'' \ts{9}{4}{6:50}. Bankers know their firms; the typical problem is not outright default
but \emph{delay}: the goods are not sold in time, so payment is late, a liquidity issue \ts{9}{4}{7:52}.

% ------------------------------------------------------------------
\section{Managing cash flow: discount and rediscount}\label{sec:bg-cashflow}

Focus on the bank: in step 2 bills come in and notes go out, in step 4 notes come in and bills go out. ``That's the banking business'' \ts{9}{5}{0:07}. The bank holds a
\emph{ladder} of bills maturing on known dates (a 90-day bill today is an 89-day bill tomorrow), so it knows the timing of its inflows; on the other side, demand deposits that
could all go tomorrow, but whose patterns it knows from experience \ts{9}{5}{0:49}. The banker's business is matching inflows and outflows: use the notes that
come in today to discount new bills, hold as little idle cash as possible \ts{9}{5}{1:52}. Two tools:
\begin{enumerate}
  \item \textbf{The discount rate.} Short of notes, a bank gives customers ``a bad price'': it raises the rate at which it discounts, which is the same as lowering the price at which
        it buys bills. Discounting slows, maturing bills bring notes back, and later it lowers the rate again. Raising the rate shifts the balance sheet from bills into cash, lowering
        it from cash into bills \ts{9}{5}{2:33}.
  \item \textbf{Rediscount at another bank.} Instead of sending the customer down the street, the bank itself takes bills it already holds to another bank: an interbank money market
        \ts{9}{5}{5:20}.
\end{enumerate}
So each bank quotes two prices, a rate to discount for customers and a rate to rediscount with other banks: a buy--sell spread. Banks are money dealers, and the prices at which they
deal are the discount and rediscount rates \ts{9}{5}{6:46}.

\begin{intuition}[Price as a valve on flows]
Raising the discount rate does not change the bank's stock of notes at once; it changes the \emph{rate of flow}: fewer bills in, the same maturities out, so the stock of notes
recovers over the following days. The interest rate here is a control on the time derivative of reserves, set by a bank watching its own cash position.
\end{intuition}

\section{The market rate of interest}\label{sec:bg-market}

Each bank looks only at its own cash flow and moves its discount rate accordingly; the effect is \emph{coordination}. Whoever quotes the lowest rate gets all the discounts, so
rates line up, some a bit off the market in one direction or the other to influence flows: a \emph{market rate of interest} emerges \ts{9}{6}{0:07}. It moves
over time with firms' demand for discounts and the banks' decisions on deposits, acceptances and how much cash to hold: all these decisions show up in one money market rate,
uniform across the economy \ts{9}{6}{1:10}. Lombard Street was the money market not only for wool merchants but for the world: firms A and B might be in two different
countries, discounting their bills in London; the London market rate was the world rate of interest, and payment imbalances anywhere showed up in it
\ts{9}{6}{1:51}.

% ------------------------------------------------------------------
\section{The central bank and bank rate}\label{sec:bg-bankrate}

The balance sheet of the Bank of England of 9 January 1924 (after Bagehot's time, but with the structure he describes) is divided, since Peel's Act of 1844, into an Issue
Department and a Banking Department \ts{9}{7}{0:07}:
\begin{taccts}
\tacct[2.8cm]{Issue Department}{Gold coin and bullion\\Government debt (fiduciary issue)}{Notes issued}\qquad
\tacct[2.8cm]{Banking Department}{Notes (reserve)\\Government securities\\Other securities (discounts, bills)}{Deposits (public, bankers, other)}
\end{taccts}
\begin{itemize}
  \item The \textbf{Issue Department} issues the currency, mostly against gold, plus a small fiduciary issue against government debt; at the margin it can issue more notes only
        against more gold, and must pay gold to foreigners who present notes \ts{9}{7}{0:49}.
  \item The \textbf{Banking Department} holds part of the notes as its reserve (the rest circulate), government securities and other securities, discounts of trade bills; its
        liabilities are deposits, not only of banks but of firms and people: unlike the Fed, the Bank was a bank for everyone \ts{9}{7}{1:30}.
\end{itemize}
The Banking Department is in the same bill market as other banks: it discounts for the market and rediscounts for other banks \ts{9}{7}{2:31}. It quotes \emph{bank rate},
the rate at which a banker who needs notes can get them from the Bank \ts{9}{7}{3:12}.

\begin{definition}[New concepts: bank rate]
\emph{Bank rate} was the Bank of England's published minimum rate for discounting eligible bills: the central bank's rediscount rate, the price at which the banking system could
obtain the best money of the system.
\end{definition}

The Bank was privately owned and paid dividends, but it held the reserve of the whole British banking system, which was at the same time the reserve of the international money
market \ts{9}{7}{3:53}. If the banking system needs money, it discounts bills at the Bank for notes, or, in a real crisis when the Bank runs out of notes, for deposits
\ts{9}{7}{4:56}. It could behave like any bank, raising or lowering its rate to manage its own cash flow; Bagehot says it does not, because as keeper of the reserve
it has public duties, whether it admits it or not \ts{9}{7}{6:03}. Usually the market rate is below bank rate and nobody goes to the Bank; when everyone scrambles
for cash the market rate rises to bank rate, and then everyone goes to the Bank, which ``disgorges'' its notes and expands its deposits, which banks treat as cash, although deposits
are promises to pay notes and notes promises to pay gold \ts{9}{7}{7:05}.

\begin{nfigure}
\begin{tikzpicture}[font=\scriptsize,x=0.8cm,y=0.8cm]
  \draw[-{Stealth[length=3pt]}] (0,0) -- (10,0) node[right]{time};
  \draw[-{Stealth[length=3pt]}] (0,0) -- (0,4) node[above]{rate};
  \draw[very thick,thmcol] (0,3) -- (9.5,3) node[right]{bank rate};
  \draw[thick,defcol] plot[smooth] coordinates {(0,1.5) (1,1.8) (2,1.4) (3,2.0) (4,2.4) (5,3.0) (5.6,3.0) (6.2,2.6) (7,1.9) (8,1.6) (9.5,1.8)};
  \node[defcol,right] at (9.5,1.8) {market rate};
  \fill[keycol!25] (4.95,0) rectangle (5.65,3);
  \node[align=center,keycol] at (5.3,3.7) {everyone goes\\to the Bank};
\end{tikzpicture}
\caption{Bagehot's two rates \ts{9}{7}{7:05}: normally the market rate is below bank rate and nobody borrows from the Bank; in a scramble for cash the market rate rises to bank rate,
which acts as a ceiling, and the Bank lends freely there.}\label{fig:bg-rates}
\end{nfigure}

% ------------------------------------------------------------------
\section{The Bagehot rule and the origin of monetary policy}\label{sec:bg-rule}

With bank rate the Bank of England is engaged in monetary policy, whether it knows it or not: ``this is the origin of monetary policy'' \ts{9}{8}{0:08}. In its minimal
form: set a maximum rate, and when the crisis comes act passively, lending to all comers at that high rate \ts{9}{8}{0:28}.

\begin{proposition}[The Bagehot rule]\label{prop:bg-rule}
In a crisis the central bank should \emph{lend freely, at a high rate, against good security} \ts{9}{8}{0:48}. That is all Bagehot says in \emph{Lombard Street}; he does not say to move
bank rate around to influence the economy \ts{9}{8}{0:48}.
\end{proposition}

\paragraph{Elasticity and discipline in one phrase.} ``Lend freely'' is elasticity: otherwise borrowers who cannot make their payments default, and liquidity kills you quick. ``At a high
rate'' is discipline: only those who really need it come, and they repay fast, borrowing in the market again as soon as the market rate falls; the lending is self-liquidating
\ts{9}{8}{3:13}.

\paragraph{From backstop to policy.} The market knows the Bank is waiting in the wings and behaves accordingly, taking more risk, lending more; this gives the Bank influence over the
economy, and once it realises it, it can move bank rate up and down rather than leave it at, say, 6\%: lower it to encourage credit growth, raise it to discourage it (``I'm not going to help
you, you're on your own''). ``That's the art of central banking'' \ts{9}{8}{1:09}. Sitting at the top of the hierarchy, the central bank can manage it,
moving bank rate, and be lender of last resort in trouble; but it can push rates around with its balance sheet only if it lets the market move onto its balance sheet: it can relax tension or
cause it \ts{9}{8}{2:31}.

% ------------------------------------------------------------------
\section{Limits on central banking: internal and external drain}\label{sec:bg-limits}

\paragraph{Looking down: internal drain.} If notes drain out of the banks into the hands of firms, the Bank can ease the drain by releasing its notes and creating good substitutes for them,
or choose not to: a policy choice, without constraint \ts{9}{9}{0:09}.

\paragraph{Looking up: external drain.} If foreigners want notes, and ultimately gold, the Bank must give them gold. London may be the centre of the world money market, but the Bank of
England is one central bank among others, and if they will not accept deposits at the Bank as equivalent to notes and demand gold, it must pay: that is the gold standard
\ts{9}{9}{1:16}. The Bank faces the same constraint as any bank, measured in gold in and gold out for the whole country \ts{9}{9}{2:18}. And it uses the same tool: it
raises its discount rate to say ``don't come to me to discount and then take the gold to France'' \ts{9}{9}{2:38}. If it does not, it can lose all its gold and have to
suspend payments \ts{9}{9}{3:20}.

\begin{finance}[Drawing gold from the moon]
A saying of the time held that a high enough bank rate would ``draw gold from the moon''. The mechanism is a contraction of credit: at a very high rate no new bills come for discount,
maturing bills are repaid, in gold, and no gold flows out again. The Bank secures its own position and maintains the gold parity by contracting credit \ts{9}{9}{4:23}.
\end{finance}

So the Bank of England is a source of elasticity looking down, and faces discipline looking up \ts{9}{9}{5:24}. Young stresses how little gold ran the whole world monetary system: only the
gold in the Banking Department moved, the rest was locked behind the notes; this was possible because of the sophistication of the money market, where gold could be borrowed, also from other
central banks \ts{9}{9}{5:48}.

\begin{intuition}[Cooperation and its limits]
Central banks willing to borrow and lend gold among themselves could jointly relax their survival constraint, as banks do in the Fed funds market. But central bank cooperation, ``easier
said than done'', involves national interests, politicians and finance ministers. Without it the survival constraint binds on central banks, which then have no choice but to raise rates,
even if that crunches all the banks below them: discipline is transmitted down the hierarchy. ``Does this sound familiar?'' \ts{9}{9}{6:49}.
\end{intuition}

\paragraph{The money market as coordinator.} The central bank is in between, not at the absolute top: it can relax its own constraint with other central banks, and impose discipline below
only if it has discipline imposed on itself, or chooses to \ts{9}{9}{8:13}. This is the world Bagehot knew, ``not so different from our own'', and the origin of all later thinking about central
banks, before IS/LM \ts{9}{9}{8:34}. The money market rate was a reliable indicator of stress, watched by everyone, coordinating the economy and, because London was the world money
market, the world: ``enough to make you think the money market is god''. Hence \emph{Lombard Street} was a bestseller \ts{9}{9}{9:14}.

\begin{history}[Bank rate in the 19th century]
Bank rate was the Bank of England's minimum discount rate; from the 1850s, as the Bank came to manage the gold reserve actively, it was changed often, and rates of 6--10\% were used in crises
(10\% in 1857 and 1866, 9\% in 1873). Peel's Act of 1844 was suspended in 1847, 1857 and 1866, allowing the Bank to issue notes beyond the legal limit. Bagehot's \emph{Lombard Street} (1873) appeared
in the year of a further crisis.\par
\emph{References:} R.~S.~Sayers, \emph{The Bank of England 1891--1944}, Cambridge UP, 1976; Bagehot, \emph{Lombard Street}, 1873.
\end{history}

% ------------------------------------------------------------------
\flushside
\section*{Key formulas and ideas}
\begin{keyformulas}
\begin{tabularx}{\linewidth}{@{}l L@{}}
Bill of exchange & 90-day trade credit tied to goods; discounted at a bank for notes: bank gains interest, takes liquidity risk\\
Economizing on notes & discount for deposits; acceptance $=$ contingent guarantee (early CDS)\\
Discount rate & raising it slows inflow of bills, notes accumulate: a cash-flow management tool\\
Banks as money dealers & discount rate (to customers) and rediscount rate (between banks): a bid--ask spread\\
Market rate & emerges from decentralised cash management; London rate $=$ world rate\\
Bank rate & central bank's rediscount rate; normally above the market rate, a ceiling in crises\\
Bagehot rule & lend freely (elasticity) at a high rate (discipline) against good security\\
Drains & internal: policy choice; external (gold): discipline from above, relaxed only by central bank cooperation\\
\end{tabularx}
\end{keyformulas}

\section*{Exercises}

\begin{exercise}[Discount arithmetic]
A 90-day bill of 100 is discounted at a discount rate of 4\% a year (bank discount basis, 360 days).
\begin{enumerate}[(a)]
  \item What does firm B receive?
  \item What is the bank's yield on the money it lays out, expressed as simple interest on a 365-day year?
\end{enumerate}
\end{exercise}

\begin{exercise}[A bank's ladder]
A bank holds bills maturing at 10 a day for the next 90 days and notes of 50; its depositors withdraw on net 8 a day.
\begin{enumerate}[(a)]
  \item How many new bills can it discount each day without running down its notes?
  \item A large withdrawal of 40 occurs today. Describe how it can use its discount rate and the rediscount market to restore its position.
\end{enumerate}
\end{exercise}

\begin{exercise}[The two rates]
Explain, with the T-accounts of a bank and of the Bank of England, what happens when the market rate reaches bank rate and a bank rediscounts 20 of bills at the Bank, first when it is paid in
notes and then when it is paid in a deposit at the Bank.
\end{exercise}

\begin{exercise}[External drain]
Foreign central banks present notes worth 10 to the Issue Department for gold. Write the T-accounts of both departments. What can the Banking Department do with bank rate, and with what
consequences for British banks?
\end{exercise}

% !TEX root = ../main.tex
\chapter{Dealers and Liquid Security Markets}\label{ch:dealers}
\begin{paracol}{2}

\begin{sources}
\begin{tabularx}{\linewidth}{@{}l l L@{}}
\toprule
\textbf{Segment} & \textbf{Length} & \textbf{Content}\\
\midrule
\seg{10}{1}{L10.1 FT: Asymmetric Credit Growth in Europe} & 6:33 & IMF: European banks deleveraging; deposits flee from periphery to core.\\
\seg{10}{2}{L10.2 Market Liquidity, Dealers, and Inventories} & 7:10 & Market versus funding liquidity; the supermarket and its inventories.\\
\seg{10}{3}{L10.3 Two-sided Dealer Basics} & 6:31 & Two inventories; dealers as smoothers in price and time.\\
\seg{10}{4}{L10.4 Economics of the Dealer Function: the Treynor Model} & 11:43 & Position limits, outside spread, inside spread.\\
\seg{10}{5}{L10.5 Leveraged Dealer Basics} & 7:22 & Dealing with borrowed money; the Volcker rule and matched book.\\
\seg{10}{6}{L10.6 Real World Dealers} & 7:56 & Gross versus net exposure; inventories ``out in the world''.\\
\seg{10}{7}{L10.7 Arbitrage and the Assumption of Perfect Liquidity} & 9:41 & Dealers move liquidity across markets; perfect liquidity as a limit case; Fischer Black's ``Noise''.\\
\bottomrule
\end{tabularx}

\smallskip
\textbf{Files.} Transcripts \file{materials/transcripts/L10-P*.txt}; Mehrling's notes \mn{10}{1}.
\textbf{Reading.} Jack Treynor, ``The Economics of the Dealer Function'', \emph{Financial Analysts Journal} (1987) (not among the course files); Fischer Black, ``Noise'' (1986) (\file{materials/readings/L10-Noise.pdf}); IMF GFSR Box~2.3 and the FT article (\file{materials/readings/L10-*.pdf}).
Recommended: Larry Harris, \emph{Trading and Exchanges: Market Microstructure for Practitioners} (2003).
\end{sources}

\section*{Motivation}
Lecture~9 showed banks as dealers in money in Bagehot's world. This lecture builds the general model of a dealer, Treynor's, which the course uses from now on:
for security dealers here, for banks as dealers in money in lecture~11, for foreign exchange and swaps later. The question it answers: where does
\emph{market liquidity} come from? Answer: from dealers, who supply it for a profit, and in doing so move prices away from fundamental values. Perfect liquidity,
the assumption behind most of finance and general equilibrium theory, is shown to be a logically impossible limit.

% ------------------------------------------------------------------
\section{FT: asymmetric credit growth in Europe}\label{sec:dl-ft}

\begin{casestudy}[FT, 9 October 2012: ``IMF warns eurozone on capital flight'']
The IMF's \emph{Global Financial Stability Report} warns that unless eurozone officials step up their response, European banks would dump \$2.8 trillion of assets, more
than 7\% of their balance sheets, by the end of next year; banks in the periphery about 10\% \ts{10}{1}{0:07}. Mehrling notes that the FT mislabelled one of the
IMF's charts, ``an F for accuracy'': pay attention to details \ts{10}{1}{0:48}.
The article (by Claire Jones and Michael Steen) adds three points:
\begin{itemize}
  \item The IMF calls the result ``extreme fragmentation'' of euro-area funding markets. Its estimates are based on the 58 largest EU banks.
  \item With additional national measures, deleveraging would be \$2.3 trillion, 6\% of balance sheets.
  \item Draghi: ``You can't have a union when you have certain countries that are permanent creditors and a set of countries that are permanent debtors.''
\end{itemize}
The chart is from Box~2.3 of the IMF \emph{Global Financial Stability Report}, October 2012. Its Figure~2.3.3 shows euro-area bank credit to the non-bank private sector (Dec 2009 = 100) rising in the core and falling in the periphery.
\end{casestudy}
Depositors move money from banks in the periphery (Spain, Italy, Portugal, Ireland) to banks in the core, fearing bankruptcy or a break-up of the euro, in which
they would end up with a liability of, say, the Irish rather than the German central bank \ts{10}{1}{1:28}. This involves movements of reserves
among central banks (TARGET2 balances, left for later) \ts{10}{1}{2:29}. For the periphery banks, losing deposits means they must find other funding, impossible, or
shrink assets: they shed securities and loans and try to raise capital, and they make no new loans \ts{10}{1}{2:50}.
\begin{taccts}
\tacct[2.3cm]{Periphery bank}{$-$Securities\\$-$Loans}{$-$Deposits\\$+$Capital}\qquad
\tacct[2.3cm]{Core bank}{$+$Reserves / securities}{$+$Deposits}
\end{taccts}
For European banks as a whole the leverage ratio, total liabilities to capital, fell from about 25 to 23 (banks have 3--4\% capital) \ts{10}{1}{3:32}. The average
hides the asymmetry: bank credit to the private sector shrinks in the periphery and keeps growing in the core \ts{10}{1}{4:37}. The IMF adds that the ECB's
pledge to buy government debt lowered sovereign yields, but it is too early to tell whether it relieves deleveraging pressure \ts{10}{1}{6:01}. The question of the lecture:
how can a mere \emph{pledge} to buy lower bond yields? \ts{10}{1}{6:01}

% ------------------------------------------------------------------
\section{Market liquidity and inventories}\label{sec:dl-liquidity}

\begin{definition}[New concepts: funding liquidity, market liquidity]\label{def:dl-liq}
\emph{Funding liquidity} is the ability to raise money, to replace funds, the problem of the periphery banks \ts{10}{2}{0:28}. \emph{Market liquidity} is the ability to buy or
sell an asset quickly, in volume, without moving the price much; it is a property of the market for a particular asset \ts{10}{2}{0:49}. In Hicks's terms,
a liquid market has a \emph{continuous} price: it may fluctuate, but without jumps or ``air pockets''; an illiquid market trades in fits and starts, at discontinuous prices
\ts{10}{2}{1:51}.
\end{definition}

Supply and demand say nothing about the market structure behind such continuity: the price simply moves from one equilibrium to another, an abstraction from the
institutions that make it continuous, the dealers \ts{10}{2}{3:01}.

\paragraph{The supermarket.} Intuition from a dealer we know, the West Side Market: in the early morning it is empty and the shelves full; in the evening ``all of New York''
is in there and the shelves are depleted; the prices are the same morning and evening \ts{10}{2}{3:42}. Supply and demand are not lined up in time, yet the price
does not move. The secret is \emph{inventories}: the shelves absorb the flow of demand and are restocked at night, from inventories in New Jersey, the wholesale market, which in
turn are restocked from farms: ``dare I say a hierarchy of inventories'' \ts{10}{2}{4:43}. The supermarket is a \emph{one-sided} dealer: it sells
retail but does not buy from you (come in with a box of mangoes and they will ask whether you stole it); it buys wholesale. An investment bank underwriting an initial public
offering is a one-sided security dealer: it buys the whole issue and distributes it, holding inventory meanwhile \ts{10}{2}{6:07}. Even Japanese ``just-in-time'' production
does not eliminate inventories, the notes observe: it pushes them back before production, into the parts and components held by suppliers \mn{10}{2}.

% ------------------------------------------------------------------
\section{Two-sided dealers}\label{sec:dl-twosided}

A two-sided dealer trades a security for money and money for the security, either way; so it needs two inventories, of cash and of securities. A simple (hypothetical)
dealer with fixed capital allocates it between the two and lets them fluctuate: when the public sells securities to it, its cash falls and its securities rise, and vice versa
\ts{10}{3}{0:28}.
\begin{taccts}
\tacct[2.5cm]{Unleveraged dealer}{Cash\\Securities}{Capital}
\end{taccts}
With large enough inventories the price would not move at all \ts{10}{3}{2:12}.

\paragraph{Life without a dealer.} A buyer must find a seller: either raise the price offered (``I know it's worth 10, but I'll pay 12'') or wait. Time and price are the two
costs \ts{10}{3}{2:33}. The housing market is the example: a broker shows you houses, and you wait; knocking on the door of an owner who is not ready to
sell, you overpay; needing to sell quickly, you cut the price \ts{10}{3}{3:56}. Dealers are \emph{smoothers}: they can flatten the price and reduce the
waiting time to nothing, because they are always ready to take the other side \ts{10}{3}{3:36}. They do it to make money, so we need the economics of the dealer
function \ts{10}{3}{4:58}.

\begin{intuition}[A missing piece of microeconomics]
Price theory, and general equilibrium, assume perfect liquidity: a market-clearing price at every moment. This assumption must be false: dealers supply market liquidity, and not
for free; if liquidity were a free good, infinite in quantity, no dealer would supply it. ``The abstraction we have in price theory is an abstraction that's logically impossible for
liquid markets'' \ts{10}{3}{5:19}.
\end{intuition}

% ------------------------------------------------------------------
\section{The economics of the dealer function: Treynor's model}\label{sec:dl-treynor}

Treynor's diagram has the dealer's \emph{inventory} on the horizontal axis and the \emph{price} on the vertical axis (\cref{fig:dl-treynor}). It has three elements.

\begin{nfigure}
\begin{tikzpicture}[font=\scriptsize,x=1cm,y=1cm]
  \draw[-{Stealth[length=4pt]}] (-4.6,0) -- (4.8,0) node[right]{inventory};
  \draw[-{Stealth[length=4pt]}] (0,-0.2) -- (0,5.2) node[above]{price};
  \draw[dashed] (-4,0) -- (-4,4.9); \draw[dashed] (4,0) -- (4,4.9);
  \node[below,align=center] at (-4,0) {dealer's\\maximum short};
  \node[below,align=center] at (4,0) {dealer's\\maximum long};
  \node[below] at (0,0) {0};
  % outside spread
  \node[defcol,rotate=-22] at (-2.0,4.05) {dealer's ask};
  \node[intcol,rotate=-22] at (-1.4,2.65) {dealer's bid};
  \draw[very thick,thmcol] (-4.3,4.4) -- (4.3,4.4) node[right,align=left]{value-based\\ask};
  \draw[very thick,thmcol] (-4.3,0.8) -- (4.3,0.8) node[right,align=left]{value-based\\bid};
  % inside spread: ask and bid lines sloping down, reaching the outside spread at the limits
  \draw[very thick,defcol] (-4,4.4) -- (4,1.15);
  \draw[very thick,intcol] (-4,4.05) -- (4,0.8);
  \draw[dotted] (0,2.6) -- (-0.15,2.6) node[left]{$p^\ast$};
  \draw[{Stealth[length=3pt]}-{Stealth[length=3pt]}] (1.5,1.74) -- (1.5,2.13);
  \node[right] at (1.55,1.95) {inside spread};
  \draw[{Stealth[length=3pt]}-{Stealth[length=3pt]}] (-4.55,0.8) -- (-4.55,4.4);
  \node[left,align=right] at (-4.6,2.6) {outside\\spread};
  \node[align=center,font=\scriptsize\itshape] at (2.6,3.6) {market liquidity\\(dealer)};
  \node[align=center,font=\scriptsize\itshape] at (0,-1.35) {funding liquidity: how far from 0\\the dealer can finance its position};
\end{tikzpicture}
\caption{Treynor's model of a dealer \ts{10}{4}{0:13}. The dealer quotes an ask and a bid (inside spread) that fall as its inventory grows; the value-based traders' prices (outside
spread) bound the market; the position limits come from the dealer's funding. $p^\ast$, the midpoint at zero inventory, can be read as fundamental value
\ts{10}{7}{4:07}.}\label{fig:dl-treynor}
\end{nfigure}

\begin{enumerate}
  \item \textbf{Position limits.} The dealer's maximum short and maximum long positions. A long position gains if the price rises, a short position gains if it falls; at zero
        inventory the dealer is not exposed \ts{10}{4}{0:36}. The limits come from capital, but more generally from the ability to \emph{borrow} to fund the inventory: a
        banker who will lend up to some amount and no more, or the dealer's own risk appetite. They are soft in reality, hard in the model, and symmetric only for convenience
        \ts{10}{4}{1:21}. This is \emph{funding liquidity}.
  \item \textbf{The outside spread.} Behind the scenes are the fundamental forces of supply and demand and deep pockets: the \emph{value-based traders}, ``Warren Buffett in the
        wings'', who buy when the price is low enough (the value-based bid) and sell or short when it is high enough (the value-based ask) \ts{10}{4}{3:07}. They are
        the dealer's ``liquidity provider of last resort'': a dealer near its long limit can unload its inventory to them. Knowing they are there, the dealer is willing to let
        inventory build up, and so to supply liquidity \ts{10}{4}{4:08}. In a crisis the value-based bid goes to zero; dealers stop making markets
        and prices are in free fall \ts{10}{4}{5:11}.
  \item \textbf{The inside spread.} The dealer's own quotes, bid and ask, are what you see in the market when your broker quotes you Apple stock; the outside spread you see only if
        inventories push the price to it \ts{10}{4}{5:51}. Two things determine it \ts{10}{4}{6:34}:
        \begin{enumerate}[(a)]
          \item its \emph{width}: the volatility of the price; holding inventory of a volatile asset is risky, and a wider spread protects the dealer \ts{10}{4}{6:55};
          \item its \emph{slope}: adverse selection; the dealer trades with anyone, and the counterparty may know more. A growing inventory means many people are selling to it: the
                order flow is information, ``maybe they know something I don't know'', so the dealer buys more only at a better price, and lowers both quotes
                \ts{10}{4}{7:38}.
        \end{enumerate}
\end{enumerate}

\begin{definition}[New concepts: inside and outside spread, adverse selection]
The \emph{inside} (dealer) spread is the difference between the dealer's ask and bid; the \emph{outside} spread is the difference between the prices at which value-based investors
will sell and buy. \emph{Adverse selection} is the risk of trading with a better-informed counterparty: the trades that happen are disproportionately those that are bad for you.
\end{definition}

\begin{proposition}[Dealers turn funding liquidity into market liquidity]
Security dealers demand funding liquidity, borrowing to finance their inventories, and supply market liquidity, quoting prices at which anyone can trade \ts{10}{4}{10:02}.
\end{proposition}

The notes describe the liquidity the dealer supplies as an \emph{option to trade}: at the quoted prices, either way, if you want to, but it is up to you. Anyone who places a limit order
supplies liquidity in this sense, and amateur day traders often do, ``until they run out of money'' (the concept is developed in Larry Harris, \emph{Trading and Exchanges}, 2003)
\mn{10}{4}. In the notes' Treynor model the value-based traders' outside spread may be ``maybe 20\% below and 20\% above'' their estimate of fundamental value; when the dealer hits his
position limits, the value-based trader becomes market maker of last resort \mn{10}{4}.

\begin{keyformulas}
\textbf{A linear Treynor dealer} (a simple parametrisation of \cref{fig:dl-treynor}): with inventory $I\in[-I_{\max},I_{\max}]$,
\[
  \text{ask}(I) = p^\ast + \frac{s}{2} - \lambda I ,\qquad \text{bid}(I) = p^\ast - \frac{s}{2} - \lambda I ,
\]
with inside spread $s$ (rising with volatility) and slope $\lambda$ (rising with adverse selection), chosen so that the quotes reach the value-based ask and bid at the position
limits: $\lambda I_{\max} + s/2 = $ half the outside spread.
\end{keyformulas}

\begin{exercise}[The pledge]
Use the Treynor diagram to explain how a central bank's credible pledge to buy a government bond at a given price can lower its yield even before any purchase: which element of the diagram does the
pledge change?
\end{exercise}

% ------------------------------------------------------------------
\section{Leveraged dealers and the Volcker rule}\label{sec:dl-leverage}

Dealing is a very competitive business, with pressure to quote narrow spreads; one way to beat the competition is to trade on less capital, with borrowed money, so that the profit on
capital is multiplied by leverage \ts{10}{5}{0:28}. Leverage is the ratio of debt to equity; deleveraging, as in the FT story, reduces it \ts{10}{5}{1:10}.
Capital-funded dealers exist, but they are the value-based traders: ask Warren Buffett what his balance sheet looks like, ``billions of dollars of cash waiting to buy a railroad''; the
inside spread is made by leveraged dealers \ts{10}{5}{1:52}. In the extreme, fully borrowed: the ``inventory of cash'' is the banker who lends \ts{10}{5}{2:33}.
\begin{taccts}
\textit{long:}\quad\tacct[2.0cm]{Leveraged dealer}{Securities}{Loan from bank}\qquad
\textit{short:}\quad\tacct[2.0cm]{Leveraged dealer}{Cash}{Securities (owed)}
\end{taccts}
Making markets now means expanding and contracting the balance sheet itself, with only a thin layer of capital to absorb price changes \ts{10}{5}{3:14}. Mehrling knows
of no regulatory leverage limit for dealers; they are usually part of larger organisations with more capital \ts{10}{5}{4:16}.

\paragraph{The Volcker rule again.} It says that banks cannot have subsidiaries that take positions: they may deal, but with a \emph{matched book}, net inventory zero. The balance sheet can
be very large, long and short positions of equal size, but there is no exposure to the price \ts{10}{5}{4:57}. The rule comes from the view that banks, or their
dealer subsidiaries, took proprietary risk that, with FDIC insurance, fell on the taxpayer \ts{10}{5}{6:18}. It is about exposure to price risk, not leverage: ``a matched-book dealer can be
infinitely leveraged'' without price risk, and will try to persuade regulators it needs no capital. But there is no perfect hedge, and matched books carry other risks, liquidity risk among
them \ts{10}{5}{5:58}.

% ------------------------------------------------------------------
\section{Real-world dealers}\label{sec:dl-real}

Real dealers hedge: long some securities and short others that are correlated \ts{10}{6}{0:07}. Their balance sheet is the repo dealer's of lecture~7: reverse repo brings securities in (and
they may be sold short), repo sends them out as collateral for borrowing cash \ts{10}{6}{0:27}.

\begin{definition}[New concepts: gross and net exposure]
\emph{Gross exposure} is the size of the balance sheet; \emph{net exposure} is long positions minus short positions. A matched book can have a huge gross and zero net exposure; a small
one-sided book can have a small gross and a large net exposure \ts{10}{6}{1:50}.
\end{definition}

Typically (apart from the present distortions from QE3) dealers' net exposure is to the \emph{term spread}: long long-term bonds, short short-term ones, borrowing short and lending long
like a bank; they profit from the failure of the expectations hypothesis \ts{10}{6}{2:57}. The dealer business is building up gross exposure while controlling
net exposure \ts{10}{6}{3:59}. The Treynor model is about net exposure, price risk; gross exposure is liquidity risk, since the positions must be funded and counterparties may fail, a topic
for the banking lectures \ts{10}{6}{3:59}. Managing price risk well allows tighter spreads, which is how dealers compete \ts{10}{6}{4:41}.

\paragraph{How dealers make money} \ts{10}{6}{5:22}--\ts{10}{6}{7:48}:
\begin{enumerate}
  \item \textbf{Access to inventories}: knowing where securities are (which long-only pension funds hold what) and having a good relationship with a banker, perhaps the parent bank, for
        cash. The inventories are not on the dealer's balance sheet, where they would need funding, but ``out in the world'', like the mangoes in New Jersey;
  \item \textbf{watching the net position} and moving prices so as not to be ``bagged'' by better-informed traders, widening the spread when volatility rises, minute by minute;
  \item \textbf{offsetting} long with short exposures, buying at the bid and selling at the ask: the more both-way business, the more spread earned.
\end{enumerate}

\begin{intuition}[A just-in-time inventory system that straddles the hierarchy]
Through repo the dealer reaches the inventory of cash, through reverse the inventory of securities: it behaves as if it had inventories that are actually out in the market, a just-in-time
system. In a crisis it matters where the inventories are: cash comes from higher in the hierarchy, from banks, securities from lower, from security holders; the dealer straddles the
layers. Troubles come when it must produce securities it reversed in and sold, and more seriously when it must produce money it repoed in and spent \mn{10}{5}\mn{10}{6}.
\end{intuition}

\begin{coursenote}
The notes give the crisis examples: ``fails'' (collateral not returned) became so widespread that the Fed lent its own Treasuries to dealers through the Term Securities Lending Facility,
and when the repo market collapsed it opened the Primary Dealer Credit Facility, a lender of last resort for dealers \mn{10}{6}. (Both were created in March 2008, around the collapse of
Bear Stearns.)
\end{coursenote}

\begin{exercise}[Gross and net]
A dealer is long \$5 billion of 10-year Treasuries and short \$4.8 billion of 2-year Treasuries, all financed by repo and reverse repo.
\begin{enumerate}[(a)]
  \item Compute gross and net exposure in dollar terms.
  \item Why is the net exposure in dollars a poor measure of its price risk? What would you use instead?
\end{enumerate}
\end{exercise}

% ------------------------------------------------------------------
\section{Arbitrage and the assumption of perfect liquidity}\label{sec:dl-arbitrage}

A dealer long Treasury bonds and short Treasury bills is \emph{transporting liquidity} from one market to another: if every security needed its own dealer with its own capital, markets would
be much less liquid \ts{10}{7}{0:11}. Arbitrage links all markets: liquidity is a \emph{macroeconomic} concept, not market by market \ts{10}{7}{0:59}. Dealers choose
where to deploy their scarce resources, seeking wide spreads and little competition; so some markets have wide spreads, some steep bid--ask curves (trade in size and the price moves), some,
like Treasury bills, are nearly flat. Beyond this cross-section, macroeconomic fluctuations of liquidity move all prices together \ts{10}{7}{1:41}.

\paragraph{Prices away from value.} In Treynor's model the price moves with the dealer's \emph{inventory}, not with the fundamental value: the dealer lowers its price not because it knows the
asset is worth less but because its inventory is large \ts{10}{7}{3:05}. If fundamental value is the midpoint at zero inventory, then whenever the dealer holds
inventory, ``the price is wrong'' \ts{10}{7}{4:07}.

\begin{proposition}[Perfect liquidity is a limit case]
Most asset pricing, and general equilibrium theory, assume \emph{perfect arbitrage}: the outside spread collapses onto the fundamental value, value-based traders buy at any price a tiny bit below
it and sell a tiny bit above, the inside spread is squeezed to zero, and liquidity is a free good \ts{10}{7}{4:49}. In Treynor's world your
ability to trade comes from trading at a price \emph{different} from fundamental value: liquidity is provided for profit, and in fluctuating quantity \ts{10}{7}{6:36}.
\end{proposition}

The notes add a ``philosophical question'' \mn{10}{6}\mn{10}{7}: finance assumes that arbitrage creates perfect liquidity by entering at the smallest deviation; but in a fully efficient market
arbitrage would not be profitable, all bets being fair, so position takers would not make money, would compete less for market making, and markets would be less liquid, with wider spreads and
more volatile prices: ``in practice it seems that we must expect liquidity to enter into asset prices''. Arbitrageurs, by taking positions against temporary distortions, spread their
impact into other markets and so counteract them: ``porters'' of liquidity, producing a result ``no part of their intention''. Arbitrage and liquidity are two sides of the same coin
\mn{10}{6}. Some markets, and some times, are close to the frictionless limit; others are not \mn{10}{7}.

\begin{history}[Treynor, Black and ``Noise'']
Jack Treynor, a friend of Fischer Black, published ``The Economics of the Dealer Function'' in the \emph{Financial Analysts Journal} (1987); in 1971 he had already written, under the pseudonym
``Walter Bagehot'', ``The Only Game in Town'' on market makers and informed traders. Black, by then a partner at Goldman Sachs, gave his presidential address to the American Finance Association,
``Noise'', in New York in December 1985 (\emph{Journal of Finance}, July 1986). There he proposed to call a market efficient if price is within a factor of 2 of value, more than half and less than
twice value, and judged that by this definition almost all markets are efficient almost all of the time, ``almost all'' meaning at least 90\%. Mehrling reads the factor of 2 as Treynor's outside
spread \ts{10}{7}{6:58}. (In class the order of the two works is given the other way round.)\par
\emph{References:} J.~L.~Treynor, \emph{Financial Analysts Journal} 43(6), 1987; W.~Bagehot (pseud.), \emph{FAJ} 27(2), 1971; F.~Black, ``Noise'', \emph{Journal of Finance} 41(3), 1986.
\end{history}

\begin{reading}[Black, ``Noise'' (1986), section I ``Finance'', pp.~529--533]
\begin{itemize}
  \item Black contrasts trading on \emph{information} with trading on \emph{noise} as if it were information.
  \item Without noise traders there would be almost no trading in individual assets. An informed trader knows that only another informed trader takes the other side, and so declines to trade. Without trading in individual shares, portfolios, index futures and options could not be priced either.
  \item Noise trading makes markets liquid and makes it pay to collect costly information. But it also puts noise into prices: ``What's needed for a liquid market causes prices to be less efficient.''
  \item Informed traders do not eliminate the noise, because larger positions mean more risk and they can never be sure their information is not already in the price.
  \item Price wanders away from value like a drunk, and is pulled back the harder the further it strays. Hence the definition of efficiency: price within a factor of 2 of value.
\end{itemize}
The connection with Treynor: Black's noise traders are the flow the dealer absorbs, and his informed traders are Treynor's value-based traders. The band within which price may wander is the outside spread.
\end{reading}

``Money and banking matters'': it is not plumbing behind the scenes to be abstracted from in theories of value; it matters for asset pricing, income, employment and business cycles, as the crisis
showed \ts{10}{7}{9:03}.

% ------------------------------------------------------------------
\flushside
\section*{Key formulas and ideas}
\begin{keyformulas}
\begin{tabularx}{\linewidth}{@{}l L@{}}
Market vs funding liquidity & trade in volume without moving the price vs ability to raise funds\\
Inventories & the secret of continuous prices (supermarket, hierarchy of inventories)\\
Treynor & position limits (funding), outside spread (value-based traders), inside spread (width $\sim$ volatility, slope $\sim$ adverse selection)\\
Dealers & demand funding liquidity, supply market liquidity\\
Leverage & dealers trade on borrowed money; matched book: no price risk, but liquidity risk\\
Gross vs net & size of balance sheet vs longs $-$ shorts; dealers build gross, control net\\
Arbitrage & transports liquidity across markets: liquidity is macroeconomic\\
Perfect liquidity & outside spread $\to$ value, inside spread $\to$ 0: a logically impossible limit; Black: price within a factor 2 of value\\
\end{tabularx}
\end{keyformulas}

\section*{Exercises}

Standalone exercises follow the section they practise; connected or integrative ones are at the end.

\begin{exercise}[A linear dealer]
A dealer has $p^\ast=100$, inside spread $s=0.2$, position limits $\pm 1000$, and value-based bid and ask at 95 and 105.
\begin{enumerate}[(a)]
  \item Compute $\lambda$ and the quotes at inventories $-1000$, $0$, $500$, $1000$.
  \item The public sells the dealer 100 units, one at a time. What is the average price received, and how does it compare with $p^\ast$?
\end{enumerate}
\end{exercise}

\begin{exercise}[Crisis in the model]
In the model of the previous exercise the value-based bid falls to 70. Redraw the diagram. What happens to the dealer's willingness to absorb sales, and to the price path when the public sells
1500 units?
\end{exercise}

% !TEX root = ../main.tex
\chapter{Banks and the Market for Liquidity}\label{ch:banksliquidity}
\begin{paracol}{2}

\begin{sources}
\begin{tabularx}{\linewidth}{@{}l l L@{}}
\toprule
\textbf{Segment} & \textbf{Length} & \textbf{Content}\\
\midrule
\seg{11}{1}{L11.1 FT: Money Market Mutual Funds} & 6:46 & The Reserve Primary Fund in court: a run on a deposit substitute.\\
\seg{11}{2}{L11.2 Banks as Money Dealers, a Puzzle} & 4:15 & Banks make markets at par, with no spread: how, and why profitably?\\
\seg{11}{3}{L11.3 Security Dealers as Money Dealers} & 11:19 & Matched book and speculative book; the dealers' repo book rearranged.\\
\seg{11}{4}{L11.4 Adapting Treynor to Liquidity Risk} & 6:33 & The Treynor model for term funding: dealers in liquidity.\\
\seg{11}{5}{L11.5 Digression: Evolution of American Banking} & 11:30 & From bank intermediation to parallel and shadow banking.\\
\seg{11}{6}{L11.6 The Fed in the Fed Funds Market} & 12:48 & The Fed's outside spread (IOER, discount rate) and daily stabilisation.\\
\seg{11}{7}{L11.7 Return to the Initial Puzzle} & 2:14 & Payments and money dealing as joint businesses.\\
\bottomrule
\end{tabularx}

\smallskip
\textbf{Files.} Transcripts \file{materials/transcripts/L11-P*.txt}; Mehrling's notes \mn{11}{1}.
\textbf{Reading.} Treynor (lecture~10); Stigum on banks.
\end{sources}

\section*{Motivation}
Lecture~10 built the Treynor model for security dealers, who turn funding liquidity into market liquidity. This lecture climbs one step up the hierarchy, to
banks \ts{11}{2}{1:09}. Banks make a market between currency and deposits at par: a fixed price, with no bid--ask spread. How can that be possible, and profitable?
\ts{11}{2}{1:49} The route is indirect: first see security dealers as dealers in \emph{money}, apply Treynor's model to liquidity risk,
then bring in the Fed as the maker of the outside spread in the money market; the answer comes at the end.

% ------------------------------------------------------------------
\section{FT: money market mutual funds}\label{sec:bl-ft}

\begin{casestudy}[FT, October 2012: ``Court drama puts focus on money funds'']
Anyone wondering why Mary Schapiro, chairman of the SEC, pushes so hard to reform money market funds should sit in courtroom 6B of Manhattan's federal courthouse,
where the SEC is reliving 2008: the oldest money market fund, the Reserve Primary Fund, suffered a run when investors discovered it was owed money by bankrupt Lehman
Brothers \ts{11}{1}{0:07}. The fund was started by Bruce Bent, and the case involves him and his son \ts{11}{1}{1:08}.
\end{casestudy}

\begin{definition}[New concepts: money market mutual fund, net asset value]
A \emph{money market mutual fund} (MMMF) holds money market assets (time deposits, Eurodollars, commercial paper, repo) and issues shares meant to look and feel like
deposits \ts{11}{1}{1:28}. Like any mutual fund, the value of its shares is the value of its assets: the \emph{net asset value} (NAV) per share
\ts{11}{1}{2:30}. Money funds use accounting that keeps the NAV at \$1, absorbing fluctuations in the interest they pay, so that a share can be treated as a
dollar deposit \ts{11}{1}{3:11}. ``Breaking the buck'' means the NAV falling below \$1.
\end{definition}

\begin{taccts}
\tacct[2.6cm]{Money market fund}{Commercial paper (incl.\ Lehman)\\Eurodollar time deposits\\Repo}{Shares (NAV kept at \$1)}
\end{taccts}
Money funds were a response to a regulatory constraint, \emph{Regulation Q}, which forbade banks to pay interest on (demand) deposits: Bent offered interest on something
very like a deposit, and ``the world beat a path to his door'' \ts{11}{1}{3:52}. In September 2008 the fund's Lehman commercial paper became worthless; the
loss could not be absorbed in the interest rate (it would have needed a negative rate), holders ran, and shares were redeemed below a dollar. The fund was diversified, so it
lost only its Lehman paper, a few cents per share \ts{11}{1}{4:33}. The SEC claims the Bents told investors they intended to support the fund with
their own money when they had no such intention; Bent replies that it was a run, a liquidity problem that no private wealth could stop \ts{11}{1}{5:35}. The run stopped when the
Treasury in effect created deposit insurance for money funds \ts{11}{1}{5:55}.

\begin{coursenote}[Facts about the Reserve Primary Fund]
In class: the first money fund, started in 1969, redeemed at about 98 cents \ts{11}{1}{1:08}. The Reserve Fund was launched by Bruce Bent and Henry Brown in
1971--1972; its NAV fell to \$0.97 on 16 September 2008. On 19 September the Treasury announced a temporary guarantee program for money market funds. Regulation Q
prohibited interest on demand deposits and capped the rates on time and savings deposits.
\end{coursenote}

% ------------------------------------------------------------------
\section{The puzzle: banks as money dealers}\label{sec:bl-puzzle}

A bank makes a market between two layers of the hierarchy, currency and deposits, at par (\cref{sec:nh-marketmakers}) \ts{11}{2}{1:49}. This seems to contradict the Treynor
model on two counts \ts{11}{2}{2:30}:
\begin{enumerate}
  \item the price is \emph{fixed}: it cannot fall as inventories grow, which is how a dealer protects itself;
  \item there is \emph{no bid--ask spread}: currency goes into deposits and deposits into currency at the same price.
\end{enumerate}
So: how is it possible, and why is it profitable? ``When you confront a puzzle like this, it's better not to hit it head on'': back up to security dealers
\ts{11}{2}{3:34}.

% ------------------------------------------------------------------
\section{Security dealers as money dealers}\label{sec:bl-moneydealers}

\paragraph{Two books.} Split a security dealer's book in two \ts{11}{3}{0:21}:
\begin{itemize}
  \item a \emph{matched book}: securities in 100, securities out 100, the same securities, long and short, netting to zero;
  \item a \emph{speculative book} (Treynor's book): a net long position of 10 in bonds financed by borrowing of 10, exposed to price risk.
\end{itemize}
In real dealers' balance sheets most of the book is matched, gross exposure, and only a little is net; they make a lot of money on the matched book too, which Treynor's model
ignores \ts{11}{3}{1:27}.

\paragraph{The gestalt switch.} Now look at the security dealer as a dealer in \emph{money}: the securities are just collateral for borrowing and lending money
\ts{11}{3}{3:11}. From the New York Fed's primary dealer statistics (the memorandum items, in billions) \ts{11}{3}{4:35}:
\begin{taccts}
\tacct[3.4cm]{Primary dealers as money dealers}{Overnight reverse\amt{854}\\Term reverse\amt{1{,}253}}{Overnight repo\amt{1{,}796}\\Term repo\amt{826}}
\end{taccts}
Reverse and repo are the same instrument from opposite sides, so there is a matched book underneath \ts{11}{3}{4:58}. Assume, for the sake of the concept, that the
term instruments have the same term \ts{11}{3}{5:39}. Carve out the matched parts \ts{11}{3}{6:00}:
\begin{taccts}
\tacct[3.4cm]{Matched book}{Overnight reverse\amt{854}\\Term reverse\amt{826}}{Overnight repo\amt{854}\\Term repo\amt{826}}\quad
\tacct[3.4cm]{The rest}{Term reverse\amt{427}\\Net financing (bonds)\amt{515}}{Overnight repo\amt{942}}
\end{taccts}
The numbers are the same, only rearranged; the ``net financing'' of 515 is bonds financed in the overnight repo market. One more step: finance the bonds with term repo instead,
adding 515 of term repo on both sides \ts{11}{3}{8:29}. The rest becomes two separate exposures, read line by line:
\begin{taccts}
\tacct[3.4cm]{Two exposures}{Term reverse\amt{942}\\Bonds\amt{515}}{Overnight repo\amt{942}\\Term repo\amt{515}}
\end{taccts}
\begin{enumerate}
  \item \textbf{Bonds 515 against term repo 515}: Treynor's \emph{price risk} \ts{11}{3}{9:10};
  \item \textbf{Term reverse 942 against overnight repo 942}: borrowing overnight, lending for a month, both secured. ``You're basically a bank'': the overnight repo is like a demand
        deposit, the term reverse like a one-month loan. This is \emph{liquidity risk}: the funding must be rolled every day, and could run, as the Reserve Primary Fund did
        \ts{11}{3}{9:31}.
\end{enumerate}

\begin{definition}[New concepts: price risk, liquidity risk (of a dealer)]
A dealer's \emph{price risk} is its exposure to changes in the prices of the securities it holds net (its net position). Its \emph{liquidity risk} is its exposure to the
need to refinance short-term borrowing that funds longer-term lending, i.e.\ to the market for funding.
\end{definition}

% ------------------------------------------------------------------
\section{Adapting Treynor to liquidity risk}\label{sec:bl-treynor}

Put \emph{liquidity risk} on the horizontal axis of Treynor's diagram: the scale of the maturity mismatch, up to some maximum \ts{11}{4}{0:07}. There are two prices,
the overnight rate and the term rate; what matters is the spread between them, the dealer's profit \ts{11}{4}{0:49}. Take the overnight rate as fixed by the Fed (Fed funds, and by
arbitrage repo) for the moment, and focus on the \emph{term} rate \ts{11}{4}{1:09}. Money markets quote yields, not prices; a yield is an inverse price, so the curves
slope \emph{upward} and bid and offer swap places \ts{11}{4}{1:52}.

\begin{nfigure}
\begin{tikzpicture}[font=\scriptsize,x=1cm,y=1cm]
  \draw[-{Stealth[length=4pt]}] (-4.4,0) -- (4.6,0) node[right,align=left]{liquidity risk\\(term lent $-$ term borrowed)};
  \draw[-{Stealth[length=4pt]}] (0,-0.2) -- (0,4.6) node[above]{term rate};
  \draw[dashed] (-4,0) -- (-4,4.3); \draw[dashed] (4,0) -- (4,4.3);
  \draw[very thick,defcol] (-4,0.9) -- (4,3.7);
  \draw[very thick,intcol] (-4,0.5) -- (4,3.3);
  \node[defcol,rotate=19] at (-1.8,2.05) {bid (rate to lend term)};
  \node[intcol,rotate=19] at (-1.6,0.55) {offer (rate to borrow term)};
  \draw[keycol,thick] (-4.2,1.3) -- (4.2,1.3) node[right]{overnight rate (Fed)};
  \draw[{Stealth[length=3pt]}-{Stealth[length=3pt]}] (1.2,1.3) -- (1.2,2.72);
  \node[right,align=left] at (1.25,1.75) {spread $=$\\dealer's\\profit};
\end{tikzpicture}
\caption{Treynor's model applied to liquidity \ts{11}{4}{1:52}: quotes are yields, so the curves slope upward. To take on more liquidity risk (lend more at term
against overnight funding) the dealer must be paid a larger spread of the term rate over the overnight rate.}\label{fig:bl-liqtreynor}
\end{nfigure}

\begin{proposition}[The price of liquidity risk]
For dealers to take on more liquidity risk they must be paid, and what they are paid is the spread between the term rate and the overnight rate. With a big enough spread,
``give it a whirl, life is short'': they expand their balance sheets, although a run is always possible \ts{11}{4}{2:57}.
\end{proposition}

\paragraph{Two margins.} The bond price is driven by the dealer's net bond position (515), the term rate by its liquidity position (942); there is no reason for them to be equal.
A dealer can change its bond position without changing its liquidity risk, by expanding both sides, and vice versa: a profit-seeking dealer adjusts both margins, dealing in
bonds and in money separately \ts{11}{4}{3:41}. The same diagram describes a bank, with deposits in place of overnight
repo and term loans to companies in place of term reverse: a bank is a dealer whose ``security'' is money \ts{11}{4}{5:24}. The other side of a bank, the payment system,
is what puzzled us \ts{11}{4}{6:06}.

% ------------------------------------------------------------------
\section{Digression: the evolution of American banking}\label{sec:bl-evolution}

Why this way of thinking is now necessary, and why textbooks, by a kind of hysteresis, still do not use it; it also helps read Stigum, who belongs to an earlier generation
\ts{11}{5}{0:07}.

\begin{enumerate}
  \item \textbf{Banks as intermediaries} (``\emph{It's a Wonderful Life}'' banking of the 1950s): households save in bank deposits, banks lend to businesses for capital
        investment; Janus-faced, giving households the short-term liquid asset they want and firms the long-term liability they need. Banks also hold government securities and
        borrow wholesale, which lets loans and deposits move separately; capital absorbs losses, reserves at the Fed serve payments \ts{11}{5}{1:12}.
  \item \textbf{The death of loans}: corporations found the commercial paper market cheaper than bank loans \ts{11}{5}{4:19}.
  \item \textbf{The death of deposits}: households and businesses moved money to money market funds, which paid interest \ts{11}{5}{4:39}.
  \item \textbf{A parallel banking system}: finance companies (General Motors' financing arm making auto loans) issue commercial paper, money funds buy it and issue deposit
        substitutes; intermediation bypasses banks \ts{11}{5}{5:20}.
  \item \textbf{The shadow banking system}: households issue mortgages, which are packaged into residential mortgage-backed securities, sliced into AAA ``quasi Treasury bills''
        held by a shadow bank (on the balance sheet of a European bank, or a structured investment vehicle of Citibank), which uses them as collateral to borrow in the wholesale
        money market, through asset-backed commercial paper and repo, bought by money funds that issue shares \ts{11}{5}{6:43}.
        None of it passes through a bank with an account at the Fed, and much happened offshore: UBS or Deutsche Bank borrowing in the dollar money market to invest in dollar RMBS
        \ts{11}{5}{9:12}.
\end{enumerate}

\begin{nfigure}
\begin{tikzpicture}[font=\scriptsize,
  b/.style={draw,rounded corners=2pt,align=center,minimum height=0.75cm,inner sep=3pt},
  arr/.style={-{Stealth[length=3pt]},semithick}]
  \node[font=\small\bfseries] at (0,3.2) {traditional};
  \node[b,fill=fincol!15] (h1) at (-1.6,2.3) {households};
  \node[b,fill=defcol!15] (b1) at (0,1.2) {bank};
  \node[b,fill=intcol!15] (f1) at (1.6,2.3) {firms};
  \draw[arr] (h1) -- node[left]{deposits} (b1);
  \draw[arr] (b1) -- node[right]{loans} (f1);
  \node[font=\small\bfseries] at (6.4,3.2) {shadow banking};
  \node[b,fill=fincol!15] (h2) at (4.0,2.3) {households\\(mortgages)};
  \node[b,fill=keycol!15] (s2) at (5.7,1.2) {shadow bank\\RMBS (AAA)};
  \node[b,fill=defcol!15] (m2) at (7.6,1.2) {money fund};
  \node[b,fill=fincol!15] (h3) at (9.0,2.3) {households\\(fund shares)};
  \draw[arr] (s2) -- node[left]{RMBS} (h2);
  \draw[arr] (m2) -- node[below]{ABCP, repo} (s2);
  \draw[arr] (h3) -- node[right]{shares} (m2);
\end{tikzpicture}
\caption{From bank intermediation to market-based credit \ts{11}{5}{1:33}: the shares of money funds fund mortgages through the wholesale money market, without passing
through a bank. (Arrows point from the funder to the funded asset.)}\label{fig:bl-shadow}
\end{nfigure}

The notes add what happened in between \mn{11}{4}\mn{11}{5}. Banks hit by the death of loans adjusted by providing backup lines of credit to commercial paper issuers. The
death of deposits came as Regulation Q ceilings eroded and money market funds competed: ``disintermediation'', ultimate borrowers and lenders taking interest rate risk without a bank
in between. As banks lost their core business they were kept out of others (the Securities Industry Association resisted bank underwriting) and subjected to capital requirements
others did not face, partly for the safety of the payment system and because of their privileged access to the Fed: this pushed them to strip the balance sheet, FRAs for deposits,
futures for actuals. The shadow banking system faces the liquidity and solvency risks of traditional banking without the government backstops, relying on the wholesale money market
and the CDS market instead; its lender of last resort is the traditional banking system, through lines of credit and liquidity commitments. (On the older practice of managing bank
interest rate and liquidity risk the notes cite M.~Stigum and R.~Branch, \emph{Managing Bank Assets and Liabilities}, 1983.)

\begin{intuition}[Two prices, both made by dealers]
A shadow bank makes money on two prices: the price of its assets, essentially a bond price, and the price of its funding, a term interest rate. Both are made in dealer markets, by
the two versions of Treynor's model of this and the previous lecture. ``The modern banking system is a capital-market-based credit system, not a bank-loan-based system'': hence
the central place of the economics of the dealer function in the course \ts{11}{5}{10:16}.
\end{intuition}

% ------------------------------------------------------------------
\section{The Fed in the Fed funds market}\label{sec:bl-fed}

Now unbracket the overnight rate \ts{11}{6}{0:07}. The Fed does two things in the Fed funds market \ts{11}{6}{1:08}:
\begin{enumerate}
  \item it makes the \emph{outside spread};
  \item in normal times it \emph{stabilises} the rate within it.
\end{enumerate}

\paragraph{The outside spread today.} Since the crisis the Fed pays \emph{interest on excess reserves} (IOER), 0.25\%, where before reserves earned nothing, like cash: a floor
\ts{11}{6}{1:49}. (Fed funds effective at 0.16\% is below the floor because some lenders cannot hold accounts at the Fed and lend to banks that can, which
pocket the spread by depositing at 0.25\% \ts{11}{6}{2:51}.) The ceiling is the discount rate, 0.75\%: the Fed lends to banks at that rate, so Fed funds cannot
really go above it \ts{11}{6}{3:32}. With a trillion dollars of excess reserves, the market is ``smack up against'' the floor: the Fed is flooding the market with liquidity.
This is unusual, maybe pathological, and the Fed is experimenting with ways to soak up reserves: reverse repos, and last month term deposits \ts{11}{6}{4:13}.
The notes add: the Fed's communication says it expects to widen this 50 basis point outside spread, with IOER some 75 basis points below the target; in a way the effective rate
should be thought of as 25, not 16, which helps explain why repo trades above Fed funds, ``a lot of distortion in the system right now''; and the first offering of the Fed's own
term deposits took place on 10 September 2012 \mn{11}{5}\mn{11}{6}\mn{11}{7}.

\paragraph{Before the crisis.} The floor was zero; the discount rate was set 100 basis points above the Fed funds \emph{target}: an outside spread that never bit, because the Fed
stabilised the rate around the target every day, anticipating fluctuations in the demand for reserves and supplying it, sometimes with repos twice a day; whole books were written on
forecasting the need for reserves \ts{11}{6}{5:34}.

\begin{nfigure}
\begin{tikzpicture}[font=\scriptsize,x=1cm,y=1cm]
  \begin{scope}
    \node[font=\small\bfseries] at (0,4.3) {before the crisis};
    \draw[-{Stealth[length=3pt]}] (-2.6,0) -- (2.8,0) node[right]{reserves};
    \draw[-{Stealth[length=3pt]}] (-2.4,-0.1) -- (-2.4,4) node[above]{rate};
    \draw[thick,thmcol] (-2.4,3.4) -- (2.5,3.4) node[right]{discount rate};
    \draw[thick,thmcol] (-2.4,0.15) -- (2.5,0.15) node[right]{0 (reserves)};
    \draw[thick,defcol] (-2,1.4) -- (2,2.9);
    \draw[dashed,keycol] (-2.4,2.15) -- (2.5,2.15) node[right]{target};
    \fill[keycol] (0,2.15) circle (2pt);
  \end{scope}
  \begin{scope}[xshift=7cm]
    \node[font=\small\bfseries] at (0,4.3) {2012};
    \draw[-{Stealth[length=3pt]}] (-2.6,0) -- (2.8,0) node[right]{reserves};
    \draw[-{Stealth[length=3pt]}] (-2.4,-0.1) -- (-2.4,4) node[above]{rate};
    \draw[thick,thmcol] (-2.4,2.0) -- (2.5,2.0) node[right]{0.75\% discount};
    \draw[thick,thmcol] (-2.4,0.8) -- (2.5,0.8) node[right]{0.25\% IOER};
    \draw[thick,defcol] (-2,1.0) -- (-0.5,1.5);
    \fill[keycol] (2.2,0.8) circle (2pt) node[above]{\$1 tn};
  \end{scope}
\end{tikzpicture}
\caption{The Fed's outside spread in the Fed funds market \ts{11}{6}{2:09}: before the crisis a wide spread with the rate stabilised at the target
inside it; in 2012 a narrow spread with the market pressed against the floor by a trillion dollars of excess reserves. (Schematic.)}\label{fig:bl-corridor}
\end{nfigure}

\paragraph{How the Fed stabilised.} If the market pushed Fed funds above the target, the Fed added reserves with a \emph{temporary open market operation}: a repo with a dealer, a
purchase and sale of a Treasury for a fixed period, overnight, three days, a week (an outright purchase would be permanent) \ts{11}{6}{8:40}. The
dealer, which has no Fed account, uses the cash to repay a loan from its bank, borrowed at Fed funds plus 50 basis points: it replaces expensive borrowing with cheap repo. The bank's
reserves rise, and it lends them in the Fed funds market, pushing the rate down \ts{11}{6}{10:08}.
\begin{taccts}
\tacct[2.0cm]{Fed}{$+$Repo (dealer)}{$+$Reserves (bank)}\quad
\tacct[2.0cm]{Dealer}{}{$+$Repo from Fed\\$-$Loan from bank}\quad
\tacct[2.0cm]{Bank}{$+$Reserves\\$-$Loan to dealer}{}
\end{taccts}
In doing so the Fed moves dealer inventories and so their prices; and it puts part of the Fed funds market on its own balance sheet, borrowing from a bank and lending to a dealer: a
``dealer of first resort'', trading not at the outside spread but inside the inside spread, since it asks the dealers to bid and takes the best bids \ts{11}{6}{11:32}.

\begin{definition}[New concepts: interest on excess reserves, corridor]
\emph{Interest on excess reserves} (IOER), paid by the Fed since October 2008, is the rate on reserve balances above the requirement; it sets a floor under the overnight rate for
institutions that can hold reserves. Together with the discount rate as ceiling, it forms a \emph{corridor} within which the policy rate moves.
\end{definition}

% ------------------------------------------------------------------
\section{Return to the initial puzzle}\label{sec:bl-answer}

Banks, like security dealers, have two faces: on one side they are money dealers, in the money market; on the other they operate the payment system \ts{11}{7}{0:07}.

\begin{proposition}[How banks can make markets at par]
\leavevmode
\begin{enumerate}[(i)]
  \item \emph{Possible}: all the slack is taken up in the money market: the ability to expand and contract the balance sheet and to get reserves when needed is what lets a bank
        hold par in the payment system \ts{11}{7}{0:49}.
  \item \emph{Profitable}: the money-market side earns a spread; the payment side may even lose money, but the two are on the same balance sheet, complementary, and the joint
        business is profitable \ts{11}{7}{1:10}.
  \item As for dealers, the size of the payment business and of the funding exposure are separate margins the bank can move independently \ts{11}{7}{1:50}.
\end{enumerate}
\end{proposition}

% ------------------------------------------------------------------
\flushside
\section*{Key formulas and ideas}
\begin{keyformulas}
\begin{tabularx}{\linewidth}{@{}l L@{}}
Money funds & deposit substitutes, NAV kept at \$1; vulnerable to runs (Reserve Primary Fund, 2008)\\
Dealer's books & matched book (gross) $+$ speculative book (net, Treynor)\\
Rearrangement & repo book $=$ matched book $+$ price risk (bonds vs term repo) $+$ liquidity risk (term reverse vs overnight repo)\\
Liquidity Treynor & term rate on vertical axis (upward-sloping quotes); dealers paid the term--overnight spread to bear liquidity risk\\
Evolution & intermediation $\to$ death of loans and deposits $\to$ parallel banking $\to$ shadow banking (capital-market-based credit)\\
Fed & outside spread: IOER (floor), discount rate (ceiling); stabilisation by temporary repos with dealers\\
Banks at par & possible thanks to the money market; profitable as a joint business with it\\
\end{tabularx}
\end{keyformulas}

\section*{Exercises}

\begin{exercise}[Rearranging a repo book]
A dealer has overnight reverse 300, term reverse 700, overnight repo 800, term repo 250 (assume equal terms).
\begin{enumerate}[(a)]
  \item Separate the matched book and the rest.
  \item What is its net bond position? Rewrite the rest as a price-risk exposure and a liquidity-risk exposure.
\end{enumerate}
\end{exercise}

\begin{exercise}[A run on a money fund]
A money fund has \$100 of assets, of which \$3 of commercial paper of a bank that fails (worth zero), and 100 shares.
\begin{enumerate}[(a)]
  \item What is the NAV? What happens if it keeps paying \$1 per share to the first 50\% of holders who redeem?
  \item Why is it rational to run? Which institution's guarantee stops the run, and why?
\end{enumerate}
\end{exercise}

\begin{exercise}[The corridor]
With IOER at 0.25\% and the discount rate at 0.75\%:
\begin{enumerate}[(a)]
  \item Explain why Fed funds can trade below 0.25\%.
  \item Describe what the Fed must do to raise the effective rate to 0.50\% with \$1 trillion of excess reserves outstanding, using the tools mentioned in class (reverse repos,
        term deposits).
\end{enumerate}
\end{exercise}

\begin{exercise}[Pricing liquidity]
A dealer can borrow overnight at 0.30\% and lends term (one month) at a rate that rises by 1 basis point for each \$10 billion of term lending it does. It must earn at least 5 basis
points over the overnight rate on its marginal lending. What is its maximum term lending if the market term rate is 0.40\% at zero dealer position?
\end{exercise}

% !TEX root = ../main.tex
\chapter{Lender/Dealer of Last Resort}\label{ch:lolr}
\begin{paracol}{2}

\begin{sources}
\begin{tabularx}{\linewidth}{@{}l l L@{}}
\toprule
\textbf{Segment} & \textbf{Length} & \textbf{Content}\\
\midrule
\seg{12}{1}{L12.1 FT: Citibank and the SIVs} & 5:10 & Pandit resigns; Citi's structured investment vehicles and the collapse of ABCP.\\
\seg{12}{2}{L12.2 The Art of Central Banking} & 3:42 & Hawtrey (1932); the Fed's crisis facilities.\\
\seg{12}{3}{L12.3 Evolution of Monetary Policy: 1951--1987} & 7:06 & Quantity and price views; operating targets from free reserves to Fed funds.\\
\seg{12}{4}{L12.4 The Taylor Rule: 1987--2007} & 7:04 & The rule, the Fisher effect, inflation targeting.\\
\seg{12}{5}{L12.5 Monetary Transmission Mechanism} & 10:00 & From Fed funds to term rates to asset prices, through three dealer markets.\\
\seg{12}{6}{L12.6 Anatomy of a Normal Crisis} & 8:18 & Dealers and banks absorb an imbalance, for a price.\\
\seg{12}{7}{L12.7 Anatomy of a Serious Crisis} & 4:32 & The Fed as backstop of the banks that backstop the dealers.\\
\seg{12}{8}{L12.8 Should the Fed Intervene or Not?} & 8:40 & Elasticity against spirals, discipline against accumulated mistakes.\\
\seg{12}{9}{L12.9 The Fed as Dealer of Last Resort: 2007--2009} & 15:50 & The stages of the crisis in balance sheets.\\
\bottomrule
\end{tabularx}

\smallskip
\textbf{Files.} Transcripts \file{materials/transcripts/L12-P*.txt}.
\textbf{Readings.} Stigum, chapter~9 on monetary policy; R.~G.~Hawtrey, \emph{The Art of Central Banking} (1932), chapter~4 (quoted).
\end{sources}

\section*{Motivation}
The last lecture before the midterm puts the pieces together: the three dealer markets of lectures 6--11 (Fed funds, term funding, bonds) form the
\emph{transmission mechanism} of monetary policy; normal fluctuations are absorbed by dealers and banks; serious crises need the Fed as lender of last
resort to the banks; and in 2007--2009 the Fed went further, becoming the money market itself and finally a \emph{dealer of last resort}. The lecture is
mostly about lender of last resort in the money market; dealer of last resort in capital markets needs derivatives and comes after the midterm
\ts{12}{2}{0:07}.

% ------------------------------------------------------------------
\section{FT: Citibank and the SIVs}\label{sec:lr-ft}

\begin{casestudy}[FT, October 2012: Vikram Pandit resigns]
Citigroup's CEO Vikram Pandit resigned under board pressure and was replaced by Michael Corbat, who had run Citi Holdings, where the toxic assets were put to
be worked out \ts{12}{1}{0:07}\ts{12}{1}{0:28}\ts{12}{1}{3:31}. Citi, built by Sandy Weill as a ``state-of-the-art financial supermarket'', came close to
nationalisation in the crisis \ts{12}{1}{0:49}\ts{12}{1}{1:09}. The FT's timeline omits an episode Mehrling would add: in November 2007, as Citi announced losses
that grew every few days (5, then 8, then 11 billion), it received a \$7.5 billion capital infusion from the Abu Dhabi Investment Authority: professional
investors still thought it was a minor blip, as many economists thought subprime too small to matter \ts{12}{1}{1:30}\ts{12}{1}{1:50}\ts{12}{1}{2:10}\ts{12}{1}{2:51}.
It was the beginning of the global financial crisis \ts{12}{1}{3:11}.
\end{casestudy}

\paragraph{The SIVs.} Citi was a famous part of the shadow banking system: it set up off-balance-sheet \emph{structured investment vehicles} that held AAA
tranches of residential mortgage-backed securities, funded with asset-backed commercial paper \ts{12}{1}{3:31}\ts{12}{1}{3:52}\ts{12}{1}{4:12}. The first thing that
happened in the crisis was the collapse of the SIVs and of the ABCP market: 90-day paper funding 30-year mortgages; at maturity the holders, money funds,
declined to roll, the SIV found no other funding and borrowed from its host, Citibank, which ended up with the assets on its own balance sheet
\ts{12}{1}{4:33}\ts{12}{1}{4:53}.

\begin{definition}[New concepts: structured investment vehicle, ABCP]
A \emph{structured investment vehicle} (SIV) is an off-balance-sheet entity sponsored by a bank that holds long-term assets (here AAA mortgage-backed securities)
funded by short-term debt. \emph{Asset-backed commercial paper} (ABCP) is commercial paper secured by a pool of assets, typically 30 to 90 days.
\end{definition}

% ------------------------------------------------------------------
\section{The art of central banking}\label{sec:lr-art}

Ralph Hawtrey, \emph{The Art of Central Banking} (London, 1932, ``an auspicious year''), chapter 4, subtitled ``the lender of last resort'': ``a central bank is
a bankers' bank'', and the central bank is the lender of last resort; Hawtrey, in the line of Bagehot and up through Keynes, Sayers and Goodhart, was an
inspiration for Mehrling \ts{12}{2}{0:48}\ts{12}{2}{1:08}\ts{12}{2}{1:29}. (His copy, withdrawn from the Kirkland House library at Harvard, cost three dollars
\ts{12}{2}{1:50}.)

\paragraph{The crisis facilities.} The New York Fed lists the Fed's lending facilities: regular open market operations and the discount window, and then a long
list that was not there before the crisis \ts{12}{2}{2:10}\ts{12}{2}{2:31}: the term discount window program, the Term Auction Facility, the Primary Dealer Credit
Facility, transitional credit extensions, reciprocal currency arrangements (the liquidity swaps), securities lending and the Term Securities Lending Facility, the
ABCP Money Market Mutual Fund Liquidity Facility, the Commercial Paper Funding Facility, the Money Market Investor Funding Facility, the Term Asset-Backed
Securities Loan Facility \ts{12}{2}{2:51}\ts{12}{2}{3:12}. ``This is what put a floor under it'' \ts{12}{2}{3:33}.

% ------------------------------------------------------------------
\section{Normal monetary policy}\label{sec:lr-normal}

\subsection{Quantity or price?}
Once central banks accepted Bagehot's responsibility, they started to think how to avoid crises in the first place, and invented monetary policy, which evolved from
1873 to what it was before the crisis \ts{12}{3}{0:07}\ts{12}{3}{0:28}. Stigum shows an ambivalence about what the Fed controls: a quantity (money, reserves, a balance
sheet measure) or a price (the Fed funds rate)? \ts{12}{3}{0:48}\ts{12}{3}{1:09} Those who emphasise quantities have the equation of exchange $MV=PT$ in mind; those who
emphasise rates, the money market view \ts{12}{3}{1:30}\ts{12}{3}{1:51}. The two can be translated into each other:
\begin{itemize}
  \item \textbf{quantity}: the Fed sits on top of the hierarchy (the pyramid of lecture~2) and controls the best money of all, its own liability, high-powered money,
        which it uses to influence the balance between discipline and elasticity \ts{12}{3}{2:12}\ts{12}{3}{2:33}\ts{12}{3}{2:54}\ts{12}{3}{3:15};
  \item \textbf{price}: the Fed pegs the shortest end of the term structure, an administered policy rate, while the market determines the long end, security prices
        \ts{12}{3}{3:38}\ts{12}{3}{3:59}.
\end{itemize}
Either way, monetary policy is about the \emph{transmission mechanism} from what the Fed controls, its own balance sheet, to the financial system and the economy
\ts{12}{3}{4:20}.

\paragraph{Operating targets} (Stigum) \ts{12}{3}{4:41}\ts{12}{3}{5:02}\ts{12}{3}{5:24}:
\begin{itemize}
  \item 1951: free reserves;
  \item 1979: non-borrowed reserves (Volcker's ``Saturday night special'');
  \item 1983: borrowed reserves;
  \item 1987: the Fed funds rate, ever since.
\end{itemize}
Mehrling is not sure the Fed changed much what it actually did, rather how it explained it \ts{12}{3}{5:45}. In 1979, facing double-digit inflation, Volcker presented
himself as a convert to monetarism; most people think it was a convenient language: saying ``I'm going to jack the Fed funds rate up to 15\%'' might have cost him his
confirmation, saying ``I'm a monetarist'' had the same consequence without saying it, and he stopped inflation in its tracks \ts{12}{3}{6:05}\ts{12}{3}{6:26}\ts{12}{3}{6:46}.

\begin{history}[The Saturday night special]
On Saturday 6 October 1979 the FOMC, chaired by Paul Volcker since August, announced that it would control the supply of bank reserves (non-borrowed reserves) instead of
the Fed funds rate; the rate rose to around 20\% in 1980--1981. The approach was relaxed in 1982.\par
\emph{References:} Federal Reserve press release, 6 Oct 1979; A.~H.~Meltzer, \emph{A History of the Federal Reserve}, vol.~2, Chicago UP, 2009.
\end{history}

\subsection{The Taylor rule (1987--2007)}
1987 is when Mehrling got his job; by the eve of the crisis almost all economists agreed on what the Fed was doing and should do: the \emph{Taylor rule}
\ts{12}{4}{0:07}\ts{12}{4}{0:28}.

\begin{keyformulas}
\textbf{Taylor rule} (in the form given in class) \ts{12}{4}{0:49}:
\[
  r = \rho + \pi^e + \alpha\,(\pi^e - \pi^\ast) + \beta\,(y - y_f),
\]
with $r$ the target (Fed funds) rate, $\rho$ the natural real rate of interest (the return on real capital), $\pi^e$ expected inflation, $\pi^\ast$ the inflation target,
$y-y_f$ the gap between output and its full-employment level \ts{12}{4}{1:09}\ts{12}{4}{1:30}\ts{12}{4}{1:51}.

\medskip
\textbf{Fisher effect}: $r = \rho + \pi^e$, nominal rate $=$ real rate $+$ expected inflation \ts{12}{4}{2:31}\ts{12}{4}{2:56}.
\end{keyformulas}

\begin{itemize}
  \item As a regression, the rule fits the behaviour of central banks around the world well in normal times (not in the crisis): a positive description; a large literature
        then made it normative \ts{12}{4}{1:51}\ts{12}{4}{2:11}\ts{12}{4}{2:31}.
  \item The first part, $\rho+\pi^e$, is the \emph{Fisher effect} (Irving Fisher, Yale): what markets do by themselves, compensating creditors for expected inflation. Fisher
        himself thought people often got inflation wrong, causing fluctuations; modern economics, with rational expectations, assumes $\pi^e$ is unbiased
        \ts{12}{4}{2:56}\ts{12}{4}{3:17}\ts{12}{4}{3:37}\ts{12}{4}{3:58}.
  \item The rest is the Fed leaning against the wind: if expected inflation exceeds the target, raise the rate by more than the difference (more discipline); if output is
        below full employment, lower it (more elasticity) \ts{12}{4}{4:18}\ts{12}{4}{4:39}.
\end{itemize}
Careers were made on the optimal $\alpha$ and $\beta$, on whether to add a financial-stability term: ``inflation targeting'' spread all over the world. Economics came to
believe it had mastered the economy to the point of arguing whether a coefficient should be 1.1 or 1.2; the Great Moderation, then the crisis. ``Nobody believes this
anymore'' \ts{12}{4}{4:59}\ts{12}{4}{5:19}\ts{12}{4}{5:40}\ts{12}{4}{6:01}. Mehrling, who was reading Hawtrey instead, was better prepared: ``in the land of the blind the
one-eyed man is king'' \ts{12}{4}{6:22}\ts{12}{4}{6:42}.

\begin{coursenote}[Taylor's original rule]
John Taylor (``Discretion versus policy rules in practice'', \emph{Carnegie--Rochester Conference Series}, 1993) wrote $i = \pi + 0.5\,y + 0.5\,(\pi-2) + 2$, with
$\pi$ inflation over the previous four quarters, $y$ the percentage output gap and 2\% for both the real rate and the inflation target. The total response of $i$ to
inflation is $1+0.5=1.5>1$ (the ``Taylor principle''). In the form above the coefficient on $\pi^e$ is $1+\alpha$, so the principle requires $\alpha>0$; in class
Mehrling says ``$\alpha$ is greater than one'' \ts{12}{4}{4:18}.
\end{coursenote}

% ------------------------------------------------------------------
\section{The monetary transmission mechanism}\label{sec:lr-transmission}

The three dealer diagrams of lectures 10--11, in reverse order, form the transmission mechanism \ts{12}{5}{0:08}\ts{12}{5}{0:28}:
\begin{enumerate}
  \item \textbf{Fed funds} (overnight): the Fed makes the outside spread (IOER, discount rate); the horizontal axis is \emph{settlement risk}, the risk of pushing payments to
        tomorrow \ts{12}{5}{1:11}\ts{12}{5}{1:32}\ts{12}{5}{1:53}\ts{12}{5}{2:16};
  \item \textbf{term funding} (three months): \emph{liquidity risk}, borrowing short and lending long; upward-sloping quotes in yields \ts{12}{5}{2:37}\ts{12}{5}{2:58};
  \item \textbf{bonds}: asset prices, downward-sloping quotes; \emph{price risk} \ts{12}{5}{2:58}\ts{12}{5}{3:19}.
\end{enumerate}

\begin{nfigure}
\begin{tikzpicture}[font=\scriptsize,x=0.64cm,y=0.64cm]
  \foreach \xs/\t/\xl/\yl in {0/{Fed funds (overnight)}/{settlement risk}/{rate},5.6/{term funding}/{liquidity risk}/{term rate},11.2/{bonds}/{price risk}/{price}} {
    \begin{scope}[xshift=\xs cm]
      \draw[-{Stealth[length=3pt]}] (0,0) -- (5.4,0) node[below left]{\xl};
      \draw[-{Stealth[length=3pt]}] (0,0) -- (0,4.4) node[above]{\yl};
      \node[font=\scriptsize\bfseries] at (2.7,5.2) {\t};
    \end{scope}}
  \begin{scope}[xshift=0cm]
    \draw[thick,thmcol] (0,3.6) -- (5.2,3.6) node[right]{discount};
    \draw[thick,thmcol] (0,0.6) -- (5.2,0.6) node[right]{IOER};
    \draw[thick,defcol] (0.5,1.2) -- (4.8,3.0);
    \draw[dashed,keycol] (0,2.1) -- (5.2,2.1) node[right]{target};
  \end{scope}
  \begin{scope}[xshift=5.6cm]
    \draw[thick,defcol] (0.4,1.6) -- (4.9,3.6);
    \draw[dashed,keycol] (0,2.1) -- (5.2,2.1);
  \end{scope}
  \begin{scope}[xshift=11.2cm]
    \draw[thick,defcol] (0.4,3.6) -- (4.9,1.4);
  \end{scope}
  \draw[-{Stealth[length=5pt]},very thick,keycol] (6.9,2.8) -- (7.8,2.8);
  \draw[-{Stealth[length=5pt]},very thick,keycol] (14.9,2.8) -- (15.8,2.8);
\end{tikzpicture}
\caption{The transmission mechanism as three dealer markets \ts{12}{5}{0:28}\ts{12}{5}{3:43}: the Fed sets the overnight rate within its outside spread; the term rate must exceed it
by enough to pay dealers for liquidity risk; bond prices depend on the cost of funding dealers' inventories. Raising the Fed funds target raises term rates and lowers bond prices.}
\label{fig:lr-transmission}
\end{nfigure}

The story \ts{12}{5}{4:04}\ts{12}{5}{6:29}\ts{12}{5}{6:49}\ts{12}{5}{7:11}:
\begin{enumerate}
  \item the Fed sets a target Fed funds rate and enforces it by trading (temporary open market operations); the level of the target is a choice between leaning toward
        discipline and leaning toward elasticity, and the same rate can mean elasticity in some circumstances and discipline in others \ts{12}{5}{4:25}\ts{12}{5}{4:46}\ts{12}{5}{5:07};
  \item raising the target (and with it the discount rate) narrows the gap between Fed funds and the term rate, which dealers do not like: they raise the term rate;
  \item funding becomes more expensive, so bond dealers want smaller inventories: downward pressure on bond prices;
  \item asset prices, mortgage-backed securities among them, affect the price of loans and so the real economy.
\end{enumerate}
This works through the price of money and the price of capital, and instantly, through arbitrage in securities markets; it does not need banks lending multiples of their
reserves (the money multiplier) or a bank-lending channel: institutions that no longer exist \ts{12}{5}{7:31}\ts{12}{5}{7:51}\ts{12}{5}{8:12}\ts{12}{5}{8:33}. The Fed may not move prices
exactly as it wants, because the market watches and anticipates it; but ``monetary policy does matter'' in a deregulated system. What it \emph{should} do has to be rethought
\ts{12}{5}{8:53}\ts{12}{5}{9:14}\ts{12}{5}{9:34}.

% ------------------------------------------------------------------
\section{Anatomy of a normal crisis}\label{sec:lr-normalcrisis}

A ``normal crisis'' is the kind of imbalance that can happen every day \ts{12}{6}{0:07}. Households (non-banks) decide, for whatever reason, maybe only for a day, that they
want to hold money instead of securities \ts{12}{6}{0:28}\ts{12}{6}{0:49}. The dealer system exists to absorb imbalances in order flow \ts{12}{6}{1:09}:
\begin{taccts}
\tacct[2.0cm]{Households}{$-$Securities\\$+$Deposits}{}\quad
\tacct[2.0cm]{Dealers}{$+$Securities}{$+$Repo (from bank)}\quad
\tacct[2.0cm]{Banks}{$+$Repo (to dealer)}{$+$Deposits}
\end{taccts}
The household sector sells securities on net; dealers buy them and finance them with repo from banks, which lend by expanding their balance sheets
\ts{12}{6}{1:32}\ts{12}{6}{1:53}\ts{12}{6}{2:14}\ts{12}{6}{2:34}. ``Sort of magical'': households think they sold to other households, but the securities were absorbed by the dealer and
banking system, and their deposits are \emph{new} deposits, indistinguishable from old ones \ts{12}{6}{2:55}\ts{12}{6}{3:16}.

\paragraph{The price of absorption.} Dealers and banks do it to make money \ts{12}{6}{3:36}. The dealer, seeing the imbalance in its order book, buys at a lower price; when
flows reverse, it sells higher. The bank lends at a slightly higher rate, and later liquidates at a lower one \ts{12}{6}{3:57}\ts{12}{6}{4:17}\ts{12}{6}{4:38}\ts{12}{6}{4:58}\ts{12}{6}{5:20}.
The households pay without seeing it, through a lower price for their bonds \ts{12}{6}{5:40}\ts{12}{6}{6:00}\ts{12}{6}{6:21}. These price distortions away from fundamental value are
possibly \emph{equilibrating}: low prices attract other buyers, the positions are liquidated \ts{12}{6}{6:43}\ts{12}{6}{7:03}. Even a seller with no information, who just needs
cash for a payment, moves the price, because the buyer fears he knows something; it comes back once the market realises he does not \ts{12}{6}{7:23}\ts{12}{6}{7:43}\ts{12}{6}{8:04}.

\section{Anatomy of a serious crisis}\label{sec:lr-seriouscrisis}

In a normal crisis there is no Fed: the private system absorbs it and the Fed funds rate stays put \ts{12}{7}{0:07}. A big shock makes dealers and banks anxious; that is
when the Fed comes in \ts{12}{7}{0:28}\ts{12}{7}{0:48}. As in Bagehot's world, the central bank sets a rate above the market rate; if the market on its own would push rates above
it, the central bank becomes lender of last resort \ts{12}{7}{1:10}\ts{12}{7}{1:32}:
\begin{taccts}
\tacct[2.0cm]{Fed}{$+$Loans to banks}{$+$Reserves}\quad
\tacct[2.0cm]{Banks}{$+$Repo to dealers\\$+$Reserves}{$+$Deposits\\$+$Borrowing from Fed}\quad
\tacct[2.0cm]{Dealers}{$+$Securities}{$+$Repo}
\end{taccts}
If the Fed backstops the banks, the banks are more willing to backstop the dealers; this keeps a large shock within bounds \ts{12}{7}{1:52}\ts{12}{7}{2:23}\ts{12}{7}{2:43}.

\begin{proposition}[Why the central bank can stop a crisis]
In a crisis everyone wants money and nobody wants securities; the securities must go somewhere. The lender of last resort works by letting the unwanted securities be absorbed
into the financial sector and creating the cash people want to fund them: expansion of dealers, of banks, ultimately of the Fed. The Fed backstops the banks' short positions
in money, i.e.\ their liquidity risk, and it cannot itself run into trouble with short positions in money, because its liability is the ultimate money, at least domestically:
it operates as a money market dealer \ts{12}{7}{3:04}\ts{12}{7}{3:24}\ts{12}{7}{3:45}\ts{12}{7}{4:06}\ts{12}{7}{4:26}.
\end{proposition}

% ------------------------------------------------------------------
\section{Should the Fed intervene or not?}\label{sec:lr-should}

Private dealers absorb fluctuations for profit; the Fed's motive is different: when should it take the pressure off, and when let it ride? The balance between discipline and
elasticity \ts{12}{8}{0:09}\ts{12}{8}{0:29}\ts{12}{8}{0:50}. The Fed can always provide elasticity, but that does not mean it should \ts{12}{8}{1:11}.

\paragraph{For elasticity.} Prices may become so distorted that people make bad decisions; or distortions may stop being equilibrating: falling asset prices make assets
unusable as collateral, holders dump them, a \emph{liquidity spiral}, and the system collapses \ts{12}{8}{1:11}\ts{12}{8}{1:32}\ts{12}{8}{4:16}\ts{12}{8}{4:36}.

\paragraph{The value investors.} A student asks: why would value-based investors not step in? In this picture the Fed backstops banks, not dealers (dealer of last resort in
the money market) \ts{12}{8}{2:13}\ts{12}{8}{2:33}. It may intervene directly in asset markets, as it did, because the value investors did not: they refused, or lost their shirts,
or wait for prices of 20 cents on the dollar, which would come only with a great depression; the Fed is not there to make money for value investors, it is concerned with
unemployment. If value investors do step in, the Fed need not \ts{12}{8}{2:54}\ts{12}{8}{3:14}\ts{12}{8}{3:34}\ts{12}{8}{3:55}.

\paragraph{For discipline.} Why not always be elastic? Minsky: what moves money market rates is the mismatch between cash commitments and cash flows; if rates never moved,
nothing would stop agents from pushing commitments into the future forever. Discipline makes mistakes get recognised: at some point you pay or go bankrupt
\ts{12}{8}{4:56}\ts{12}{8}{5:16}\ts{12}{8}{5:37}\ts{12}{8}{5:57}. Otherwise mistakes build up like dry tinder, and ``Smokey the Bear'' produces larger forest fires
\ts{12}{8}{6:18}. In stress people are forced to sell at prices they would never accept, and so to recognise their mistakes and start afresh; but too much discipline ``kills the
patient'' \ts{12}{8}{6:59}\ts{12}{8}{7:20}.

\begin{intuition}[The art of central banking]
Central bankers ``feel this in their bones'': not every hiccup, not even every significant stress, deserves elasticity. In this crisis the Fed did not flood the market at once:
it let some bad mortgages ``walk'' while preventing collapse; then things got worse than expected and it did more, and more \ts{12}{8}{6:38}\ts{12}{8}{7:40}\ts{12}{8}{8:00}\ts{12}{8}{8:21}.
\end{intuition}

% ------------------------------------------------------------------
\section{The Fed as dealer of last resort, 2007--2009}\label{sec:lr-dolr}

The Fed's balance sheet in the crisis (lecture~3) shows the stages \ts{12}{9}{0:07}\ts{12}{9}{0:28}:
\begin{enumerate}
  \item \textbf{Interest rates} (August 2007 to Bear Stearns): Fed funds from 5\% to 2\%, to make it more profitable for banks to take on liquidity risk and support asset prices
        indirectly \ts{12}{9}{0:49}\ts{12}{9}{1:09};
  \item \textbf{composition} (after Bear Stearns): of about a trillion of Treasury bills, \$500 billion sold and the proceeds lent to brokers and dealers \ts{12}{9}{1:09}\ts{12}{9}{1:30};
  \item \textbf{expansion} (after Lehman and AIG): about two trillion more of deposits and of lending; the interbank market broke down and ``the Fed became the interbank market'':
        surplus institutions deposited at the Fed, deficit institutions borrowed from it \ts{12}{9}{1:30}\ts{12}{9}{1:51};
  \item \textbf{capital market}: purchases of mortgage-backed securities \ts{12}{9}{2:11}.
\end{enumerate}

\paragraph{The crisis in balance sheets} \ts{12}{9}{2:32}--\ts{12}{9}{10:15}:
\begin{enumerate}
  \item \textbf{Before}: Citi's SIV holds RMBS funded by ABCP, held by a money fund that issues shares to (international) investors; nothing is on Citi's or the Fed's balance sheet.
  \item \textbf{First stage}, which looked like a normal crisis: the money fund declines to roll the ABCP; Citi replaces the funding with a loan to the SIV, funded on its own
        balance sheet, say with repo. The fund now holds a liability of Citibank with recourse to its whole balance sheet. Citi did this as a business: take assets off customers who
        no longer want them, and charge for it \ts{12}{9}{3:56}\ts{12}{9}{4:17}\ts{12}{9}{4:38}\ts{12}{9}{4:58}\ts{12}{9}{5:19}.
  \item \textbf{Second stage}: the fund prefers Citi's unsecured commercial paper, based on its whole book (Lehman too issued such paper, held by the Reserve Primary Fund), or
        Eurodollar deposits at foreign banks \ts{12}{9}{5:39}\ts{12}{9}{6:00}\ts{12}{9}{6:20}\ts{12}{9}{7:00}. The shadow banking system collapses onto the traditional banking system,
        backstopped by the Fed and the FDIC \ts{12}{9}{7:23}\ts{12}{9}{7:43}.
  \item \textbf{Third stage}: European banks, with little access to the Fed, cannot roll their dollar funding when money funds stop renewing Eurodollar deposits
        \ts{12}{9}{8:04}\ts{12}{9}{8:25}. The Fed lends dollars to foreign central banks through liquidity swaps (about \$600 billion at the peak), and the Treasury issues bills
        through a supplementary account at the Fed, so that money funds can hold what they want, Treasury bills \ts{12}{9}{8:48}\ts{12}{9}{9:10}\ts{12}{9}{9:34}\ts{12}{9}{9:55}\ts{12}{9}{10:15}.
\end{enumerate}
\begin{taccts}
\tacct[2.0cm]{Fed $+$ Treasury}{$+$Swap (loan to ECB)}{$+$Treasury bills}\quad
\tacct[2.0cm]{ECB}{$+$Dollar loans to European banks}{$+$Swap (dollars owed to Fed)}\quad
\tacct[2.0cm]{Money fund}{$-$Eurodollar deposits\\$+$Treasury bills}{}
\end{taccts}

\paragraph{From the best money down the hierarchy.} The sequence: first the normal absorption, letting prices move to give the banking system an incentive; then lower Fed funds;
then direct entry into the term funding market (liquidity swaps, lending to dealers, the term facilities), backstopping banks; finally, as value investor, the mortgage-backed
securities market \ts{12}{9}{10:35}\ts{12}{9}{10:56}\ts{12}{9}{11:16}\ts{12}{9}{11:37}. The Fed moved down the hierarchy, putting more and more of the financial system on its own
balance sheet. At no point did it want to: central bankers are conservative, they do not want credit risk or interest rate risk, which is why their money is the best money; but in
a crisis, ``lend freely at a high rate against what would be good security in a normal time'' \ts{12}{9}{11:57}\ts{12}{9}{12:18}.

\begin{intuition}[A floor under the system]
Normally the Fed acts at a distance, setting the overnight rate and letting profit-seeking dealers and arbitrage transmit it \ts{12}{9}{12:39}\ts{12}{9}{12:59}. When imbalances built up
over years are too large, the crisis is not temporary and keeps getting worse, and the Fed's responsibility as ultimate backstop forces it to do more, until it has put a floor under
the system. It can, because its liabilities are the best money: ultimately it could buy every risky security for cash, though that would be the end of the market system
\ts{12}{9}{13:19}\ts{12}{9}{13:40}\ts{12}{9}{14:01}\ts{12}{9}{14:21}. It always hopes to do just enough; the ECB is doing that now, hoping the pressure will make politicians act: perhaps
``a slow-motion collapse'' of the European system, in which long periods of nothing alternate with bursts of excitement \ts{12}{9}{14:41}\ts{12}{9}{15:02}\ts{12}{9}{15:23}\ts{12}{9}{15:43}.
\end{intuition}

\begin{definition}[New concepts: dealer of last resort]
A central bank acts as \emph{dealer of last resort} when, instead of only lending to banks against collateral (lender of last resort), it stands ready to buy and sell, or to
backstop the prices of, securities that private dealers will no longer absorb, putting a floor under their prices.
\end{definition}

% ------------------------------------------------------------------
\flushside
\section*{Key formulas and ideas}
\begin{keyformulas}
\begin{tabularx}{\linewidth}{@{}l L@{}}
Quantity vs price & controlling the best money (top of the hierarchy) $\equiv$ pegging the short end of the term structure\\
Taylor rule & $r=\rho+\pi^e+\alpha(\pi^e-\pi^\ast)+\beta(y-y_f)$; Fisher effect $r=\rho+\pi^e$\\
Transmission & Fed funds (settlement risk) $\to$ term rate (liquidity risk) $\to$ bond prices (price risk) $\to$ economy\\
Normal crisis & dealers and banks absorb order imbalances for a price (temporary distortions)\\
Serious crisis & Fed backstops banks, banks backstop dealers; the Fed cannot fail on short positions in its own money\\
Intervene? & elasticity against liquidity spirals; discipline to make mistakes recognised\\
2007--2009 & rates $\to$ composition $\to$ expansion (Fed $=$ interbank market) $\to$ swaps, facilities $\to$ MBS: dealer of last resort\\
\end{tabularx}
\end{keyformulas}

\section*{Exercises}

\begin{exercise}[Taylor rule]
With $\rho=2\%$, $\pi^\ast=2\%$, $\alpha=0.5$, $\beta=0.5$:
\begin{enumerate}[(a)]
  \item compute the prescribed rate for $\pi^e=3\%$ and an output gap of $-2\%$;
  \item what rate does the rule prescribe for $\pi^e=1\%$ and a gap of $-6\%$? What happens at the zero lower bound?
\end{enumerate}
\end{exercise}

\begin{exercise}[Transmission]
Using the three diagrams of \cref{fig:lr-transmission}, explain step by step what happens to the three-month rate and to bond prices when the Fed raises its target by 25 basis points,
and why the effect may be smaller if the market had already anticipated it.
\end{exercise}

\begin{exercise}[A normal crisis]
Households sell \$100 of bonds. Dealers buy them at a price 1\% below fundamental value, financed by repo at the term rate. Three days later the flow reverses and dealers sell at
fundamental value. Write the T-accounts at each step and compute the dealers' profit if repo costs 0.5\% a year.
\end{exercise}

\begin{exercise}[The SIV unwinds]
An SIV holds RMBS of 100 funded by ABCP of 100 held by a money fund. Write the T-accounts (SIV, Citi, money fund, Fed) for each stage of \cref{sec:lr-dolr}, ending with the money fund
holding Treasury bills.
\end{exercise}


% Part II: lectures 13--22
\part{Advanced: foreign exchange and capital markets}\label{part:advanced}
% !TEX root = ../main.tex
\chapter{Chartalism, Metallism and Key Currencies}\label{ch:chartalism}
\begin{paracol}{2}

\begin{sources}
\begin{tabularx}{\linewidth}{@{}l l L@{}}
\toprule
\textbf{Segment} & \textbf{Length} & \textbf{Content}\\
\midrule
\seg{13}{1}{L13.1 FT: Autonomy of Bank of Japan} & 2:14 & A timid Bank of Japan, a stronger yen.\\
\seg{13}{2}{L13.2 Key Currencies as a Hierarchical System} & 8:38 & Quoting conventions; the hierarchy of currencies.\\
\seg{13}{3}{L13.3 What is Money? Chartalism versus Metallism} & 8:37 & Schumpeter's two traditions; king's money and international money.\\
\seg{13}{4}{L13.4 Chartalism as a Theory of Money} & 2:10 & Currency as zero-interest government funding; taxes.\\
\seg{13}{5}{L13.5 Quantity Theory of Money} & 4:39 & $MV=PT$ as a way to tame the power of the state; velocity and bank money.\\
\seg{13}{6}{L13.6 Purchasing Power Parity} & 3:02 & The chartalist theory of the exchange rate.\\
\seg{13}{7}{L13.7 Metallism as a Theory of Money} & 5:07 & Mint pars, gold points, forward parity.\\
\seg{13}{8}{L13.8 A Money View of International Payments, FX Dealers} & 11:36 & A deficit country pays in dollars it does not have: matched-book and speculative FX dealers.\\
\seg{13}{9}{L13.9 Chartalism, Metallism, and the Money View Compared} & 4:29 & Goods, assets, money: three readings of the exchange rate.\\
\seg{13}{10}{L13.10 Private and Public Money: A Hybrid System} & 7:30 & How war finance merged the king's money and the bankers' money.\\
\seg{13}{11}{L13.11 Hybridity in FX Market-making} & 5:33 & Private dealers and central banks along the currency hierarchy.\\
\bottomrule
\end{tabularx}

\smallskip
\textbf{Files.} Transcripts \file{materials/transcripts/L13-P*.txt}.
\end{sources}

\section*{Motivation}
The second part of the course begins with four lectures on the \emph{exchange rate}, one of the four prices of money \ts{13}{1}{0:07}. Foreign exchange is
the relative price of monies, so it is a \emph{money} market, not a capital market \ts{13}{2}{2:05}. To talk about the price of one money in terms of another
we must face the question avoided so far, by focusing on credit: what kind of thing is money? \ts{13}{2}{3:07} Two traditions answer it,
chartalism and metallism; each yields a theory of the exchange rate; the money view offers a third, and the modern system turns out to be a hybrid of the
two kinds of money.

% ------------------------------------------------------------------
\section{FT: autonomy of the Bank of Japan}\label{sec:cm-ft}

\begin{casestudy}[FT, October 2012: ``Autonomy of BoJ called into question'']
The Bank of Japan was less aggressively expansionary than the market expected, and the yen rose about 0.5\% against the dollar within 30 minutes of the
announcement: bad for a big export economy that has long struggled with deflation \ts{13}{1}{0:29}. ``Set alongside the
spectacular easing of the US Federal Reserve or the European Central Bank, the more subdued Bank of Japan often struggles to compete''; and if prices keep
falling, raising the real cost of borrowing, executives will not lever up \ts{13}{1}{1:31}. The article links central bank balance sheets to the price of
money relative to foreign monies \ts{13}{1}{1:53}.
\end{casestudy}

% ------------------------------------------------------------------
\section{Key currencies as a hierarchical system}\label{sec:cm-hierarchy}

\paragraph{Quoting conventions.} A FRED chart shows the dollar price of gold, rising from 2003 to about \$1800 an ounce, and the dollar price of a euro, about
\$1.30 \ts{13}{2}{0:00}. There are two ways to quote an exchange rate \ts{13}{2}{1:01}:

\begin{definition}[New concepts: quoting conventions]
In the \emph{American convention} the exchange rate is the number of dollars per unit of foreign currency (\$1.30 per euro); in the \emph{European convention}
it is the number of units of foreign currency per dollar (0.77 euro per dollar). The theory is the same, but every formula changes into its inverse
\ts{13}{2}{1:01}. In this course, so far, the European convention has been used.
\end{definition}

\paragraph{The international hierarchy.} So far the best money was reserves at the Fed (and currency), then Fed funds, then Eurodollars, all promises to pay
reserves and trading at about the same price \ts{13}{2}{2:26}. What lies above? A stylised description of what one sees
\ts{13}{2}{4:08}:
\begin{enumerate}
  \item the \textbf{international dollar}, the world reserve currency, in which most trade is denominated and settled \ts{13}{2}{4:28};
  \item the \textbf{domestic dollar}, at par with it (with difficulty in the crisis, as seen) \ts{13}{2}{4:55};
  \item \textbf{key currencies}: the yen in Asia (eventually perhaps the renminbi), the pound and the euro in Europe: with the dollar the four most traded
        \ts{13}{2}{5:15};
  \item other \textbf{majors}: the Swiss franc, the Canadian and Australian dollars: deep, liquid markets against the dollar, where very large trades do not move
        the price \ts{13}{2}{5:35};
  \item the \textbf{minor} currencies of the hundreds of other countries, which do not trade against each other: to trade Hungarian forints against Mexican pesos
        one goes through dollars \ts{13}{2}{7:20}.
\end{enumerate}
In trading volume (as quoted in class), 51\% involves only the dollar, euro, yen and sterling, and 85\% of trades have the dollar on one side \ts{13}{2}{6:39}.
``It really is a dollar system.'' Being at the centre makes it hard for Americans to think about foreign exchange, which Europeans think about all the time
\ts{13}{2}{8:00}.

\begin{nfigure}
\begin{tikzpicture}[font=\scriptsize]
  \fill[keycol!30] (0,4) -- (-0.7,3.2) -- (0.7,3.2) -- cycle;
  \fill[defcol!25] (-0.7,3.2) -- (-1.5,2.2) -- (1.5,2.2) -- (0.7,3.2) -- cycle;
  \fill[intcol!20] (-1.5,2.2) -- (-2.3,1.2) -- (2.3,1.2) -- (1.5,2.2) -- cycle;
  \fill[fincol!15] (-2.3,1.2) -- (-3.2,0) -- (3.2,0) -- (2.3,1.2) -- cycle;
  \draw[thick] (0,4) -- (-3.2,0) -- (3.2,0) -- cycle;
  \node[align=center] at (0,3.45) {int'l\\dollar};
  \node at (0,2.65) {domestic \$, euro, yen, pound};
  \node at (0,1.65) {Swiss franc, CAD, AUD};
  \node at (0,0.55) {minor currencies};
  \node[right,align=left] at (3.4,3.4) {world reserve currency:\\85\% of trades have \$ on one side};
  \node[right,align=left] at (3.4,1.7) {majors: deep, liquid\\markets against the dollar};
  \node[right,align=left] at (3.4,0.5) {traded through the dollar;\\markets made by central banks};
\end{tikzpicture}
\caption{The hierarchy of currencies \ts{13}{2}{4:28} (a stylised description, not a theory).}\label{fig:cm-currencies}
\end{nfigure}

% ------------------------------------------------------------------
\section{What is money? Chartalism versus metallism}\label{sec:cm-what}

The two words come from Joseph Schumpeter's treatise on the history of economic analysis, which classifies monetary theories into chartalist and metallist
\ts{13}{3}{0:20}.

\begin{definition}[New concepts: chartalism, metallism]\label{def:cm-chartalism}
\emph{Chartalism}: a piece of money is an emblem of the power of the state that issues it; the quintessential money is fiat currency, ``the king's money''. Its
classic theorist is Georg Friedrich Knapp \ts{13}{3}{0:41}. \emph{Metallism}: the essence of money is a physical, metallic value,
gold or silver \ts{13}{3}{1:43}. (Chartalism from Latin \emph{charta}, a token or ticket, the state-stamped token; \emph{fiat}, ``let it be done''.)
\end{definition}\begin{history}[Knapp and Schumpeter]
Georg Friedrich Knapp's \emph{Staatliche Theorie des Geldes} (1905; English \emph{The State Theory of Money}, 1924) argued that money is a creature of law,
defined by what the state accepts in payment. Schumpeter's posthumous \emph{History of Economic Analysis} (1954) contrasts ``theoretical metallism'' with
``theoretical cartalism''.\par
\emph{References:} G.~F.~Knapp, \emph{The State Theory of Money}, Macmillan, 1924; J.~A.~Schumpeter, \emph{History of Economic Analysis}, Oxford UP, 1954, part~II ch.~6.
\end{history}

Both traditions allow a large superstructure of credit, so most of the course is consistent with either; but the debate can get very heated
\ts{13}{3}{2:04}. For metallists inconvertible money is a kind of robbery, unanchored, anxiety-inducing; for chartalists the state is the ultimate
authority, and who is gold to tell it what to do? It is discipline versus elasticity again: state money against the discipline of gold and of the market
\ts{13}{3}{2:45}.

\paragraph{Two moneys in history.} When two traditions are at loggerheads forever, look for a higher point of view, or for a historical origin; here both apply
\ts{13}{3}{3:45}. Before the modern nation-state there were two parallel monetary systems in the same country \ts{13}{3}{4:06}:
\begin{itemize}
  \item the \textbf{king's money}: local, domestic, used in retail trade and to pay taxes; the king's face on it matters if he is your king;
  \item \textbf{international money}: gold, which you simply weigh, used in wholesale and international business.
\end{itemize}
Both circulated, with an exchange rate between them. Mehrling learned this from Fernand Braudel, \emph{The Wheels of Commerce}, volume 2 of his trilogy on the
15th to 18th centuries, which describes the wheels of trade at the lowest and at the highest level \ts{13}{3}{5:29}. It is no surprise that there are
two traditions in monetary theory: there were always two kinds of money \ts{13}{3}{6:10}.

\begin{taccts}
\tacct[2.4cm]{King's money system}{}{King's money\\\quad domestic credit superstructure}\qquad
\tacct[2.4cm]{International money system}{}{Gold\\\quad business credit superstructure}
\end{taccts}
Each money is at the top of its own hierarchy \ts{13}{3}{6:31}. The systems touch: merchants who trade internationally and sell retail move between them,
and kings fighting foreign wars need gold to pay soldiers and buy war material abroad \ts{13}{3}{7:13}. Today there are many hierarchies: each currency has its
domestic credit hierarchy, and the international dollar a whole international one; the world funding markets are dollar markets \ts{13}{3}{8:15}.

% ------------------------------------------------------------------
\section{Chartalism: the state, the quantity theory, PPP}\label{sec:cm-chartalist}

\paragraph{Currency as cheap funding.} If money is a creature of the state, think of the central bank as a government bank: the Treasury issues bills paying, say, 5\%;
the central bank buys them and issues currency paying 0\%. Looking through the central bank, currency is zero-interest funding for the government, which is why states
like to monopolise its issue \ts{13}{4}{0:00}.
\begin{taccts}
\tacct[2.2cm]{Treasury}{(authority to tax)}{Treasury bills (5\%)}\qquad
\tacct[2.2cm]{Central bank}{Treasury bills (5\%)}{Currency (0\%)}
\end{taccts}
What supports the value of such a liability? The government's authority, especially its taxing authority: it insists that taxes be paid in its currency, and if you
do not pay, it throws you in jail \ts{13}{4}{1:26}.

\paragraph{The quantity theory.} Economists turn this into a proper theory with the quantity theory of money \ts{13}{5}{0:00}:
\begin{keyformulas}
\textbf{Quantity theory} \ts{13}{5}{0:21}: $MV = PT$ (money $\times$ velocity $=$ price level $\times$ transactions). With $V$ fixed by institutions and $T$ by the
real economy, $M$ determines $P$ \ts{13}{5}{0:43}.

\medskip
\textbf{Purchasing power parity} \ts{13}{6}{0:23}: $E\,P = P^\ast$, with $E$ the exchange rate (foreign currency per dollar), $P$ and $P^\ast$ the domestic
and foreign price levels: the exchange rate is the relative price of goods in the two countries.
\end{keyformulas}
The quantity theory tells the government: you can say what money is, but not what it is worth; print too much and it is inflated away \ts{13}{5}{1:25}. People get
heated about it because of a deep anxiety about the power of the state and about green pieces of paper that could become worthless; the equation is ``the way economists
regulate their anxiety'' \ts{13}{5}{2:06}. But the government is not the only issuer: if quantity causes prices, bank money should worry us too; it
seems to worry people less because it looks like credit \ts{13}{5}{2:48}. And since the hierarchy expands and contracts, velocity is not constant: in an
expansion private money grows with transactions even if government money does not, which in the equation appears as rising $V$ \ts{13}{5}{3:28}.

\paragraph{PPP.} If money is fiat and its quantity sets the price level, the natural theory of the exchange rate is \emph{purchasing power parity}: the quantity theory works
behind the scenes in each country, and the relative price of monies is just the relative price of goods \ts{13}{6}{0:02}. Two prices of money, the
exchange rate and the price level, connected, and no interest rate anywhere \ts{13}{6}{1:27}. It suits economists who look through the veil of money to a barter economy of
goods for goods \ts{13}{6}{1:50}. But it does not hold empirically except perhaps in the long run: price levels move slowly, exchange rates ``bop all around''
\ts{13}{6}{2:30}.

% ------------------------------------------------------------------
\section{Metallism: mint pars and gold points}\label{sec:cm-metallist}

If money is gold, the exchange rate depends only on how much gold each money contains \ts{13}{7}{0:03}. If the dollar is $x$ ounces of gold and the pound $y$
ounces, the exchange rate is the ratio of the \emph{mint pars}, $x/y$ \ts{13}{7}{0:45}. Arbitrage enforces it: if the rate deviates, melt the heavier coin into the
lighter one \ts{13}{7}{1:29}. But that means shipping gold, which is costly, so the rate can deviate within a band \ts{13}{7}{2:10}.

\begin{definition}[New concepts: mint par, gold points]
The \emph{mint par} is the quantity of gold a currency unit represents at the mint; the ratio of two mint pars is the \emph{mint parity} exchange rate. The \emph{gold
points} are the exchange rates above and below mint parity at which it pays to ship gold, given the costs of shipping and insurance: outside them arbitrage moves gold,
inside them it does not \ts{13}{7}{2:54}.
\end{definition}

Mint parity with gold points is a finance view: the forward parity condition of lecture~8 (in class here, with the spot rate $E$ related to the forward rate by the ratio of
the two interest rates) holds under the gold standard too, because the gold points bound the spot rate between a floor and a cap \ts{13}{7}{3:20}.
A floor and a cap suggest Treynor's outside spread, with dealers trading inside it and determining the exchange rate: that is where the money view enters
\ts{13}{7}{4:24}.

\begin{history}[The classical gold standard]
Under the classical gold standard (roughly 1880--1914) the dollar was defined as 23.22 grains of pure gold and the pound as 113.0016 grains, giving a mint parity of
about \$4.8665 per pound; the gold points between New York and London were typically within about 0.5--1\% of parity.\par
\emph{References:} L.~Officer, \emph{Between the Dollar--Sterling Gold Points}, Cambridge UP, 1996.
\end{history}

% ------------------------------------------------------------------
\section{A money view of international payments}\label{sec:cm-fxdealers}

A \emph{deficit} country has bought from a \emph{surplus} country and owes \$10, in dollars, the international reserve currency (neither is the US)
\ts{13}{8}{0:03}. If it has \$10, no problem; but suppose it does not \ts{13}{8}{1:25}. The point of a payment system is that willing buyers
and sellers are not stopped by not having the means of payment at hand; domestically the answer was credit; internationally it is harder \ts{13}{8}{1:47}.

\paragraph{The FX dealer.} A foreign exchange dealer in the middle buys $10E$ of the deficit country's currency and sells it \$10 spot, which settles the debt: the surplus
country has been paid, the deficit country thinks it is off the hook \ts{13}{8}{2:49}. The dealer made the
payment, with dollars ``created from nothing'', promises to pay dollars: a short position in dollars against a long position in foreign exchange, just as a bank makes the
payments its customers cannot \ts{13}{8}{5:00}.

\paragraph{Hedging in the term market.} The dealer is exposed to the exchange rate, and to the surplus country demanding currency for its dollar claim \ts{13}{8}{6:33}. It hedges by
trading term deposits with another dealer: a \$10 term deposit asset and a $10E$ foreign-currency term deposit liability; whatever the exchange rate does, it gains on one side
what it loses on the other \ts{13}{8}{6:54}. This uses forward interest parity: in effect it hedges its spot position at the forward rate. It is now a
\emph{matched-book} FX dealer \ts{13}{8}{7:55}. Someone must take the other side of those term deposits: a \emph{speculative} dealer, exposed to exchange risk and
compensated for it \ts{13}{8}{8:58}.
\begin{taccts}
\tacct[1.9cm]{Deficit country}{$-$FX $10E$}{$-$Due to surplus \$10}\quad
\tacct[1.9cm]{Matched-book FX dealer}{FX $10E$ (spot)\\\$10 term}{\$10 spot\\FX $10E$ term}\quad
\tacct[1.9cm]{Speculative FX dealer}{FX $10E$ term}{\$10 term}

\medskip
\tacct[1.9cm]{Surplus country}{$-$Due from \$10\\$+$\$10 spot}{}
\end{taccts}
Foreign exchange trading is about \$4 trillion a day, one of the largest markets \ts{13}{8}{8:37}. Every asset is someone's liability, but the functions differ: the matched-book dealer
bears another kind of risk (later), and can be called a forward-interest-parity dealer; the speculative dealer, recognisable from Treynor's model, bears exchange risk and must expect a
profit, which is what moves the forward rate away from the expected spot rate: an uncovered-interest-parity dealer \ts{13}{8}{9:39}.

\begin{definition}[New concepts: spot and forward FX, matched-book and speculative FX dealers]
A \emph{spot} foreign exchange transaction is settled now (in practice in two business days); a \emph{forward} one at a future date at a rate fixed today. A
\emph{matched-book} FX dealer balances its currency positions (spot against term), bearing no exchange-rate risk; a \emph{speculative} FX dealer holds an open position in a
currency and bears the exchange risk.
\end{definition}

\section{Three readings of the exchange rate}\label{sec:cm-compared}

\begin{nfigure}
\renewcommand{\arraystretch}{1.35}
\begin{tabularx}{\linewidth}{l|L|L}
\textbf{Tradition} & \textbf{The exchange rate is} & \textbf{Theory}\\\hline
Chartalism & relative price of goods & purchasing power parity; trade (current) account\\
Metallism (finance) & relative price of assets & interest parity; capital account\\
Money view & price of money in terms of money & dealers' balance sheets\\
\end{tabularx}
\caption{What kind of a thing is an exchange rate? \ts{13}{9}{0:04}}\label{fig:cm-three}
\end{nfigure}

The chartalist road leads to the exchange rate as a relative price of tradeable goods (PPP); the metallist road to a relative price of financial assets, interest rates. Some theories
stress the trade or current account, others the capital account \ts{13}{9}{0:04}. The money view sees the exchange rate as the price of money in terms of
money, on the dealers' balance sheets \ts{13}{9}{1:06}. Economics abstracts from liquidity to price commodities, finance to price assets; the course is about liquidity, so it cannot.
PPP and interest parity are not wrong, but they abstract from liquidity, so they cannot be fully right \ts{13}{9}{1:27}.

\begin{intuition}[Pushing exchange risk into the future]
In the example nobody lends: not the surplus country, not the FX dealer; they swap currencies. The dealer takes risk in the spot market and hedges it in the forward market, where
someone else takes it: exchange rate risk located today is pushed into the future, exactly as the money market pushes a payment one cannot make today to tomorrow or 90 days. The
speculative dealer's position is a \emph{carry trade}, and it is compensated because the payment system needs it \ts{13}{9}{2:08}.
\end{intuition}

% ------------------------------------------------------------------
\section{Private and public money: a hybrid system}\label{sec:cm-hybrid}

How did the two parallel systems of the 15th century become today's? Mostly through war \ts{13}{10}{0:05}. Start with two separate systems
\ts{13}{10}{1:34}:
\begin{taccts}
\tacct[2.4cm]{Government bank}{Treasury bills}{King's money (currency)}\qquad
\tacct[2.4cm]{Private bankers' bank}{Gold\\Private credit}{Deposits (promises to pay gold)}
\end{taccts}
The government bank needs no reserves: the king can always print more pieces of paper with his head on them, a cheap way to fund himself \ts{13}{10}{2:36}. Then comes
war, and, as in the Civil War (\cref{sec:st-civilwar}), the Treasury asks the private bank for a loan, which it makes by creating deposits, and then withdraws them in gold
\ts{13}{10}{3:17}:
\begin{taccts}
\tacct[2.4cm]{Government}{$+$Gold}{$+$Loan from bankers' bank}\qquad
\tacct[2.4cm]{Private bankers' bank}{$+$Loan to government\\$-$Gold}{($+$Deposits, then $-$Deposits withdrawn in gold)}
\end{taccts}
``Here's hybridity'': the government bank has all the gold, and the bankers' bank starts using domestic currency as its reserve \ts{13}{10}{4:35}. It happened in the
Bank of England, in the US, in most countries: the state used the private and international bank for war finance, twisting its arm \ts{13}{10}{5:17}.

\begin{proposition}[Why hybridity lasted]
After the war both sides kept the arrangement, ``a marriage made in heaven'': the state blessed private money, giving bank deposits legal protection as money (political cover for
private money issuers); the private market accepted Treasury bills as promises to pay money, giving the state access to capital markets. A government that funds itself by printing
money cannot fund much; one that sells Treasury bills to the private sector can fund a huge war. Modern monetary systems are always hybrid \ts{13}{10}{5:37}.
\end{proposition}

\section{Hybridity in FX market-making}\label{sec:cm-fxhybrid}

Are the FX dealers private banks or central banks? ``It depends'' on the place in the currency hierarchy \ts{13}{11}{0:03}:
\begin{itemize}
  \item in the dollar and the majors, a huge number of private dealers, highly liquid, with very thin margins;
  \item in the Hungarian forint, the central bank: there is no deep and liquid private market; the weaker the currency, the more it is all central bank \ts{13}{11}{1:05}.
\end{itemize}
So the theory of exchange rates must be about the relationship between private, profit-seeking FX dealers and the public, welfare-seeking central bank \ts{13}{11}{1:26}.
(Not that central banks are the good guys: in many countries the central bank is the agent of a rapacious dictator, and the private dealers who refuse to do business with it are
the good guys \ts{13}{11}{2:12}.) For some currencies the exchange rate is mostly politics; for others the central bank is a dealer of last resort in the wings
\ts{13}{11}{2:33}. Under the gold standard the outside spread was the gold points; today there are none, but central banks, by controlling short-term interest rates,
create an outside spread for the FX market, as the next lectures show \ts{13}{11}{3:14}.

\begin{finance}[Why so much FX trading?]
Foreign exchange is more ``mind-blowing'' than domestic money markets: all of a money market in one country, all of it in another, and a set of dealers in between
\ts{13}{11}{3:34}. Economists and finance professors simplify it to a relative price of goods or assets, and wonder why \$4 trillion a day is traded. But behind a
\$10 trade there are many trades ``below the line'' that separate the risk into categories and put each in the right place; today it is electronic and cheap, so there is a lot of
positioning without burning real resources \ts{13}{11}{4:16}.
\end{finance}

% ------------------------------------------------------------------
\flushside
\section*{Key formulas and ideas}
\begin{keyformulas}
\begin{tabularx}{\linewidth}{@{}l L@{}}
Conventions & American: \$ per foreign unit; European: foreign units per \$ (inverse formulas)\\
Currency hierarchy & international \$ $>$ domestic \$, euro, yen, pound $>$ other majors $>$ minors (via \$)\\
Chartalism & money $=$ state token; currency $=$ zero-interest funding; value from taxes; $MV=PT$; PPP $EP=P^\ast$\\
Metallism & money $=$ gold; exchange rate $=$ ratio of mint pars within gold points; interest parity\\
Money view & exchange rate $=$ price of money in money, set by FX dealers (matched-book and speculative)\\
Hybridity & war finance merged king's money and bankers' money; modern systems always hybrid\\
FX dealers & private in major currencies, central banks in weak ones; central banks set the outside spread via interest rates\\
\end{tabularx}
\end{keyformulas}

\section*{Exercises}

\begin{exercise}[Conventions]
The euro trades at \$1.30. Write the rate in the European convention. If the dollar depreciates by 10\% against the euro, what are the new quotes in both conventions? Rewrite
covered interest parity (lecture~8) in the American convention.
\end{exercise}

\begin{exercise}[Gold points]
Mint parity is \$4.8665 per pound; shipping gold from New York to London costs 0.5\% of its value. Between which exchange rates is gold not shipped? In which direction does gold
flow if the pound trades at \$4.92 in New York?
\end{exercise}

\begin{exercise}[PPP]
A basket costs 100 dollars in the US and 12{,}000 yen in Japan, and the yen trades at 80 per dollar. What does PPP predict? US inflation is 2\% and Japanese inflation 0\% for a
year, while the yen goes from 80 to 78 per dollar. Has PPP moved closer or further?
\end{exercise}

\begin{exercise}[A payment through FX dealers]
Redo the deficit-country example with numbers: $E=1.25$ units per dollar, six-month dollar rate 1\%, foreign rate 3\%. Write all the T-accounts, compute the forward rate the matched-book
dealer locks in, and state what the speculative dealer gains or loses if the spot rate in six months is 1.30.
\end{exercise}

% !TEX root = ../main.tex
\chapter{Money and the State: International}\label{ch:intl}
\begin{paracol}{2}

\begin{sources}
\begin{tabularx}{\linewidth}{@{}l l L@{}}
\toprule
\textbf{Segment} & \textbf{Length} & \textbf{Content}\\
\midrule
\seg{14}{1}{L14.1 FT: Costs of Japan's Monetary Policy} & 2:49 & Zombie firms and a strong yen.\\
\seg{14}{2}{L14.2 Reading: Robert Mundell} & 10:27 & Mundell's 20th century in the language of the hierarchy: central banks avoiding the discipline of gold.\\
\seg{14}{3}{L14.3 Confrontation} & 11:48 & Act one: the Fed and the gold standard, 1931.\\
\seg{14}{4}{L14.4 Contradiction} & 14:08 & Act two: national management versus fixed rates; Keynes's bancor plan versus the IMF.\\
\seg{14}{5}{L14.5 The Dollar System} & 7:04 & Kindleberger: the US as a bank to the world; the end of Bretton Woods.\\
\seg{14}{6}{L14.6 Flexible exchange} & 8:28 & Act three: floating rates, inflation, inflation targeting, speculative FX markets.\\
\seg{14}{7}{L14.7 Global Financial Crisis} & 17:58 & European banks funding US mortgages in dollars; central bank swap lines; return of home bias.\\
\bottomrule
\end{tabularx}

\smallskip
\textbf{Files.} Transcripts \file{materials/transcripts/L14-P*.txt}; Mehrling's notes \mn{14}{1}.
\textbf{Reading.} R.~A.~Mundell, ``A Reconsideration of the Twentieth Century'', \emph{American Economic Review} 90(3), 2000, 327--340.
\end{sources}

\section*{Motivation}
Lecture~\ref{ch:chartalism} placed the dollar at the top of a hierarchy of currencies. This lecture tells the history of the international monetary system in
the 20th century, following Mundell's Nobel lecture, and translates it into the language of the course: discipline versus elasticity at the top of the
hierarchy \ts{14}{2}{3:12}. It ends with the global financial crisis, where the same tension reappears as dollar funding of non-US banks and
central bank swap lines \ts{14}{7}{0:03}.

% ------------------------------------------------------------------
\section{FT: the costs of Japan's monetary policy}\label{sec:intl-ft}

\begin{casestudy}[FT, 2012: John Plender, ``Japan counts zombie cost of ultra-loose monetary moves'']
About 20 years earlier Japan had a financial crisis; its responses included quantitative easing, the first experiment with this kind of unconventional
policy, and fiscal expansion. The economy never really got going again and has had deflation ever since \ts{14}{1}{0:07}.
Plender worries about the structural damage of long periods of low rates: uncompetitive firms do not go out of business, because rolling over their loans
costs them almost nothing, and by competing they keep the productive firms from making profits \ts{14}{1}{0:50}. The
discipline avoided for so long now comes in through the back door: a yen at 79--80 per dollar makes Japanese exports uncompetitive
\ts{14}{1}{1:53}.
\end{casestudy}

\begin{definition}[New concepts: zombie firm]
A firm that is not viable but survives because cheap credit lets it roll over its debts (the living dead of horror films: neither alive nor buried)
\ts{14}{1}{1:11}.
\end{definition}

% ------------------------------------------------------------------
\section{Mundell's reconsideration of the twentieth century}\label{sec:intl-mundell}

\paragraph{Still no global currency.} Mundell wrote the article in 1999, published in 2000 when he received the Nobel prize; it ends optimistically: after a bad
century we now have stable units, the euro, the dollar and the yen, and perhaps we are on our way to an international monetary system
\ts{14}{2}{0:03}. Twelve years later the yen and the euro still fluctuate a lot against the dollar and against each other, there is no global
currency, the euro (born in 1999) is in crisis, and after the financial crisis there has been some retreat from financial globalization towards ``national
capitalism'' \ts{14}{2}{0:44}.

\begin{coursenote}
On the FRED chart shown, the euro is quoted as dollars per euro (American convention) and the yen as yen per dollar (European convention), so a rise of the
first line and a fall of the second both mean appreciation against the dollar. With no movement between yen and euro the two lines would move in opposite
directions \ts{14}{2}{1:04}.
\end{coursenote}

\paragraph{The good old days of gold.} For Mundell the international gold standard is the ``good old days'' because the payment system disciplined central banks
and finance ministries, and, since everyone adjusted to the same discipline, the whole world was integrated and stable \ts{14}{2}{3:32}.
There were problems, notably a tendency to long-run deflation when the value of gold changed, but in retrospect the system looks good. It broke down in World
War~I \ts{14}{2}{4:55}.

\paragraph{Central banks avoiding discipline.} The 20th-century story is of countries, central banks and finance ministers who found the discipline too harsh and tried
to get out from under it \ts{14}{2}{5:36}. In the hierarchy gold $>$ currency $>$ deposits, the discipline sits in the link between gold and
currency: as the economy grows, currency expands but gold does not \ts{14}{2}{6:20}. Mundell keeps in mind two clean ways out:
\begin{enumerate}
  \item \textbf{revalue gold}: a higher price of gold is like more gold; but all countries must do it together, and since gold holdings differ it is a wealth
        transfer, which international politics could not agree on \ts{14}{2}{6:46};
  \item \textbf{create artificial gold}: increase the quantity instead of the value, as attempted with the SDR at the IMF, but too little too late
        \ts{14}{2}{7:47}.
\end{enumerate}
Neither was found, so there was ``a lot of bumbling around''. The main culprit is the Fed: a new central bank (1913), with no history, constantly making errors;
but the US is dominant, everyone pegs to the dollar, and what the Fed does the world does. Instead of a gold standard, a dollar standard with an immature Fed
\ts{14}{2}{8:27}. Then floating, until in 1999 Europe decided it no longer wanted floating rates among its currencies
\ts{14}{2}{9:32}.

\begin{nfigure}
\begin{tikzpicture}[font=\scriptsize]
  \draw[thick,->] (0,0) -- (14.5,0);
  \foreach \x/\y in {0/1900,4.8/1933--34,10.2/1971--72,14/1999}{\draw (\x,0.1) -- (\x,-0.1) node[below]{\y};}
  \fill[thmcol!20] (0,0.25) rectangle (4.6,1.35);
  \fill[fincol!20] (5.0,0.25) rectangle (10.0,1.35);
  \fill[defcol!20] (10.4,0.25) rectangle (14,1.35);
  \node[align=center] at (2.3,0.8) {Act 1: confrontation\\Fed vs gold standard};
  \node[align=center] at (7.5,0.8) {Act 2: contradiction\\national management vs fixed rates};
  \node[align=center] at (12.2,0.8) {Act 3: flexible rates\\learning from experience};
  \node[below,align=center] at (2.3,-0.5) {1913 Fed; 1914--18 war;\\1931 Fed raises rates};
  \node[below,align=center] at (7.5,-0.5) {1944 Bretton Woods; \$35/oz;\\1971 gold window closed};
  \node[below,align=center] at (12.2,-0.5) {1970s stagflation;\\inflation targeting; euro 1999};
\end{tikzpicture}
\caption{Mundell's 20th century in three acts, with the dates used in the lecture \ts{14}{3}{0:03}\ts{14}{4}{0:45}\ts{14}{6}{1:12}.}\label{fig:intl-acts}
\end{nfigure}

\begin{reading}[R.~A.~Mundell, A Reconsideration of the Twentieth Century (2000)]
Mundell's Nobel lecture (delivered December 1999). Read with the hierarchy of money in mind: every episode is a central bank, at the top of a national
hierarchy, facing an international standard above it \ts{14}{2}{9:52}.

Some facts from the article (\file{materials/readings/L14-Mundell.pdf}), checked against the scan:
\begin{itemize}
  \item Cassel had warned in 1925 and 1928 that restoring gold would bring a scarcity of gold and falling prices (p.~329).
  \item Wholesale prices fell from their 1929 peak to September 1931 by between 22\% (Germany) and 40.5\% (Japan); the United States figure is 29.5\% (p.~329).
  \item In October 1931 the Fed raised the rediscount rate in two steps from 1\textonehalf\ to 3\textonehalf\%; US wholesale prices fell 35\% between 1929 and 1933 (p.~330).
  \item Bank failures rose from about 500 a year in the 1920s to 1{,}350 in 1930, 2{,}293 in 1931 and 1{,}453 in 1932 (p.~330).
  \item In 1934 the dollar's gold value was cut by 40.94\%, raising the official gold price by 69.33\% to \$35 an ounce (p.~331).
\end{itemize}
Both quotations used below (p.~330 on ``pumping liquidity'', p.~331 on the counterfactual) are verbatim in the article.
\end{reading}

\begin{history}[Mundell]
Robert A.~Mundell (1932--2021), Canadian economist, received the 1999 Nobel memorial prize for his analysis of monetary and fiscal policy under different
exchange rate regimes and of optimum currency areas; he is often called an intellectual father of the euro.\par
\emph{References:} NobelPrize.org, ``Robert A.~Mundell, Facts'', 1999; R.~A.~Mundell, \emph{AER} 90(3), 2000.
\end{history}

\begin{exercise}[Two ways to more gold]
World gold reserves are fixed while nominal incomes must double. Show the two adjustments Mundell has in mind (price of gold, price level) and explain who gains
and who loses from each when gold holdings are unevenly distributed.
\end{exercise}

% ------------------------------------------------------------------
\section{Act one: confrontation (1900--1933)}\label{sec:intl-act1}

Mundell calls the first act the confrontation of the Fed and the gold standard; it really begins only after World War~I \ts{14}{3}{0:03}.
The Fed is created in 1913; almost immediately there is war, and every country goes off gold except the US. That was no feat: gold was flowing to the US to pay
for war material \ts{14}{3}{1:05}. So gold did not discipline the US during the war, and there was some wartime inflation; the discipline came back after the
war, with deflation in 1920--21, though not enough to return to pre-war price levels \ts{14}{3}{1:47}. One by one the mark,
the pound and the franc returned to gold, not all at the pre-war parities (the French devalued) \ts{14}{3}{2:28}.

\paragraph{Gold becomes scarce.} Mundell's point: each currency that re-links to gold adds to the demand for gold reserves, so gold becomes increasingly scarce,
undervalued \ts{14}{3}{3:13}. A first chance to revalue gold was missed. There are two ways to have ``more gold'': raise the price of gold, or keep
it and lower the price of everything else, which is deflation, bad for the economies and resisted by national central banks \ts{14}{3}{3:54}.

\paragraph{The 1931 mistake.} The key mistake, for Mundell, was the Fed's in 1931: to defend the gold standard it raised interest rates in the middle of a
deflation, from 1.5\% to 3.5\%, which in deflation are very high real rates \ts{14}{3}{4:40}. Gold flowed to the US, forcing the
other countries off gold, eventually the US too, and the depression became worldwide. Defending the parity of the dollar destroyed the gold standard, and at
home it sacrificed the US banking system, which collapsed; Roosevelt was elected and the banking legislation of the 1930s followed \ts{14}{3}{5:42}.

\begin{history}[The Fed in October 1931]
After Britain left gold (21 September 1931), the Federal Reserve Bank of New York raised its discount rate from 1.5\% to 2.5\% on 9 October and to 3.5\% on
16 October 1931 to stop the gold outflow, while US banks were failing in large numbers.\par
\emph{References:} M.~Friedman and A.~J.~Schwartz, \emph{A Monetary History of the United States, 1867--1960}, Princeton UP, 1963, ch.~7; B.~Eichengreen,
\emph{Golden Fetters}, Oxford UP, 1992, ch.~9.
\end{history}

\begin{intuition}[Two pars, one choice]
There are two par relationships: the mint par (dollar to gold) and the deposit par (deposits to currency). Defending the first, the Fed sacrificed the second:
people preferred currency to deposits, banks were run and closed. Alternatively it could have stopped paying gold and supported the banks as lender of last
resort. ``You can't do both'': lender of last resort to your banks, or defender of the dollar's value in world markets \ts{14}{3}{8:48}.
\end{intuition}

The notes quote Mundell (p.~330): ``Instead of pumping liquidity into the system, it chose to defend the gold standard'', compounding the secular deflationary tendency (every return to gold
raised the monetary demand for gold, so lowered gold prices of everything else; France, which devalued, got a fresh start) with explicitly deflationary policy. Falling farm prices meant loan
defaults, tight policy added liquidity pressure, and bank failures added a further deflationary impulse (on which the notes cite Irving Fisher, ``The Debt-Deflation Theory of Great Depressions'',
1933) \mn{14}{1}\mn{14}{2}.

\paragraph{A monetary theory of history.} On p.~331 Mundell writes that had the price of gold been raised in the late 1920s, or had the major central banks pursued
price stability instead of adhering to the gold standard, there would have been no Great Depression, no Nazi revolution and no World War~II
\ts{14}{3}{6:44}. Mehrling finds this overboard but dramatic: the point is that instability of the world polity and of the world
monetary system are tied together. As a student of politics he would say that both are symptoms of the same failure: we could not get together to do the right
thing \ts{14}{3}{7:45}.

\paragraph{Why did the Fed err?} A further story: Benjamin Strong, who ran the New York Fed and knew all the world's central bankers, had died, and so had his adviser
Allyn Young; the young Fed had no deep bench, and in the crisis the knee-jerk conservative reaction was to defend gold. Some say that had Strong lived the
depression might not have happened; ``we'll never know'' \ts{14}{3}{9:53}.

\begin{coursenote}
In class the date of Strong's death is given as ``1929 I think''. Benjamin Strong died on 16 October 1928, in New York (he had long suffered from
tuberculosis). Allyn Young died in London on 7 March 1929, of pneumonia during an influenza epidemic, as said. The biography mentioned is L.~V.~Chandler,
\emph{Benjamin Strong, Central Banker}, Brookings, 1958 \ts{14}{3}{10:14}.
\end{coursenote}

% ------------------------------------------------------------------
\section{Act two: contradiction (1934--1971)}\label{sec:intl-act2}

The second act is the confrontation of national macroeconomic management with a fixed exchange rate system, set up at Bretton Woods but starting already in
the late 1930s \ts{14}{4}{0:03}. Every finance minister and central bank wants policies good for its own country, and at the same
time commits to fixed rates with other countries: this cannot be squared unless someone absorbs all the contradiction, which would have been the US, and it
did not \ts{14}{4}{1:06}.

\paragraph{Keynes's bancor plan.} Before Bretton Woods (a resort in New Hampshire where the post-war plan was made), Keynes proposed a world central bank and a true
global currency, bancor \ts{14}{4}{1:48}. Deficit countries would borrow bancor, created out of thin air by the
international bank and credited to the surplus countries; so the quantity of international money expands with the imbalances. It is the classic banking
solution, a lender of last resort at the international level, just as inside the US a deficit agent ultimately borrows from a bank that borrows from the central
bank \ts{14}{4}{3:19}.

\begin{definition}[New concepts: bancor, asymmetry of adjustment]
\emph{Bancor}: Keynes's proposed international bank money (from French \emph{banque} and \emph{or}, ``bank gold''). \emph{Asymmetry of adjustment}: under gold
(and under the IMF) the deficit country must adjust, because it has to acquire the international reserve, while the surplus country need not: the survival
constraint at the level of nations \ts{14}{4}{6:07}.
\end{definition}

\begin{nfigure}
\begin{taccts}
\tacct[2.4cm]{Deficit country (UK)}{}{$+$Loan from int'l bank (bancor)}\quad
\tacct[2.4cm]{International bank}{$+$Loan to deficit country}{$+$Bancor deposit of surplus country}\quad
\tacct[2.4cm]{Surplus country (US)}{$+$Bancor deposit}{}
\end{taccts}
\caption{The bancor plan: imbalances expand the international bank's balance sheet \ts{14}{4}{3:19}.}\label{fig:intl-bancor}
\end{nfigure}

\begin{history}[Bretton Woods and Keynes]
The United Nations Monetary and Financial Conference met at Bretton Woods, New Hampshire, 1--22 July 1944; the US plan (Harry Dexter White) prevailed over
Keynes's Clearing Union. Keynes died on 21 April 1946, a month after the inaugural meeting of the IMF and World Bank governors at Savannah.\par
\emph{References:} B.~Steil, \emph{The Battle of Bretton Woods}, Princeton UP, 2013; R.~Skidelsky, \emph{John Maynard Keynes}, vol.~3, Macmillan, 2000.
\end{history}

\paragraph{Why the US said no.} The US asked: which of these agents is us? Clearly the surplus country, facing the UK, all of Europe, everyone else. The plan looked
like a mechanism forcing the US to lend to other countries, through an intermediary it did not control; true after the war, perhaps not forever
\ts{14}{4}{4:44}. Worse, to remove the asymmetry Keynes proposed charging interest on the surplus country's deposit as
well as on the loan, to push the US to spend its bancor on European goods: a forced loan at a negative rate \ts{14}{4}{6:07}.
It was an attempt to eliminate the survival constraint for the deficit countries, which needed US goods to rebuild \ts{14}{4}{7:29}.

\begin{coursenote}
Keynes's 1942--43 Clearing Union proposals did charge a fee on large \emph{credit} balances as well as on debit balances (at rates rising with the size of the
balance, of the order of 1--2\% a year); the 3\% in class is an illustrative number.\par
\emph{Reference:} ``Proposals for an International Clearing Union'', Cmd.~6437, HMSO, April 1943.
\end{coursenote}

\paragraph{What we got: the IMF.} The US plan. In stylised form (not the IMF's actual, more complicated accounting, seen in problem set~1), member countries
deposit gold and their own currencies and receive deposit-like claims, now called special drawing rights \ts{14}{4}{7:51}.
The Fund ends up with all kinds of currencies, the important one being the dollar \ts{14}{4}{10:44}. Deficit countries pay surplus countries in
SDRs: the quantity of international money does not change, it is only transferred, unlike bancor \ts{14}{4}{11:25}. When you hit zero you
borrow from the IMF, at a positive interest rate and with policy conditionality: exactly the asymmetry Keynes wanted to avoid \ts{14}{4}{12:06}.
Keynes died shortly after; perhaps he accepted the IMF thinking it could be fixed later (``the famous economist death theory of international finance'',
Mehrling's ``pet theory'') \ts{14}{4}{12:49}.

\begin{definition}[New concepts: special drawing right (SDR)]
An international reserve asset created by the IMF and held by member countries, which can use it to make payments to one another; ``paper gold'' of a fixed
quantity \ts{14}{4}{9:12}. (A \emph{drawing right} is a right to draw, to obtain other currencies, from the Fund.)
\end{definition}

\begin{coursenote}
The lecture's IMF balance sheet is a deliberate simplification. Historically, members paid their quota partly in gold and partly in their own currency and
could draw on the Fund up to limits; SDRs were created only later, by the First Amendment to the Articles (agreed 1967, in force 1969), with a first allocation
of about SDR 9.3 billion in 1970--72. They are allocated to members, not issued against deposits of gold.\par
\emph{Reference:} IMF, ``Special Drawing Rights (SDR)'' factsheet; M.~G.~de~Vries, \emph{The International Monetary Fund 1966--1971}, IMF, 1976.
\end{coursenote}

\paragraph{Discipline versus elasticity.} The IMF is a discipline system, bancor an elasticity system; maybe bancor would have been too elastic and the IMF too
disciplined. For a while the shortage of ultimate international reserves was mitigated by the expansion of dollars further down the hierarchy
\ts{14}{4}{13:10}.

\begin{exercise}[Bancor accounting]
The UK runs a deficit of 100 bancor with the US for three years. Write the balance sheets of the UK, the international bank and the US at the end, with
1\% interest charged on both debit and credit balances. Do the same under the stylised IMF, with initial SDR holdings of 150 each.
\end{exercise}


% ------------------------------------------------------------------
\section{The dollar system}\label{sec:intl-dollar}

Under Bretton Woods the hierarchy was: gold plus SDRs (a fixed quantity of paper gold) at the top; the dollar, pegged to gold at \$35 an ounce; the other
currencies (mark, pound, franc) pegged to the dollar \ts{14}{5}{0:03}. The extra reserves that allowed post-war growth came
neither from new gold, nor SDRs, nor revaluation, but from more dollars \ts{14}{5}{0:44}.

\begin{history}[Gold at \$35 and the closing of the gold window]
The Gold Reserve Act of January 1934 set the dollar price of gold at \$35 an ounce (from \$20.67). On 15 August 1971 President Nixon suspended the
convertibility of dollars held by foreign official institutions into gold. The view of the US as ``bank of the world'' was set out by E.~Despres,
C.~P.~Kindleberger and W.~S.~Salant, ``The Dollar and World Liquidity: A Minority View'', \emph{The Economist}, 5 February 1966.\par
\emph{References:} Federal Reserve History, ``Nixon Ends Convertibility of U.S. Dollars to Gold'' and ``Gold Reserve Act of 1934''.
\end{history}

\paragraph{The US as a bank (Kindleberger).} Mundell emphasizes trade deficits as the way dollars went abroad. Charles Kindleberger, Mundell's professor at MIT,
stressed a different mechanism, which needs no trade deficit at all: the US lends long term to the rest of the world (financing European reconstruction by
buying foreign bonds), while the rest of the world, which wants to grow, accumulates short-term dollar balances as monetary reserves
\ts{14}{5}{1:06}.

\begin{nfigure}
\begin{taccts}
\tacct[2.8cm]{United States (as a bank)}{Gold\\$+$Foreign bonds (long term)}{$+$Dollar deposits of the rest of the world (short term)}\qquad
\tacct[2.8cm]{Rest of the world}{$+$Dollar reserves}{$+$Bonds issued}
\end{taccts}
\caption{Kindleberger's view: the US as a bank to the world, borrowing short and lending long, with no net capital or trade flow \ts{14}{5}{2:42}.}\label{fig:intl-usbank}
\end{nfigure}

The US is a bank holding the world's international reserves, as long as the world accepts dollars and does not demand gold: a sort of bancor plan, but on the
US balance sheet; and the US earns the spread, because dollar reserves pay little or no interest \ts{14}{5}{3:03}.
The IMF refused to expand reserves at the top of the hierarchy; the US expanded them one step down \ts{14}{5}{4:04}.

\paragraph{The end of Bretton Woods.} Dollars expanded, gold did not, so gold and SDRs became undervalued and stress built up \ts{14}{5}{4:24}.
Again: revalue gold or create more SDRs. In 1967 there was an increase in SDRs, about \$10 billion, not enough. Mundell tells, for the first time in Mehrling's
reading, that Arthur Burns sounded out European central bankers about revaluing gold and suggested to Nixon to do it right after the 1968 election, which did
not happen \ts{14}{5}{4:44}. By 1971 there were too many dollars for the gold, there was a run on the dollar, and
instead of paying out the gold Nixon closed the gold window \ts{14}{5}{5:48}. Kindleberger argued all his life that it would have
been fine had foreigners accepted dollars as the reserve and not insisted on the ``fiction'' of gold \ts{14}{5}{6:15}.

\begin{exercise}[The US as a bank]
Starting from the T-accounts of figure~\ref{fig:intl-usbank}, show what happens when foreign central banks convert \$10 of their dollar balances into gold. Why is
the system a bank-run problem, and what are the options of the US?
\end{exercise}


% ------------------------------------------------------------------
\section{Act three: flexible exchange rates (1972--1999)}\label{sec:intl-act3}

The contradiction destroyed fixed rates; central bankers and finance ministers were now free to run their economies their own way, and the 1970s were not a good
time: stagflation, worldwide \ts{14}{6}{0:03}. Mundell calls 1972--1999 the experiment with flexible exchange rates and learning
from experience \ts{14}{6}{1:12}.

\begin{definition}[New concepts: stagflation]
Inflation and unemployment (stagnation) at the same time; the word joins \emph{stagnation} and \emph{inflation} \ts{14}{6}{0:23}.
\end{definition}

\paragraph{What was learned.} Flexible rates are very volatile, and without international discipline national policy tends to inflation: the greatest peacetime
inflation, because no central bank, least of all the Fed, faced any discipline \ts{14}{6}{1:40}. The stagnation came from
uncertainty: with exchange rates moving 100\% back and forth you cannot tell where your comparative advantage lies, and you invest in the wrong things
\ts{14}{6}{2:44}. Two consequences:
\begin{enumerate}
  \item \textbf{central bank self-discipline}: gold was too much discipline, but some is needed; central banks impose price discipline on themselves, inflation
        targeting \ts{14}{6}{3:27}. Friedman and the advocates of floating thought this would also stabilize exchange rates (PPP,
        lecture~\ref{ch:chartalism}); it did not: inflation differentials were small, yet exchange rates moved all over the place
        \ts{14}{6}{4:09}. Mundell wants central bank discipline \emph{and} fixed rates, central banks cooperating \ts{14}{6}{4:50};
  \item \textbf{private speculative markets}: firms selling abroad must hedge their exchange risk, so forward FX markets develop, very sophisticated; but the same
        markets allow speculation, such as the carry trade, which pushes currencies apart and then, when reversed, snaps them back
        \ts{14}{6}{4:50}. The dealers take risk and want to be paid for it, which distorts prices
        (next lecture) \ts{14}{6}{6:34}.
\end{enumerate}

\begin{definition}[New concepts: carry trade]
Borrowing in a low-interest currency and lending in a high-interest one (borrow yen, lend dollars) to pick up the interest differential, the \emph{carry},
bearing the exchange risk \ts{14}{6}{5:53}.
\end{definition}

\paragraph{Mundell's hope.} In 1999 Mundell foresaw three key currencies, each anchoring a large area (Asian currencies pegged to the yen, the dollar spreading
to Latin America, the euro taking over Europe), linked to each other into a unified world system. That is not what happened \ts{14}{6}{7:38}.

The notes add \mn{14}{4}\mn{14}{5}: a system of national currencies with floating rates is typical in wartime, when commerce is restricted, ``quite anomalous in peace time''. Friedman persuaded
himself that market prices beat administered ones and that domestic price stability would bring exchange stability through PPP; Mundell disagrees, valuing fixed rates as par clearing knits
together the states of the US. Inflation targeting did bring price stability, but not exchange stability: volatility of dollar, euro and yen continued. A lasting legacy of the period of disorder is
the currency futures market for hedging exchange risk (lecture~\ref{ch:futures}). And Mundell's three-currency picture is not symmetric: the world funding market is a dollar market, as the crisis showed.
In the 1960s, too, Mundell suggests that more SDRs, substituted for gold in reserves, would have lowered the demand for gold and taken the pressure off; one response to the recent crisis was indeed
a new issue of SDRs by the IMF \mn{14}{4}.

\begin{coursenote}
The IMF made a general SDR allocation of about SDR 161 billion (some \$250 billion) in August 2009, plus a special allocation in September 2009. The notes also mention, as an alternative
proposal of the 1940s, Benjamin Graham's commodity-reserve currency (\emph{World Commodities and World Currency}, 1944) \mn{14}{2}.\par
\emph{Reference:} IMF, ``Special Drawing Right (SDR)'' factsheet.
\end{coursenote}

% ------------------------------------------------------------------
\section{The global financial crisis}\label{sec:intl-gfc}

\paragraph{Dollar funding outside the US.} In the build-up to the crisis, big European banks (UBS, for example) bought US residential mortgage-backed securities and
funded them short term in the international dollar funding markets: shadow banking, though on balance sheet \ts{14}{7}{0:23}.
The money came, say, from money market mutual funds whose shares, a deposit substitute, were held by non-US residents, including central banks and countries that
prefer to keep dollars offshore, out of reach of the US legal system \ts{14}{7}{1:52}. Dollar borrowing and lending outside
the US, by institutions without access to the US banking system or to the Fed \ts{14}{7}{2:54}.

\paragraph{The run.} After Lehman and AIG, money funds that had lost on Lehman commercial paper wanted only Treasury bills and asked for their money back. The
European bank went to the ECB, which could lend euros but had no dollars: ``I'm not the Fed'' \ts{14}{7}{3:15}.

\paragraph{The swap-line operation.} The engineered solution \ts{14}{7}{5:20}:
\begin{enumerate}
  \item the Fed and the ECB swap deposits: the Fed creates dollars, the ECB creates euros (a liquidity swap);
  \item the ECB lends the dollars to the European bank, which rolls its funding and need not sell its mortgage securities;
  \item the dollars end up as a deposit of the US Treasury at the Fed, which issues additional T-bills to the money funds: the funds wanted T-bills, could not
        hold deposits at the Fed, and the Treasury acted as the intermediary.
\end{enumerate}

\begin{nfigure}
\begin{taccts}
\tacct[1.9cm]{European bank}{RMBS}{$-$MMF funding\\$+$Dollar loan from ECB}\quad
\tacct[1.9cm]{ECB}{$+$Dollar deposit at Fed\\$+$Dollar loan to bank}{$+$Euro deposit of Fed\\$-$Dollar deposit (lent)}\quad
\tacct[1.9cm]{Fed}{$+$Euro deposit at ECB}{$+$Dollar deposit of ECB\\(then Treasury deposit)}
\end{taccts}
\begin{taccts}
\tacct[1.9cm]{Treasury}{$+$Deposit at Fed}{$+$T-bills}\quad
\tacct[1.9cm]{Money market fund}{$-$Loan to European bank\\$+$T-bills}{Shares}
\end{taccts}
\caption{The central bank liquidity swap in the crisis, as drawn in class (changes only) \ts{14}{7}{4:37}.}\label{fig:intl-swap}
\end{nfigure}

\begin{definition}[New concepts: central bank liquidity swap line]
An agreement by which two central banks exchange their own currencies, to be swapped back later at the same exchange rate; the receiving central bank lends the
foreign currency on to its own banks \ts{14}{7}{6:27}.
\end{definition}

\begin{coursenote}
The Treasury deposit is the Supplementary Financing Program, announced in September 2008: the Treasury issued extra bills and left the proceeds on deposit at the
Fed. The Fed's swap lines peaked at about \$580 billion in December 2008 (``600 billion'' in class).\par
\emph{References:} Federal Reserve, H.4.1 releases 2008; Federal Reserve History, ``Central Bank Liquidity Swaps''.
\end{coursenote}

\paragraph{Temporary cooperation, then repatriation.} This is the central bank cooperation Mundell asks for, but only a holding action, to keep the shadow banks from
dumping their securities and collapsing the market; eventually the Fed bought mortgage securities itself \ts{14}{7}{7:54}. The
permanent solution was repatriation: bringing the securities back home, undoing the internationalization of finance; a return of \emph{home bias} in
investment, because lender-of-last-resort support is national \ts{14}{7}{8:56}. In the notes: the money fund still funds the shadow bank's MBS, but through a chain, money fund to
Treasury to Fed to ECB to shadow bank; afterwards the construction was dismantled by taking the MBS back to the US, where dollar funding was easier to cobble together: ``it is no accident
that there is now over a trillion MBS on the balance sheet of the Fed'' \mn{14}{5}\mn{14}{6}. The dollar is a national currency,
the liability of the Fed; there is no global currency \ts{14}{7}{9:58}.

\begin{definition}[New concepts: home bias]
The tendency of investors to hold domestic assets more than international diversification would suggest \ts{14}{7}{9:38}.
\end{definition}

The reversal of globalization worries Mehrling: we saw it at the end of the 1920s, and it is a headwind against which national policies do not work; countries
try fiscal or monetary stimulus and only become more indebted and vulnerable \ts{14}{7}{10:40}. Meanwhile the system is held together
by international central bank cooperation and by speculative trading \ts{14}{7}{11:42}.

\paragraph{Is the Fed taking risk?} A student asks; Congress asked Bernanke the same about the roughly \$600 billion of swaps \ts{14}{7}{12:03}.
Two answers:
\begin{enumerate}
  \item the swap is reversed at the same amounts: the Fed gets back its dollars, the ECB its euros. Each books in its own currency, so whatever the exchange rate
        does nobody books a loss; commercially, with both parties keeping their books in the same currency, one side would show a loss
        \ts{14}{7}{13:07};
  \item it is international lender-of-last-resort action: with private funding gone, the central banks had no real choice; they were doing their job (``try to
        explain this to Congress'') \ts{14}{7}{14:49}.
\end{enumerate}

\paragraph{Standing swap lines.} The big central banks (Japan, ECB, Switzerland, Bank of England, Fed) have agreed swap lines in all currencies, of unlimited amount:
they are aware of the history and of the stakes, that monetary instability means economic and then social instability \ts{14}{7}{15:50}.
Maybe they overestimate their influence: without political will, a lender of last resort cannot prevent social instability; wars happen when people cannot
figure out what to do next, and in war the central bank becomes the financing agent of the government. ``I'd rather not go down that route''
\ts{14}{7}{17:12}.

\begin{coursenote}
The coordinated action of 30 November 2011 involved six central banks (the five named plus the Bank of Canada). The swap lines became standing, unlimited
arrangements among these six only on 31 October 2013, after this lecture.\par
\emph{Reference:} Federal Reserve press releases of 30 November 2011 and 31 October 2013.
\end{coursenote}

% ------------------------------------------------------------------
\flushside

\begin{exercise}[Swap gains and losses]
The Fed swaps \$100 for euros at \$1.25 per euro; three months later the euro trades at \$1.40. Compute the gain or loss of each central bank if both kept their
books in dollars, and then if each kept its books in its own currency.
\end{exercise}
\section*{Key formulas and ideas}
\begin{keyformulas}
\begin{tabularx}{\linewidth}{@{}l L@{}}
Mundell's thesis & integration and stability of the world economy depend on the discipline of an international standard\\
Escapes from gold & revalue gold (price) or create paper gold (quantity); neither coordinated in time\\
1931 & defending the mint par sacrificed the deposit par: LLR to banks or defend the currency, not both\\
Bancor vs IMF & elasticity (reserves grow with imbalances) vs discipline (fixed SDRs, deficit country adjusts)\\
Dollar system & US as a bank: short-term dollar liabilities, long-term foreign assets, gold reserve\\
Floating & inflation, volatility; answers: inflation targeting, private hedging/speculative FX markets\\
Crisis & offshore dollar funding; swap lines; repatriation and home bias; no global currency\\
\end{tabularx}
\end{keyformulas}

% !TEX root = ../main.tex
\chapter{Banks and Global Liquidity}\label{ch:globalliq}
\begin{paracol}{2}

\begin{sources}
\begin{tabularx}{\linewidth}{@{}l l L@{}}
\toprule
\textbf{Segment} & \textbf{Length} & \textbf{Content}\\
\midrule
\seg{15}{1}{L15.1 FT: European money market funds} & 2:37 & EUR~100 billion moved across currency areas.\\
\seg{15}{2}{L15.2 International transactions} & 11:42 & Bagehot's world, international: bills discounted in London, two FX dealers, the Bank of England.\\
\seg{15}{3}{L15.3 Dealer model for foreign exchange} & 10:05 & Gold points as the outside spread of FX dealers.\\
\seg{15}{4}{L15.4 Central banking, defense of domestic exchange} & 8:43 & The deficit country's central bank as dealer: pay out, sell assets, borrow reserves.\\
\seg{15}{5}{L15.5 Bank of England, defense against external drain} & 12:08 & Bank rate as outside spread; external drain; raising Bank rate; suspension.\\
\seg{15}{6}{L15.6 Toward a theory of exchange} & 9:37 & Without gold both outside spreads vanish; Keynes's \emph{Tract} (1923) and covered interest parity.\\
\bottomrule
\end{tabularx}

\smallskip
\textbf{Files.} Transcripts \file{materials/transcripts/L15-P*.txt}; Mehrling's notes \mn{15}{1} (headed ``18.'' in the PDF, from an earlier numbering of the course).
\end{sources}

\section*{Motivation}
Today's world is one of flexible, volatile exchange rates (next lecture). This lecture goes back a hundred years, to the gold standard, to introduce concepts
needed next time \ts{15}{1}{2:10}. It is ``the world that Bagehot knew'' (lecture~\ref{ch:bagehot}) again, now with its international piece: firms in different
countries trading by bills of exchange discounted in London, foreign exchange dealers, and the Bank of England on top \ts{15}{2}{0:03}. The tool is
Treynor's dealer model (lecture~\ref{ch:dealers}), applied twice: to exchange rates and to interest rates.

% ------------------------------------------------------------------
\section{FT: European money market funds}\label{sec:gl-ft}

\begin{casestudy}[FT, 2012: ``Europe's money market funds look further afield'']
Citing a Fitch Ratings report: rating downgrades, worries over the eurozone and lower bank demand for short-term funding are pushing European money market
funds towards Asia and the Middle East. Over two years they moved almost a fifth of their assets, about EUR~100 billion, from the UK, the US and the eurozone
periphery to Germany, France, the Netherlands, the Nordics, Asia and the Middle East \ts{15}{1}{0:48}. Moving EUR~100 billion
across currency areas has exchange-rate consequences: part of the volatility of flexible rates comes from such flows, and from people trying to avoid
currencies that are going to depreciate \ts{15}{1}{1:29}.
\end{casestudy}

\begin{aside}[Coming attractions]
The same day's FT reported China lifting foreign access to its markets, linked to the week's reading by McCauley (BIS) on the internationalization of the renminbi,
discussed next lecture \ts{15}{1}{0:07}.
\end{aside}

% ------------------------------------------------------------------
\section{International transactions in Bagehot's world}\label{sec:gl-transactions}

\paragraph{The bill discounted in London.} A surplus firm (or country) sells goods to a deficit firm and receives a bill of exchange, a promise to pay in 90 days;
neither is in England \ts{15}{2}{0:44}. The bill is not money, so the seller takes it to a private bank in the City of London,
the centre of the world money market in 1873, which discounts it, accepting the bill and paying out sterling notes \ts{15}{2}{1:50}.
(Notes rather than deposits, because the intuition is sharper \ts{15}{2}{2:54}.) The surplus firm is paid before the deficit firm has paid. In 90 days the
deficit firm pays the bill with notes, which flow back to the bank \ts{15}{2}{3:14}. The banks manage the outflow and inflow of
notes, holding bills that mature every day: classic 19th-century bill banking \ts{15}{2}{3:56}.

\paragraph{The foreign exchange dimension.} Sterling notes are high-powered money for international banking, but they cannot pay the surplus firm's workers;
and the deficit firm sells its goods for its own currency and must turn it into sterling to pay the bill \ts{15}{2}{4:38}. So
there are two foreign exchange dealers (or banks: in this course a bank is a special kind of dealer and a dealer a special kind of bank)
\ts{15}{2}{5:39}:
\begin{itemize}
  \item the \textbf{surplus dealer} buys sterling notes from the surplus firm and pays out domestic currency \ts{15}{2}{6:01};
  \item the \textbf{deficit dealer}, 90 days later, sells sterling notes to the deficit firm for its domestic currency \ts{15}{2}{7:06}.
\end{itemize}
The two domestic currencies may differ; what the dealers have in common is that both exchange against the international currency, sterling \ts{15}{2}{8:10}.

\paragraph{The Bank of England.} In 1873 sterling notes are liabilities of the Issue Department of the Bank of England, which holds gold and keeps the notes at par
with it \ts{15}{2}{8:31}. Both dealers use Bank of England notes as international reserves, and may cash them for gold or vice versa: the Bank makes a two-way
market between notes and gold \ts{15}{2}{8:52}. Three kinds of dealers, in a hierarchy: gold at the top, Bank of
England notes promising gold, then the other currencies kept at par with those notes \ts{15}{2}{9:58}.

\begin{nfigure}
\begin{taccts}
\tacct[2.2cm]{Surplus firm}{$-$Goods\\$+$Bill, then\\$-$Bill $+$Notes}{}\quad
\tacct[2.2cm]{City bank}{$+$Bill\\(90 days later: $-$Bill)}{$-$Notes (paid out)\\(90 days later: $+$Notes)}\quad
\tacct[2.2cm]{Deficit firm}{$+$Goods}{$+$Bill (due in 90 days)}
\end{taccts}
\begin{taccts}
\tacct[2.2cm]{Surplus dealer}{$+$Sterling notes}{$+$Domestic currency}\quad
\tacct[2.2cm]{Deficit dealer}{$-$Sterling notes}{$-$Domestic currency}\quad
\tacct[2.2cm]{Bank of England (Issue Dept.)}{$\pm$Gold}{$\pm$Notes}
\end{taccts}
\caption{International trade finance in 1873, as drawn in class \ts{15}{2}{1:08}.}\label{fig:gl-taccts}
\end{nfigure}

This is how it worked, and how it works today with the dollar: most world trade is invoiced in dollars, though much of it does not involve the US; back then it
was sterling \ts{15}{2}{10:59}.

% ------------------------------------------------------------------
\section{The dealer model for foreign exchange}\label{sec:gl-fxdealer}

Recall from lecture~\ref{ch:chartalism}: if the dollar is $x$ ounces of gold and the pound $y$ ounces, the mint parity exchange rate, quoted in pounds per dollar
(the convention of this course), is $x/y$ \ts{15}{3}{0:44}. The actual rate can deviate, because moving gold is costly:
\begin{keyformulas}
\textbf{Gold points} \ts{15}{3}{1:46}: with $\delta$ set by the cost of shipping gold,
\[
  \frac{x}{y}-\delta \;\le\; E \;\le\; \frac{x}{y}+\delta \qquad (E \text{ in pounds per dollar}).
\]
\end{keyformulas}

Inside the band there is no arbitrage to snap the rate back. This is the Treynor model hinted at two lectures ago: the gold points are an \emph{outside spread},
and the FX dealers do business inside it \ts{15}{3}{2:26}.

\paragraph{The deficit country's dealer.} Let the deficit country be the US (deliberately: the world currency is the pound, not the dollar; ``God'' did not plan
the dollar to be the reserve currency, certainly not in 1873) \ts{15}{3}{0:24}. Its dealers buy dollars and sell sterling: long the
domestic currency, short the international one, losing reserves, so long liquidity risk \ts{15}{3}{3:50}. They do it only at a lower price of the
dollar: with pounds per dollar, a lower number is a depreciated dollar \ts{15}{3}{5:13}. They take exchange risk only if paid in
expectation, hoping the dollar snaps back so they can sell high \ts{15}{3}{6:14}.

\paragraph{Why the dealer dares.} What if the dollar keeps falling? It cannot fall far: dollars can be taken to the central bank and converted into gold, which can
be shipped to England for sterling notes. The gold point is a floor, the outside spread; the dealers' quotes are the inside spread
\ts{15}{3}{6:55}. The quotes seen in the market are the dealers'; everyone knows the mint
pars, and a large volume of business pushes the price to the gold points, where gold flows keep it within bounds as long as central banks can supply the
gold \ts{15}{3}{8:37}.

\begin{coursenote}
The US central bank in this example is hypothetical: in 1873 the US had no central bank (the Fed dates from 1913), as Mehrling remarks \ts{15}{3}{7:36}.
\end{coursenote}

\begin{nfigure}
\begin{tikzpicture}[font=\scriptsize,x=1cm,y=1cm]
  % left panel: FX dealer
  \begin{scope}
  \draw[-{Stealth[length=4pt]}] (-3.2,0) -- (3.4,0) node[below left,align=right]{dealer's dollar\\inventory};
  \draw[-{Stealth[length=4pt]}] (0,-0.2) -- (0,4.4) node[above,align=center]{$E$ (\pounds/\$)};
  \draw[very thick,thmcol] (-3.2,3.4) -- (3.2,3.4) node[above left]{$x/y+\delta$ (gold point)};
  \draw[very thick,thmcol] (-3.2,0.8) -- (3.2,0.8) node[below left]{$x/y-\delta$ (gold point)};
  \draw[dotted] (-3.2,2.1) -- (3.2,2.1); \node[left] at (-3.2,2.1) {$x/y$};
  \draw[very thick,defcol] (-3,3.4) -- (3,1.1);
  \draw[very thick,intcol] (-3,3.1) -- (3,0.8);
  \fill (1.8,1.4) circle (1.5pt);
  \draw[thin] (1.8,1.4) -- (2.2,2.6) node[above]{deficit country};
  \node[align=center] at (0,-1.25) {FX dealer: exchange rate space};
  \end{scope}
  % right panel: bank as dealer in yield space
  \begin{scope}[xshift=7.3cm]
  \draw[-{Stealth[length=4pt]}] (-3.2,0) -- (3.4,0) node[below left,align=right]{bank's bill\\holdings};
  \draw[-{Stealth[length=4pt]}] (0,-0.2) -- (0,4.4) node[above]{discount rate};
  \draw[very thick,thmcol] (-3.2,3.4) -- (3.2,3.4) node[above left]{Bank rate};
  \draw[very thick,defcol] (-3,1.1) -- (3,3.4);
  \draw[very thick,intcol] (-3,0.8) -- (3,3.1);
  \node[align=center] at (1.6,1.3) {market rate};
  \node[align=center] at (0,-1.25) {City bank: interest rate (yield) space};
  \end{scope}
\end{tikzpicture}
\caption{Two outside spreads of the gold standard: the gold points for FX dealers \ts{15}{3}{3:28} and Bank rate for the discounting banks
\ts{15}{5}{1:26}. In yield space the curves slope up. Schematic.}\label{fig:gl-spreads}
\end{nfigure}

% ------------------------------------------------------------------
\section{The central bank defends the domestic exchange}\label{sec:gl-defense}

Suppose the dollar is pushed to the gold point: private dealers, who have as much risk as they want, bring domestic currency to the central bank and ask for
international currency, gold or sterling notes \ts{15}{4}{0:04}. Domestic currency is the central bank's own liability, so it is
simply cancelled: the central bank's balance sheet shrinks, it loses reserves and high-powered money contracts \ts{15}{4}{1:10}.
Textbooks say a currency is defended by contractionary monetary policy; here that follows literally from acting as a dealer, not from an intention
\ts{15}{4}{1:51}.

Its gold and notes are limited, so the central bank must think ahead. Three logical options, as in the flow-of-funds accounting
\ts{15}{4}{2:33}:
\begin{enumerate}
  \item \textbf{pay out reserves} (dishoarding);
  \item \textbf{sell an asset} for reserves, for example Treasury bills, which pushes bill prices down and interest rates up \ts{15}{4}{3:17};
  \item \textbf{borrow reserves} from another central bank, possibly the Bank of England itself \ts{15}{4}{3:58}.
\end{enumerate}

\begin{nfigure}
\begin{taccts}
\tacct[2.6cm]{Deficit country's central bank}{$-$Gold or sterling notes\\(1) $-$Reserves\\(2) $-$T-bills, $+$Reserves\\(3) $+$Reserves}{$-$Domestic currency\\\\(3) $+$Borrowed reserves}
\end{taccts}
\caption{Defending the domestic exchange at the gold point \ts{15}{4}{0:50}.}\label{fig:gl-cbdefense}
\end{nfigure}

The central bank is a backstop not so much because it holds large reserves, but because it can replenish them: it has assets (its special relationship with the
sovereign) and relationships with other central banks, which a private dealer lacks \ts{15}{4}{4:42}.

\begin{intuition}[What ``raising rates to defend the currency'' really means]
The notes read the three channels in balance sheets \mn{15}{3}\mn{15}{4}. Upward pressure on the currency is no problem: issue currency to buy reserves. Downward pressure means buying the
currency, the central bank's own liability, so its balance sheet contracts: paying out reserves works only while they last; selling Treasury bills pushes their price down and interest rates up,
perhaps a lot if speculators fear devaluation. ``We often hear loose talk about central banks defending their currency by raising interest rates'': what really happens is that the central bank
offers interest-bearing securities in exchange for currency, hoping the interest is enough to keep holders from demanding international reserves. Borrowing reserves from other central banks
(or the IMF, ``somewhat like an international central bank'') leaves the balance sheet's size unchanged: part of the monetary base in private hands shifts to other central banks. In every
case, contraction of the money supply supports the currency because the central bank buys its own liability. And the deficit entity that must adjust need not be below the Bank of England: it
can be the Bank itself \mn{15}{4}.
\end{intuition}

\begin{intuition}[What counts as reserves depends on the layer]
A student asks what the difference is between domestic currency and reserves. For the US central bank in a sterling world, reserves are gold or Bank of England
notes; for a bank inside the US, reserves would be domestic currency. What counts as reserves and what as credit depends on where you are in the hierarchy.
Foreign exchange is everything learned so far, one layer higher, with one difference: inside a country the relations are par relations, between countries the
rates move between the gold points \ts{15}{4}{5:44}.
\end{intuition}

\paragraph{The abstraction that removes liquidity.} Economists often treat the gold points as so narrow that the rate is just the mint par ratio, $\delta=0$, a
horizontal line. That abstracts from the essence of monetary economics: with no price movement no dealer would make a market, and without dealers there is no
liquidity. The course puts back the market mechanics that give liquidity \ts{15}{4}{7:28}.

% ------------------------------------------------------------------
\section{The Bank of England and the external drain}\label{sec:gl-boe}

\paragraph{City banks as dealers in yield space.} The City banks discount bills and wait 90 days for the money: they lend long and borrow short, take liquidity
risk, and so are dealers too, in interest-rate space; in yield space the bid and offer curves slope up, and the outside spread is \emph{Bank rate}
\ts{15}{5}{0:03}. A bank losing cash faster than it comes in raises its discount rate, so fewer
bills come in and the inflow exceeds the outflow \ts{15}{5}{1:47}. If all banks do it, the market rate of interest rises,
eventually to Bank rate, at which the Bank of England, as lender of last resort, rediscounts bills for notes \ts{15}{5}{2:28}.
The banks dare to discount because the backstop is there \ts{15}{5}{4:17}.

\paragraph{Internal drain.} The Banking Department pays out the notes it holds as its reserve, or it expands its balance sheet on both sides, crediting deposits:
as long as banks accept deposits as notes it economizes on its note issue and keeps rates from rising above Bank rate (lecture~\ref{ch:bagehot})
\ts{15}{5}{4:39}.

\paragraph{External drain.} New today: what if the dollar is bid \emph{up} to the other gold point? That outside spread is the Bank of England's market: surplus
dealers with too many notes ask for gold, and the Issue Department pays, $-$notes, $-$gold, defending the mint par of sterling
\ts{15}{5}{5:40}. Its gold is limited, so like any central bank it seeks more reserves. Its equivalent
of selling Treasury bills is to \textbf{raise Bank rate}, at a time when everyone is bringing bills for discount: the world rate of interest rises, borrowers
repay and do not renew, and gold flows in \ts{15}{5}{7:26}. But a high Bank rate can bankrupt people and cause
further problems \ts{15}{5}{8:50}.

\begin{definition}[New concepts: external drain, suspension of specie payments]
An \emph{internal drain} is a demand for the central bank's notes from the domestic system; an \emph{external drain} is a demand for gold to pay abroad
\ts{15}{5}{5:40}. \emph{Suspending specie payments}: the central bank stops converting its notes into gold (``specie'', coined metal, from Latin
\emph{in specie}, ``in kind'') while still issuing notes and deposits \ts{15}{5}{8:50}.
\end{definition}

\paragraph{The nuclear option.} The third option is suspension: ``I'm no longer going to pay notes in gold'', the system becomes a pure credit system for the time
being \ts{15}{5}{9:11}. Other central banks that suspend drop out of the international system; when the Bank of England does it, it
\emph{is} the international system: people keep using sterling as reserves, now unattached to gold, and perhaps decades later it returns to gold
\ts{15}{5}{9:51}.

\begin{keyformulas}
\textbf{The gold standard as discipline} \ts{15}{5}{10:53}. Domestically the central bank is at the apex and can
always print: no internal run can trouble it. Internationally an external drain forces it to defend its gold by raising Bank rate, at the cost of the domestic
economy, because foreigners did their banking business in London; the only alternative is suspension.
\end{keyformulas}

\begin{coursenote}
In class the comparison is said to be between ``lecture 10 and lecture 18''; it is between lecture~9 of these notes (``The World that Bagehot Knew'', which
Mehrling calls lecture~10) and this lecture~15 \ts{15}{5}{10:32}.
\end{coursenote}

\begin{history}[Suspensions of the Bank Charter Act]
Under the Bank Charter Act of 1844 the Issue Department's notes beyond a fixed fiduciary issue had to be backed by gold. In the crises of 1847, 1857 and 1866 the
government authorised the Bank to exceed that limit (a ``suspension'' of the Act, though convertibility was kept). Convertibility itself was suspended from 1797 to
1821 (the Restriction period) and, in effect, at the outbreak of the First World War; Britain returned to gold in 1925 and left it in 1931.\par
\emph{References:} J.~Clapham, \emph{The Bank of England: A History}, vol.~2, Cambridge UP, 1944; Bank of England, ``A brief history of banknotes''.
\end{history}

% ------------------------------------------------------------------
\section{Toward a theory of exchange}\label{sec:gl-theory}

\paragraph{After suspension.} A student: if payments are suspended, does the model not break down, and should sterling not be strongly devalued? Two answers
\ts{15}{6}{0:03}:
\begin{enumerate}
  \item yes, the story of mint par ratios and outside spreads goes away; it is replaced by something else, next lecture \ts{15}{6}{0:24};
  \item ``depreciate against what?'' Certainly against gold; but sterling may still be the best currency in the world, so not necessarily against other
        currencies: the whole currency system may depreciate against gold, but maybe not sterling against the rest \ts{15}{6}{1:05}.
\end{enumerate}

\paragraph{Two outside spreads gone.} The gold standard had two central bank outside spreads: the gold points, giving FX dealers room to work, and Bank rate,
giving discounting banks room to work \ts{15}{6}{2:08}. Moving away from gold removes both. Bank rate was a discount rate for 90-day bills; today
the Fed sets an overnight rate, the Fed funds rate, and does not set 90-day bill prices except in a crisis (when it walked up the yield curve)
\ts{15}{6}{2:49}. Next time the same two diagrams return, but the outside spreads will not be central banks
\ts{15}{6}{3:51}. Why explain the gold standard with the dealer model? Because, unlike the usual textbook explanation, it extends to fiat money: going off gold
does not mean ``la la land'' with no economics \ts{15}{6}{4:12}.

\paragraph{Keynes, 1923.} In \emph{A Tract on Monetary Reform}, his first real money book, written after the war, while countries were trying to get back to gold,
Keynes calls the gold standard ``already a barbarous relic'' (p.~172): central bankers now care about the stability of prices and employment and will not
sacrifice them to the ``outworn dogma'' of a gold price of \pounds3 17s 10\textonehalf d per ounce \ts{15}{6}{4:53}.
War finance had accustomed the world to paper money not backed by gold \ts{15}{6}{6:14}. His hope: each central bank stabilizes prices at home,
exchange rates float, and through purchasing power parity they are stable \ts{15}{6}{6:55}. It did not work: the world went back to gold, which
broke apart, and there was a world depression; this is the dream Mundell talked about \ts{15}{6}{7:17}.

But Keynes also started to discuss forward exchange, and wrote an equation equivalent to covered interest parity (lecture~\ref{ch:eurodollar}), the key
relationship for translating the two diagrams to a world without gold \ts{15}{6}{7:58}:
\begin{keyformulas}
\textbf{Covered interest parity} \ts{15}{6}{8:39}: with $S$ the spot and $F$ the forward rate (pounds per dollar), $R$ and $R^\ast$ the dollar and sterling term
rates of interest,
\[
  \frac{F}{S}=\frac{1+R^\ast}{1+R}.
\]
\end{keyformulas}

\begin{coursenote}
The auto-captions give the gold price as ``\pounds3 17 shillings 102''; the British mint price of standard gold was \pounds3 17s 10\textonehalf d per troy
ounce, as Keynes himself writes it, ``\pounds3:17:10\textonehalf\ per ounce''.

Checked against the scan (\file{materials/readings/L15-Keynes_Tract_on_Monetary_Reform.pdf}):
\begin{itemize}
  \item ``In truth, the gold standard is already a barbarous relic'' is on p.~172.
  \item ``Outworn dogma'' follows on p.~173, with the remark that a regulated non-metallic standard ``has slipped in unnoticed. \emph{It exists.}''
\end{itemize}
The equation is written on the board in class and is garbled in the captions; it is reproduced here in the notation of lecture~\ref{ch:eurodollar}.
\par\emph{Reference:} J.~M.~Keynes, \emph{A Tract on Monetary Reform}, Macmillan, 1923, ch.~3 (forward exchange) and ch.~4 (``barbarous relic'').
\end{coursenote}

% ------------------------------------------------------------------
\flushside
\section*{Key formulas and ideas}
\begin{keyformulas}
\begin{tabularx}{\linewidth}{@{}l L@{}}
Trade finance & bill discounted in London for sterling notes; FX dealers in each country exchange local currency for sterling\\
Hierarchy & gold $>$ Bank of England notes $>$ other currencies (kept within the gold points)\\
Gold points & $x/y-\delta \le E \le x/y+\delta$: outside spread for FX dealers\\
Defense & deficit central bank: pay out reserves, sell assets, borrow reserves; balance sheet shrinks\\
Bank rate & outside spread for discounting banks; internal drain absorbed by expanding the balance sheet\\
External drain & raise Bank rate (world rate of interest) or suspend specie payments\\
Without gold & both outside spreads vanish; CIP $F/S=(1+R^\ast)/(1+R)$ is the bridge to the modern system\\
\end{tabularx}
\end{keyformulas}

\section*{Exercises}

\begin{exercise}[A bill from Buenos Aires]
An Argentine exporter sells wheat to a German importer against a 90-day bill for \pounds1000, discounted in London at 4\% a year. Write all the T-accounts on the day
of discount and 90 days later, including the two FX dealers. How much sterling does the exporter receive?
\end{exercise}

\begin{exercise}[Gold points]
The dollar is 0.04838 ounces of gold and the pound 0.2354 ounces; shipping gold costs 0.6\% of its value. Compute the mint parity in pounds per dollar and the gold
points. Draw the FX dealer's diagram when the US is a deficit country.
\end{exercise}

\begin{exercise}[Three ways to defend]
A central bank holds 50 of gold, 100 of T-bills, and has 150 of currency outstanding. Dealers bring 60 of currency for gold. Show the balance sheet after each of
the three defenses, and say which ones raise domestic interest rates.
\end{exercise}

\begin{exercise}[Internal and external drain]
Explain why the Bank of England can stop an internal drain by lending freely at Bank rate but must raise Bank rate to stop an external drain. What does suspension
change for a country that is not at the top of the hierarchy?
\end{exercise}

% !TEX root = ../main.tex
\chapter{Foreign Exchange}\label{ch:fx}
\begin{paracol}{2}

\begin{sources}
\begin{tabularx}{\linewidth}{@{}l l L@{}}
\toprule
\textbf{Segment} & \textbf{Length} & \textbf{Content}\\
\midrule
\seg{16}{1}{L16.1 FT: High frequency trading} & 4:41 & EBS curbs ``predatory'' high-frequency traders.\\
\seg{16}{2}{L16.2 Uncovered interest parity} & 2:22 & Three FX facts: CIP, UIP, EH; UIP and EH fail for the same reason.\\
\seg{16}{3}{L16.3 FX dealers under the gold standard, redux} & 5:03 & The two gold-standard diagrams and their exogenous bounds.\\
\seg{16}{4}{L16.4 Private FX dealing system} & 10:41 & The \$10 payment: matched-book and speculative dealers.\\
\seg{16}{5}{L16.5 Economics of the dealer function, speculative dealer} & 5:57 & Forward rate below the expected spot rate: expected profit.\\
\seg{16}{6}{L16.6 Economics of the dealer function, matched-book dealer} & 6:15 & $F/S$ rises with the book; through CIP, foreign term rates are bid up.\\
\seg{16}{7}{L16.7 Digression: Why do UIP and EH fail?} & 9:45 & Dealers must be paid: liquidity premia, time-varying with the pattern of payments.\\
\seg{16}{8}{L16.8 Central bank as FX dealer of last resort} & 16:09 & Central bank swaps; the deficit central bank as speculative dealer; when intervention helps.\\
\seg{16}{9}{L16.9 Reading: McCauley on internationalization of renminbi} & 10:29 & Three tracks of China's financial system; FX as hybrid of state and market.\\
\bottomrule
\end{tabularx}

\smallskip
\textbf{Files.} Transcripts \file{materials/transcripts/L16-P*.txt}; Mehrling's notes \mn{16}{1}.
\textbf{Reading.} R.~N.~McCauley, ``Renminbi internationalisation and China's financial development'', \emph{BIS Quarterly Review}, December 2011; Baba and Packer, BIS WP~267 (2008); Mehrling, ``Essential hybridity (2013) (both in ile{materials/readings/}).
\end{sources}

\section*{Motivation}
Under the gold standard the central bank ``put its arms around the system'': the gold points and Bank rate were explicit outside spreads (lecture~\ref{ch:globalliq}).
The modern system has no such explicit bounds; the challenge is to understand how it works \ts{16}{3}{4:12}. Mehrling calls this ``in a way the hardest
lecture of the whole semester'': more wheels are turning than in any other \ts{16}{9}{8:46}. The payoff: the failures of uncovered interest parity and of the
expectations hypothesis turn out to be the same fact \ts{16}{2}{1:25}.

% ------------------------------------------------------------------
\section{FT: high frequency trading}\label{sec:fx-ft}

\begin{casestudy}[FT, 2012: ``EBS hails success of HFT curbs'']
High frequency trading (``cyborg trading''): algorithms, often written by computer scientists and engineers, run on servers co-located next to the exchange
\ts{16}{1}{0:06}. HFT firms make money, but it is not clear they make markets: they take money from slower regular dealers, who may
protect themselves by widening the bid--ask spread, making markets less liquid \ts{16}{1}{1:09}. EBS, one of the largest currency trading platforms, took measures
against ``predatory practices'': the fifth decimal of a quote must end in 0 or 5, and limits on orders cancelled before execution. Volumes fell, but the chief
executive says the real traders, and liquidity, stayed \ts{16}{1}{2:11}.
\end{casestudy}

Why does it matter especially in FX? Dealers hedge a spot position by a forward position; if a human takes a second between the two trades, a computer can take
both first and sell the hedge back at a profit. Banks that discover this back off \ts{16}{1}{1:30}. Platforms protect their other
customers because otherwise those customers go elsewhere: an arms race, the ``war of the cyborgs versus the humans'', in a market of \$4 trillion a day
\ts{16}{1}{3:32}.

\begin{definition}[New concepts: high frequency trading]
Automated trading by algorithms that place and cancel orders in fractions of a second, exploiting speed (\emph{high frequency}: many trades per unit of time)
\ts{16}{1}{0:06}.
\end{definition}

% ------------------------------------------------------------------
\section{Three FX facts}\label{sec:fx-facts}

\begin{keyformulas}
With $S$ the spot and $F$ the forward exchange rate (FX per dollar), $R$ and $R^\ast$ the dollar and foreign term rates \ts{16}{2}{0:03}:
\begin{enumerate}
  \item \textbf{covered interest parity} (CIP), a strict arbitrage condition that more or less defines the forward rate, since a synthetic forward can be built
        from deposits:\quad $\dfrac{F}{S}=\dfrac{1+R^\ast}{1+R}$;
  \item \textbf{uncovered interest parity} (UIP): the forward rate equals the expected spot rate,\quad $F=\E[S_T]$; empirically false;
  \item \textbf{expectations hypothesis} (EH): the long rate equals the expected rollover of short rates,\quad $(1+R_{0,T}) = (1+R_{0,N})\,\E[1+R_{N,T}]$; empirically
        false.
\end{enumerate}
\end{keyformulas}

The claim of the lecture: taking a money view of foreign exchange explains why UIP and EH fail, and shows that they fail for exactly the same reason:
``these are not different facts, they're actually the same fact'' \ts{16}{2}{1:25}.

% ------------------------------------------------------------------
\section{The gold standard, redux}\label{sec:fx-redux}

The two diagrams of last lecture \ts{16}{3}{0:24}:
\begin{itemize}
  \item exchange rate (pounds per dollar) against the dealers' liquidity risk: as the dollar depreciates the private dealers' quotes move down until the gold
        point $x/y-\delta$, where the US central bank must defend the dollar; the rate fluctuates between $x/y-\delta$ and $x/y+\delta$ even though there is a mint
        par \ts{16}{3}{1:47};
  \item the rate of interest at the centre (the UK), an upward-sloping curve: stress moves the market rate until Bank rate, where the Bank of England backstops
        the system, subject to convertibility; under an external drain it must raise Bank rate or suspend \ts{16}{3}{2:27}.
\end{itemize}
These were exogenous bounds set by the central bank. The modern system does not have them, at least not explicitly \ts{16}{3}{4:12}.

% ------------------------------------------------------------------
\section{The private FX dealing system}\label{sec:fx-private}

Back to the balance sheets of lecture~\ref{ch:chartalism}. A surplus country has a \$10 ``due from'' item on a deficit country; neither is the US, but trade is
invoiced and settled in dollars, the world reserve currency, and the deficit country has no dollars \ts{16}{4}{0:24}.

\paragraph{Step 1: the payment.} An FX dealer buys $10S$ of foreign exchange (FX, the deficit country's domestic currency) from the deficit country and pays with a new
liability, \$10 spot, which the deficit country transfers to the surplus country \ts{16}{4}{1:29}.
Both countries are done; the dealer is not happy: short dollars, long FX, exposed to the spot rate \ts{16}{4}{3:18}.

\paragraph{Step 2: the hedge.} The spot dealer hedges in the forward market, which is like a matched pair of term deposits, short and long, in the two currencies:
it adds a $10S$ FX term deposit liability (at $R^\ast$) and a \$10 term deposit asset (at $R$). This is CIP at work
\ts{16}{4}{4:19}. If spot and forward rates move together it loses on one side what it gains on
the other; some basis risk remains \ts{16}{4}{6:02}. One payment required two trades: this is why FX trading is \$4 trillion a day \ts{16}{4}{6:22}.

\paragraph{Step 3: the counterparty.} Whoever is on the other side of the hedge, neither country, takes the opposite position: a \$10 term deposit liability at $R$
and a $10S$ FX term deposit asset at $R^\ast$ \ts{16}{4}{6:44}.

\begin{nfigure}
\begin{taccts}
\tacct[2.0cm]{Deficit country}{$-$FX $10S$}{$-$Due to \$10}\quad
\tacct[2.0cm]{Surplus country}{$-$Due from \$10\\$+$\$10 spot}{}\quad
\tacct[2.0cm]{Matched-book dealer}{FX $10S$ spot\\\$10 term ($R$)}{\$10 spot\\FX $10S$ term ($R^\ast$)}
\end{taccts}
\begin{taccts}
\tacct[2.0cm]{Speculative dealer}{FX $10S$ term ($R^\ast$)}{\$10 term ($R$)}
\end{taccts}
\caption{The private FX dealing system: one \$10 payment, two dealers \ts{16}{4}{1:50}.}\label{fig:fx-private}
\end{nfigure}

\begin{definition}[New concepts: matched-book and speculative FX dealer (recalled)]
The \emph{matched-book dealer} hedges a long FX spot position with a short FX forward position: not exactly matched, but hedged. The \emph{speculative dealer} is
naked, speculating in the forward market, not in spot: it does not need to come up with spot dollars on demand. Its position is equivalent to borrowing at the
dollar rate and lending at the foreign rate: a carry trade \ts{16}{4}{7:53}.
\end{definition}

This is what finance people are paid for: making mutually beneficial trade possible across currencies, and shaving off their share. One dealer may take all of
these on its books; or no dealer will, and a central bank does it \ts{16}{4}{9:35}.

% ------------------------------------------------------------------
\section{The economics of the dealer function}\label{sec:fx-dealers}

\paragraph{The speculative dealer} is exposed to price risk, the subject of Treynor's model \ts{16}{5}{0:03}. On the vertical axis put the \emph{value} of FX (the
inverse of the course's usual quote, so that going down means depreciation of the deficit country's currency) \ts{16}{5}{0:24}.
What matters is the forward rate locked in today relative to the expected spot rate at maturity: being long FX, the dealer wants to buy it cheap. If the forward
price of FX is below the expected spot, the difference is expected profit \ts{16}{5}{1:46}.
Prices fluctuate and the dealer may be wrong, but right more than half the time it makes money. The speculative dealer is paid to take the exchange risk off the
matched-book dealer's hands \ts{16}{5}{4:37}.

\paragraph{The matched-book dealer} has hedged price risk but holds liquidity risk: it borrows short and lends long in the world reserve currency, with dollar
liabilities that can be withdrawn \ts{16}{6}{0:03}. It expands its book only if it buys spot more cheaply than it sells forward:
in $F/S$ space it has an upward-sloping bid--ask curve, higher the larger the book; in a crisis the spread also widens a lot \ts{16}{6}{1:08}.

\begin{nfigure}
\begin{tikzpicture}[font=\scriptsize,x=0.82cm,y=1cm]
  \begin{scope}
  \draw[-{Stealth[length=4pt]}] (-0.2,0) -- (6.4,0) node[below left,align=right]{long FX forward\\position};
  \draw[-{Stealth[length=4pt]}] (0,-0.2) -- (0,4.4) node[above]{value of FX (\$ per FX)};
  \draw[dashed,thmcol] (0,3.2) -- (6.2,3.2) node[above left]{$\E[S_T]$ (expected spot)};
  \draw[very thick,defcol] (0,3.1) -- (6,1.1);
  \draw[very thick,intcol] (0,2.8) -- (6,0.8);
  \fill (4,1.77) circle (1.5pt);
  \draw[{Stealth[length=3pt]}-{Stealth[length=3pt]}] (4,1.8) -- (4,3.17);
  \node[right,align=left] at (4.05,2.5) {expected\\profit};
  \node[left] at (4,1.55) {$F$};
  \node[align=center] at (3,-1.0) {speculative dealer};
  \end{scope}
  \begin{scope}[xshift=8cm]
  \draw[-{Stealth[length=4pt]}] (-0.2,0) -- (6.4,0) node[below left,align=right]{size of book\\(liquidity exposure)};
  \draw[-{Stealth[length=4pt]}] (0,-0.2) -- (0,4.4) node[above]{$F/S$, i.e.\ $(1+R^\ast)/(1+R)$};
  \draw[very thick,defcol] (0,1.2) -- (6,3.4);
  \draw[very thick,intcol] (0,0.9) -- (6,3.1);
  \node[align=center] at (3,-1.0) {matched-book dealer};
  \node[align=center,font=\scriptsize\itshape] at (2.2,3.6) {outside spread: ?};
  \end{scope}
\end{tikzpicture}
\caption{The two private FX dealers \ts{16}{5}{3:35}\ts{16}{6}{1:30}. Unlike the gold standard, no outside spread is drawn: who sets it is the question
\ts{16}{6}{2:58}. Schematic.}\label{fig:fx-dealers}
\end{nfigure}

\begin{intuition}[Expected spot as the new mint par]
In the notes, the expected spot rate plays, under floating, the role the mint par ratio played under gold: deviations create dealer profit as long as there is some expectation of return
\mn{16}{3}. But under floating there is no long-run anchor; and if people believe UIP, they read the forward depreciation that arises to induce dealers to absorb imbalances as a forecast of
future spot depreciation: a source of volatility. ``If that forecast leaks into expectations of actual future spot, then the whole system can get unhinged'' \mn{16}{3}\mn{16}{7}.
\end{intuition}

\paragraph{Where is the outside spread?} Under gold the central bank set the bounds; here nobody seems to. The key is CIP: $F/S=(1+R^\ast)/(1+R)$. If $F/S$ is
driven up by the matched-book dealer seeking profit, the term rates must move too; holding the US rate $R$ constant, the foreign term rate $R^\ast$ is bid up, an
upward-sloping curve in interest-rate space \ts{16}{6}{3:39}. The deficit country's central bank cares about
those term rates. Central banks set overnight rates, not term rates; the conceptual link between the two is the expectations hypothesis: with no liquidity
premia, term rates are expected rollovers of overnight rates, so a credible move in the overnight rate moves term rates \ts{16}{6}{5:10}.

% ------------------------------------------------------------------
\section{Why do UIP and EH fail?}\label{sec:fx-fail}

\begin{intuition}[One fact, two failures]
\begin{itemize}
  \item \textbf{EH fails}: to absorb the payment imbalance, term rates are bid up above what EH predicts. That is what we observe: on average the term structure is
        not flat but upward sloping. Dealers must make a profit to make markets, and this is where one kind of dealer's profit comes from \ts{16}{7}{0:03};
  \item \textbf{UIP fails}: if the forward rate equalled the expected spot rate, no speculative dealer would take the risky trade with zero expected profit:
        ``uncovered interest parity must be false'', or dealers will not make the market \ts{16}{7}{1:05}.
\end{itemize}
Both failures come from the need to mobilize profit-making dealers of two kinds to handle international payment imbalances \ts{16}{7}{1:46}.
\end{intuition}

\paragraph{An inverted yield curve?} A student saw an inverted repo yield curve on a terminal. Inverted curves usually signal monetary tightening, the central
bank pulling up the short rate; in repo just then, probably collateral shortage, unrelated to this argument \ts{16}{7}{2:27}.
An inverted curve is a sign of severe money-market stress: people bid for spot money because they have payments to make. Central banks usually cap overnight rates
unless they want stress; overnight rates pull the one-week and one-month rates through term-structure arbitrage, while at long horizons something other than
liquidity sets prices \ts{16}{7}{3:49}.

\begin{finance}[When CIP itself broke down]
In the crisis the need to roll dollar funding was so strong that those who could not borrow Eurodollars borrowed yen and swapped them into dollars: the dollar
funding crisis became an FX crisis, and covered interest parity broke down, because under stress people stopped doing the arbitrage and looked after themselves.
The BIS documented it \ts{16}{7}{5:11}. The lecture's argument is about normal times \ts{16}{7}{5:52}.
\end{finance}

\begin{coursenote}
The BIS article is not named in class. Representative BIS studies of the 2007--08 dislocation of FX swap markets are N.~Baba, F.~Packer and T.~Nagano, ``The
spillover of money market turbulence to FX swap and cross-currency swap markets'', \emph{BIS Quarterly Review}, March 2008, and N.~Baba and F.~Packer,
``Interpreting deviations from covered interest parity during the financial market turmoil of 2007--08'', BIS Working Paper 267, 2008.

The second paper is on the course's reading list (\file{materials/readings/L16-BIS\_work267\_Baba\_Packer.pdf}). Its abstract makes two points:
\begin{itemize}
  \item The sharp and persistent deviations from CIP between dollar and euro were significantly associated with the difference in counterparty risk between European and US institutions.
  \item The ECB's dollar term auctions, funded by the Fed's swap lines, stabilised the FX swap market by lowering the volatility of those deviations.
\end{itemize}
\end{coursenote}

\paragraph{What risk premium?} Abstracting from liquidity, standard economics has no story for either failure; it can invoke ``time-varying risk premia'', but what is
the risk? Here it is clear: exposure to liquidity, funding, turnover risk \ts{16}{7}{6:13}. The premia vary over time not because
utility functions bounce around but because the pattern of payments does: with small imbalances the premia should nearly vanish, since dealers need no
incentive \ts{16}{7}{7:14}. Inefficiencies appear even in the most liquid market in the world, \$4 trillion a day: it is not about
volume but about the incentive for dealers \ts{16}{7}{7:54}.

\paragraph{Pushing the survival constraint forward.} The dealers move the deficit country's survival constraint on into the next day: ``I don't have any dollars today,
I'm going to have dollars tomorrow, I think'', and I pay dealers meanwhile. Since they do not know when the flow reverses, they must price as if they will hold to
maturity; ``that's where the prices come from'' \ts{16}{7}{8:36}.

% ------------------------------------------------------------------
\section{The central bank as FX dealer of last resort}\label{sec:fx-cb}

The central bank comes in when private dealers are done: prices have moved, and they will take no more risk even if prices move dramatically
\ts{16}{8}{0:03}. Add the two countries' central banks; for simplicity they deal with each other, not with a third country or the IMF
\ts{16}{8}{0:50}. The private agent in the deficit country still needs to sell FX for dollars \ts{16}{8}{2:12}:
\begin{enumerate}
  \item the deficit central bank borrows term dollars from the surplus central bank at a rate $R^\ast_G$ (``G'' for government, not necessarily a market rate),
        and the surplus central bank opens a \$10 spot deposit for it: a swap of IOUs, like all banking \ts{16}{8}{2:33};
  \item the deficit central bank sells the \$10 to its resident for $10S_G$ of FX, at a rate that need not be the market rate (``we like you, you are an important
        industry, or a connected fellow'') \ts{16}{8}{4:46};
  \item the resident pays the surplus country: the dollar deposit is transferred \ts{16}{8}{5:28}.
\end{enumerate}

\paragraph{Sterilization.} The surplus central bank's balance sheet, and its money supply, expand; if it dislikes that it sterilizes by selling a Treasury bill.
Net, it has sold an asset and lent the proceeds to another central bank: its government may be unhappy about lending to foreigners instead of holding Treasury
bills, which ``might ring a bell'' from the crisis \ts{16}{8}{6:14}. The deficit
central bank's sterilization is the mirror, an expansionary operation adding an FX term asset \ts{16}{8}{8:20}.

\begin{nfigure}
\begin{taccts}
\tacct[2.6cm]{Surplus central bank}{$+$\$10 term loan to deficit CB ($R^\ast_G$)\\$-$T-bill (sterilization)}{($+$\$10 spot deposit of deficit CB,\\then transferred; net zero)}\qquad
\tacct[2.6cm]{Deficit central bank (net)}{$+$FX $10S_G$ term}{$+$\$10 term borrowing ($R^\ast_G$)}
\end{taccts}
\caption{The central bank as FX dealer of last resort, net positions after sterilization \ts{16}{8}{7:18}.}\label{fig:fx-cb}
\end{nfigure}

\paragraph{The deficit central bank as speculative dealer.} After the cancellations the deficit central bank borrows term dollars and lends term FX: it is a
speculative dealer, with a naked forward position in its own currency, but without a profit motive; if there were expected profit, private dealers would be doing
it \ts{16}{8}{9:04}. It ends up as FX dealer of last resort as a consequence of its commitment to peg overnight
interest rates \ts{16}{8}{10:32}.

\begin{definition}[New concepts: FX dealer of last resort, sterilization]
\emph{FX dealer of last resort}: the central bank taking the naked forward position private speculative dealers will no longer take, as a backstop for the FX
market \ts{16}{8}{10:32}. \emph{Sterilization}: an offsetting operation (for example selling a Treasury bill) that cancels the effect of an
FX operation on the domestic money supply (from Latin \emph{sterilis}, ``barren'': the operation is made to produce no monetary effect) \ts{16}{8}{6:34}.
\end{definition}

\paragraph{When is intervention efficient?} Knowing it is the dealer of last resort, the central bank may intervene earlier: if private speculative dealers charge
``an arm and a leg'' for the naked exposure, pushing the forward rate away from its efficient level, the central bank can provide the forward hedge to the
matched-book dealers itself \ts{16}{8}{10:53}. Typically these are central banks of poor countries, or countries without
well-developed FX dealer markets, where premia can be very large: ``I know I'm a minor currency, but my people are just as important as your people''
\ts{16}{8}{11:54}. Prices are signals: of competitiveness, and, for term rates, of whether saving is scarce. But here it is liquidity, not
saving, that is scarce, and the central bank can make it less scarce, accepting zero profit, or even losses for a national purpose such as industrialization
\ts{16}{8}{12:36}.

\begin{coursenote}
Not a blanket endorsement of currency intervention. A central bank offering FX hedges to the whole world would cause bubbles and ``all kinds of problems'';
with small distortions, why bother? Deep, liquid private markets need to make money, and a central bank that replaces them must decide politically who gets
foreign exchange, or end up ``offering Goldman Sachs their next bonuses''. The aim is to facilitate the underlying trade \ts{16}{8}{14:19}.

The notes widen the options \mn{16}{5}\mn{16}{6}. Central banks that target overnight rates ``are inevitably drawn into serving as FX dealers of last resort'': first as speculative dealers selling
hedges to matched-book dealers, then as matched-book dealers themselves through swaps with other central banks. The surplus central bank lending to a foreign central bank rather than its own
government is a limiting case; other sources of reserves work too, regional liquidity pools or the IMF. And no other central bank is needed: the deficit central bank can serve as speculative dealer
for private hedgers, or offer forward cover directly to its citizens; in the limit, forward cover for all comers at the current spot rate amounts to fixing the exchange rate. The rates involved
($R^\ast_G$, $S_G$) are policy rates, reflecting non-commercial relationships. The danger of trading with all comers at non-commercial prices is arbitrage by speculators; the notes quote DeRosa
(2009) on currency crises, all preceded by large positions short dollars and long local currency that the whole market, ``not counting the central bank'', had to hedge at once. Hence: ``the
exchange rate is not a free variable left open for state choice, but neither is it a market price that always fully reflects fundamental valuation'' \mn{16}{7}.
\end{coursenote}

% ------------------------------------------------------------------
\section{McCauley on the internationalization of the renminbi}\label{sec:fx-rmb}

\begin{reading}[R.~N.~McCauley, Renminbi internationalisation and China's financial development (BIS, 2011)]
McCauley describes China's monetary, credit and bond markets as having three tracks \ts{16}{9}{0:03}:
\begin{enumerate}
  \item the \textbf{old track}, the controlled economy, still large: banks are told whom to lend to and at what rate \ts{16}{9}{0:24};
  \item \textbf{peripheral markets}: loosened at the edges, a money market, a small interbank market, bond and loan markets with market rates; some market
        signals are better than none, in a controlled transition \ts{16}{9}{1:11};
  \item a new \textbf{offshore track}, from the internationalization of the renminbi \ts{16}{9}{2:16}.
\end{enumerate}
\end{reading}

The analogy in McCauley's mind is the Eurodollar market of the late 1950s and early 1960s: deposits and loans in dollars offshore, outside US regulation, often a
way of evading capital controls and interest-rate ceilings (Regulation~Q) \ts{16}{9}{2:16}. Offshore and onshore
renminbi prices are not on top of each other, a sign of obstacles to arbitrage (as Eurodollar and Fed funds rates normally are, within ``gold points'' of a few
basis points); in stress, the third and fourth quarters of 2011, they move farther apart, as Eurodollar and Fed funds did in the crisis
\ts{16}{9}{3:38}. Offshore renminbi bonds yield less than onshore ones, even with lower credit quality, because
holders of offshore renminbi deposits can switch only into offshore bonds \ts{16}{9}{5:01}. The Chinese authorities seem to allow it, as they
allowed track two; but this is harder to turn off once developed \ts{16}{9}{5:41}.

The notes place the offshore track in Hong Kong, money and capital markets both; the first step will be breaking down the barriers between the tracks, and in the process the central bank's
role will shift ``from a gold-standard like price setting function to a modern price backstopping function'' \mn{16}{7}.

\paragraph{A reserve currency?} Behind the discussion is the renminbi as a world reserve currency. China is a surplus country: everyone has to make payments to it,
and the survival constraint is asymmetric, so its liabilities are money, as US liabilities were after the two world wars. Today the Eurodollar market has little to
do with the US trade balance \ts{16}{9}{6:03}.

\begin{intuition}[FX is inherently hybrid]
Foreign exchange is a negotiated relationship between the state and the market, different in different countries. The Fed basically lets the private market
work: there are deep spot and forward markets with the dollar on one side. Elsewhere there are no dealers and central banks make the markets; most countries are
in between. These are not ideological choices: near the top of the hierarchy you look like the US, near the bottom like the central-bank case, in the middle a
hybrid. ``It's where states meet states and markets meet markets'' \ts{16}{9}{7:04}. There is no
general rule that intervention or capital controls are always bad; prices need not be at their efficient levels. Economists of all political stripes abstract
from liquidity, ``but liquidity doesn't abstract from them'' \ts{16}{9}{9:08}.
\end{intuition}

\begin{reading}[Mehrling, ``Essential hybridity: a money view of FX'' (\emph{J.~Comparative Economics} 41, 2013)]
The written version of this lecture (\file{materials/readings/L16-Essential\_hybridity.pdf}). Economics treats the exchange rate as the relative price of goods, finance as the relative price of assets. The paper treats it instead as the \emph{relative price of moneys}, set in dealer markets by order flow.
\begin{itemize}
  \item \textbf{A spectrum of regimes.} At one extreme the central bank quotes buy and sell prices and absorbs the order flow on its own balance sheet. At the other extreme private dealers set prices. In the general, hybrid case private dealers make the market and the central bank acts as ``dealer of last resort''.
  \item \textbf{The central bank as speculative dealer.} The deficit country's central bank borrows term dollars, for instance from the Fed through a swap line, and sells them spot to its citizens. It ends up with a naked forward position, the same position a private speculative dealer would take for profit; it shoulders FX risk that no private speculator will take. In the limit it offers forward cover to all comers at the current spot rate, which amounts to fixing the exchange rate (pp.~360--361).
  \item \textbf{The exchange rate is ``essentially hybrid''.} It is neither a market price that always reflects fundamentals, nor a free policy instrument.
  \item \textbf{The hierarchy of currencies.} According to the BIS figures quoted, 84.9\% of FX volume has the dollar as one leg. Of \$4 trillion a day, \$1.5 trillion is spot and \$1.8 trillion FX swaps. ``The FX market is fundamentally a money market, not a capital market'' (p.~361).
  \item \textbf{Majors, minors and crosses.} The majors all have the dollar as one leg; the minors trade against a major; most minor cross pairs do not trade at all. Every currency is ultimately a promise to pay dollars, through a regional key currency (Fig.~5, p.~362).
\end{itemize}
\end{reading}

% ------------------------------------------------------------------
\flushside
\section*{Key formulas and ideas}
\begin{keyformulas}
\begin{tabularx}{\linewidth}{@{}l L@{}}
CIP & $F/S=(1+R^\ast)/(1+R)$: arbitrage, defines the forward rate (breaks down in crises)\\
UIP & $F=\E[S_T]$: fails, since speculative dealers need $F\neq\E[S_T]$ to earn expected profit\\
EH & long rate $=$ expected rollover of short rates: fails, since term rates are bid up to pay matched-book dealers\\
Private system & matched-book dealer hedges spot with forward; speculative dealer holds the naked forward (carry trade)\\
Outside spread & no explicit bound; via CIP the deficit country's term rate, linked by EH to the overnight rate the central bank sets\\
Last resort & deficit central bank borrows term dollars, lends term FX: a speculative dealer without profit motive\\
Hybridity & FX regime reflects place in the hierarchy: markets at the top, central banks at the bottom\\
\end{tabularx}
\end{keyformulas}

\section*{Exercises}

\begin{exercise}[The two dealers with numbers]
The spot rate is $S=1.25$ FX per dollar, the six-month rates are $R=1\%$ and $R^\ast=3\%$ (not annualized). Write the T-accounts of the four agents for the \$10
payment, compute $F$ from CIP, and the speculative dealer's profit in dollars if the spot rate in six months is 1.25, 1.27 or 1.30.
\end{exercise}

\begin{exercise}[UIP and expected profit]
Show that if $F=\E[S_T]$ the speculative dealer's expected profit is zero. If its required expected return is 0.5\% over six months, how far must $F$ be from
$\E[S_T]$? Translate this into a premium in $R^\ast$ via CIP.
\end{exercise}

\begin{exercise}[Central bank swap]
Redo the central bank case with numbers: the deficit central bank borrows \$10 at $R^\ast_G=2\%$ and sells it to a resident at $S_G=1.20$ when the market rate is
1.25. Write all T-accounts, including both sterilizations, and identify who bears the exchange risk and who receives the subsidy.
\end{exercise}

\begin{exercise}[Offshore and onshore]
Offshore renminbi bonds yield 3\% and comparable onshore bonds 4\%. Explain why arbitrage does not close the gap, what happens to the gap as controls are relaxed,
and compare with the Eurodollar--Fed funds spread before and during 2008.
\end{exercise}

% !TEX root = ../main.tex
\chapter{Direct and Indirect Finance}\label{ch:directindirect}
\begin{paracol}{2}

\begin{sources}
\begin{tabularx}{\linewidth}{@{}l l L@{}}
\toprule
\textbf{Segment} & \textbf{Length} & \textbf{Content}\\
\midrule
\seg{17}{1}{L17.1 FT: Shadow banking} & 4:15 & Adair Turner: shadow banking is like cholesterol; traditional versus shadow bank.\\
\seg{17}{2}{L17.2 Bagehot's World} & 8:51 & Money market and capital market as separate worlds.\\
\seg{17}{3}{L17.3 The New World} & 10:55 & American banks: long-term assets, interbank repo, the origin of rating agencies.\\
\seg{17}{4}{L17.4 Funding liquidity versus market liquidity} & 2:33 & Moulton's shiftability (1918).\\
\seg{17}{5}{L17.5 Digression: Schumpeter on banking} & 4:05 & Banks as engines of development.\\
\seg{17}{6}{L17.6 Payment versus funding} & 5:31 & A loan funded by deposits, then taken out by a bond.\\
\seg{17}{7}{L17.7 Reading: Gurley and Shaw} & 2:04 & \emph{Money in a Theory of Finance} (1960): direct and indirect finance.\\
\seg{17}{8}{L17.8 Financial evolution} & 13:12 & Pension funds and insurers; from intermediation to mutual funds.\\
\seg{17}{9}{L17.9 Banking evolution} & 11:12 & Deposits as cheap permanent funding; from Jimmy Stewart banks to shadow banks.\\
\seg{17}{10}{L17.10 Preview: Central banking and shadow banking} & 8:28 & Liquidity puts, solvency versus liquidity, market liquidity provided by dealers.\\
\bottomrule
\end{tabularx}

\smallskip
\textbf{Files.} Transcripts \file{materials/transcripts/L17-P*.txt}; Mehrling's notes \mn{17}{1}.
\textbf{Reading.} J.~G.~Gurley and E.~S.~Shaw, \emph{Money in a Theory of Finance}, Brookings Institution, 1960 (passages).
\end{sources}

\section*{Motivation}
This lecture opens the fourth and last part of the course, which extends the money view from the money market to the capital market, to asset prices and the
prices of risk \ts{17}{1}{0:49}\ts{17}{2}{0:03}. The two subjects are usually taught separately: capital markets and CAPM in finance courses in the
business school, money and banking (perhaps IS--LM) in economics departments \ts{17}{2}{0:03}. Shadow banking shows that this
intellectual separation does not match the real world, where the two are more intertwined than at any time since World War~II, perhaps World War~I
\ts{17}{2}{1:04}. The separation has a historical origin, which the lecture reconstructs.

% ------------------------------------------------------------------
\section{FT: regulators to tackle shadow banking}\label{sec:di-ft}

\begin{casestudy}[FT, November 2012: ``Regulators to tackle shadow banking'']
A report led by Lord Adair Turner, the UK regulator, calls for regulation of shadow banking; Turner: shadow banking ``is like cholesterol'', there is good and
bad \ts{17}{1}{0:07}. Mehrling, who was working on the subject, thinks shadow banking is in our future and hopes
it will be the good kind \ts{17}{1}{0:28}.
\end{casestudy}

\begin{nfigure}
\begin{taccts}
\tacct[2.4cm]{Traditional (``Jimmy Stewart'') bank}{Loans}{Deposits}\qquad
\tacct[3.2cm]{Shadow bank}{RMBS (tranched)\\Interest rate swap\\Credit default swap}{Repo (RP)\\ABCP held by MMMF}
\end{taccts}
\caption{Traditional and shadow bank, as on Mehrling's slide \ts{17}{1}{1:51}. The traditional bank has a solvency backstop
(FDIC) and a liquidity backstop (Fed); the shadow bank has neither directly, only liquidity puts to traditional banks.}\label{fig:di-banks}
\end{nfigure}

The \textbf{traditional bank}, deposits against loans, is the ``Jimmy Stewart bank'' of the holiday film with a run on his bank. It is backstopped by the FDIC,
which insures the deposits (a solvency backstop), and by the Fed (a liquidity backstop); it is the balance sheet most people have in mind when they say ``bank''
\ts{17}{1}{1:51}. The world is increasingly like the \textbf{shadow bank}: mortgages packaged into residential mortgage-backed
securities, tranched into risky and risk-free tranches, insured by interest rate and credit default swaps, and funded in the money market by repo or asset-backed
commercial paper held by money market mutual funds \ts{17}{1}{2:53}. On both sides of its balance sheet the prices are market prices: the money
market funding rate and the price of risk; hence the coming lectures on swaps \ts{17}{1}{3:34}.

\begin{definition}[New concepts: shadow banking]
\emph{Shadow banking}: money market funding of capital market lending, outside the traditional banking system and its backstops \ts{17}{1}{0:49}
(``shadow'' because it does banking business outside the light of bank regulation).
\end{definition}

\begin{history}[Jimmy Stewart's bank]
In Frank Capra's \emph{It's a Wonderful Life} (1946) James Stewart plays George Bailey, who runs the Bailey Building and Loan in Bedford Falls; during a run he
tells the depositors that their money is not in a vault but in each other's houses.\par
\emph{Reference:} F.~Capra (dir.), \emph{It's a Wonderful Life}, Liberty Films/RKO, 1946.
\end{history}

% ------------------------------------------------------------------
\section{Bagehot's world: two separate markets}\label{sec:di-bagehot}

\paragraph{The money market: indirect finance.} In Bagehot's world banks discounted short-term bills that firms issued to finance goods on their way to final sale;
a portfolio of bills maturing at various dates assured the bank's liquidity \ts{17}{2}{2:06}. New in the course: the bank seen as an
\emph{intermediary}, between a primary lender (the household, for the first time in the course, holding deposits to make payments) and a primary borrower (a
business financing working capital with bills): deposits fund loans. This is \emph{indirect finance} \ts{17}{2}{2:47}.

\paragraph{The capital market: direct finance.} There was also a capital market: government bonds (including perpetual bonds paying a coupon forever), some
corporate bonds for large projects, and equity, held directly by rich people as wealth: \emph{direct finance}, no intermediary
\ts{17}{2}{4:30}. The two markets touch when bonds are issued, since the household pays with deposits that the firm
spends; afterwards the bond market seems separate. Bonds promise small coupons and principal only at maturity, perhaps 30 years later; the staggered-timing
liquidity issue of bills does not arise \ts{17}{2}{5:54}.

\begin{nfigure}
\begin{taccts}
\tacct[2.0cm]{Household (primary lender)}{Deposits\\Bonds, equity}{}\quad
\tacct[2.0cm]{Bagehot bank}{Bills}{Deposits}\quad
\tacct[2.0cm]{Business (primary borrower)}{Working capital\\Fixed capital}{Bills\\Bonds, equity}
\end{taccts}
\caption{Bagehot's world: indirect finance in the money market, direct finance in the capital market \ts{17}{2}{2:47}.}\label{fig:di-bagehot}
\end{nfigure}

\begin{aside}[A cult of liquidity?]
Some economic historians (among them some of Mehrling's professors) argue that Britain's hyper-developed short-term money market, the centre of world trade
finance, came with underdeveloped long-term finance: banks were forbidden by a ``cult of liquidity'' from long-term lending, which may help explain Britain
lagging after the industrial revolution. Mehrling does not vouch for it: ``I'm not an expert on Britain'' \ts{17}{2}{7:15}.
\end{aside}

% ------------------------------------------------------------------
\section{The new world}\label{sec:di-newworld}

The US was a developing country, with an extreme need for long-term finance: railroads, buildings, factories, homes. American banks held corporate bonds and
long-term loans, including ``six-month'' loans that were always rolled over, and few short-term bills \ts{17}{3}{0:03}.
Before there was a central bank, how did they manage? By cash reserves, by bankers' (correspondent) balances, and by interbank borrowing against bonds as
collateral: essentially repo, a paired sale and repurchase \ts{17}{3}{1:26}. A deficit bank
at the clearing pays in cash, in bankers' balances, or borrows against a bond.

\begin{nfigure}
\begin{taccts}
\tacct[2.4cm]{Deficit bank}{Bonds (collateral)\\Loans\\$-$Cash \emph{or} $-$Bankers' balances}{Deposits\\\emph{or} $+$Repo borrowing}\qquad
\tacct[2.4cm]{Surplus bank}{Bonds, loans\\Cash reserves\\$+$Cash \emph{or} $+$Repo lending}{Deposits\\Bankers' balances}
\end{taccts}
\caption{Interbank settlement in the US before the Fed \ts{17}{3}{1:50}.}\label{fig:di-us}
\end{nfigure}

\paragraph{Deficit at the clearing.} A student asks whether, in a developing country with many borrowers, the central bank acted as the surplus bank. There was no
central bank; and ``deficit'' refers to the clearing, not the whole system. Banks lend, as always, by creating deposits; when the borrower spends the deposit
in New York, the bank must pay the New York bank, with reserves, with bankers' balances at its correspondent, or by borrowing against a bond
\ts{17}{3}{3:55}.

\paragraph{The origin of rating agencies.} In Mehrling's account rating agencies arose in this period: banks paying each other did not want to analyse every bond
offered as collateral, so they bought Moody's book of ratings, which said a bond was unlikely to default in the next year, good enough for short-term collateral;
the rating entered the terms of the repo \ts{17}{3}{5:38}. Unsecured lending through bankers'
balances needs a relationship; secured lending against a rated railroad bond can be done with a New York bank that has never heard of you \ts{17}{3}{7:20}.
Capital market and money market are completely intertwined: ``19th-century shadow banking'' \ts{17}{3}{7:41}.

\begin{history}[Moody's, 1909]
John Moody published \emph{Moody's Analyses of Railroad Investments} in 1909, assigning letter ratings to railroad bonds; Poor's (1916), Standard Statistics
(1922) and Fitch (1924) followed. The link to interbank collateral is Mehrling's interpretation.\par
\emph{References:} R.~Sylla, ``An historical primer on the business of credit rating'', in R.~M.~Levich et al.\ (eds), \emph{Ratings, Rating Agencies and the
Global Financial System}, Kluwer, 2002.
\end{history}

\paragraph{The Fed and the shadow system.} After the 1907 crisis the reformers tried to replace this system with discounting of short-term bills at a new Fed, but
not enough bills were issued: the US needed long-term finance \ts{17}{3}{7:41}. Mehrling's argument: the Fed's mistake around 1929
was to let the shadow banking system of its day collapse, refusing to support mortgage credit and bonds as improper for banks, and so it brought down the US
banking system and the world's \ts{17}{3}{8:22}. In 2007 we were fortunate that the Fed did backstop the shadow banking system,
much of it in Europe, through the liquidity puts linking shadow banks to Citibank and other traditional banks \ts{17}{3}{9:03}. The
founders of the Fed ``pasted a kind of central bank onto'' a system it did not fit; today's question is how to have a central bank that supports a system like
this \ts{17}{3}{10:04}. The notes put it sharply: the banking collapse of the 1930s can be read as a run on the shadow banking system of the time, private funding liquidity through repo on
non-Treasury securities; the Fed refused to support it, it collapsed, and the government had to take on long-term lending itself, most famously through the Reconstruction Finance Corporation:
``this mistake was not, thankfully, repeated in 2007--8!'' The transmission from funding liquidity to market liquidity and asset prices is, in the US, ``not recent, but in fact ancient''
\mn{17}{2}\mn{17}{3}.

% ------------------------------------------------------------------
\section{Funding liquidity versus market liquidity}\label{sec:di-shiftability}

The British concept of liquidity was \emph{funding liquidity}: guaranteed cash inflows. In the new world a bank's liquidity was having bonds pledgeable or
shiftable to someone else: \emph{shiftability}, which the course calls market liquidity, the ability to buy or sell an asset quickly without moving the price much
\ts{17}{4}{0:04}.

\begin{definition}[New concepts: shiftability]
\emph{Shiftability}: the liquidity of an asset that can be shifted to another holder (sold or pledged) without loss, as opposed to liquidity from self-liquidating
short-term assets. Term of Harold Moulton, 1918 \ts{17}{4}{0:45}.
\end{definition}

Moulton argued that shiftability was the right concept for the US, and that the Fed should backstop it. The Fed did, not for all assets but for Treasury bills:
that is where open market operations came from, supporting market liquidity in the capital market \ts{17}{4}{1:06}. This course
uses both kinds of liquidity and their interlink, dealers who create market liquidity must fund their inventories, an understanding not yet there in 1918
\ts{17}{4}{1:46}.

\begin{history}[Harold G.~Moulton]
Harold Glenn Moulton (1883--1965) published ``Commercial Banking and Capital Formation'' in four parts in the \emph{Journal of Political Economy} in 1918; he
was the first president of the Brookings Institution (1927--1952).\par
\emph{References:} H.~G.~Moulton, \emph{JPE} vol.~26, 1918 (four parts); Brookings Institution, ``History''.
\end{history}

% ------------------------------------------------------------------
\section{Schumpeter on banking}\label{sec:di-schumpeter}

Joseph Schumpeter, sympathetic to the American scene, emphasized all his life the importance of banks willing to do long-term development finance: money market
funding of capital market lending \ts{17}{5}{0:04}. Banks create new money for entrepreneurs, who need not find
someone who has saved: ``created from nothing'', the alchemy of banking. To say so risks sounding like a monetary crank, and Schumpeter was not one
\ts{17}{5}{1:08}. Expanding the balance sheet is not ``bread from heaven'': something must discipline it (next section) \ts{17}{5}{1:50}.

Schumpeter admired H.~D.~Macleod, who stressed this kind of banking in the development of Scotland. The Lewis model of development, with surplus labour in
subsistence agriculture, also gives banking a role in mobilizing it towards cities and factories \ts{17}{5}{2:10}.
Development economists are usually ``real side'' people; Mehrling would give half of a development course to finance, ``but hey, I would say that''
\ts{17}{5}{3:32}.

\begin{intuition}[There is such a thing as a free lunch]
The notes warn that most self-styled monetary alchemists are cranks: ``the ability to create money from nothing is not the same as the ability to create bread from nothing'', and economics has
come close to making ``there ain't no such thing as a free lunch'' a creed. ``But there is.'' Adam Smith (1776) already urged paper money to economize on gold, trading a sterile commodity for
productive capital; the main source of the free lunch, though, is that banking can \emph{mobilize unused resources}: if entrepreneurs spend new deposits on things that would otherwise go unsold,
they raise activity, not prices \mn{17}{3}.
\end{intuition}

\begin{coursenote}
In class Schumpeter's ``thesis'' is dated 1912. The book meant is \emph{Theorie der wirtschaftlichen Entwicklung} (published late 1911, dated 1912; English
\emph{The Theory of Economic Development}, 1934), not his doctoral thesis. Henry Dunning Macleod (1821--1902) was a Scottish economist, author of \emph{The Theory
and Practice of Banking} (1855--56). The Lewis model is W.~A.~Lewis, ``Economic Development with Unlimited Supplies of Labour'', \emph{Manchester School}, 1954,
which discusses capital formation financed by credit creation \ts{17}{5}{0:04}.
\end{coursenote}

% ------------------------------------------------------------------
\section{Payment versus funding}\label{sec:di-payment}

A business wants to buy a machine from ``society''. The bank lends by a swap of IOUs, $+$loan against $+$deposit, and the business pays with the deposit
\ts{17}{6}{0:05}. The question: what \emph{funds} the loan? If society is happy to add the
deposit to its portfolio, the deposit funds the machine, even though it is a demand deposit and the machine a loom with a ten-year life
\ts{17}{6}{1:53}. If society does not want to hold it, the business must fund the machine otherwise, by issuing a bond that
society buys with the deposit, and repaying the loan \ts{17}{6}{2:55}.

\begin{nfigure}
\begin{taccts}
\tacct[2.0cm]{Bank}{$+$Loan\\(later $-$Loan)}{$+$Deposit\\(later $-$Deposit)}\quad
\tacct[2.0cm]{Business}{$+$Machine}{$+$Loan\\(later $-$Loan, $+$Bond)}\quad
\tacct[2.0cm]{Society}{$-$Machine, $+$Deposit\\(later $-$Deposit, $+$Bond)}{}
\end{taccts}
\caption{Payment versus funding: bank credit as a bridge loan, the bond as take-out financing \ts{17}{6}{1:10}.}\label{fig:di-payment}
\end{nfigure}

The money market intermediation is a bridge loan: it gets the ball rolling, using the elasticity of the payment system; once the machine is running, the capital
market provides permanent, take-out financing. First intermediation, as in Bagehot but long; then direct finance, bond for bond, the intermediary drops out
\ts{17}{6}{4:17}.

% ------------------------------------------------------------------
\section{Gurley and Shaw: intermediaries in the capital market}\label{sec:di-gurleyshaw}

\begin{reading}[J.~G.~Gurley and E.~S.~Shaw, Money in a Theory of Finance (1960)]
Gurley and Shaw (Brookings, 1960) introduced the terms \emph{direct finance} (the primary lender lends directly to the primary borrower) and \emph{indirect
finance} (an intermediary in between). They stressed intermediation not just in the money market but in the capital market, for long-term finance, following
Moulton; and not only by banks but by the other intermediaries that had grown since his day \ts{17}{7}{0:04}.
\end{reading}

\begin{definition}[New concepts: direct and indirect finance, financial intermediary]
\emph{Direct finance}: the primary borrower's securities are held by the primary lender. \emph{Indirect finance}: a \emph{financial intermediary} (from Latin
\emph{inter}, between, and \emph{medius}, middle) issues its own liabilities to lenders and holds the borrowers' securities \ts{17}{7}{0:31}.
\end{definition}

% ------------------------------------------------------------------
\section{Financial evolution}\label{sec:di-finevol}

The two main capital market intermediaries in the US are pension funds and insurance companies \ts{17}{8}{0:05}:
\begin{itemize}
  \item a \textbf{pension fund}, typically a defined benefit plan in 1960, promises workers a fraction of final income for life, a long-term, annuity-like promise,
        and hedges it with long-term assets, especially equities \ts{17}{8}{0:46};
  \item an \textbf{insurance company} has contingent liabilities (fire, life policies) and holds bonds \ts{17}{8}{1:30}.
\end{itemize}
Both stand between households and businesses, like Bagehot's bank, though their liabilities are not money; they became bigger than banks in funding American
development \ts{17}{8}{2:11}.

\begin{intuition}[Gurley and Shaw's vision: intermediaries face both ways]
Households want to hold liquid assets, money, insurance, pensions; businesses want to borrow long. How can the capital development of the nation be financed?
An intermediary faces both ways: to households it says ``I'll give you what you want'', deposits, policies, pensions; to businesses, ten-year loans, or it buys
their bonds and equity. Without it, the mismatch between what households want and what businesses issue could block development
\ts{17}{8}{3:33}.
\end{intuition}

Their liquidity: insurance has an actuarial quality (not every house burns at once), and in a hurricane the insurer sells bonds; pension funds turn equities into
bonds (annuities) when a worker retires. Both depend on market liquidity, shiftability, as Moulton said for banks; Moulton became president of Brookings, which
may be why Brookings hired Gurley and Shaw \ts{17}{8}{5:17}.

\paragraph{The flaw: risk does not go away.} The risk of the assets does not disappear because you promise something else: it falls on the liabilities, or on
whoever guarantees them, typically the government. With risky assets, guaranteed benefits may not be paid; the risk is borne by the sponsor (often a corporation),
by the government through the Pension Benefit Guaranty Corporation, or by the beneficiaries when promises are not kept. Finance and arbitrage will make you
appreciate it, through mispricing \ts{17}{8}{7:20}.

\paragraph{From indirect to direct finance.} Since 1960 indirect finance has been hollowed out in favour of mutual funds: stock funds hold equities and issue shares
in the portfolio (Fidelity's Magellan), bond funds hold bonds (PIMCO): the risk passes right through to the holder, ``this is not investment advice''
\ts{17}{8}{9:07}. Pensions moved from defined benefit to defined contribution plans such as the
401(k), where you see your risk \ts{17}{8}{10:51}. In insurance, term policies replaced whole-life policies with their savings component: both trends reflect mutual funds competing
with traditional indirect finance. A mutual fund looks like an intermediary, but it mainly pools risk through diversification, with little transformation: its shares have the risk of the
underlying pool; the liquidity of buying back at NAV requires cash or credit lines, paid for by the shareholders. Bond funds replaced the savings in whole-life policies and bank time deposits,
stock funds replaced defined-benefit plans \mn{17}{6}\mn{17}{7}. Gurley and Shaw ``overdid it'', the notes say: modern finance stresses that intermediation eliminates no risk, only transfers
it, sometimes opaquely, so capital markets equilibrate by price, and attempts to avoid this create arbitrage opportunities that have transformed the system \mn{17}{6}.

\begin{keyformulas}
\textbf{Quantities versus prices} \ts{17}{8}{11:32}. Indirect finance solves the mismatch between what lenders want and what
borrowers want with \emph{quantities}, on its own balance sheet. Modern finance solves it with \emph{prices}: if you don't want to hold bonds, bonds fall in
price until you do. Price does much more work than in 1960, which is why everyone watches the stock market.
\end{keyformulas}

The same evolution, from intermediation to direct finance, has happened in banking and the money market: shadow banking is its natural next step \ts{17}{8}{12:33}.

% ------------------------------------------------------------------
\section{Banking evolution}\label{sec:di-bankevol}

Banks as intermediaries are a new conception in the course, which has treated banking as a payment system and as market making, not as funding
\ts{17}{9}{0:06}. Why can banks fund things? Because people want deposits not just to pay but as an asset, a liquidity reserve in their wealth
\ts{17}{9}{0:26}:
\begin{itemize}
  \item deposits are \textbf{permanent} funding for the system as a whole: money demand fluctuates but does not go to zero, so reserves need cover only the
        fluctuations; a banker who sees deposits never fall below 80\% of the balance sheet can hold 80\% in earning loans \ts{17}{9}{2:10};
  \item deposits are \textbf{cheap} funding, close substitutes for cash that pays no interest; even when banks compete by paying interest, it is less than on bonds
        \ts{17}{9}{1:30}.
\end{itemize}

\paragraph{Liquidity at the margin.} In the notes: because households and firms want to \emph{hold} means of payment, not only make payments, the banking system can keep a permanent
short position in cash and use it to fund long positions in non-cash assets. Liquid liabilities need liquid assets only at the margin: cash reserves plus the ability to replenish them (Fed funds,
the discount window, ``secondary reserve'' assets that can be sold); the rest of the deposits can go into less liquid assets. Once the focus is intermediation, solvency risk, interest rate risk and
credit risk, comes to the fore \mn{17}{7}\mn{17}{8}.

\paragraph{Sharing out the spoils.} Everyone wants this cheap funding \ts{17}{9}{3:52}: the government through the central bank (Treasury bills
against cash), business through commercial banks (industrial loans against deposits), households through savings and loans, the Jimmy Stewart banks (mortgages
against ``deposits'', legally shares but used as money) \ts{17}{9}{4:33}.

\begin{nfigure}
\begin{taccts}
\tacct[1.9cm]{Central bank}{Treasury bills}{Cash}\quad
\tacct[1.9cm]{Commercial bank}{Business loans}{Deposits}\quad
\tacct[1.9cm]{Savings and loan}{Household mortgages}{``Deposits'' (shares)}
\end{taccts}
\caption{Cheap permanent funding shared out among government, business and households \ts{17}{9}{4:33}.}\label{fig:di-spoils}
\end{nfigure}

\paragraph{The risk passes through, to the taxpayer.} The banks face both ways, but the risk does not vanish: in the Jimmy Stewart bank the FDIC guarantees
solvency and the Fed liquidity, so the taxpayer bears it. Hence regulation (no risky loans, no interest on deposits), and hence a regulatory arbitrage to create
non-banks doing the same thing: that is where shadow banking came from \ts{17}{9}{6:17}. In
the Jimmy Stewart bank the risk stays in the community: ``your money is in your houses'' \ts{17}{9}{7:40}.

\paragraph{Shadow banking strips the risks.} In shadow banking the loans become traded securities whose price is the price of all their risks: a mortgage has
duration risk (30 years) and default risk (household income). The interest rate swap and the credit default swap strip them off, leaving a synthetic Treasury
bill funded in the money market: risk is not eliminated but hedged, sold off separately, rather than borne by a government that pretends it is not there, with the
moral hazard that implies \ts{17}{9}{8:22}. As pension funds evolved into stock funds and insurers into bond funds,
banks evolved into money market mutual funds, many outside the US, holding dollars for sovereign wealth funds and central banks: ``a global dollar system funding
subprime mortgages in Cleveland'' \ts{17}{9}{9:44}. From indirect finance (Gurley and Shaw, 1960) to direct finance (2012),
with all risks priced in markets and money and capital markets completely integrated \ts{17}{9}{10:45}.

% ------------------------------------------------------------------
\section{Preview: central banking and shadow banking}\label{sec:di-preview}

\paragraph{The central bank's exposure.} A student asks how the central bank becomes exposed to shadow banks \ts{17}{10}{0:10}. To the traditional
bank the Fed gives a liquidity backstop only; if it is insolvent the FDIC shuts it down and pays depositors from accumulated premia, and beyond that the
government backstops solvency \ts{17}{10}{0:52}. Nobody directly backstops the shadow bank, but there is the liquidity put: when the quality of the
collateral came into question, money funds refused to roll their funding (``you told me these were just like Treasury bills with a yield advantage''), and the
shadow banks drew on their parents' promises, say Citibank's, to roll their funding; so they came back onto traditional banks' balance sheets
\ts{17}{10}{1:33}. Regulators had endorsed the system to get risk out of banks: it removed most solvency risk but not liquidity risk
\ts{17}{10}{2:35}.

\begin{definition}[New concepts: liquidity put]
A commitment by a bank to provide funding to an off-balance-sheet vehicle if it cannot roll its money market funding: in effect an option to ``put'' the funding
problem back to the bank \ts{17}{10}{1:33}.
\end{definition}

\paragraph{Where is solvency?} Most banking courses are about solvency; this course stresses liquidity. In the shadow bank the solvency risk sits with the holders of
the swaps and of the mezzanine and junior tranches. For AAA credit default swaps it was AIG, an insurance company with a lot of capital and bonds, the ``FDIC
equivalent'' of the system, whose capital was sucked out very quickly by the counterparties, Soci\'et\'e G\'en\'erale and Goldman Sachs and their clients, which were
acting as market makers \ts{17}{10}{2:55}.

\begin{intuition}[The fantasy of modern finance]
The shadow banking system had no explicit lender of last resort, on the belief that removing solvency risk removes liquidity risk: a risk-free asset is like a
Treasury bill, hence liquid by design. False: someone must still buy it. Market liquidity is created by dealers quoting bid--ask spreads, and dealers must fund
their inventories; if they cannot, they stop making markets, prices go away, nobody knows what the collateral is worth, and the ability to fund the portfolio goes
away. That is what happened in 2007--09 \ts{17}{10}{5:03}.
\end{intuition}

A traditional bank has no market prices on either side: deposits were forbidden to pay interest, and loans are illiquid, specific to each borrower. A shadow bank
lives on two prices, the market price of the RMBS and the world dollar funding rate \ts{17}{10}{6:26}. The system had moved off
the Fed's balance sheet, ``out of sight, out of mind'', Europe's problem; wrong: a global funding system came under threat almost immediately, and the expansion of
the Fed's balance sheet shown in the first lecture bailed it out \ts{17}{10}{7:49}.

% ------------------------------------------------------------------
\flushside
\section*{Key formulas and ideas}
\begin{keyformulas}
\begin{tabularx}{\linewidth}{@{}l L@{}}
Shadow banking & money market funding of capital market lending; prices on both sides\\
Bagehot's world & money market (bills, indirect) and capital market (bonds, equity, direct) separate\\
New world & long-term bank assets, interbank repo against rated bonds: capital and money markets intertwined\\
Liquidity & funding liquidity (British) versus shiftability $=$ market liquidity (Moulton 1918)\\
Payment vs funding & bank credit as bridge loan; bonds as take-out finance\\
Gurley--Shaw & intermediaries face both ways; but risk passes through to someone\\
Evolution & intermediation $\to$ mutual funds; banks $\to$ MMMFs and shadow banks; quantities $\to$ prices\\
Backstops & traditional: FDIC (solvency), Fed (liquidity); shadow: liquidity puts, AIG as solvency\\
\end{tabularx}
\end{keyformulas}

\section*{Exercises}

\begin{exercise}[Bridge and take-out]
A firm borrows 100 from a bank to buy a machine and, a year later, issues a 10-year bond of 100 bought by a household. Write the T-accounts of bank, firm and
household at each step, and say which step is indirect and which direct finance.
\end{exercise}

\begin{exercise}[Permanent funding]
A bank's deposits have fluctuated between 800 and 1000 over ten years. How much of its balance sheet can it hold in long-term loans, and how much in reserves?
What changes if a money market fund offers a close substitute for deposits?
\end{exercise}

\begin{exercise}[Who bears the risk?]
For a defined benefit pension fund, a defined contribution 401(k), a Jimmy Stewart bank and a shadow bank holding an RMBS hedged with swaps, say who bears the
default risk and who the liquidity risk.
\end{exercise}

\begin{exercise}[Market liquidity needs funding]
Using Treynor's model, explain why a dealer in AAA mortgage securities stops quoting when its repo funding is cut, and what happens to the value of the same
securities as repo collateral.
\end{exercise}

% !TEX root = ../main.tex
\chapter{Forwards and Futures}\label{ch:futures}
\begin{paracol}{2}

\begin{sources}
\begin{tabularx}{\linewidth}{@{}l l L@{}}
\toprule
\textbf{Segment} & \textbf{Length} & \textbf{Content}\\
\midrule
\seg{18}{1}{L18.1 Preview: FT: Argentina in court to fight debt ruling} & 4:56 & Holdouts, pari passu and collective action clauses.\\
\seg{18}{2}{L18.2 Banking as advance clearing} & 5:55 & Changed expectations cause cash flows today: the survival constraint again.\\
\seg{18}{3}{L18.3 Forwards versus futures} & 13:41 & Forward rate agreements as matched book; speculators and futures; $f>$ futures $>\E[r]$.\\
\seg{18}{4}{L18.4 Futures contracts, fluctuations II} & 14:07 & Convergence; the value of a forward contract over its life.\\
\seg{18}{5}{L18.5 Futures contracts, fluctuations II} & 6:07 & Futures marked to market daily: cash flows all along.\\
\seg{18}{6}{L18.6 Cash and carry arbitrage, defined} & 7:28 & Stigum's example; full carry pricing; implied repo rate.\\
\seg{18}{7}{L18.7 Cash and carry arbitrage, explained as liquidity risk} & 5:26 & Fully hedged, but not in cash flow.\\
\seg{18}{8}{L18.8 Cash and carry arbitrage as counterparty risk} & 2:49 & An alternative explanation.\\
\seg{18}{9}{L18.9 Cash and carry arbitrage, as natural banking business} & 3:22 & Banks' comparative advantage in liquidity risk.\\
\bottomrule
\end{tabularx}

\smallskip
\textbf{Files.} Transcripts \file{materials/transcripts/L18-P*.txt}. (The two segments 4 and 5 carry the same title on the playlist, but their content differs.) Mehrling's notes \mn{18}{1}.
\textbf{Reading.} M.~Stigum, \emph{The Money Market}, pp.~718--722 (cash and carry), as cited in class; J.~R.~Hicks, \emph{Value and Capital} (1939), ch.~XI; FT press cutting on Argentina (all in \file{materials/readings/}).
\end{sources}

\section*{Motivation}
The last section of the course is called ``banking as advance clearing'' \ts{18}{2}{0:03}. The first section showed the payment system: cash inflows and
outflows balanced each day, deficits pushed into the future, stress showing up as an interest rate, and defaults when it gets big enough \ts{18}{2}{0:26}.
The economy is coherent not because it is an intertemporal general equilibrium but because the clearing constraint forces everyone to pay attention to their place in
the system \ts{18}{2}{1:49}. Forward-looking financial markets do the same: changed expectations change asset prices (standard finance), and the
money view adds that \emph{price changes cause cash flows today}, tightening or loosening the survival constraint \ts{18}{2}{2:30}.
The dramatic case: the second round of funding is refused, or the bond meant to take out short-term funding cannot be sold, ``then I'm dead''; the everyday case:
pressure on money market prices \ts{18}{2}{4:39}.

% ------------------------------------------------------------------
\section{FT: Argentina in court}\label{sec:ff-ft}

\begin{casestudy}[FT, November 2012: ``Argentina in court to fight debt ruling'']
In 2001 Argentina defaulted. Eventually 93\% of bondholders by value accepted a restructuring (a reduction of the debt) and got new bonds; the other 7\%,
holdouts, mostly hedge funds, sued for full payment \ts{18}{1}{0:07}. A New York court ruled that unless the
holdouts are paid in full, Argentina may not pay the restructured bonds, since all bondholders must be treated equally; this undermines restructuring itself
\ts{18}{1}{1:29}. With a payment on the restructured bonds due the following month, Argentina faces a ``monetary cliff'': pay the
holdouts too, or default on everything \ts{18}{1}{2:51}. The FT's editorial underlines the importance of collective action clauses: one holdout
blocking agreement is a social loss, and taking a loss can be in the bondholders' collective interest \ts{18}{1}{3:12}.

The article in the course file (Jude Webber, \emph{FT}, 26 November 2012, p.~6; \file{materials/readings/L18-Financial_Times_-_Presscuttings.pdf}) adds the details:
\begin{itemize}
  \item The 2001 default was on nearly \$100bn. The restructurings were in 2005 and 2010.
  \item Judge Thomas Griesa ordered Argentina to pay \$1.3bn to the holdouts when it paid the restructured bonds on 15 December, and to deposit the money in escrow pending appeal. Third parties helping it evade the order would count as aiders and abettors.
  \item The dispute turns on the ``pari passu'' (``equal step'') clause.
  \item Argentina's own law forbade reopening the restructuring. Failing to comply could mean a technical default ``perhaps as early as mid-January''.
\end{itemize}
\end{casestudy}

\begin{definition}[New concepts: restructuring, holdout, collective action clause]
\emph{Restructuring}: renegotiating a debt to less (longer terms, lower rates, lower principal). \emph{Holdouts}: creditors who refuse the deal. A \emph{collective
action clause} (CAC) lets a qualified majority of bondholders bind the minority, so there can be no holdouts; the 2001 Argentine bonds, issued under New York law,
had none \ts{18}{1}{0:48}.
\end{definition}

\begin{history}[The end of the Argentine holdout saga]
The ruling was by Judge Thomas Griesa (S.D.N.Y.), interpreting the bonds' \emph{pari passu} clause; it was upheld on appeal, and in June 2014 the US Supreme Court
declined to hear Argentina's appeal. Argentina defaulted again in July 2014 and settled with most holdouts in 2016.\par
\emph{References:} \emph{NML Capital, Ltd.\ v.\ Republic of Argentina}, 699 F.3d 246 (2d Cir.\ 2012); IMF, ``Strengthening the Contractual Framework to Address
Collective Action Problems in Sovereign Debt Restructuring'', 2014.
\end{history}

% ------------------------------------------------------------------
\section{Forwards versus futures}\label{sec:ff-fwdfut}

\paragraph{A forward is advance clearing.} Recall the forward rate agreement as a swap of IOUs between a bank and firm A: a six-month loan against a three-month
deposit, notional amounts, no money changing hands, locking in a three-month rate three months ahead, $f_{3,6}$, defined by forward interest parity
\ts{18}{3}{0:24}:
\begin{keyformulas}
\textbf{Forward interest parity} \ts{18}{3}{1:50}: $(1+r_{0,3})(1+f_{3,6}) = 1+r_{0,6}$, where $r_{0,3}$ and $r_{0,6}$ are today's three- and six-month
rates: the forward rate makes the long rate equal to rolling over two short ones.
\end{keyformulas}
The bank, having promised to lend in three months, hedges with a firm B wanting the opposite, to lock in a deposit rate: the bank's positions offset. In practice
these are forward contracts listed off balance sheet, not parallel loans \ts{18}{3}{2:54}. The
bank has \emph{cleared in advance}: in three months A will be a deficit agent and B a surplus agent, and the flows are already agreed; one way to evade the survival
constraint is to take care of it ahead of time \ts{18}{3}{4:40}. The bank runs a matched book and earns the bid--ask spread
\ts{18}{3}{5:42}.

\paragraph{From matched book to speculators.} There is no reason for equal numbers of fundamental customers on both sides; imbalance pushes prices until dealers
take the other side onto their balance sheets. When the banking system as a whole has as much risk as it wants, banks find hedge funds and speculators in the
\emph{futures} market: a speculator goes long futures to help the bank square its book with a short futures position
\ts{18}{3}{6:22}. The chain: client $\to$ bank $\to$ other banks (FRAs) $\to$ futures exchange
(speculators) $\to$ spot money market, where a rate is realized in three months \ts{18}{3}{8:07}.

\begin{nfigure}
\begin{tikzpicture}[font=\scriptsize,node distance=1.2cm,box/.style={draw,rounded corners,align=center,minimum width=2.1cm,minimum height=0.8cm}]
  \node[box] (c) {client (firm A)\\locks in a loan};
  \node[box,right=0.6cm of c] (b) {bank\\(matched book?)};
  \node[box,right=0.6cm of b] (o) {other banks\\(FRAs)};
  \node[box,right=0.6cm of o] (f) {futures exchange\\(speculators)};
  \node[box,right=0.6cm of f] (s) {spot money\\market at $t=3$};
  \draw[-{Stealth}] (c) -- (b); \draw[-{Stealth}] (b) -- (o); \draw[-{Stealth}] (o) -- (f); \draw[-{Stealth}] (f) -- (s);
  \node[below=0.35cm of o,align=center] {$f_{3,6} \;>\; \text{futures rate} \;>\; \E[r_{3,6}]$: each hedge takes a piece of the premium};
\end{tikzpicture}
\caption{A sequence of hedging transactions \ts{18}{3}{8:31}.}\label{fig:ff-chain}
\end{nfigure}

\paragraph{A stylized fact.} The forward rate is not an unbiased forecast of the future spot rate (the expectations hypothesis fails), and one reason, in this
course, is a mismatched book in the forward market: dealers must be paid to take the risk \ts{18}{3}{9:32}. The sequence of hedges
makes sense only if the forward rate exceeds the expected spot rate, each taking a piece; so one also expects a gap between forward and futures rates
\ts{18}{3}{10:33}:
\begin{keyformulas}
\textbf{Stylized fact (as stated in class)} \ts{18}{3}{11:14}: \quad $f_{3,6} \;>\; \text{futures rate} \;>\; \E[r_{3,6}]$. The forward--futures gap
is small compared with the forward--expected spot gap; finance usually attributes it to something other than liquidity.
\end{keyformulas}

\begin{reading}[Hicks, \emph{Value and Capital} (1939), ch.~XI ``Interest'', pp.~141--152]
The course reading (\file{materials/readings/L18-Hicks_Value_and_Capital.pdf}) is the classic source of this stylized fact.
\begin{itemize}
  \item A loan is a transaction with only one side executed now. Any loan can be reduced to a money loan plus spot and forward transactions. With forward markets, ``own'' rates of interest are implicitly fixed: with money at 5\% and coffee futures 3\% above spot, the coffee rate is about 2\%.
  \item Long rates are built up from \emph{forward} short rates, determined like commodity futures prices by hedgers and speculators.
  \item The forward market for loans has a ``constitutional weakness'' on one side: most people prefer to lend short. Only speculators fill the gap, so the forward rate must exceed the rate they expect, by a risk premium which ``corresponds exactly to the `normal backwardation' of the commodity markets'' (p.~147).
  \item Hence, if short rates are not expected to change, forward rates exceed current short rates by this premium.
\end{itemize}
This is the course's mismatched forward book, with the dealer's required compensation as Hicks's premium.
\end{reading}

\paragraph{The difference is cash flow.} A forward causes no cash flow until maturity; futures are \emph{marked to market} and cause cash flows every day, in or
out, as prices move. So do most derivatives, through margin and collateral calls, ``and those collateral calls cause people to go out of business'': the survival
constraint \ts{18}{3}{12:16}.

\begin{definition}[New concepts: futures contract, marking to market]
A \emph{futures contract} is a standardized forward traded on an exchange, whose price is reset every day; \emph{marking to market} means settling the change in
value in cash each day (through margin), so that the contract's value returns to zero \ts{18}{3}{12:38}\ts{18}{5}{0:04}. (\emph{Futures}: short for contracts for future delivery.)
\end{definition}

% ------------------------------------------------------------------
\section{The value of a forward over its life}\label{sec:ff-value}

\paragraph{Convergence.} Rewrite forward interest parity as \ts{18}{4}{0:03}
\[
  \frac{1}{1+f^{0}_{3,6}} = \frac{1+r_{0,3}}{1+r_{0,6}}, \qquad \frac{1}{1+f^{1}_{3,6}} = \frac{1+r_{1,3}}{1+r_{1,6}},\ \ldots,\ \frac{1}{1+f^{3}_{3,6}} = \frac{1}{1+r_{3,6}},
\]
where $f^{t}_{3,6}$ is the forward rate quoted at time $t$; in general they all differ. At maturity the forward rate equals the spot rate
\ts{18}{4}{1:11}. So the contract between firm A and the bank, a zero-value swap of
IOUs at inception, changes value: a zero-sum game in which, if rates behave as usual, the bank wins, having locked in a lending rate above the expected spot rate;
the firm pays a premium for peace of mind about a serious survival constraint \ts{18}{4}{3:22}.

\begin{intuition}[Coordinating views of the future]
In the notes' framing \mn{18}{1}: in finance the future determines the present, but nobody knows the future, so there are many views of it. How is one path chosen and diverse views
coordinated? By market pricing of the views, and by the effect of prices on behaviour through the survival constraint. At one extreme, if the market changes its mind about your view you
cannot roll your funding; more subtly, changing ideas about the future change cash flows today, loosening the constraint for some and tightening it for others, as futures prices do.
\end{intuition}

The classic forward is the wheat farmer, naturally long wheat, and the baker, naturally short: by a forward contract both lock in the price, each hedging a natural exposure; at delivery one
side ``wins'', but in a forward the winnings are only notional until the end \mn{18}{3}. The forward price is an arbitrage condition: if $K>S_0e^{rT}$, buy the bond spot and sell it forward,
locking in a return above the borrowing rate $r$; if $K<S_0e^{rT}$, sell spot and buy forward \mn{18}{4}. (The notes follow J.~Hull, \emph{Options, Futures, and Other Derivatives}, 5th ed.,
equations 3.5 and 3.9.)

\paragraph{Continuous-time version.} In price terms, with $K$ the delivery (forward) price, $S_t$ the spot price of the underlying (here the six-month bill) and $r$
the interest rate \ts{18}{4}{5:46}:
\begin{keyformulas}
\textbf{Forward price and value} \ts{18}{4}{7:10}:
\[
  K = S_0\,e^{rT} \quad\text{(1)}, \qquad f_t = S_t - K\,e^{-r(T-t)} \quad\text{(2)}, \qquad F_t = S_t\,e^{r(T-t)} \quad\text{(3)}.
\]
$K$ is fixed at inception; the value $f_t$ of the forward to the long side moves with the spot price and with the shrinking discount as maturity approaches:
$f_0=0$ and $f_T = S_T - K$. The futures price $F_t$ (3) is $K$ reset every day so that the contract's value is again zero \ts{18}{5}{0:04}.
\end{keyformulas}

Banking deals constantly with zero-value contracts that later acquire value, ``a little mind-blowing'' \ts{18}{4}{10:00}. The
same apparatus applies to commodities with storage costs and to coupon bonds with more complicated formulas; here we keep discount bonds of the money market
\ts{18}{4}{11:28}. With a forward there are no cash flows until the end, when the shorts pay the longs $S_T-K$ (in financial
contracts by cash netting) or the reverse \ts{18}{4}{12:30}.

\begin{nfigure}
\begin{tikzpicture}[font=\scriptsize,x=1cm,y=1cm]
  \draw[-{Stealth[length=4pt]}] (0,0) -- (8.3,0) node[right]{$t$};
  \draw[-{Stealth[length=4pt]}] (0,-1.6) -- (0,2.2) node[above]{$f_t$ (value to the long side)};
  \draw[dashed] (7.5,-1.6) -- (7.5,2.0) node[above]{$T$};
  \draw[very thick,defcol] (0,0) .. controls (1,1.2) and (1.8,1.4) .. (2.5,0.6) .. controls (3.3,-0.4) and (4,-1.3) .. (4.8,-0.6) .. controls (5.6,0.2) and (6.6,1.2) .. (7.5,1.0);
  \fill (7.5,1.0) circle (1.5pt) node[right,align=left]{$S_T-K$: the only\\cash flow of a forward};
  \node[align=left] at (1.5,1.75) {futures: cash to longs};
  \node[align=left] at (4.0,-1.45) {futures: cash to shorts};
  \node[align=left] at (6.4,1.7) {to longs};
\end{tikzpicture}
\caption{The same price history, two cash-flow profiles: a forward pays once at maturity, a futures contract pays every day as the value changes
\ts{18}{4}{11:01}\ts{18}{5}{1:47}. Schematic.}\label{fig:ff-paths}
\end{nfigure}

\paragraph{Margin.} The payments go through margin accounts at the futures clearinghouse, and both sides post margin, since at inception both are equally likely to lose; the clearinghouse
accepts liquid securities as margin, repriced every day, so collateral and contract are both marked to market. The cumulative futures payments equal the forward's final payment, but they come
every day: ``a very concrete way in which views about the future are settled today'' \mn{18}{6}. (The notes' formula abstracts from interest on the fluctuating margin balances \mn{18}{5}.)

\paragraph{Futures cash flows.} Whenever the futures price changes, cash moves from losers to winners: while the forward's value rises, from the shorts to the
longs; when it falls, back \ts{18}{5}{1:06}. With a forward, who was winning in the middle does
not matter; with futures the largest swings may come in the middle, and you must be ready to absorb them. Futures carry liquidity risk that forwards do not, which
may be why their prices differ, though the formulas treat them as interchangeable \ts{18}{5}{2:50}.

% ------------------------------------------------------------------
\section{Cash and carry arbitrage}\label{sec:ff-cashcarry}

\begin{aside}[Ten years of worrying]
Stigum's pages 718--722 show real cash and carry trades making money, which, by finance theory, should not be possible. ``This lecture is a consequence of 10 years
of me figuring this stuff out'': find something that does not seem right and keep worrying at it \ts{18}{5}{4:12}.
\end{aside}

\paragraph{The trade.} Buy a six-month Treasury bill and finance it with a three-month repo: borrowing short, lending long; the price risk (the value of the then
three-month bill in three months) is hedged by a short futures position \ts{18}{6}{0:03}.
Long bill plus short repo is a synthetic forward (the six-month loan against the three-month deposit of the balance sheets): so the position is \emph{long forward,
short futures} for the same term. Why should it make any profit? \ts{18}{6}{2:08}

\begin{nfigure}
\begin{taccts}
\tacct[2.6cm]{Cash and carry position}{6-month Treasury bill}{3-month repo\\Short futures}
\end{taccts}
\caption{Stigum's cash and carry: long positions above, short positions below the line in class; long bill $+$ short repo $=$ synthetic long forward
\ts{18}{6}{1:48}.}\label{fig:ff-cc}
\end{nfigure}

\begin{keyformulas}
\textbf{Full carry pricing and the implied repo rate} \ts{18}{6}{3:51}:
\[
  F_0 = S_0\,e^{rT}\ \text{(full carry)}, \quad F_0 > S_0\,e^{rT}\ \text{(cash and carry)}, \quad F_0 < S_0\,e^{rT}\ \text{(reverse).}
\]
The \emph{implied repo rate} $\rho$ is defined by $F_0 = S_0\,e^{\rho T}$: the rate that would make full carry pricing true. Cash and carry ($\rho>r$) is like
borrowing at the actual repo rate $r$ and lending at $\rho$: an interest rate spread.
\end{keyformulas}

\begin{reading}[Stigum's example, \emph{The Money Market}, pp.~718--722]
Stigum's own numbers (\file{materials/readings/L18-Pages_from_Stigum.pdf}) are dated 28 October 1982, ``when liquidity in T-bills was at its peak''.
\begin{enumerate}
  \item Buy \$1 million of the 24 March 1983 bill, 147 days to maturity, at a discount rate of 8.26: price \$966{,}271.67.
  \item Sell the December 1982 bill future, which expires in 56 days, at 91.72 (rate 8.28). Delivering the bill, then a 91-day bill, brings in \$979{,}070.00.
  \item The holding-period yield over the 56 days, on a 360-day basis, is the \emph{implied repo rate}: 8.51\%.
  \item With 56-day term repo at 8.25, the trade earns 26\,bp. On \$10 million that is \$3{,}908.03.
\end{enumerate}
Stigum lists the ``slips twixt the cup and the lip'': commission, margin calls, convergence-price risk and back-office costs. He judges margin calls a small threat. Even a 100\,bp rally on the first day would cost only 2\textonequarter\,bp of the spread. In Mehrling's reading this cash-flow timing is the whole point: what is left after the price risk is hedged is liquidity risk.

(Table 15.4 dates the futures sale ``March 28, 1982''; the context shows it is 28 October 1982.)
\end{reading}

\paragraph{Explained as liquidity risk.} Look at the cash flows \ts{18}{7}{0:03}:
\begin{itemize}
  \item \textbf{long forward}: since the expectations hypothesis fails, the long side probably wins, but receives cash only at $T$;
  \item \textbf{short futures}: probably loses, and pays all along, with margin calls until $T$; if you cannot keep up, the clearinghouse liquidates your
        position \ts{18}{7}{1:45}.
\end{itemize}
All price risk is gone; what is left is liquidity risk: the inflows will cover the outflows, ``but not at the same time''. The position is risk-free only if one could
somehow use the certain final inflow to pay the earlier outflows \ts{18}{7}{2:26}. The same banking view that explains the failure
of EH (an imbalance in the forward market) explains Stigum's confusing pages: forwards and futures do not trade at the same price \ts{18}{7}{3:28}. The notes stress that the negative
flows might all come at once: the volatility of the bill's spot price is what creates liquidity risk, and the two anomalies, the failure of EH and the profit of cash and carry, have the same origin,
a mismatch in the natural forward interest rate positions of the real economy. A former teaching assistant of the course, Daniel Neilson, showed how forward--futures deviations can serve as
an observable proxy for the unobservable deviation of forward rates from expected spot rates, a test of the liquidity premium theory of the term structure \mn{18}{8}.

\begin{coursenote}
Mehrling notes that the standard explanation, reinvestment of daily futures cash flows, ``goes in the wrong direction'' and may even hide the size of the liquidity
premium \ts{18}{7}{4:12}. In the literature this is the forward--futures price difference under stochastic interest
rates (J.~C.~Cox, J.~E.~Ingersoll and S.~A.~Ross, ``The relation between forward prices and futures prices'', \emph{Journal of Financial Economics} 9, 1981); for
interest-rate futures it implies a futures rate \emph{above} the forward rate (the ``convexity adjustment''), the opposite of the class's stylized fact
$f>$ futures.
\end{coursenote}

\paragraph{Or counterparty risk?} A student suggests credit risk: a forward with a repo counterparty is one big payment at the end, never sure until paid, while
futures are margined by an exchange, with warning from many small payments. More credit risk in the forward would call for compensation in the right direction
\ts{18}{8}{0:03}. Typical of these debates: a spread exists, and one must list all candidate causes
(cash-flow timing, counterparty risk, reinvestment) and do the empirical work to separate them; probably traders have done it and are making money, but ``it's not
been published in the academic literature'' \ts{18}{8}{1:46}.

\paragraph{A natural banking business.} Another student: futures need cash at a moment's notice, so how does that affect banks? Banks use futures to hedge forward
positions, so value fluctuations offset; the naked speculator on the other side must post margin. What the bank is not hedged against is the liquidity risk, but a
bank has the discount window and interbank relationships, a comparative advantage in liquidity risk, so cash and carry is natural banking business; speculators bear
liquidity risk and are paid the term premium, the difference between the forward rate and the expected spot rate
\ts{18}{9}{0:01}.

% ------------------------------------------------------------------
\flushside
\section*{Key formulas and ideas}
\begin{keyformulas}
\begin{tabularx}{\linewidth}{@{}l L@{}}
Advance clearing & forward contracts settle future payments today; asset price changes cause cash flows today\\
Forward parity & $(1+r_{0,3})(1+f_{3,6})=1+r_{0,6}$; at maturity the forward rate equals the spot rate\\
Forward & $K=S_0e^{rT}$, $f_t=S_t-Ke^{-r(T-t)}$; one cash flow $S_T-K$ at $T$\\
Futures & $F_t=S_te^{r(T-t)}$ reset daily; marked to market, cash flows all along\\
Stylized fact & $f_{3,6}>$ futures rate $>\E[r_{3,6}]$: each hedger takes part of the liquidity premium\\
Cash and carry & long bill $+$ short repo $+$ short futures; profit if $\rho>r$; the residual risk is liquidity risk\\
\end{tabularx}
\end{keyformulas}

\section*{Exercises}

\begin{exercise}[Forward interest parity]
Three-month and six-month rates are 1.0\% and 2.2\% (not annualized). Compute $f_{3,6}$. One month later the two-month and five-month rates are 0.7\% and 1.8\%: compute
the new forward rate and say whether the bank that lent forward has gained or lost.
\end{exercise}

\begin{exercise}[Forward versus futures cash flows]
A forward and a futures contract on the same bill have $K=F_0=98$. The futures price then goes 98.5, 97.0, 96.5, 98.0, 99.0 at maturity. Compute the daily cash
flows of a short futures position and the single cash flow of a short forward. What is the largest cumulative cash outflow of the futures position?
\end{exercise}

\begin{exercise}[Implied repo rate]
A six-month bill trades at $S_0=99$ and the three-month futures price on it is $F_0=99.6$; the three-month repo rate is 2\% a year (continuous compounding). Compute the
implied repo rate. Is a cash and carry trade profitable, and what liquidity must the trader hold?
\end{exercise}

\begin{exercise}[Who should do it?]
Explain why a bank with access to the discount window, rather than a hedge fund, is the natural holder of a cash and carry position, and what happens to the
forward--futures spread if banks' access to liquidity is restricted.
\end{exercise}

% !TEX root = ../main.tex
\chapter{Interest Rate Swaps}\label{ch:irs}
\begin{paracol}{2}

\begin{sources}
\begin{tabularx}{\linewidth}{@{}l l L@{}}
\toprule
\textbf{Segment} & \textbf{Length} & \textbf{Content}\\
\midrule
\seg{19}{1}{L19.1 FT: Sovereign debt crises} & 3:21 & Argentina, Greece, Draghi; pension funds leaving long sovereign debt.\\
\seg{19}{2}{L19.2 Reading: FOMC Report (1952)} & 7:40 & \emph{No captions available; not covered in these notes} (see the coursenote in \cref{sec:irs-arbland}).\\
\seg{19}{3}{L19.3 Treasury-swap spread, a puzzle} & 11:04 & Negative swap spreads at the long end; what a swap is; swap lingo.\\
\seg{19}{4}{L19.4 What is a swap} & 10:35 & Same captions as L19.5.\\
\seg{19}{5}{L19.5 Why swap? An example from Stigum} & 10:37 & AA and BBB share 35 basis points; exposure in a swap.\\
\seg{19}{6}{L19.6 Market making in swaps} & 8:42 & The swap dealer's matched book; where capital market prices come from.\\
\seg{19}{7}{L19.7 Money market swaps, example} & 5:57 & Stigum's one-year swap through J.P.~Morgan.\\
\seg{19}{8}{L19.8 Life in Arbitrage Land} & 5:04 & Stigum's ``arbitrage land''; from corporate bonds to mortgages.\\
\seg{19}{9}{L19.9 Treasury-swap spread} & 7:21 & Why nobody arbitrages the negative swap spread: liquidity risk.\\
\bottomrule
\end{tabularx}

\smallskip
\textbf{Files.} Transcripts \file{materials/transcripts/L19-P*.txt}; no transcript exists for L19.2 (YouTube has no captions for it). The video published as L19.4
(``What is a swap'') contains the same lecture as L19.5 (the two auto-caption files are independent but match word for word with the same timings): an error in the
playlist, so L19.5 is cited throughout. The material that title promises is covered at the end of L19.3.
\textbf{Reading.} Federal Open Market Committee, report of the Ad Hoc Subcommittee on the Government Securities Market (1952); M.~Stigum, \emph{The Money Market}.
\end{sources}

\section*{Motivation}
Interest rate swaps are a private-sector yield curve, a kind of corporate bond term structure, usually seen as a spread over the Treasury curve
\ts{19}{3}{0:02}. Since the crisis the spread has gone ``a little crazy'', negative at the long end: an apparent arbitrage nobody takes
\ts{19}{3}{0:46}. To understand why, the lecture translates swaps into the course's balance sheets, shows the
swap dealers who make the market, and finds that the prices of the capital market are made by dealers who fund themselves in the money market.

% ------------------------------------------------------------------
\section{FT: sovereign debt crises}\label{sec:irs-ft}

\begin{casestudy}[FT, November 2012: sovereign debt]
A full page on Argentina's ``unforgiving debt'' after the New York ruling on the holdouts (``vulture'' bondholders), quoting people from a research project Mehrling
was involved in; a full page on Greece, where, as Gavyn Davies writes, the eurozone has embarked on official debt forgiveness while avoiding outright
cancellation of nominal debt: most Greek debt is now held by official European entities, so the renegotiation is with euro partners, not bondholders
\ts{19}{1}{0:07}. Spanish and Italian yields have fallen after Draghi's announcement of Outright
Monetary Transactions, though nothing has been bought yet, as the theory of the course would predict (a midterm question) \ts{19}{1}{1:52}.
``This party is over'': some pension funds are giving up long-dated sovereign debt, expecting yields to rise, possibly related to the Treasury and swap curves
\ts{19}{1}{2:32}.
\end{casestudy}

% ------------------------------------------------------------------
\section{The Treasury--swap spread: a puzzle}\label{sec:irs-puzzle}

On the day of the lecture: at one year the swap rate is a little above the Treasury rate; at ten years they are on top of each other; at 30 years the swap spread is
negative \ts{19}{3}{1:07}. That suggests an arbitrage, short the ``corporate'' rate and long Treasuries. Why does nobody do it? The
negative long-end spread has persisted throughout and since the crisis: something is seriously skewed in the structure of the capital market
\ts{19}{3}{1:48}.

\begin{nfigure}
\begin{tikzpicture}[font=\scriptsize,x=0.36cm,y=1.6cm]
  \draw[-{Stealth[length=4pt]}] (0,0) -- (31.5,0) node[right]{maturity (years)};
  \draw[-{Stealth[length=4pt]}] (0,0) -- (0,3.2) node[above]{yield};
  \foreach \x in {1,5,10,30} \draw (\x,0.05) -- (\x,-0.05) node[below]{\x};
  \draw[very thick,defcol] (1,0.25) -- (5,0.75) -- (10,1.7) -- (30,2.85) node[right]{Treasury};
  \draw[very thick,thmcol] (1,0.4) -- (5,0.85) -- (10,1.68) -- (30,2.55) node[right]{swap};
  \node[align=left] at (21,0.9) {schematic shape: swap above Treasury at 1 year,\\equal at 10 years, below at 30 years};
\end{tikzpicture}
\caption{Treasury and swap yield curves, late 2012, as shown in class (four maturities joined by straight lines) \ts{19}{3}{1:07}. Shape only, not the actual numbers.}\label{fig:irs-curves}
\end{nfigure}

% ------------------------------------------------------------------
\section{What is a swap}\label{sec:irs-what}

Stigum's example: an AA firm can borrow relatively cheaply at a fixed long-term rate (a ten-year bond), a BBB firm most cheaply at a short-term floating rate (say
three-month LIBOR); BBB worries about rolling over and would like to lock in long-term funding, AA is happy with short-term borrowing. They swap exposures
\ts{19}{3}{3:11}.

\paragraph{A swap is a parallel loan.} As always, ask what IOUs are notionally swapped \ts{19}{3}{4:55}:
\begin{itemize}
  \item AA, wanting floating: lends fixed and borrows at LIBOR, i.e.\ \emph{receives fixed, pays floating};
  \item BBB, wanting fixed: \emph{pays fixed, receives floating}.
\end{itemize}
In reality there is no loan and no principal: the parties net the interest payments, so the swap is off balance sheet and does not show up as leverage
\ts{19}{3}{7:42}.

\begin{definition}[New concepts: interest rate swap, swap rate]
An \emph{interest rate swap} is an agreement to exchange, on a notional principal, fixed-rate interest payments for floating-rate payments (e.g.\ six-month LIBOR)
for the life of the swap; only the net payment changes hands. The \emph{swap rate} is the fixed rate quoted against the floating leg \ts{19}{3}{5:59}.
(\emph{Swap}: exchange, from the old sense ``to strike hands'' on a bargain.)
\end{definition}

\begin{finance}[Reading the FT quotes]
The FT's second section quotes swaps daily for euros, pounds, Swiss francs, dollars and yen, from one to 30 years, with a bid and an ask: for the ten-year dollar
swap, 1.66--1.69, against six-month LIBOR as the footnotes say \ts{19}{3}{6:19}.
\end{finance}

\paragraph{Street lingo.} AA, receiving fixed and paying floating, is the \emph{seller} of the swap, ``short'' the swap; BBB, paying fixed, is the \emph{buyer}, ``long''
the swap. Following the lingo, the course books a short swap on the liability side and a long swap on the asset side, though actual balance sheets have no such
convention \ts{19}{3}{8:23}.

\begin{intuition}[A short swap is like a repo'd bond]
A short swap position (receive fixed, pay floating) is like owning a fixed-rate bond and financing it short term in the repo market: a long-term repo in which a
corporate bond is the collateral \ts{19}{3}{9:49}.
\end{intuition}

% ------------------------------------------------------------------
\section{Why swap? An example from Stigum}\label{sec:irs-why}

AA and BBB ask their bankers for a floating-rate term loan and a five-year fixed-rate Eurobond \ts{19}{5}{0:01}:

\begin{nfigure}
\renewcommand{\arraystretch}{1.35}
\begin{tabular}{lccc}
\toprule
 & AA & BBB & difference\\
\midrule
floating (term loan) & six-month LIBOR $+\tfrac18$ & six-month LIBOR $+\tfrac14$ & 12.5 bp\\
fixed (five-year Eurobond) & 5.375\% & 5.85\% & 47.5 bp\\
\midrule
room for a deal & \multicolumn{3}{c}{$47.5-12.5 = 35$ bp}\\
\bottomrule
\end{tabular}
\caption{Stigum's quotes \ts{19}{5}{1:05}.}\label{tab:irs-stigum}
\end{nfigure}

\paragraph{Comparative advantage.} AA has an absolute advantage at every maturity, but its advantage is much larger in the long market: BBB has a comparative
advantage in short-term borrowing. Mehrling is ``not sure that is really a very good way to understand this, but the math works that way''
\ts{19}{5}{2:51}.

\paragraph{The deal.} AA issues fixed at 5.375\% and does a swap with BBB at 5.50\% against LIBOR flat; BBB borrows at six-month LIBOR $+\tfrac14$
\ts{19}{5}{4:14}:
\begin{itemize}
  \item BBB pays $(\text{LIBOR}+0.25) + 5.50 - \text{LIBOR} = 5.75\%$ fixed: 10 bp less than its 5.85\% \ts{19}{5}{5:17};
  \item AA pays $5.375 + \text{LIBOR} - 5.50 = \text{LIBOR}-\tfrac18$: 25 bp less than its LIBOR $+\tfrac18$ \ts{19}{5}{5:59}.
\end{itemize}
The 35 bp are shared 25 to AA and 10 to BBB; the split is negotiated, and one would expect the better borrower to get most of it \ts{19}{5}{2:30}.

\begin{nfigure}
\begin{taccts}
\tacct[2.6cm]{AA}{(Fixed loan to BBB, 5.50) $=$ short swap}{Fixed bond 5.375\%\\(LIBOR loan from BBB)}\qquad
\tacct[2.6cm]{BBB}{(LIBOR loan to AA) $=$ long swap}{LIBOR $+\tfrac14$ bank loan\\(Fixed loan from AA, 5.50)}
\end{taccts}
\caption{Stigum's swap as a notional parallel loan (in parentheses: not actual loans; only net payments are made) \ts{19}{5}{4:14}.}\label{fig:irs-parallel}
\end{nfigure}

\paragraph{Exposure.} Why can AA get funding below LIBOR? It has an exposure to BBB. BBB has hedged LIBOR movements, not its own credit spread: if in six months its
banker wants LIBOR $+1$, it still receives only LIBOR flat on the swap, may get into trouble rolling its loan, and then paying AA \ts{19}{5}{6:40}.
But it is not credit risk on a principal: if BBB stops paying, AA stops paying and tears up the swap, losing only the change in value of net payments, far less than
on an actual parallel loan \ts{19}{5}{8:27}.

\paragraph{Market imperfections.} Many swaps exist because one borrower has cheap local funding (the government likes it, capital controls, an immature bond market)
and swaps it to others, breaking down inequalities across markets until the bond market grows up \ts{19}{5}{9:28}.

% ------------------------------------------------------------------
\section{Market making in swaps}\label{sec:irs-dealers}

The 35 bp are an incentive for a dealer to set the deal up and keep a couple of basis points. In the FT the bid--ask spread is 0.32--0.35 at one year, three basis
points, and three basis points at two and 30 years too: the dealer's spread \ts{19}{6}{0:01}. An investment bank
stands between AA (short a swap) and BBB (long a swap), doing a swap with each, and runs a swap book, here shown as a matched book earning the spread
\ts{19}{6}{1:24}.

\begin{nfigure}
\begin{taccts}
\tacct[1.9cm]{AA}{}{Short swap}\quad
\tacct[1.9cm]{Swap dealer (investment bank)}{Swap with BBB (bid)}{Swap with AA (ask)}\quad
\tacct[1.9cm]{BBB}{Long swap}{}
\end{taccts}
\caption{The swap dealer's matched book \ts{19}{6}{1:24}.}\label{fig:irs-dealer}
\end{nfigure}

\paragraph{Where prices come from.} The swap yield curve comes from dealers quoting bid and ask; customers decide whether to buy or sell and need not find a
counterparty, which is the dealer's problem; the dealer hedges imbalances with Treasuries, forwards or futures \ts{19}{6}{2:46}. With an investment
bank in the middle, AA's exposure is to Goldman Sachs, not to BBB; the dealer holds a diversified portfolio of low-rated counterparties on one side and high-rated
on the other, in effect a diversified corporate bond portfolio in swap space that takes up no balance sheet room \ts{19}{6}{5:31}.
Usually the fixed rate exceeds the floating, so BBB pays AA the difference over the swap's life, while still paying its bank; if its six-month loan is not renewed,
it is in crisis: that is why the market lent to it only for six months \ts{19}{6}{4:09}.

\begin{keyformulas}
\textbf{The swap market is the government securities dealer market for corporate bonds} \ts{19}{6}{6:55}.
Interest rate swaps are where forward rates all along the yield curve come from; the Fed, by acting on the short end, also influences the swap curve, since the two
curves trade off each other and the same dealers often run both books, hedging one with the other. The dealers are highly leveraged and fund themselves in the
money market: capital market prices are much closer to the money market than one thought.
\end{keyformulas}

% ------------------------------------------------------------------
\section{Money market swaps}\label{sec:irs-mmswap}

Stigum's second example, a money market swap, mostly an interbank market \ts{19}{7}{0:02}. An AA bank receives a one-year deposit
from a client but wants a floating-rate liability, to match its assets: it swaps with J.P.~Morgan, one year fixed against three-month LIBOR, at 4.44\%
\ts{19}{7}{1:05}. Morgan, now exposed to interest rate risk, may hedge short term in futures or
FRAs until it finds the opposite customer: a leveraged buyout firm (LBO) with a three-month floating bank loan that wants to swap into one year, at 4.47\%
\ts{19}{7}{2:56}.

\begin{nfigure}
\begin{taccts}
\tacct[2.2cm]{AA bank}{(3-month LIBOR loan to Morgan)}{1-year deposit\\(1-year loan from Morgan, 4.44)}\quad
\tacct[2.2cm]{J.P.~Morgan (dealer)}{(1-year, 4.44, to AA bank)\\(3-month LIBOR to LBO)}{(3-month LIBOR from AA bank)\\(1-year, 4.47, from LBO)}\quad
\tacct[2.2cm]{LBO}{(1-year, 4.47, to Morgan)}{3-month bank loan\\(3-month LIBOR from Morgan)}
\end{taccts}
\caption{A money market swap as notional parallel loans: Morgan receives 4.47 and pays 4.44, the LIBOR legs net out \ts{19}{7}{4:25}.}\label{fig:irs-mm}
\end{nfigure}

\begin{coursenote}
Students often object that 4.47 on a liability exceeds 4.44 on an asset, which does not look like profit. The lingo trips you up: think of the parallel loans
behind the swaps. Morgan is paid 4.47 by the LBO and pays 4.44 to the AA bank, and the two LIBOR legs net: three basis points. Booking a ``short swap'' as a
liability does not mean subtracting it from assets to get net worth \ts{19}{7}{4:25}.
\end{coursenote}

% ------------------------------------------------------------------
\section{Life in arbitrage land}\label{sec:irs-arbland}

Swap market making happens at all maturities, one to 30 years. Liquidity lives in derivatives markets because they use little capital: no loans, ``swaps of IOUs
all the way down'', profit from selling one IOU for more than you buy the other, cash flows netted. ``This is banking, but without the actual loans'': the core of
the modern system is not Jimmy Stewart banking but swap dealers quoting two-sided markets and so making prices \ts{19}{8}{0:02}.

\begin{reading}[M.~Stigum, \emph{The Money Market}, p.~900: arbitrage land]
Stigum describes money market swaps as living in what could be called ``arbitrage land'': traders arbitrage swaps against futures, against cash, against FRAs, FRAs
against futures, and so on. An event that moves one rate sends a ripple through the related markets, creating basis points for every player, except perhaps the
unhedged speculator in the futures pit, who usually is the one who loses \ts{19}{8}{1:24}.
\end{reading}

\paragraph{A hierarchy of arbitrages.} The chain ends in the futures market, where price risk is hedged and liquidity risk is borne: hence last lecture on forwards
and futures \ts{19}{8}{2:26}. The same idea is in the 1952 FOMC report: government securities dealers doing term structure arbitrage, buying the
five-year when cheap, financing it in a primitive repo market, selling the ten-year. ``1952 is the future'', extended from government securities to corporate
securities and then to mortgages: an RMBS is a 30-year bond with option features (early repayment), and the swap apparatus built for corporate securities was moved
into mortgage space, producing a global market in household bonds: ``AA, LBO, subprime'' \ts{19}{8}{2:47}.

\begin{coursenote}
The segment devoted to the 1952 FOMC report (L19.2) has no captions on YouTube, so its content is not reported here; only what is said about it in other segments
\ts{19}{6}{7:15}\ts{19}{8}{2:47}. The report is that of the FOMC's Ad Hoc Subcommittee on the Government Securities Market (chaired by William
McChesney Martin, November 1952), which recommended that open market operations be confined to the short end of the market, normally Treasury bills
(``bills only''), leaving the rest of the yield curve to dealers' arbitrage.\par
\emph{Reference:} Federal Reserve, \emph{Federal Reserve Bulletin}, 1953; FOMC Minutes, March 1953.
\end{coursenote}

% ------------------------------------------------------------------
\section{The Treasury--swap spread explained}\label{sec:irs-spread}

To exploit a swap curve below the Treasury curve, one would borrow at the ``corporate'' rate and lend at the Treasury rate: be long a Treasury, financed in repo
with the Treasury as collateral, and long a swap (pay fixed, receive floating) \ts{19}{9}{0:22}.
Corporations that can borrow are in effect doing it on their own balance sheets, issuing 30-year debt and holding cash or Treasuries; it is not enough to close
the gap \ts{19}{9}{2:27}.

\begin{nfigure}
\begin{taccts}
\tacct[2.6cm]{Arbitrageur}{30-year Treasury\\Long swap (pay fixed, receive floating)}{Repo (short-term funding)}
\end{taccts}
\caption{The arbitrage that would close a negative swap spread: no credit risk, but borrowing short to lend long \ts{19}{9}{1:24}.}\label{fig:irs-arb}
\end{nfigure}

\begin{intuition}[Why the gap stays open]
The position borrows short and lends long: no credit risk (it is the Treasury), but liquidity risk, since one does not know whether, or at what price, the funding
can be rolled. Even with the market flooded by QE, people do not want that risk. Other stories (capital constraints coming with Basel~III) are possible; any story
must explain why nobody has done the arbitrage for four years \ts{19}{9}{3:09}.
\end{intuition}

\begin{aside}[The fiscal cliff]
A manufactured crisis: unless Congress agreed by a date, taxes would rise and spending fall, with large demand effects. Mehrling sees political brinkmanship; it may
lead people to stay long cash, and firms to pay out cash as dividends ahead of tax rises, but it cannot explain a four-year anomaly
\ts{19}{9}{4:11}.
\end{aside}

\paragraph{Where the mortgage rate comes from.} The interest rate swap market is enormous, ``trillions, trillions, trillions''. When a Jimmy Stewart banker quotes you
a mortgage rate, it comes from here: the mortgage will be packaged and sold as a bond, so the banker looks at the swap curve for mortgage-backed securities. Prices
are made in dealer markets in derivatives, with supply and demand for loans behind the scenes: ``that's the way modern banking works''
\ts{19}{9}{5:54}.

% ------------------------------------------------------------------
\flushside
\section*{Key formulas and ideas}
\begin{keyformulas}
\begin{tabularx}{\linewidth}{@{}l L@{}}
Swap & notional parallel loan, fixed against floating; only net payments; off balance sheet\\
Lingo & receive fixed $=$ seller, short the swap; pay fixed $=$ buyer, long the swap\\
Analogy & short swap $\approx$ fixed-rate bond financed by term repo\\
Stigum's gains & $(5.85-5.375)-(\tfrac14-\tfrac18) = 0.35$: AA gets 25 bp (LIBOR $-\tfrac18$), BBB 10 bp (5.75\%)\\
Dealers & swap book earns the bid--ask (about 3 bp); swap curve made by leveraged dealers funded in the money market\\
Arbitrage land & arbitrages chain to the futures market, where liquidity risk ends up\\
Negative spread & long Treasury in repo $+$ long swap: no credit risk, but funding liquidity risk\\
\end{tabularx}
\end{keyformulas}

\section*{Exercises}

\begin{exercise}[Splitting the gains]
With Stigum's quotes, find the swap fixed rate at which AA and BBB split the 35 bp equally, and show each firm's all-in cost. What limits the range of feasible
swap rates?
\end{exercise}

\begin{exercise}[Swap cash flows]
A five-year swap on \$100 million pays 5.50\% fixed against six-month LIBOR, settled semi-annually. LIBOR fixes at 5.0\%, 5.6\%, 6.2\%. Compute the net payment
each period and who pays it.
\end{exercise}

\begin{exercise}[A dealer's mismatched book]
Morgan has done the AA bank's swap but finds no LBO for a week. Write its exposure and describe how it can hedge with three-month Eurodollar futures. What does it
earn when the LBO arrives at 4.47\%?
\end{exercise}

\begin{exercise}[The negative swap spread]
The 30-year Treasury yields 2.85\% and the 30-year swap rate is 2.55\%. Write the balance sheet of the arbitrage, the annual carry if three-month repo costs 0.20\%
and three-month LIBOR is 0.30\%, and the risks that remain.
\end{exercise}

% !TEX root = ../main.tex
\chapter{Credit Default Swaps}\label{ch:cds}
\begin{paracol}{2}

\begin{sources}
\begin{tabularx}{\linewidth}{@{}l l L@{}}
\toprule
\textbf{Segment} & \textbf{Length} & \textbf{Content}\\
\midrule
\seg{20}{1}{L20.1 FT: Internationalization of the Euro} & 5:07 & Offshore euro trading in London; bonds replacing syndicated loans.\\
\seg{20}{2}{L20.2 Credit indices} & 3:40 & The Markit CDX high-yield index.\\
\seg{20}{3}{L20.3 Fischer Black (1970), risk-free security} & 2:14 & A risky bond $=$ risk-free bond $+$ interest rate risk $+$ default risk.\\
\seg{20}{4}{L20.4 What is a Credit Default Swap (CDS)} & 5:12 & Buyer and seller of protection on the balance sheet.\\
\seg{20}{5}{L20.5 Corporate bonds} & 3:43 & Price as present value; spreads and ratings.\\
\seg{20}{6}{L20.6 CDS pricing} & 11:24 & Bond $+$ IRS $=$ LIBOR $+s$; CDS premium $u$; arbitrage $s=u$; is it insurance?\\
\seg{20}{7}{L20.7 Market making in CDS} & 4:03 & Dealers hedge single names with the index.\\
\seg{20}{8}{L20.8 Example: Negative basis trade and Liquidity Risk} & 10:20 & UBS: AAA CDO tranches, insured, funded in the money market.\\
\seg{20}{9}{L20.9 Example: Private backstop of marketmaking in CDS} & 15:15 & AIG, Goldman Sachs and collateral calls: ``liquidity kills you quick''.\\
\seg{20}{10}{L20.10 Example: Synthetic CDO as Collateral Prepayment} & 6:43 & Paulson, Abacus and IKB.\\
\bottomrule
\end{tabularx}

\smallskip
\textbf{Files.} Transcripts \file{materials/transcripts/L20-P*.txt}; L20.1 has no YouTube captions and was transcribed locally from the audio (faster-whisper). Mehrling's notes \mn{20}{1}.
\end{sources}

\section*{Motivation}
All semester Mehrling promised that the course's balance sheets would handle credit default swaps; now he delivers \ts{20}{2}{0:06}. Think of a corporate bond with a lot
of credit risk: a credit default swap peels that risk off and sells it separately, just as an interest rate swap peels off interest rate exposure
\ts{20}{2}{3:11}. Three examples from the crisis (UBS, AIG, Abacus) then show that the trouble came not from defaults but from liquidity.

% ------------------------------------------------------------------
\section{FT: the internationalization of the euro}\label{sec:cds-ft}

Mehrling had just come back from Rome, where Italian central bankers and securities regulators asked his shadow banking team whether shadow banking could channel
credit to small and medium-sized enterprises: they depend on banks, and banks, having bought government bonds, now need to build up capital and are not lending
\ts{20}{1}{0:00}. He reads two articles of the day as symptoms of sea changes in the European financial system
\ts{20}{1}{1:24}.

\begin{casestudy}[FT, November 2012: two European evolutions]
\begin{enumerate}
  \item ``France targets UK euro trade supremacy'': 40\% of euro trading takes place in London, and France wants it inside the eurozone, under control
        \ts{20}{1}{1:35}. Mehrling recalls the rise of the Eurodollar market in London and the old American worry about losing control of the dollar;
        the US concluded it could not stop it, and the same holds for the euro: people who want to trade euro derivatives will find somewhere to do it, London or
        Hong Kong. A ``little naive'', but a sign of the maturation of the euro as an international currency \ts{20}{1}{2:03};
  \item ``Syndicated bank loans in eurozone slump to ten-year low'': top-rated corporations no longer take syndicated loans (a bank making a loan much bigger than it
        could hold and parcelling out pieces to other banks) but issue bonds directly, in big amounts: the capital market replaces bank lending, as in the US
        evolution from indirect to direct finance (lecture~\ref{ch:directindirect}) \ts{20}{1}{3:21}.
\end{enumerate}
One evolution is in the money market (London), the other in the capital market; link the two and you have shadow banking, money market funding of capital market
lending, driven by market forces, while ``none of the regulations are right for this'' \ts{20}{1}{4:15}.
\end{casestudy}

\begin{coursenote}
The figure of 40\% is as said in class, quoting the FT. The ``30 years ago'' of the Eurodollar debate is approximate: the Eurodollar market grew in London from the
late 1950s, and the US debate on controlling it was intense in the late 1960s and 1970s \ts{20}{1}{1:46}.
\end{coursenote}

% ------------------------------------------------------------------
\section{Credit indices}\label{sec:cds-index}

\begin{finance}[The Markit CDX family]
Markit's CDX indices are, in its words, the standard tradable family of credit default swap indices for North America and emerging markets: they offer exposure to
a diversified portfolio of credits, a hedge of credit risk, and liquidity that flows into the single-name CDS market \ts{20}{2}{0:27}.
The CDX.NA.HY (North America high yield, i.e.\ junk) index: spreads of 450--580 basis points over a riskless rate between late September and the day of the lecture;
the index is set at 100 and moves inversely to the spread (about 530 bp at that point). Its constituents are 100 companies, each with 1\% weight, with ratings
such as BB, B, CCC \ts{20}{2}{1:29}.
\end{finance}

% ------------------------------------------------------------------
\section{Fischer Black's vision (1970)}\label{sec:cds-black}

In an unpublished 1970 memo, which Mehrling read while writing Black's biography, the young Fischer Black imagined a long-term corporate bond sold to three
separate persons: one supplying the money, one bearing the interest rate risk and one the default risk, the last two putting up no capital but perhaps posting
collateral. ``That's the world that came to be'' \ts{20}{3}{0:05}.

\begin{keyformulas}
\textbf{Black's decomposition} \ts{20}{3}{0:48}:
\[
  \text{risk-free security} = \text{risky (corporate) bond} + \text{interest rate hedge (IRS)} + \text{credit hedge (CDS)}.
\]
\end{keyformulas}

\begin{history}[Fischer Black]
Fischer Black (1938--1995), co-author of the Black--Scholes option pricing formula (1973), worked at Arthur D.~Little, the University of Chicago, MIT and Goldman
Sachs. Mehrling's biography draws on his unpublished papers, including the 1970 memo ``Fundamentals of Liquidity''.\par
\emph{Reference:} P.~Mehrling, \emph{Fischer Black and the Revolutionary Idea of Finance}, Wiley, 2005.
\end{history}

% ------------------------------------------------------------------
\section{What is a credit default swap}\label{sec:cds-what}

There is a buyer and a seller of ``insurance''. The buyer is \emph{long the swap} and \emph{short credit risk} \ts{20}{4}{0:04}. Balance sheet
construction \ts{20}{4}{0:56}:
\begin{enumerate}
  \item start from a risky corporate bond;
  \item add (in brackets: notional) a long Treasury bond and a short position in the same corporate bond, of the same maturity: the net exposure is now a
        Treasury bond. The pair is the CDS: it carves off the credit risk, the 530 bp;
  \item add a short Treasury bond and a long Treasury bill: the net exposure is a Treasury bill. This pair, long fixed and short floating, is an interest rate swap.
\end{enumerate}

\begin{nfigure}
\begin{taccts}
\tacct[2.6cm]{Buyer of protection}{Corporate bond\\{}[$+$Treasury bond, $-$corporate bond] $=$ CDS (long swap)\\{}[$+$T-bill, $-$Treasury bond] $=$ IRS}{}\qquad
\tacct[2.6cm]{Seller of protection}{}{CDS (short swap, long credit risk)\\IRS}
\end{taccts}
\caption{A CDS and an IRS as notional swaps of IOUs: the hedges are booked as contingent assets of the buyer and contingent liabilities of the seller
\ts{20}{4}{2:50}.}\label{fig:cds-bs}
\end{nfigure}

\begin{definition}[New concepts: credit default swap, protection buyer and seller]
A \emph{credit default swap} (CDS) is a contract in which the \emph{protection buyer} pays a periodic premium and, if the reference bond defaults, receives
from the \emph{protection seller} the difference between face value and the bond's post-default value \ts{20}{4}{0:04}\ts{20}{6}{4:28}.
(\emph{Default}: from Old French \emph{defaute}, a failing or lack.)
\end{definition}

The course books any hedge as a contingent asset, so the buyer lists both the long credit exposure (the bond) and the short one (the CDS) as assets; real-world
practice is not consistent, often off balance sheet \ts{20}{4}{3:12}. The seller is long credit risk and short the swap: the
lingo is consistent between CDS and IRS. ``Credit default swaps are not some weird exotic creature of Wall Street'' \ts{20}{4}{4:33}.

% ------------------------------------------------------------------
\section{Corporate bonds}\label{sec:cds-bonds}

For the first time in the course, a corporate bond: a fixed coupon every six months or year for, say, ten years, and the principal at the end; unlike a mortgage it
is not amortized \ts{20}{5}{0:04}.
\begin{keyformulas}
\textbf{Bond price} \ts{20}{5}{1:05}: with coupons $C_t$, face value $F$ at maturity $T$, and a risk-weighted discount rate
(the risk-free rate plus a credit spread $\delta$, e.g.\ 530 bp),
\[
  P_0 = \sum_{t=1}^{T}\frac{C_t}{(1+r+\delta)^t} + \frac{F}{(1+r+\delta)^T}.
\]
\end{keyformulas}
Bonds have ratings, and at any moment there are credit spreads between rating classes, which change over time. So a bond's price changes when the spread for its
rating changes and when its rating changes, a ``double whammy'' if a downgrade comes with wider spreads \ts{20}{5}{2:12}.

% ------------------------------------------------------------------
\section{CDS pricing}\label{sec:cds-pricing}

\paragraph{Step 1: strip the interest rate risk.} Take a ten-year bond with constant coupon and swap fixed for LIBOR: what remains is a ten-year floating-rate bond paying
LIBOR $+\,s$, with the credit spread $s$ locked in \ts{20}{6}{0:02}.

\paragraph{Step 2: the CDS as a swap.} Each year the bond survives, one side pays LIBOR times face value and the other LIBOR $+\,u$: the protection buyer pays the net
premium $u$ times face value \ts{20}{6}{1:51}. If the bond defaults, say in period five, the buyer
swaps it for a Treasury bond: it receives face value minus the bond's price at that time \ts{20}{6}{4:05}. The contract follows the
two parts of a bond: the annual (coupon) business and the terminal payment if it defaults \ts{20}{6}{5:35}.

\begin{nfigure}
\begin{tikzpicture}[font=\scriptsize,x=1.3cm,y=1cm]
  \draw[-{Stealth[length=4pt]}] (0,0) -- (6.4,0) node[right]{period};
  \foreach \t in {1,...,4} {\draw[thick,-{Stealth[length=3pt]}] (\t,0) -- (\t,-0.7); \node[below] at (\t,-0.7) {$-u F$};}
  \draw[very thick,-{Stealth[length=4pt]}] (5,0) -- (5,1.8);
  \node[above,align=center] at (5,1.8) {default: $+\,(F - P_5)$};
  \foreach \t in {1,...,5} \node[above] at (\t,0.05) {\t};
\end{tikzpicture}
\caption{Cash flows to the protection buyer: the premium $u$ times face value $F$ while the bond survives, face value minus the bond's price $P_5$ at default
\ts{20}{6}{3:45}.}\label{fig:cds-flows}
\end{nfigure}

\begin{keyformulas}
\textbf{CDS pricing by arbitrage} \ts{20}{6}{6:36}: the CDS premium $u$ should equal the bond's credit spread $s$;
otherwise go long the bond and insure it (if $u<s$), or short the bond and sell insurance (if $u>s$). This is how CDS are priced, ``assuming that there's no
liquidity premium anywhere in this picture''. The difference $u-s$ is the \emph{CDS--bond basis}.
\end{keyformulas}

The notes call the protection seller's IOU a \emph{mirror bond}: it promises LIBOR $+U$ on face value as long as the issuer pays, and the liquidation value on default; the buyer's IOU promises
LIBOR, and face value on default. On default the buyer recovers $F$ in full: $P_5$ from the bond, $F-P_5$ from the swap. $U$ is the premium that makes the two IOUs equal in value at inception; but
since any change in the spread $S$ moves its value, ``CDS is not so much insurance against eventual default as it is insurance against change in the credit spread'' \mn{20}{3}\mn{20}{4}. Pooling
CDS reduces idiosyncratic risk (the less correlated, the better), letting sellers charge less: a flourishing swap market should lower the price of credit risk overall, ``and that is exactly what
happened''; for subprime the index used for hedging was ABX, the index of top-tranche mortgage CDOs \mn{20}{1}\mn{20}{4}\mn{20}{5}.

\paragraph{Not just insurance.} A CDS is a promise to pay a fixed premium for its life, but its value fluctuates, inversely to the bond, whenever the warranted
discount rate changes; so CDS are used not only as insurance against default but as a way to get exposure to credit risk, traded and reversed like any liquid
derivative. Unlike fire insurance, you are not just waiting for the house to burn \ts{20}{6}{8:00}.

\begin{intuition}[Is a CDS insurance?]
``Insurance'' is used to remember which side of the balance sheet to book; but many people dislike the word, on both sides of the political spectrum. If it were
insurance it would be regulated as insurance, with reserves and suitability standards; and insurance works by pooling independent risks (fires), while credit risk
is systemic: corporate bonds and their spreads move together. The mathematicians and theorists say: it is not insurance, just a swap
\ts{20}{6}{9:43}.
\end{intuition}

% ------------------------------------------------------------------
\section{Market making in CDS}\label{sec:cds-dealers}

Opposite the protection buyer is typically a dealer writing CDS on bonds $i$, $j$, $k$\ldots: a portfolio of CDS is like a diversified portfolio of corporate bonds
\ts{20}{7}{0:03}. The dealer cannot find someone wanting exactly the opposite single-name positions, so it hedges with a CDS on the
CDX index, which is correlated with its portfolio; on the other side may be an asset manager seeking exposure to the bond market as a whole
\ts{20}{7}{1:08}.

\begin{nfigure}
\begin{taccts}
\tacct[1.9cm]{Hedger (single name)}{CDS on bond $i$ (long)}{}\quad
\tacct[1.9cm]{CDS dealer (investment bank)}{CDS on CDX index}{CDS on bonds $i, j, k, \ldots$}\quad
\tacct[1.9cm]{Asset manager}{}{CDS on CDX index (long credit risk)}
\end{taccts}
\caption{A CDS dealer hedging a portfolio of single names with the index \ts{20}{7}{1:31}.}\label{fig:cds-dealer}
\end{nfigure}

With a dealer in the middle, order flow pushes CDS prices around without necessarily moving bond prices; someone must arbitrage to bring them in line. So
arbitrage opportunities appear whenever CDS prices are out of line with bond prices, ``and sometimes it is'' one \ts{20}{7}{2:55}.

% ------------------------------------------------------------------
\section{UBS: the negative basis trade and liquidity risk}\label{sec:cds-ubs}

UBS lost a lot of money on what it thought was a riskless trade, and its regulator made it write a public report explaining how; Mehrling wishes US regulators had
required the same of Citibank and the TARP banks \ts{20}{8}{0:03}. UBS noticed that the spread $s$ on a security exceeded the
cost $u$ of insuring it against default, at least from AIG \ts{20}{8}{0:45}:
\begin{enumerate}
  \item it bought AAA tranches of CDOs on mortgage-backed securities, some on subprime mortgages, with the credit risk supposedly carved off into lower tranches
        \ts{20}{8}{1:26};
  \item its European regulators required insurance to use them as bank collateral, so it bought CDS, from AIG among others: ``double insurance''
        \ts{20}{8}{2:09};
  \item it funded them in the money market, asset-backed commercial paper or repo, with the tranche as collateral: money market funding of capital market lending,
        shadow banking, but right on UBS's balance sheet, unlike Citibank's structured investment vehicles off balance sheet \ts{20}{8}{3:17}.
\end{enumerate}

\begin{nfigure}
\begin{taccts}
\tacct[3.0cm]{UBS (negative basis trade)}{AAA CDO tranche (yield LIBOR $+\,s$)\\CDS bought from AIG (premium $u<s$)}{Money market funding (ABCP, repo), collateral: the tranche}
\end{taccts}
\caption{The negative basis trade \ts{20}{8}{1:48}.}\label{fig:cds-ubs}
\end{nfigure}

\paragraph{The numbers.} The notes give them from UBS's report: UBS paid AIG 11 basis points for 100\% protection on a super-senior CDO tranche, financed the tranche in the wholesale money
market, and netted an apparently riskless 20 basis points; ``because it was apparently riskfree, they did massive amounts of it'' \mn{20}{5}. In the notes' notation, if the CDS premium $U$ is below
the spread $S$, the buyer of protection swaps out the credit risk and keeps $S-U$ over LIBOR.

\paragraph{Booking 20 years of profit today.} The \emph{negative basis trade} left money over each period for the life of the tranche, 10 or 20 years; UBS booked the
whole stream as profit today and paid bonuses on it \ts{20}{8}{4:42}.

\paragraph{The liquidity spiral.} Then small doubts about the AAA tranche: its value fluctuated a little; the CDS moved the other way, fine; but the tranche, not the
CDS, was the collateral, so the funding could not be rolled, the tranche had to be sold, pushing its price down further, and others doing the same trade made a
downward liquidity spiral \ts{20}{8}{5:47}. Moreover, AIG's insurance was expensive, so late in the game UBS insured
only against a 2\% fall (from 100 to 98) and kept the tail risk itself \ts{20}{8}{6:48}. They felt they took no risk at all and lost
billions, because they assumed they could always roll their funding \ts{20}{8}{7:50}.

\begin{coursenote}
A student asks why the lenders did not accept the CDS as collateral together with the tranche. Perhaps by analogy with Bagehot's bills, where the acceptance was
attached to the bill; here they are separate securities, and money market funds are restricted in the collateral they may accept. Mehrling recalls that when AIG
was failing he proposed that the government write CDS itself, since a CDS is like an acceptance, which is what a liquidity problem needs
\ts{20}{8}{8:11}.
\end{coursenote}

\begin{history}[The UBS report]
At the request of the Swiss Federal Banking Commission, UBS published in April 2008 a ``Shareholder Report on UBS's Write-Downs'', describing among other things its
negative basis trades and the ``Amplified Mortgage Portfolio'' (AMPS), in which only a small first slice of losses on super-senior positions was hedged.\par
\emph{Reference:} UBS AG, \emph{Shareholder Report on UBS's Write-Downs}, 18 April 2008.
\end{history}

% ------------------------------------------------------------------
\section{AIG and Goldman Sachs: a private backstop}\label{sec:cds-aig}

Goldman Sachs, a dealer in CDS selling protection to its clients (not completely matched book: short credit risk when it mattered), kept buying CDS from AIG, which held
CDS liabilities and capital \ts{20}{9}{0:00}. When the CDS came into the money, AIG's contingent liability, zero when
written, became, say, \$30 billion owed to Goldman \ts{20}{9}{1:36}. Goldman had written the contracts so that if AIG's rating fell below AAA it
would post margin; Moody's downgraded AIG, and it had to \emph{pay} the \$30 billion, not just owe it \ts{20}{9}{2:19}.

\begin{nfigure}
\begin{taccts}
\tacct[2.4cm]{Goldman Sachs (dealer)}{CDS bought from AIG\\($+$\$30 bn when in the money: margin received)}{CDS sold to clients\\($+$\$30 bn owed to clients)}\qquad
\tacct[2.4cm]{AIG}{Capital (bonds)\\($-$cash paid as margin)}{CDS on AAA securities\\($+$\$30 bn contingent liability)}
\end{taccts}
\caption{AIG as private backstop of CDS market making; ``say \$30 billion'' in class; the notes give about \$30 billion of collateral posted to Goldman's segregated account by the
time AIG failed \ts{20}{9}{1:16}\mn{20}{6}.}\label{fig:cds-aig}
\end{nfigure}

\begin{intuition}[Liquidity kills you quick]
AIG thought it was insuring against defaults that would never happen; but you can go bankrupt without any default if the CDS fluctuate in value and you must pay
margin. The contractual promise to post collateral killed AIG; monoline insurers that wrote CDS without margin saw the marks go against them but paid nothing.
``Insolvency you can survive for a long time\ldots{} but liquidity kills you quick'': AIG, the anchor of the system, fell two or three days after Lehman in September
2008, and the system went into free fall \ts{20}{9}{2:40}.
\end{intuition}

\begin{coursenote}
The AIG executive is not named in class. Joseph Cassano, head of AIG Financial Products, told investors in August 2007 that it was hard to see a scenario in which
the unit would lose ``one dollar'' on these transactions (``a nickel'' in class). The US government took a 79.9\% equity stake in AIG in September 2008 (``80'' in
class).\par
\emph{References:} Financial Crisis Inquiry Commission, \emph{Final Report}, 2011, ch.~13; SIGTARP, \emph{Factors Affecting Efforts to Limit Payments to AIG Counterparties},
November 2009.
\end{coursenote}

\begin{intuition}[Was Goldman paid at par by the government?]
The notes correct ``a lot of loose talk'' \mn{20}{6}: since the CDS were marked to market, Goldman already held the collateral, it had already been paid. When AIG could not meet further collateral
calls the government took over, lending \$85 billion; that money was used to acquire the referenced securities at liquidation value to end the swaps. In the algebra of the pricing example, AIG had
already paid $F-P_5$; the government lent AIG $P_5$ to buy the bonds, which went onto the Fed's balance sheet as Maiden Lane II and III.
\end{intuition}

\paragraph{Should the government have clawed it back?} Neil Barofsky (SIGTARP), whose book is on the final-exam reading, wrote a report on AIG: Goldman and, even bigger,
Soci\'et\'e G\'en\'erale sucked the capital out of AIG, the government took over AIG, and the question was what to do about the money already paid to Goldman
\ts{20}{9}{4:03}. But to the extent Goldman was a market maker with offsetting positions, taking back the \$30 billion would have left it
unable to pay its clients and the whole thing unravels; ``I don't say that Geithner necessarily did the right thing'', but it was not simply unearned profit: the
point of a CDS is that when the corporate bond goes south you replace it with a Treasury \ts{20}{9}{5:05}.

\paragraph{No fraud, a misunderstanding of risk.} Mehrling is not sure there was fraud, at AIG or at UBS: a systematic misunderstanding of risk, a supposedly risk-free
arbitrage on which you bet the bank, and do more of it. ``If it looks like free money, look again'', with the money view: is the arbitrage just a liquidity premium
you think should be zero? Then you are becoming the lender of last resort to the whole system without being a central bank that can print money
\ts{20}{9}{7:28}.

\paragraph{Moral hazard?} A student suggests they took the risk expecting a bailout. There is always some such calculation, but Mehrling thinks they were surprised.
The one really relying on the government was Goldman: it chose AIG though it was more expensive (UBS found it too expensive), because AIG promised margin and had a
huge capital base; and AIG, reserving nothing and paying bonuses on each sale, could not be allowed to fail, with fire insurance policies worldwide. But a
counterparty that dies is not the right counterparty: Goldman surely did not expect to destroy AIG \ts{20}{9}{9:18}.

\paragraph{What triggered payment?} Without margin, the premium flows and a large payment comes only at actual default; for AAA tranches, with many tranches beneath,
defaults rarely reached them. It was not default but the fluctuation in CDS values that caused the problem. That is why many underestimated the crisis: ``there
just aren't a lot of subprime mortgages'', ignoring leverage, liquidity and the embedding in world money and capital markets. Those who took the low tranches and
knew they bore risk mostly survived; those who thought they took no risk failed \ts{20}{9}{12:24}.

% ------------------------------------------------------------------
\section{Abacus: the synthetic CDO as collateral prepayment}\label{sec:cds-abacus}

John Paulson wanted to short subprime, expecting the housing bubble to end. He bought CDS on CDO tranches of mortgage-backed securities, insurance without owning the
securities \ts{20}{10}{0:02}. Goldman Sachs set up an entity, Abacus, which took the other side of the CDS, credit exposure as if owning
RMBS, ``just a side bet'' \ts{20}{10}{1:04}. Abacus sold the German bank IKB a CDO tranche paying a spread over the risk-free rate; IKB
borrowed to buy it, and Abacus invested the proceeds in Treasury bills \ts{20}{10}{2:08}.

\begin{nfigure}
\begin{taccts}
\tacct[1.9cm]{Paulson}{CDS on CDO tranches (long protection)}{}\quad
\tacct[1.9cm]{Abacus (synthetic CDO)}{Treasury bills}{CDS sold to Paulson\\CDO tranche sold to IKB}\quad
\tacct[1.9cm]{IKB}{CDO tranche}{Borrowing}
\end{taccts}
\caption{The synthetic CDO: T-bills plus a short CDS replicate a long position in mortgage-backed securities \ts{20}{10}{2:28}.}\label{fig:cds-abacus}
\end{nfigure}

\begin{definition}[New concepts: synthetic CDO]
A \emph{collateralized debt obligation} (CDO) issues tranches backed by a pool of debt; a \emph{synthetic} CDO owns no such debt but replicates the exposure with
risk-free securities plus CDS sold on the reference assets \ts{20}{10}{2:49}.
\end{definition}

When the tranches fell, the CDS rose and the T-bills went to Paulson as margin. From a liquidity perspective that is the genius of it: IKB had paid 100\% upfront,
so there was no problem getting paid, no bill to send to IKB that it could dispute, no need to liquidate RMBS exposure
\ts{20}{10}{3:30}. There were issues about Paulson's role in choosing the tranches most
likely to default and about Goldman's representations, and a fine was paid; but no issue of payment. IKB was bailed out by the German government; AIG's counterparties
included Soci\'et\'e G\'en\'erale: the international aspects again \ts{20}{10}{4:55}.

\begin{history}[Abacus 2007-AC1]
In April 2010 the SEC charged Goldman Sachs with misleading investors in the synthetic CDO Abacus 2007-AC1, whose reference portfolio Paulson \& Co.\ had helped select
while betting against it; Goldman settled in July 2010 for \$550 million. IKB, an early casualty of the crisis in 2007, was rescued with the support of the state-owned
KfW.\par
\emph{References:} SEC press releases 2010-59 (16 April 2010) and 2010-123 (15 July 2010); G.~Zuckerman, \emph{The Greatest Trade Ever}, 2009 (cited in Mehrling's notes, which also
recall, from G.~Tett, \emph{Fool's Gold}, 2009, that the first CDS were built in this synthetic way for J.~P.~Morgan to hedge tail risk on corporate loans) \mn{20}{6}.
\end{history}

% ------------------------------------------------------------------
\flushside
\section*{Key formulas and ideas}
\begin{keyformulas}
\begin{tabularx}{\linewidth}{@{}l L@{}}
Black (1970) & risky bond $+$ IRS $+$ CDS $=$ risk-free bond\\
CDS & notional long Treasury $+$ short corporate bond of the same maturity; buyer long the swap, short credit risk\\
Bond price & $P_0=\sum C_t/(1+r+\delta)^t+F/(1+r+\delta)^T$; moves with spreads and ratings\\
Pricing & bond $+$ IRS $=$ LIBOR $+s$; CDS premium $u$; no-arbitrage $u=s$ if no liquidity premium\\
Negative basis & $u<s$: buy bond, buy protection, fund in money market; liquidity risk in the funding\\
AIG & margin calls on in-the-money CDS: ``liquidity kills you quick''\\
Abacus & synthetic CDO $=$ T-bills $+$ short CDS: collateral prepaid\\
\end{tabularx}
\end{keyformulas}

\section*{Exercises}

\begin{exercise}[Black's decomposition]
Write the balance sheet of an investor who holds a ten-year BBB bond, an interest rate swap and a CDS on it, and show that its net exposure is a Treasury bill. Who
bears the interest rate risk and who the default risk?
\end{exercise}

\begin{exercise}[CDS cash flows]
A protection buyer on \$10 million face pays $u=400$ bp a year. The bond defaults at the end of year 3 and then trades at 35. Write the cash flows each year. What is
the buyer's position worth just before default if spreads have widened to 900 bp?
\end{exercise}

\begin{exercise}[Negative basis]
A AAA tranche yields LIBOR $+$ 60 bp, protection costs 20 bp, and repo funding costs LIBOR $+$ 10 bp. Compute the annual carry on \$1 billion. The tranche's price falls
from 100 to 96 and the repo haircut rises from 2\% to 10\%: how much cash must be found, and from where?
\end{exercise}

\begin{exercise}[Margin or no margin]
Compare an AIG-style CDS seller that posts margin when downgraded with a monoline that posts none, if the marked-to-market value of the CDS rises by \$5 billion but no
default occurs. Which one fails, and why?
\end{exercise}

% !TEX root = ../main.tex
\chapter{Central Banking for Shadow Banking}\label{ch:cbshadow}
\begin{paracol}{2}

\begin{sources}
\begin{tabularx}{\linewidth}{@{}l l L@{}}
\toprule
\textbf{Segment} & \textbf{Length} & \textbf{Content}\\
\midrule
\seg{21}{1}{L21.1 FT: Regulation of shadow banking} & 3:34 & Bail-in debt for banks; central counterparties for swaps.\\
\seg{21}{2}{L21.2 Shadow banking vs traditional banking} & 6:55 & Pozsar's diagram; a stylized shadow banking system without government backstops.\\
\seg{21}{3}{L21.3 Liquidity and solvency backstops} & 7:18 & Immature private backstops: liquidity puts and CDS; collapse onto the Fed.\\
\seg{21}{4}{L21.4 Global dimension} & 5:16 & Dollar funding of Korean banks and of shadow banks; a mature funding system, an immature risk transfer system.\\
\seg{21}{5}{L21.5 Evolution of modern finance} & 3:19 & Securitization, globalization, derivatives: risk transfer.\\
\seg{21}{6}{L21.6 What is shadow banking?} & 12:22 & Money market funding of capital market lending; a world of only shadow banking.\\
\seg{21}{7}{L21.7 Backstopping the market makers} & 7:46 & The Fed's balance sheet in December 2011 as a dealer in swaps.\\
\seg{21}{8}{L21.8 Regulation of systemic risk} & 6:27 & Inherent instability of credit: boom and bust through dealers' balance sheets.\\
\seg{21}{9}{L21.9 Regulation of Collateral and Payment} & 10:29 & Normal fluctuations: collateral flows and payment flows.\\
\seg{21}{10}{L21.10 Private backstop, and public} & 8:20 & Four principles; the central bank as outside spread.\\
\bottomrule
\end{tabularx}

\smallskip
\textbf{Files.} Transcripts \file{materials/transcripts/L21-P*.txt}; L21.4 has no YouTube captions and was transcribed locally from the audio (faster-whisper). Mehrling's slides \mn{21}{1}.
\end{sources}

\section*{Motivation}
This is the last lecture about banking: it draws together forwards and futures, FX swaps, interest rate swaps and credit default swaps, and asks what they mean for
banking and central banking; the next lecture connects the course to economics \ts{21}{2}{0:03}. The question: why is shadow banking
in our future, ``given how desperately it has failed us'', and how should it be managed? \ts{21}{2}{2:27}

% ------------------------------------------------------------------
\section{FT: old banking and new banking}\label{sec:cbs-ft}

\begin{casestudy}[FT, December 2012: two articles on reform]
``Regulators in long-term debt plan for big banks'': more long-term debt, subordinated and ``bail-inable'' bonds that convert to equity when a bank runs into
trouble, could lessen reliance on cheaper short-term funding. This is the traditional Basel approach (capital adequacy), now extended from the equity tranche to
subordinated debt, for traditional banks \ts{21}{1}{0:07}. Three pages later, ``Regulators wrestle
with hidden risks in swaps shake-up'': standardized interest rate swaps and credit derivatives to be traded on exchanges and cleared through central
counterparties, to increase price transparency, ensure adequate collateral and keep taxpayers from paying next time \ts{21}{1}{1:53}.
Old banking and new banking, a few pages apart \ts{21}{1}{2:55}.
\end{casestudy}

\begin{definition}[New concepts: central counterparty, bail-in]
A \emph{central counterparty} (CCP) interposes itself between the two sides of every trade, becoming buyer to every seller and seller to every buyer, and manages
margin and netting. \emph{Bail-in}: imposing losses on a bank's creditors (by conversion to equity or write-down) instead of a taxpayer \emph{bail-out}
\ts{21}{1}{1:32}.
\end{definition}

% ------------------------------------------------------------------
\section{Shadow banking versus traditional banking}\label{sec:cbs-vs}

Zoltan Pozsar's famous diagram, made at the New York Fed, shows loans funded by deposits through seven stages of securitization, slicing and dicing, wrapped by the
Fed's U-shaped embrace of liquidity facilities during the crisis: a Rube Goldberg machine, built piece by piece; much of it fell apart, but the next stage can be
simpler \ts{21}{2}{1:06}.

\begin{history}[The shadow banking map]
Z.~Pozsar, T.~Adrian, A.~Ashcraft and H.~Boesky, ``Shadow Banking'', Federal Reserve Bank of New York Staff Report 458, July 2010, includes a large fold-out map of
the shadow banking system and of the Fed's crisis facilities.\par
\emph{Reference:} FRBNY Staff Report 458, 2010 (revised 2012).
\end{history}

\paragraph{Stylized comparison.} The traditional (Jimmy Stewart) bank takes retail deposits, makes retail loans, holds cash reserves and a capital buffer; beyond
them the Fed is the liquidity backstop and the FDIC the capital backstop \ts{21}{2}{2:48}. The stylized shadow banking system, three
entities instead of seven: RMBS tranched into high, mid and low tranches; the high tranche held by a shadow bank as collateral for repo funding; the repo bought by
a money market mutual fund issuing ``deposits'' \ts{21}{2}{3:49}. Differences: it is wholesale (pension funds, foreign central banks,
corporations; a ``small'' deposit is \$10 million), its assets are securities in the capital market, and it may be offshore, on a computer in the Cayman Islands
\ts{21}{2}{4:53}. Seen through traditional banking it looks similar, deposits funding loans, but no government backstop
is anywhere to be seen; the Fed and the FDIC said ``good riddance'' to offshore risk \ts{21}{2}{6:14}.

\begin{nfigure}
\begin{taccts}
\tacct[1.8cm]{Traditional bank}{Loans\\Reserves}{Deposits\\Capital}\quad
\tacct[1.8cm]{Shadow bank}{High tranche RMBS}{Repo (RP)}\quad
\tacct[1.8cm]{MMMF}{Repo (RP)}{``Deposits'' (shares)}
\end{taccts}
\caption{Traditional bank, backstopped by Fed and FDIC, versus a stylized shadow banking chain with no government backstop \ts{21}{2}{3:08}.}\label{fig:cbs-chain}
\end{nfigure}

% ------------------------------------------------------------------
\section{Immature backstops}\label{sec:cbs-backstops}

\paragraph{A private liquidity backstop.} Invented over the last 30 years or so, the system had an immature, private liquidity backstop: liquidity puts held by the
shadow bank and by the money fund, liabilities of traditional banks. Citibank wrote in contracts that if its offshore entities could not roll their repo, for any
reason, it would lend them what they needed; other backstops were only reputational \ts{21}{3}{0:02}.
All money funds were run, but only the Reserve Primary Fund failed, because it had no larger parent to back it \ts{21}{3}{1:43}. A private backstop works
only if the backstopper is healthy; in a large crisis it becomes the traditional bank's crisis, and then the government backstops kick in: the government was
indirectly on the hook, perhaps without realizing it \ts{21}{3}{2:25}.

\paragraph{A private solvency backstop.} Even more immature: the shadow bank bought CDS on its high tranche from AIG; mid tranches went to pension funds and insurers,
low tranches to credit hedge funds, which might buy CDS from an investment bank (think Lehman or Bear Stearns) making markets in CDS on all tranches; AIG, writing
naked CDS on the highest tranche, ``picking up nickels'', was the ultimate support \ts{21}{3}{3:06}.
Lehman went, AIG went, prices went into free fall; through the liquidity backstop everything dropped onto the Fed's balance sheet, which tripled more or less
overnight \ts{21}{3}{6:10}.

\begin{nfigure}
\begin{taccts}
\tacct[1.8cm]{Shadow bank}{High tranche\\CDS on high tranche\\Liquidity put}{Repo}\quad
\tacct[1.8cm]{AIG}{Capital}{CDS on high tranche}\quad
\tacct[1.8cm]{Traditional bank}{}{Liquidity put}
\end{taccts}
\begin{taccts}
\tacct[1.8cm]{Pension fund / insurer}{Mid tranche\\CDS on mid}{DB pensions, annuities}\quad
\tacct[1.8cm]{Credit hedge fund}{Low tranche\\CDS on low}{Capital}\quad
\tacct[1.8cm]{Investment bank (CDS dealer)}{CDS on high}{CDS on mid, low}
\end{taccts}
\caption{Immature private liquidity and solvency backstops of the shadow banking system \ts{21}{3}{0:22}.}\label{fig:cbs-immature}
\end{nfigure}

\begin{history}[The Reserve Primary Fund]
On 16 September 2008, the day after Lehman's bankruptcy, the Reserve Primary Fund, which held \$785 million of Lehman commercial paper, ``broke the buck'' (net asset
value below \$1 per share, \$0.97); on 19 September the US Treasury announced a temporary guarantee of money market funds.\par
\emph{References:} Financial Crisis Inquiry Commission, \emph{Final Report}, 2011, ch.~18; US Treasury press release HP-1147, 19 September 2008.
\end{history}

% ------------------------------------------------------------------
\section{The global dimension}\label{sec:cbs-global}

What Mehrling had not adequately appreciated: shadow banking tapped a \emph{global} funding market, the dollar funding markets, which are mature and have been
tested many times \ts{21}{4}{0:00}. His slides call this ``shadow banking as global banking'' \mn{21}{7}.

\paragraph{Development lending in dollars.} A Korean bank borrows short term in dollars from a global bank (say Deutsche Bank), which issues short-term money market
liabilities held by a ``dollar bank'', perhaps a money market mutual fund: since the dollar is the world reserve currency, a country that wants to invest more than its
domestic savings can tap these markets \ts{21}{4}{0:27}. Lending in won and borrowing in dollars means
FX risk, hedged with FX swaps or dollar reserves, perhaps by the central bank as counterparty; someone must bear it \ts{21}{4}{1:18}. When
the money stops there is a banking crisis and probably a currency crisis: the Asian financial crisis, from which ``we learned something'', though it did not turn out
so well for the Tigers \ts{21}{4}{1:50}.

\paragraph{Shadow banking used the same system for dollar loans.} RMBS tranched, the high tranche held by a shadow bank funded by ABCP, the ABCP held by a money fund:
borrowing in dollars to lend in dollars, so no FX risk, the risk behind the Asian crisis; credit and duration risk remained, but those existed in Korea too
\ts{21}{4}{2:20}.

\begin{nfigure}
\begin{taccts}
\tacct[1.9cm]{Korean bank}{Domestic currency loans\\FX swaps\\\$ reserves}{Dollar funding}\quad
\tacct[1.9cm]{Global bank}{Dollar lending}{Wholesale money market}\quad
\tacct[1.9cm]{Dollar bank}{Wholesale money market}{``Deposits''}
\end{taccts}
\begin{taccts}
\tacct[1.9cm]{RMBS}{Mortgage loans}{Hi tranche\\Mid tranche\\Lo tranche}\quad
\tacct[1.9cm]{Shadow bank}{Hi tranche}{ABCP}\quad
\tacct[1.9cm]{MMMF}{ABCP}{``Deposits''}
\end{taccts}
\caption{Shadow banking as global banking: the same dollar funding chain for Korean development lending and for US mortgages \ts{21}{4}{0:35}\mn{21}{7}.}\label{fig:cbs-global}
\end{nfigure}

\begin{keyformulas}
\textbf{A mature funding system, an immature risk transfer system} \ts{21}{4}{3:08}: shadow banking tapped a mature global
funding system for a new purpose, while the risk transfer piece (the CDS) was new, ``invented on the fly''. It was easy to borrow, hard to transfer risk.
\end{keyformulas}

This is why the crisis was global: it hit the global money market at once. The spread between LIBOR (Eurodollar rates) and term Fed funds, which should be about equal
for the same term, blew out to about 100 basis points before Bear Stearns and 300 or more after Lehman and AIG, with the Fed apparently stabilizing it in between but
nowhere near zero \ts{21}{4}{3:44}. The shadow banks were in Europe, funded through the Eurodollar market,
and the crisis touched every major money market, Japan's too \ts{21}{4}{4:36}. The next slide is titled ``From domestic to international
LOLR'' \mn{21}{8}: the central bank swap lines of lecture~\ref{ch:intl}. For financial globalization the slides cite H.~S.~Shin, ``Global Banking Glut and Loan Risk
Premium'' \mn{21}{11}.

% ------------------------------------------------------------------
\section{The evolution of modern finance: risk transfer}\label{sec:cbs-risk}

Shadow banking used recent developments of modern finance: securitization, which came with globalization (individual small loans are not transferred across borders,
but packaged bonds with a market price are bought by foreign asset managers); global capital markets; derivatives markets for moving risk, of which tranching was an
awkward approximation \ts{21}{5}{0:04}.

\paragraph{An idealized shadow banking system.} A shadow bank holds RMBS funded in the money market and strips out all the risk with derivatives (interest rate swaps, FX
swaps, CDS): what is left is a quasi-Treasury bill, which can be financed in the money market without price risk. On the other side, an asset manager gets its
customers' risk exposure entirely through derivatives, without holding risky assets, equivalent to holding RMBS outright: Black's 1970 vision. That is what shadow
banking was trying to do, and probably what its future looks like \ts{21}{5}{1:06}.

% ------------------------------------------------------------------
\section{What is shadow banking?}\label{sec:cbs-what}

\begin{definition}[New concepts: shadow banking as market-based credit]
Shadow banking is \emph{money market funding of capital market lending}, on whatever balance sheet it sits (a bank's, as at UBS, a non-bank's, or a special
vehicle's), with three features \ts{21}{6}{0:03}:
\begin{enumerate}
  \item global (dollar) wholesale funding of local lending;
  \item market prices on both sides, for money and for capital, visible and changing over time;
  \item market makers, dealers quoting bid--ask spreads according to inventories and risk tolerance: that is where the prices come from.
\end{enumerate}
\end{definition}

The existing literature starts from traditional banking and stresses the ``shadow'': banking by non-banks, evading regulation, without a direct government backstop,
something a little fishy; so it would move the same regulations and backstops over, make it a bank \ts{21}{6}{2:06}.
Mehrling disagrees: there was plenty of regulatory arbitrage and fraud, but the fundamental forces are financial globalization and the revolution in modern finance;
removing the fraud will not remove shadow banking \ts{21}{6}{3:28}.

\paragraph{A world of only shadow banking.} His project since \emph{The New Lombard Street}: imagine a world with only shadow banking and ask how to regulate it. It
will never be literally true, but it frees us from Jimmy Stewart habits of thought; extending old regulations creates new obstacles that modern finance will get
around, ``a disastrous way of thinking about the problem'' \ts{21}{6}{4:10}.

\begin{nfigure}
\begin{taccts}
\tacct[1.8cm]{Capital funding bank}{RMBS\\CDS, IRS, FXS}{Money market funding}\quad
\tacct[1.8cm]{Global money dealer}{Money market loans}{Deposits}\quad
\tacct[1.8cm]{Asset manager}{Deposits}{Capital\\CDS, IRS, FXS}
\end{taccts}
\begin{taccts}
\tacct[2.6cm]{Derivative dealer}{CDS, IRS, FXS (with asset manager)}{CDS, IRS, FXS (with capital funding bank)}
\end{taccts}
\caption{Mehrling's idealized market-based credit system: the dealers stand between the asset manager and the capital funding bank in the funding market and in the
risk market \ts{21}{6}{5:37}.}\label{fig:cbs-ideal}
\end{nfigure}

\paragraph{The dealers are the key.} The global money dealer, borrowing and lending three-month deposits, is where the price of money is set; the derivative dealer, on the
other side of all the derivatives, is the risk market \ts{21}{6}{6:18}. Today these balance sheets are the big investment banks' swap
books; central counterparties would move those positions (Lehman's, Bear Stearns's) into one place where they can be seen and netted. The money funding system is
relatively mature; the risk transfer system is very immature, and this is the part to clean up \ts{21}{6}{7:20}. If the
dealers stop making markets the whole thing collapses; so regulation focused on the capital funding bank (``let's have long-term funding'') looks in the wrong place
\ts{21}{6}{8:44}.

\paragraph{Where is the capital?} In the idealization the only capital is the asset manager's: the dealers are perfectly matched, the hedges perfect, no liquidity
reserves anywhere. If risk is properly transferred, the capital backing the capital funding bank sits on the asset manager's balance sheet, not the bank's
\ts{21}{6}{9:45}. The model abstracts from retail depositors, traditional banks, bad underwriting, reserves and
speculative dealers, to focus on system interlinkages, normal liquidity risk and tail backstops \ts{21}{6}{11:07}.

% ------------------------------------------------------------------
\section{Backstopping the market makers}\label{sec:cbs-backstop}

In a crisis, the backstop: a liquidity put to the global money dealer (lender of last resort, familiar from Bagehot, 1873) and one to the derivative dealer, a
backstop for capital market prices (dealer of last resort, the new bit), both contingent liabilities of a central bank or a consortium of central banks
\ts{21}{7}{0:03}.

\paragraph{The Fed as a dealer in swaps.} The Fed's H.4.1 for 15 December 2011 \mn{21}{16}: about \$1.7 trillion of long-term Treasuries and almost \$1 trillion of MBS and GSE
(Fannie and Freddie) securities, funded by about \$1 trillion of currency and \$1.6 trillion of excess reserves (against \$50 billion of reserves in normal times)
\ts{21}{7}{1:47}. Now add the same items to both sides \ts{21}{7}{2:50}:

\begin{nfigure}
\begin{taccts}
\tacct[3.0cm]{Federal Reserve (restated, \$ trillion)}{MBS, GSE \amt{0.9}\\T-bonds ($1.7+0.9$) \amt{2.6}\\{}[T-bills] \amt{2.6}}{{}[T-bonds] \amt{0.9}\\{}[T-bills] \amt{2.6}\\Currency and reserves \amt{2.6}}
\end{taccts}
\caption{The Fed's balance sheet, 15 December 2011, restated as exposures: risky securities against T-bonds (a CDS), T-bonds against T-bills (an IRS), T-bills
against overnight money (a funding swap) \ts{21}{7}{3:10}.}\label{fig:cbs-fed}
\end{nfigure}

The bracketed pairs are swaps, parallel loans: risky securities against T-bonds is a CDS; long against short bonds an interest rate swap; T-bills against currency and
reserves, three-month against overnight money, like the overnight index swap that blew out in the crisis \ts{21}{7}{3:53}. The
out-of-the-money liquidity puts came into the money and onto the Fed's balance sheet: in parallel-loan terms about \$6.4 trillion of notional swap exposure, the Fed
acting both as global money dealer and as global derivative dealer \ts{21}{7}{4:36}.

\begin{intuition}[A Bagehot moment]
What a central bank would have to do to backstop a collapsing shadow banking system is exactly what the Fed did: evidence that shadow banking is the centre of the
system and traditional banking the edges. ``We're living in a Bagehot moment'': the role of the central bank is shifting, and monetary policy, bank regulation and
inflation targeting, designed for a previous banking system, must be rethought from scratch. As Bagehot said, money does not manage itself
\ts{21}{7}{5:37}.
\end{intuition}

% ------------------------------------------------------------------
\section{Regulating systemic risk}\label{sec:cbs-systemic}

Start from the inherent instability of credit, not from a self-equilibrating system \ts{21}{8}{0:02}\mn{21}{18}. Earlier authors located its endogenous sources:
Hawtrey coined ``the inherent instability of credit'' for the British banking system (1919); Minsky, \emph{Stabilizing an Unstable Economy}, focused on business
credit; Adrian and Shin, ``Liquidity and Leverage'', on dealer balance sheets \ts{21}{8}{0:43}. Mehrling uses Treynor's model: dealers
create market liquidity by moving order-flow imbalances onto their own balance sheets, positions that must be funded \ts{21}{8}{1:24}.

\paragraph{Boom.} Asset managers pile up cash and seek credit risk exposure faster than capital funding banks absorb it; dealers absorb the imbalance, which pushes
money market rates down and risk premia down: exactly what makes setting up capital funding banks very profitable. Prices are distorted and more shadow banking is
built than matched-book equilibrium models would suggest \ts{21}{8}{2:05}.

\paragraph{Bust.} Once overbuilt, the asset managers' flows fall short of the shadow banks' funding needs; the dealers' imbalance reverses, rates and risk premia rise,
shadow banks look unprofitable; losses shrink the dealers' capital and financing limits, leaving excess exposure to liquidate, and in Treynor's model liquidation
happens at the outside spread: the central bank \ts{21}{8}{4:08}. A Hawtrey--Minsky story updated for
modern times, ``a starting place'' \ts{21}{8}{5:53}.

\begin{history}[Hawtrey and the instability of credit]
Ralph Hawtrey (1879--1975), a Treasury economist, argued in \emph{Currency and Credit} (1919) that credit is inherently unstable, a cumulative process that the
banking system must restrain. Hyman Minsky's \emph{Stabilizing an Unstable Economy} appeared in 1986.\par
\emph{References:} R.~G.~Hawtrey, \emph{Currency and Credit}, Longmans, 1919; H.~P.~Minsky, \emph{Stabilizing an Unstable Economy}, Yale UP, 1986; T.~Adrian and
H.~S.~Shin, ``Liquidity and Leverage'', \emph{Journal of Financial Intermediation} 19(3), 2010.
\end{history}

\begin{reading}[Adrian and Shin, ``Liquidity and Leverage'' (FRBNY Staff Report 328, 2008, rev.~2010)]
The course file (\file{materials/readings/L21-sr328.pdf}) is the dealer-balance-sheet version of the Hawtrey--Minsky story.
\begin{itemize}
  \item With balance sheets marked to market, asset price changes show up at once as changes in net worth, and intermediaries respond by adjusting the size of their balance sheets.
  \item For US security dealers and brokers, \emph{leverage is procyclical}: it is high in booms and low in busts. Leverage and total assets grow together, and the margin of adjustment is repo.
  \item Changes in dealer repos forecast changes in market risk, measured by innovations in the VIX.
  \item ``Aggregate liquidity can be seen as the rate of change of the aggregate balance sheet of the financial intermediaries.''
\end{itemize}
\end{reading}

% ------------------------------------------------------------------
\section{Normal times: collateral and payments}\label{sec:cbs-collateral}

In normal times the issue is the payment system: the survival constraint, ``liquidity kills you quick''. But besides payment flows there are collateral flows, the
subject of the paper by Singh and Aitken \ts{21}{9}{0:03}\mn{21}{25}.

\paragraph{The model with numbers.} The capital funding bank holds RMBS of 100, funded by 100 of secured money market funding, the collateral perhaps passing through
the money dealer to secure the asset manager's deposit; CDS and IRS are worth 0 at inception everywhere; the asset manager has plenty of deposits as cash collateral
\ts{21}{9}{1:06}. No counterparty risk by construction: the bank is exactly hedged, and the asset
manager's capital (the pensioners') can absorb even a total loss \ts{21}{9}{2:52}.

\paragraph{A temporary fall: collateral.} RMBS fall to 90, CDS rise to 10. Collateral must flow from the asset manager through the derivative dealer; and the money market
funding of 100 is now backed by only 90 of RMBS, so the CDS, or the collateral behind it, must be usable for the funding. Otherwise the asset manager may pull its deposits
from a money dealer lending against inadequate collateral, and a liquidity spiral can start from collateral problems alone
\ts{21}{9}{4:14}. The central bank can help as lender of last resort to the dealers, or
as dealer of last resort in RMBS or CDS, putting a floor under the collateral market \ts{21}{9}{6:20}.

\begin{definition}[New concepts: rehypothecation]
\emph{Rehypothecation}: reusing collateral received from one party as collateral for one's own borrowing (from \emph{hypothecation}, Greek \emph{hypoth\=ek\=e}, a
pledge) \ts{21}{9}{7:02}. Making it illegal, as some propose, would block the collateral flows through the dealers and push the work onto the central bank
\ts{21}{9}{7:23}.
\end{definition}

\begin{coursenote}
The paper is given in the captions as by ``Aiken Singh''. It is M.~Singh and J.~Aitken, ``The (sizable) role of rehypothecation in the shadow banking system'', IMF
Working Paper 10/172, July 2010 \ts{21}{9}{0:44}.
\end{coursenote}

\begin{reading}[Singh and Aitken (2010), IMF WP/10/172]
What the paper says (\file{materials/readings/L21-Singh_Aitken_wp10172.pdf}):
\begin{itemize}
  \item \textbf{The rule.} A hedge fund's collateral posted to its prime broker can be reused by the broker for its own funding.
  \item \textbf{US cap.} In the US, Regulation~T caps rehypothecation at 140\% of the customer's debit balance. Owing \$200, a customer lets the broker reuse up to \$280.
  \item \textbf{UK.} In the UK there is no cap and there are no customer protection rules, so the UK ``provides a platform for higher leveraging''.
  \item \textbf{Collapse after Lehman.} Collateral that the large US broker-dealers were permitted to repledge fell from about \$4.5 trillion at end-2007 to \$2.1 trillion at end-2009.
  \item \textbf{Churning factor.} About \$1 trillion of hedge fund securities was rehypothecated at end-2007. About 40\% of the \$10 trillion of pledgeable collateral received by the ten largest global banks came from hedge funds. The ``churning factor'', or re-use, of collateral is therefore about 4.
  \item \textbf{Conclusion.} Counting rehypothecation, the shadow banking system was ``at least 50 percent bigger than documented so far''.
\end{itemize}
\end{reading}

\paragraph{A permanent fall: payments.} If the loss is permanent, the system deleverages: RMBS 90, funding 90, CDS back to 0. The asset manager pays 10 to the derivative
dealer, who pays the capital funding bank, which repays 10 of money market funding: a circle of payments, and if someone says ``I'm not paying until I get paid'',
everything stops \ts{21}{9}{7:43}.

\begin{nfigure}
\begin{taccts}
\tacct[1.9cm]{Capital funding bank}{RMBS: 100, 90\\CDS: 0, 10, 0}{Funding: 100, 90}\quad
\tacct[1.9cm]{Derivative dealer}{CDS: 0, 10, 0}{CDS: 0, 10, 0}\quad
\tacct[1.9cm]{Asset manager}{Deposits: 100, 90}{Capital\\CDS: 0, 10, 0}
\end{taccts}
\caption{A fall in RMBS from 100 to 90: first collateral flows, then, if permanent, a circle of payments of 10 \ts{21}{9}{4:35}.}\label{fig:cbs-numbers}
\end{nfigure}

\paragraph{Outside spread, not inside.} Collateral flows and payment flows are both key to keeping normal fluctuations from becoming crises. Having become dealer of last
resort, the central bank's job, as in Bagehot's day, is to set up the system so that it is not called upon every day, because of moral hazard: ``I'm willing to be
outside spread, but not inside spread'' \ts{21}{9}{9:27}.

% ------------------------------------------------------------------
\section{Private backstop, and public}\label{sec:cbs-principles}

To keep the central bank from acting daily, make the dealers robust, especially the new one, the derivative dealer: matched-book dealers need liquidity reserves,
speculative (proprietary) dealers need capital to absorb losses \ts{21}{10}{0:04}.

\begin{keyformulas}
\textbf{Four ideas for regulating the shadow banking system of the future} (``they might change tomorrow'') \ts{21}{10}{1:27}\mn{21}{30}:
\begin{enumerate}
  \item focus on the dealers, money dealers and derivative dealers, not on the shadow banks;
  \item distinguish matched-book from proprietary (speculative) dealers;
  \item liquidity for the matched-book dealers, capital for the proprietary ones (``liquidity kills you quick'');
  \item the survival constraint concerns collateral flows as well as payment flows: a system of secured funding; a collateral call by Goldman and Soci\'et\'e G\'en\'erale
        killed AIG, perhaps not insolvent on fundamental valuations \ts{21}{10}{2:50}.
\end{enumerate}
\end{keyformulas}

``I think we're at a Bagehot moment, 1873. We're building it from scratch'' \ts{21}{10}{3:53}.

\paragraph{How to keep dealers from leaning on the Fed?} A student asks. First, adequate capital for the speculative dealers, so their financing limits do not shift
(and, since the line between matched-book and speculative dealing is not sharp, probably for matched-book dealers too) \ts{21}{10}{4:34}.
Second, the central bank should create the outside spread, not the inside spread: ``insure freely but at a high premium'', put a floor under capital markets ``at 80 cents
on the dollar'', uneconomic except in need. Some programs, the self-liquidating ones, did this; mostly the Fed played catch-up, forced to act on weekends
\ts{21}{10}{5:35}. The modern analogue of Bagehot's ``lend freely at a high rate against good security'' is the outside spread: be away
from the market, but act boldly when you are; make it a known principle, expensive for those who come, who must come ``begging''
\ts{21}{10}{6:36}. There will be another crisis: the inherent instability of credit \ts{21}{10}{7:58}.

% ------------------------------------------------------------------
\flushside
\section*{Key formulas and ideas}
\begin{keyformulas}
\begin{tabularx}{\linewidth}{@{}l L@{}}
Shadow banking & money market funding of capital market lending: global funding, market prices on both sides, dealers\\
Immature backstops & private liquidity puts (traditional banks), private solvency backstop (AIG CDS)\\
Ideal system & capital funding bank $+$ asset manager linked by a money dealer and a derivative dealer\\
Fed, Dec.\ 2011 & restated as CDS $+$ IRS $+$ funding swap: dealer of last resort in money and risk\\
Instability & order-flow imbalances on dealers' balance sheets: boom (low rates, low premia), bust (outside spread $=$ central bank)\\
Normal times & collateral flows and payment flows must not stop\\
Principles & regulate dealers; liquidity for matched book, capital for proprietary; central bank as outside spread\\
\end{tabularx}
\end{keyformulas}

\section*{Exercises}

\begin{exercise}[Liquidity put]
A sponsored vehicle holds \$50 billion of AAA RMBS funded by three-month ABCP, with a liquidity put from its parent bank. Write the balance sheets before and after the
ABCP fails to roll, and then after the bank borrows from the Fed.
\end{exercise}

\begin{exercise}[Restating the Fed]
Using the December 2011 figures, redo the restatement of the Fed's balance sheet and check the \$6.4 trillion of notional swap exposure. Which swap corresponds to QE in
mortgage securities and which to QE in Treasuries?
\end{exercise}

\begin{exercise}[Collateral and payments]
In the 100/90 example, write all balance sheets after the temporary fall, assuming collateral flows work, and after the permanent fall. Which payment, if refused, stops
the circle?
\end{exercise}

\begin{exercise}[Outside spread]
Compare Bagehot's rule with ``insure freely at a high premium''. Design a central bank backstop bid for AAA RMBS at a price that is an outside spread, and explain why it
should rarely be used.
\end{exercise}


\end{document}
